%T-14
Este formulario define los entregables del contrato, su diccionario, trazabilidad y programa de 56 meses. SD7 explica las decisiones y presenta la síntesis; T-15 desarrolla actividades, esfuerzo, recursos, nivelación y CPM/PERT; T-18 detalla implantación y cambio de autoridad. T-13 y T-17 contienen las pruebas y el protocolo formal de recepción. Cada referencia conecta esos documentos sin reproducir su desarrollo completo.

La fecha de inicio utilizada para el escenario es el 1 de diciembre de 2026. Los meses contractuales prevalecen; un inicio distinto exige trasladar fechas y revisar estaciones, calendarios y ventanas de faena. Las ventanas de las fichas son intervalos de planificación, no jornadas continuas ni recepciones ya obtenidas.

\Control{\textbf{Control de integración.} El programa y la carta Gantt se concilian con las ventanas del diccionario vigente, incluidas las dependencias de IN4 y la actualización de la transferencia final. Las barras por cuenta resumen ventanas de entregables; T-15 desarrolla y verifica las actividades que los producen, su esfuerzo, disponibilidad y red CPM/PERT. T-18 individualiza las olas y cortes y la secuencia de revisión/decisión. Las cantidades activas de formación, selección regional y referencias marcadas siguen pendientes de validación documental. El cierre de la línea base exige resolver esas entradas y demostrar la cabida de recursos y recepciones, sin alterar las fases contractuales.}

\section{EDT y cobertura del alcance}\label{sec:t14-edt}
\subsection{Criterio de descomposición y fronteras}
La raíz \texttt{1} representa la solución y el servicio del contrato. Las doce cuentas agrupan resultados; son una organización de Synaptix y no un número impuesto por las bases. Cada paquete terminal tiene entregable, frontera, responsable y criterio de aceptación. Los padres agregan resultados y no reciben esfuerzo adicional al de sus hijos.

Las capacidades funcionales de los núcleos se entregan una sola vez: especificación, lógica, persistencia, contratos, vistas, pruebas unitarias y documentación pertenecen al mismo resultado. Sus actividades se programan en T-15. Identidad, auditoría, canales, eventos, telemetría e intercambio de datos compartidos se construyen una vez; cada dominio incorpora solamente sus reglas y adaptación particular. Pruebas integrales y expedientes se localizan en \texttt{1.10}, sin imputar nuevamente las pruebas unitarias. Las innovaciones incluyen únicamente su aporte incremental.

E1, E2 y Operación son atributos de alcance y tiempo. Los nueve dominios conservan sus nombres y propiedad: Operación marítima, Escuela náutica e Infraestructura; Varadero y contratistas, Combustible y Ambiente y consumos; Clientes y contratos y Finanzas y cobranza; Auditoría y seguridad transversal. Los grupos RF de T-12 no crean dominios adicionales.

SD3 excluye otorgar zarpes oficiales, búsqueda y salvamento, emisión fiscal del ERP contable, concesión del restaurante y bloqueos físicos de embarcaciones en agua por morosidad sin autorización. Obras y adquisiciones atribuidas al cliente mantienen ese responsable; Synaptix especifica, coordina y verifica interfaces. La dependencia no autoriza ampliar el plazo ni acredita provisión recibida.

\FuenteTabla{Bases Administrativas, T-7/SD7 y T-14; pauta disponible, capítulo 7; SD3 v0.9 y T-12 v0.9.}
\subsection{Cuentas y resultados}
\begin{figure}[htbp]\centering
\begin{tikzpicture}[font=\fontsize{11}{13}\selectfont,
caja/.style={draw=SynLine,fill=SynSurface,text width=40mm,minimum height=15mm,align=center,inner sep=2mm}]
\node[caja,fill=SynNavy,text=white,text width=138mm] (r) at (0,1.7) {\textbf{1 · Solución y servicio del Puerto Deportivo}};

\node[caja] (n0) at (-5,-0.0) {\textbf{1.1}\\Dirección y control};
\node[caja] (n1) at (0,-0.0) {\textbf{1.2}\\Plataforma híbrida};
\node[caja] (n2) at (5,-0.0) {\textbf{1.3}\\Auditoría y seguridad};
\node[caja] (n3) at (-5,-1.95) {\textbf{1.4}\\Núcleo Operacional};
\node[caja] (n4) at (0,-1.95) {\textbf{1.5}\\Núcleo de Servicios};
\node[caja] (n5) at (5,-1.95) {\textbf{1.6}\\Núcleo Administrativo};
\node[caja] (n6) at (-5,-3.9) {\textbf{1.7}\\Canales, integración y datos};
\node[caja] (n7) at (0,-3.9) {\textbf{1.8}\\Migración e historia};
\node[caja] (n8) at (5,-3.9) {\textbf{1.9}\\Innovaciones};
\node[caja] (n9) at (-5,-5.85) {\textbf{1.10}\\Verificación y aceptación};
\node[caja] (n10) at (0,-5.85) {\textbf{1.11}\\Formación e implantación};
\node[caja] (n11) at (5,-5.85) {\textbf{1.12}\\Servicio y salida};
\draw[SynBlue] (r.south) -- (0,0.9);\draw[SynBlue] (-7.3,0.9)--(2.7,0.9);
\draw[SynBlue] (-7.3,0.9)--(-7.3,-5.85);
\draw[SynBlue] (-7.3,-0.0)--(n0.west);
\draw[SynBlue] (-7.3,-1.95)--(n3.west);
\draw[SynBlue] (-7.3,-3.9)--(n6.west);
\draw[SynBlue] (-7.3,-5.85)--(n9.west);
\draw[SynBlue] (-2.3,0.9)--(-2.3,-5.85);
\draw[SynBlue] (-2.3,-0.0)--(n1.west);
\draw[SynBlue] (-2.3,-1.95)--(n4.west);
\draw[SynBlue] (-2.3,-3.9)--(n7.west);
\draw[SynBlue] (-2.3,-5.85)--(n10.west);
\draw[SynBlue] (2.7,0.9)--(2.7,-5.85);
\draw[SynBlue] (2.7,-0.0)--(n2.west);
\draw[SynBlue] (2.7,-1.95)--(n5.west);
\draw[SynBlue] (2.7,-3.9)--(n8.west);
\draw[SynBlue] (2.7,-5.85)--(n11.west);
\end{tikzpicture}
\caption{Doce cuentas de entregables de la EDT}\label{fig:t14-edt}
\end{figure}
\begingroup\setlength{\tabcolsep}{3pt}
\begin{longtable}{L{\dimexpr0.250000\textwidth-2\tabcolsep\relax}L{\dimexpr0.650000\textwidth-2\tabcolsep\relax}L{\dimexpr0.100000\textwidth-2\tabcolsep\relax}}
\caption{Cuentas de entregables y paquetes terminales}\label{tab:t14-cuentas}\\\hline
\EncabezadoTabla{Cuenta} & \EncabezadoTabla{Resultado} & \EncabezadoTabla{N.º}\\\hline\endfirsthead
\multicolumn{3}{l}{\textit{Continuación de la tabla \ref{tab:t14-cuentas}}}\\\hline
\EncabezadoTabla{Cuenta} & \EncabezadoTabla{Resultado} & \EncabezadoTabla{N.º}\\\hline\endhead
\hline\endfoot
\hline\endlastfoot
\texttt{1.\allowbreak{}1}\newline{}\textbf{Dirección y control} & Línea base, decisiones, riesgos, cambios y custodia documental. & 56\\\hline
\texttt{1.\allowbreak{}2}\newline{}\textbf{Plataforma híbrida} & Infraestructura híbrida, cinco ambientes, borde, conectividad y observabilidad. & 28\\\hline
\texttt{1.\allowbreak{}3}\newline{}\textbf{Auditoría y seguridad} & Identidad, permisos, autoridad local, cifrado, consentimiento y auditoría. & 11\\\hline
\texttt{1.\allowbreak{}4}\newline{}\textbf{Núcleo Operacional} & Operación marítima, Escuela náutica e Infraestructura; capacidades funcionales únicas. & 10\\\hline
\texttt{1.\allowbreak{}5}\newline{}\textbf{Núcleo de Servicios} & Varadero/contratistas, Combustible y Ambiente/consumos; capacidades funcionales únicas. & 12\\\hline
\texttt{1.\allowbreak{}6}\newline{}\textbf{Núcleo Administrativo} & Clientes/contratos y Finanzas/cobranza; capacidades funcionales únicas. & 8\\\hline
\texttt{1.\allowbreak{}7}\newline{}\textbf{Canales, integración y datos compartidos} & Canales, contratos, diseño UX y arquitectura integrada; telemetría, interfaces, ANAL-5, MASS-5, ENV-5 y actualidad. & 39\\\hline
\texttt{1.\allowbreak{}8}\newline{}\textbf{Migración e historia} & Diagnóstico, mitigación del legado mes 6, transformación, saneamiento, cargas y consulta histórica. & 103\\\hline
\texttt{1.\allowbreak{}9}\newline{}\textbf{Innovaciones} & Cinco aportes incrementales, pruebas, pilotos y evaluación de beneficio. & 89\\\hline
\texttt{1.\allowbreak{}10}\newline{}\textbf{Verificación y aceptación} & Ensayos, pruebas integrales y expedientes de recepción; excluye pruebas unitarias ya imputadas. & 48\\\hline
\texttt{1.\allowbreak{}11}\newline{}\textbf{Formación e implantación} & Formación, activación, convivencia, acompañamiento/estabilización por etapa y soporte de implementación. & $61+N_F$\\\hline
\texttt{1.\allowbreak{}12}\newline{}\textbf{Servicio y salida} & Prestación de atención/guardia, mantenimiento, ingeniería, formación estacional, controles, mentoría y salida. & $695+N_E$\\\hline
\end{longtable}\endgroup
En la cuenta 1.11, $N_F=\sum_{e=1}^{2}\sum_{p=1}^{5}N_{p,e}$ suma las cohortes iniciales activadas. En la cuenta 1.12, $N_E=\sum_{m}\sum_{p=1}^{5}N^{est}_{p,m}$ suma las cohortes estacionales efectivamente activadas durante el contrato. El diccionario contiene 349 fichas y 1.449 códigos nominales únicos: 214 reservas iniciales y 75 estacionales no constituyen ejecución confirmada. El universo activo es $1.160+N_F+N_E$, sujeto a manifiestos; son resultados de proyecto y períodos de prestación, no un conteo de personas o turnos. Si la demanda supera los sufijos disponibles se amplía el registro versionado sin limitar la formación. Los 560 códigos diarios anteriores se conservan como referencias históricas de actividad/evidencia de cuatro paquetes agregados; no se vuelven a contar como entregables terminales.

\textbf{Línea base provisional de cantidades de formación.} No se dispone en las fuentes de esta versión de una nómina aprobada por perfil, turno y sitio que permita fijar $N_{p,e}$. Los valores quedan pendientes de validación; no equivalen a cero ni a todos los códigos reservados. José Lara Arce solicita la población a los responsables del cliente, calcula y versiona el manifiesto; Ignacio Silva incorpora su estado a la línea base y la contraparte valida universo y calendario. La solicitud y dimensionamiento preliminar se controlan en H1 y se actualizan para E2 en H8. Antes de cerrar cantidades y HH de formación en T-15 debe recibirse el manifiesto; el plan de ejecución se confirma en 1.11.3.1.1/1.11.3.2.1 antes de activar las cohortes. Una cantidad todavía no validada mantiene ese cierre pendiente.

Para cada perfil y etapa se parte la población en grupos compatibles por turno, sitio y modalidad. Si $P_{p,e,g}$ es la cantidad de personas del grupo compatible $g$ y $K_{p,e,g}$ su capacidad efectiva, la cantidad es $N_{p,e}=\sum_g\lceil P_{p,e,g}/K_{p,e,g}\rceil$, con $1\leq K_{p,e,g}\leq12$. El máximo de doce asistentes y las dos sesiones de dos horas son decisiones de diseño de Synaptix, sujetas a contraste de disponibilidad; no se presentan como cifras obligatorias de las bases. Una persona se cuenta una sola vez por perfil, etapa y objetivo formativo; su participación justificada en otro perfil o incremento tiene registro distinto. Recuperaciones se controlan aparte, sin crear una segunda cohorte para inflar cobertura. FALTA REFERENCIA A SD3; FALTA REFERENCIA A SD10.

\begingroup\setlength{\tabcolsep}{3pt}
\begin{longtable}{L{\dimexpr.24\textwidth-2\tabcolsep\relax}L{\dimexpr.36\textwidth-2\tabcolsep\relax}L{\dimexpr.40\textwidth-2\tabcolsep\relax}}
\caption{Control de activación de cohortes; cantidad provisional}\label{tab:t14-cohortes}\\\hline
\EncabezadoTabla{Perfil} & \EncabezadoTabla{Códigos reservados por etapa} & \EncabezadoTabla{Cantidad activa y fuente}\\\hline\endfirsthead
\multicolumn{3}{l}{\textit{Continuación de la tabla \ref{tab:t14-cohortes}}}\\\hline
\EncabezadoTabla{Perfil} & \EncabezadoTabla{Códigos reservados por etapa} & \EncabezadoTabla{Cantidad activa y fuente}\\\hline\endhead
\hline\endfoot
\hline\endlastfoot
1: Personas usuarias finales & \texttt{1.11.7.e.1.1.\{1--100\}} & $N_{1,E1}$ y $N_{1,E2}$ pendientes; población nominal, turno, sitio y modalidad.\\\hline
2: Personas usuarias avanzadas & \texttt{1.11.7.e.2.1.\{1,2\}} & $N_{2,E1}$ y $N_{2,E2}$ pendientes; referentes funcionales designados y cobertura de turnos.\\\hline
3: Administración & \texttt{1.11.7.e.3.1.1} & $N_{3,E1}$ y $N_{3,E2}$ pendientes; administradores/referentes autorizados y suplencias.\\\hline
4: Equipo técnico & \texttt{1.11.7.e.4.1.1} & $N_{4,E1}$ y $N_{4,E2}$ pendientes; personas que reciben transferencia, sin presumir un área TI del cliente.\\\hline
5: Soporte niveles 1 y 2 & \texttt{1.11.7.e.5.1.\{1--3\}} & $N_{5,E1}$ y $N_{5,E2}$ pendientes; dotación asignada y turnos de soporte.\\\hline
\end{longtable}\endgroup
En los códigos de la tabla, \texttt{e=1} identifica E1 y \texttt{e=2} E2. Son reservas de identificación; la cantidad de sufijos no es una estimación de cohortes. Cada manifiesto identifica documento de origen, versión, estado de validación, cálculo por perfil y códigos activados, con responsable y fecha. Se enlaza desde esta tabla al registro recibido de 1.11.3.1.1/1.11.3.2.1; no se incorporan aquí nóminas personales ni el calendario detallado de T-18.

\textbf{Control de formación estacional.} Los 75 sufijos de \texttt{1.12.10} son reservas de tres campañas de referencia, no cohortes aprobadas. Se calcula $N^{est}_{p,m}=\sum_g\lceil P^{est}_{p,m,g}/K^{est}_{p,m,g}\rceil$, con población nueva confirmada y grupos compatibles; $1\leq K^{est}_{p,m,g}\leq12$. Las 300 plazas anuales ofrecidas representan capacidad, no un reparto fijo de personas por perfil. José Lara Arce obtiene y versiona el manifiesto de cada campaña; la contraparte valida población/turnos/calendario antes de fijar cantidades/HH y activar códigos. El perfil 4 no presume un área TI del cliente. Recuperaciones y altas fuera de las tres campañas se identifican sin duplicar asistentes y sin excluirlas del derecho de formación. Las jornadas de actualización de \texttt{1.12.6} son distintas de la incorporación de personal nuevo. FALTA REFERENCIA A SD10; FALTA REFERENCIA A SD7.

\textbf{Unidad de implantación.} \texttt{1.11.4.1/2} recibe el acompañamiento de E1/E2; \texttt{1.11.6.1/2} recibe su estabilización. Cada manifiesto de olas identifica objetivo, población, sitio, versión y grupo de datos; los tres grupos de \texttt{1.11.1} no determinan por sí solos su cantidad. Las cuatro posiciones históricas por etapa dejan de fijar cuatro olas ejecutadas. T-18 define su universo y secuencia antes de cerrar T-15. La agrupación conserva el trabajo y evidencia diaria, sin reducir horas por cambiar códigos. Soporte continuo de implementación se recibe por período en \texttt{1.11.8}; desde mes 21 la prestación de ambos alcances se recibe en \texttt{1.12.14}.

\subsection{Correspondencia con el catálogo anterior}
La revisión reúne 306 registros funcionales anteriores en 30 entregables. Los 216 suplementos y las actividades de especificación, implementación y vistas no se eliminan como trabajo: se enlazan con su resultado único y se concilian una sola vez en T-15. El código terminado en \texttt{.1} se conserva como identificador del entregable completo; cambia su significado respecto de la ficha anterior de especificación. Los demás códigos se conservan como identificadores históricos de actividad, no como paquetes terminales adicionales de T-14. Esta correspondencia debe propagarse a T-12, T-15, T-18 y a sus redes y registros.

\begingroup\setlength{\tabcolsep}{3pt}
\begin{longtable}{L{\dimexpr0.420000\textwidth-2\tabcolsep\relax}L{\dimexpr0.580000\textwidth-2\tabcolsep\relax}}
\caption{Correspondencia de resultados funcionales}\label{tab:t14-equivalencias}\\\hline
\EncabezadoTabla{Entregable actual} & \EncabezadoTabla{Registros anteriores asociados}\\\hline\endfirsthead
\multicolumn{2}{l}{\textit{Continuación de la tabla \ref{tab:t14-equivalencias}}}\\\hline
\EncabezadoTabla{Entregable actual} & \EncabezadoTabla{Registros anteriores asociados}\\\hline\endhead
\hline\endfoot
\hline\endlastfoot
\texttt{1.\allowbreak{}4.\allowbreak{}1.\allowbreak{}1.\allowbreak{}1}\newline{}Movimientos de salida y regreso identificados & \texttt{1.\allowbreak{}4.\allowbreak{}1.\allowbreak{}1.\allowbreak{}\{\{1--3\};\allowbreak{}\{4,\allowbreak{}5\}.\allowbreak{}\{1--4\};\allowbreak{}\{6,\allowbreak{}7\}.\allowbreak{}\{1,\allowbreak{}2\}\}}\\\hline
\texttt{1.\allowbreak{}4.\allowbreak{}1.\allowbreak{}2.\allowbreak{}1}\newline{}Planes voluntarios y discrepancias operacionales & \texttt{1.\allowbreak{}4.\allowbreak{}1.\allowbreak{}2.\allowbreak{}\{\{1--3\};\allowbreak{}4.\allowbreak{}\{1--3\};\allowbreak{}\{5--7\}.\allowbreak{}1\}}\\\hline
\texttt{1.\allowbreak{}4.\allowbreak{}1.\allowbreak{}3.\allowbreak{}1}\newline{}Atención y escalamiento de planes vencidos & \texttt{1.\allowbreak{}4.\allowbreak{}1.\allowbreak{}3.\allowbreak{}\{\{1--3\};\allowbreak{}4.\allowbreak{}\{1--3\};\allowbreak{}\{5--7\}.\allowbreak{}1\}}\\\hline
\texttt{1.\allowbreak{}4.\allowbreak{}1.\allowbreak{}4.\allowbreak{}1}\newline{}Autorización operacional de amarra y compatibilidad & \texttt{1.\allowbreak{}4.\allowbreak{}1.\allowbreak{}4.\allowbreak{}\{\{1--3\};\allowbreak{}4.\allowbreak{}\{1--3\};\allowbreak{}\{5--7\}.\allowbreak{}1\}}\\\hline
\texttt{1.\allowbreak{}4.\allowbreak{}2.\allowbreak{}1.\allowbreak{}1}\newline{}Clases, nóminas y participación por intervalo & \texttt{1.\allowbreak{}4.\allowbreak{}2.\allowbreak{}1.\allowbreak{}\{\{1--3\};\allowbreak{}\{4,\allowbreak{}5\}.\allowbreak{}\{1--4\};\allowbreak{}\{6,\allowbreak{}7\}.\allowbreak{}\{1,\allowbreak{}2\}\}}\\\hline
\texttt{1.\allowbreak{}4.\allowbreak{}2.\allowbreak{}2.\allowbreak{}1}\newline{}Retiros coordinados y entrega única de alumno & \texttt{1.\allowbreak{}4.\allowbreak{}2.\allowbreak{}2.\allowbreak{}\{\{1--3\};\allowbreak{}4.\allowbreak{}\{1--3\};\allowbreak{}\{5--7\}.\allowbreak{}1\}}\\\hline
\texttt{1.\allowbreak{}4.\allowbreak{}2.\allowbreak{}3.\allowbreak{}1}\newline{}Retorno de participantes y vista mínima de salud & \texttt{1.\allowbreak{}4.\allowbreak{}2.\allowbreak{}3.\allowbreak{}\{\{1--3\};\allowbreak{}\{4,\allowbreak{}5\}.\allowbreak{}\{1--4\};\allowbreak{}\{6,\allowbreak{}7\}.\allowbreak{}\{1,\allowbreak{}2\}\}}\\\hline
\texttt{1.\allowbreak{}4.\allowbreak{}3.\allowbreak{}1.\allowbreak{}1}\newline{}Inventario y condiciones de puestos y fondeos & \texttt{1.\allowbreak{}4.\allowbreak{}3.\allowbreak{}1.\allowbreak{}\{\{1--3\};\allowbreak{}4.\allowbreak{}\{1--3\};\allowbreak{}\{5--7\}.\allowbreak{}1\}}\\\hline
\texttt{1.\allowbreak{}4.\allowbreak{}3.\allowbreak{}2.\allowbreak{}1}\newline{}Inspecciones de boyas y evidencia desconectada & \texttt{1.\allowbreak{}4.\allowbreak{}3.\allowbreak{}2.\allowbreak{}\{\{1--3\};\allowbreak{}\{4,\allowbreak{}5\}.\allowbreak{}\{1--4\};\allowbreak{}\{6,\allowbreak{}7\}.\allowbreak{}\{1,\allowbreak{}2\}\}}\\\hline
\texttt{1.\allowbreak{}4.\allowbreak{}3.\allowbreak{}3.\allowbreak{}1}\newline{}Asociaciones temporales de medidores e instalaciones & \texttt{1.\allowbreak{}4.\allowbreak{}3.\allowbreak{}3.\allowbreak{}\{\{1--3\};\allowbreak{}4.\allowbreak{}\{1--3\};\allowbreak{}\{5--7\}.\allowbreak{}1\}}\\\hline
\texttt{1.\allowbreak{}5.\allowbreak{}1.\allowbreak{}1.\allowbreak{}1}\newline{}Maestros de faena y habilitaciones de consulta inicial & \texttt{1.\allowbreak{}5.\allowbreak{}1.\allowbreak{}1.\allowbreak{}\{\{1--3\};\allowbreak{}4.\allowbreak{}\{1--3\};\allowbreak{}\{5--7\}.\allowbreak{}1\}}\\\hline
\texttt{1.\allowbreak{}5.\allowbreak{}1.\allowbreak{}2.\allowbreak{}1}\newline{}Agenda de varadero y posiciones compatibles & \texttt{1.\allowbreak{}5.\allowbreak{}1.\allowbreak{}2.\allowbreak{}\{\{1--3\};\allowbreak{}4.\allowbreak{}\{1--3\};\allowbreak{}\{5--7\}.\allowbreak{}1\}}\\\hline
\texttt{1.\allowbreak{}5.\allowbreak{}1.\allowbreak{}3.\allowbreak{}1}\newline{}Plan de izaje, validación y suspensión de faena & \texttt{1.\allowbreak{}5.\allowbreak{}1.\allowbreak{}3.\allowbreak{}\{\{1--3\};\allowbreak{}4.\allowbreak{}\{1--3\};\allowbreak{}\{5--7\}.\allowbreak{}1\}}\\\hline
\texttt{1.\allowbreak{}5.\allowbreak{}1.\allowbreak{}4.\allowbreak{}1}\newline{}Órdenes, ejecución y reprogramación de trabajos & \texttt{1.\allowbreak{}5.\allowbreak{}1.\allowbreak{}4.\allowbreak{}\{\{1--3\};\allowbreak{}4.\allowbreak{}\{1--3\};\allowbreak{}\{5--7\}.\allowbreak{}1\}}\\\hline
\texttt{1.\allowbreak{}5.\allowbreak{}1.\allowbreak{}5.\allowbreak{}1}\newline{}Acreditación, inducción y acceso a intervención & \texttt{1.\allowbreak{}5.\allowbreak{}1.\allowbreak{}5.\allowbreak{}\{\{1--3\};\allowbreak{}4.\allowbreak{}\{1--3\};\allowbreak{}\{5--7\}.\allowbreak{}1\}}\\\hline
\texttt{1.\allowbreak{}5.\allowbreak{}2.\allowbreak{}1.\allowbreak{}1}\newline{}Despacho identificado y captura segura de combustible & \texttt{1.\allowbreak{}5.\allowbreak{}2.\allowbreak{}1.\allowbreak{}\{\{1--3\};\allowbreak{}4.\allowbreak{}\{1--3\};\allowbreak{}\{5--7\}.\allowbreak{}1\}}\\\hline
\texttt{1.\allowbreak{}5.\allowbreak{}2.\allowbreak{}2.\allowbreak{}1}\newline{}Conciliación de tramos y rectificación de despachos & \texttt{1.\allowbreak{}5.\allowbreak{}2.\allowbreak{}2.\allowbreak{}\{\{1--3\};\allowbreak{}4.\allowbreak{}\{1--3\};\allowbreak{}\{5--7\}.\allowbreak{}1\}}\\\hline
\texttt{1.\allowbreak{}5.\allowbreak{}3.\allowbreak{}1.\allowbreak{}1}\newline{}Intervalos y series conciliadas de agua y energía & \texttt{1.\allowbreak{}5.\allowbreak{}3.\allowbreak{}1.\allowbreak{}\{\{1--3\};\allowbreak{}4.\allowbreak{}\{1--3\};\allowbreak{}\{5--7\}.\allowbreak{}1\}}\\\hline
\texttt{1.\allowbreak{}5.\allowbreak{}3.\allowbreak{}2.\allowbreak{}1}\newline{}Consultas de consumo y evidencia ambiental inicial & \texttt{1.\allowbreak{}5.\allowbreak{}3.\allowbreak{}2.\allowbreak{}\{\{1--3\};\allowbreak{}4.\allowbreak{}\{1--3\};\allowbreak{}\{5--7\}.\allowbreak{}1\}}\\\hline
\texttt{1.\allowbreak{}5.\allowbreak{}3.\allowbreak{}3.\allowbreak{}1}\newline{}Recepción de residuos y discrepancias de origen & \texttt{1.\allowbreak{}5.\allowbreak{}3.\allowbreak{}3.\allowbreak{}\{\{1--3\};\allowbreak{}4.\allowbreak{}\{1--3\};\allowbreak{}\{5--7\}.\allowbreak{}1\}}\\\hline
\texttt{1.\allowbreak{}5.\allowbreak{}3.\allowbreak{}4.\allowbreak{}1}\newline{}Entregas RESPEL parciales y certificados vinculados & \texttt{1.\allowbreak{}5.\allowbreak{}3.\allowbreak{}4.\allowbreak{}\{\{1--3\};\allowbreak{}\{4,\allowbreak{}5\}.\allowbreak{}\{1--4\};\allowbreak{}\{6,\allowbreak{}7\}.\allowbreak{}\{1,\allowbreak{}2\}\}}\\\hline
\texttt{1.\allowbreak{}5.\allowbreak{}3.\allowbreak{}5.\allowbreak{}1}\newline{}Indicadores ambientales e imputación comercial de consumo & \texttt{1.\allowbreak{}5.\allowbreak{}3.\allowbreak{}5.\allowbreak{}\{\{1--3\};\allowbreak{}4.\allowbreak{}\{1--3\};\allowbreak{}\{5--7\}.\allowbreak{}1\}}\\\hline
\texttt{1.\allowbreak{}6.\allowbreak{}1.\allowbreak{}1.\allowbreak{}1}\newline{}Maestros de personas, organizaciones y embarcaciones & \texttt{1.\allowbreak{}6.\allowbreak{}1.\allowbreak{}1.\allowbreak{}\{\{1--3\};\allowbreak{}4.\allowbreak{}\{1--3\};\allowbreak{}\{5--7\}.\allowbreak{}1\}}\\\hline
\texttt{1.\allowbreak{}6.\allowbreak{}1.\allowbreak{}2.\allowbreak{}1}\newline{}Contratos y documentos vigentes de consulta E1 & \texttt{1.\allowbreak{}6.\allowbreak{}1.\allowbreak{}2.\allowbreak{}\{\{1--3\};\allowbreak{}4.\allowbreak{}\{1--3\};\allowbreak{}\{5--7\}.\allowbreak{}1\}}\\\hline
\texttt{1.\allowbreak{}6.\allowbreak{}1.\allowbreak{}3.\allowbreak{}1}\newline{}Atención a visitas y asignación transitoria inicial & \texttt{1.\allowbreak{}6.\allowbreak{}1.\allowbreak{}3.\allowbreak{}\{\{1--3\};\allowbreak{}4.\allowbreak{}\{1--3\};\allowbreak{}\{5--7\}.\allowbreak{}1\}}\\\hline
\texttt{1.\allowbreak{}6.\allowbreak{}1.\allowbreak{}4.\allowbreak{}1}\newline{}Espera y gestión contractual completa & \texttt{1.\allowbreak{}6.\allowbreak{}1.\allowbreak{}4.\allowbreak{}\{\{1--3\};\allowbreak{}\{4,\allowbreak{}5\}.\allowbreak{}\{1--4\};\allowbreak{}\{6,\allowbreak{}7\}.\allowbreak{}\{1,\allowbreak{}2\}\}}\\\hline
\texttt{1.\allowbreak{}6.\allowbreak{}2.\allowbreak{}1.\allowbreak{}1}\newline{}Hechos facturables, tarifas y lotes de envío & \texttt{1.\allowbreak{}6.\allowbreak{}2.\allowbreak{}1.\allowbreak{}\{\{1--3\};\allowbreak{}4.\allowbreak{}\{1--3\};\allowbreak{}\{5--7\}.\allowbreak{}1\}}\\\hline
\texttt{1.\allowbreak{}6.\allowbreak{}2.\allowbreak{}2.\allowbreak{}1}\newline{}Respuesta contable, pagos y saldos informados & \texttt{1.\allowbreak{}6.\allowbreak{}2.\allowbreak{}2.\allowbreak{}\{\{1--3\};\allowbreak{}4.\allowbreak{}\{1--3\};\allowbreak{}\{5--7\}.\allowbreak{}1\}}\\\hline
\texttt{1.\allowbreak{}6.\allowbreak{}2.\allowbreak{}3.\allowbreak{}1}\newline{}Cobranza y medidas con decisión competente & \texttt{1.\allowbreak{}6.\allowbreak{}2.\allowbreak{}3.\allowbreak{}\{\{1--3\};\allowbreak{}4.\allowbreak{}\{1--3\};\allowbreak{}\{5--7\}.\allowbreak{}1\}}\\\hline
\texttt{1.\allowbreak{}6.\allowbreak{}2.\allowbreak{}4.\allowbreak{}1}\newline{}Expedientes de custodia y exportación íntegra & \texttt{1.\allowbreak{}6.\allowbreak{}2.\allowbreak{}4.\allowbreak{}\{\{1--3\};\allowbreak{}4.\allowbreak{}\{1--3\};\allowbreak{}\{5--7\}.\allowbreak{}1\}}\\\hline
\end{longtable}\endgroup
El registro de estimación anterior sumaba 4.200 HH funcionales originales y 8.508 HH suplementarias, equivalentes a 12.708 HH. Es una conciliación aritmética histórica, no la estimación aprobada del modelo agrupado. T-15 revisará fronteras, actividades y método UCP antes de cerrar cifras. El T-14 no añade HH a los padres ni asigna un esfuerzo ficticio a los resultados nuevos de datos y legado. La revisión de esfuerzo puede modificar recursos y ventanas internas, conservando obligaciones y meses contractuales.

\begingroup\setlength{\tabcolsep}{3pt}
\begin{longtable}{L{\dimexpr.52\textwidth-2\tabcolsep\relax}L{\dimexpr.48\textwidth-2\tabcolsep\relax}}
\caption{Correspondencia de registros diarios y paquetes agregados}\label{tab:t14-diarios}\\\hline
\EncabezadoTabla{Registros históricos de actividad/evidencia} & \EncabezadoTabla{Paquete terminal vigente}\\\hline\endfirsthead
\multicolumn{2}{l}{\textit{Continuación de la tabla \ref{tab:t14-diarios}}}\\\hline
\EncabezadoTabla{Registros históricos de actividad/evidencia} & \EncabezadoTabla{Paquete terminal vigente}\\\hline\endhead
\hline\endfoot
\hline\endlastfoot
\texttt{1.\allowbreak{}11.\allowbreak{}4.\allowbreak{}1.\allowbreak{}\{1--4\}.\allowbreak{}\{1--28\}.\allowbreak{}\{1,\allowbreak{}2\}} & \texttt{1.\allowbreak{}11.\allowbreak{}4.\allowbreak{}1} Acompañamiento E1\\\hline
\texttt{1.\allowbreak{}11.\allowbreak{}6.\allowbreak{}1.\allowbreak{}\{1--28\}.\allowbreak{}\{1,\allowbreak{}2\}} & \texttt{1.\allowbreak{}11.\allowbreak{}6.\allowbreak{}1} Estabilización E1\\\hline
\texttt{1.\allowbreak{}11.\allowbreak{}4.\allowbreak{}2.\allowbreak{}\{1--4\}.\allowbreak{}\{1--28\}.\allowbreak{}\{1,\allowbreak{}2\}} & \texttt{1.\allowbreak{}11.\allowbreak{}4.\allowbreak{}2} Acompañamiento E2\\\hline
\texttt{1.\allowbreak{}11.\allowbreak{}6.\allowbreak{}2.\allowbreak{}\{1--28\}.\allowbreak{}\{1,\allowbreak{}2\}} & \texttt{1.\allowbreak{}11.\allowbreak{}6.\allowbreak{}2} Estabilización E2\\\hline
\end{longtable}\endgroup
Los códigos diarios se relacionan con un solo paquete. Turnos, olas activadas, zonas y evidencia se conservan en registros de ejecución; nuevas olas se identifican en el manifiesto, sin reutilizar códigos de otra población. Esta correspondencia y los nuevos productores de prestación deben propagarse a T-12, T-15, T-18, RACI y registros de avance.

\section{Diccionario de entregables y familias}\label{sec:t14-dic}
\subsection{Reglas de aceptación y lectura}
Cada ficha identifica códigos, resultado, responsable técnico, ventana, frontera, aceptación, insumos y condiciones. El responsable integra la entrega; no equivale al único ejecutor ni acredita contratación o dedicación exclusiva. T-15 distribuye ejecución, validación y revisión por perfil. La recepción contractual corresponde a la Contraparte Técnica autorizada; la firma de un responsable interno no la sustituye.

Todo resultado conserva código, versión, origen, universo, evidencia, revisión y observaciones. Fuentes y configuración permiten reproducción; informes identifican período, población y método. Condiciones comunes y específicas se aplican conjuntamente. Una decisión bloqueante, un insumo no recibido o una prueba pendiente impiden declarar el resultado recibido. HH consumidas no constituyen aceptación.

Los rangos enteros son inclusivos; los conjuntos representan exactamente sus códigos. Cada instancia tiene registro nominal, ventana y decisión propios. Se distingue aceptación individual del cierre agregado: recibir un lote, cohorte o período no recibe los restantes. Las reservas se activan contra trabajo autorizado; el saldo sin usar no es un resultado ejecutado.

Las actividades de análisis usan alcance, requisitos y antecedentes autorizados. Diseño, contratos y ambientes recibidos son requisitos de las actividades que los consumen, no insumos anticipados del análisis que debe producirlos. En cada resultado funcional, T-15 conserva el orden análisis/diseño, construcción, verificación y documentación, con sus esperas y recursos efectivos.

Las capacidades funcionales reciben también sus adaptaciones de administración, parámetros y flujos, documentos, búsqueda y exportación cuando la cláusula y el alcance de su dominio las exigen. La consola común de 1.7.3.3 produce la interfaz reutilizable; cada resultado propietario de 1.4--1.6 produce y verifica su lógica y sus datos. 1.7.4.1 produce mensajería, 1.3.3.1 conserva auditoría y 1.7.6.6 produce firma/evidencia, sin reconstruir esos motores en cada dominio. Los eventos de consultas sensibles y exportaciones se comprueban en el dominio emisor y en el registro receptor. La versión de diseño UX recibida precede la construcción definitiva; UAT y WCAG posteriores no sustituyen investigación ni validación previa.

Las referencias pendientes se identifican en mayúsculas, con el número del subdocumento correspondiente; no son evidencia de cumplimiento. Se sustituyen por documento, versión y apartado efectivos al integrar las fuentes. Los criterios particulares se enlazan con la cláusula y prueba de T-13/T-17, sin copiar esos protocolos completos.
\subsection{Fichas por cuenta}
\begingroup\setlength{\tabcolsep}{3pt}
\begin{longtable}{L{\dimexpr0.320000\textwidth-2\tabcolsep\relax}L{\dimexpr0.680000\textwidth-2\tabcolsep\relax}}
\caption{Diccionario: Dirección y control}\label{tab:t14-dic-1-1}\\\hline
\EncabezadoTabla{Paquete y resultado} & \EncabezadoTabla{Frontera, aceptación e insumos}\\\hline\endfirsthead
\multicolumn{2}{l}{\textit{Continuación de la tabla \ref{tab:t14-dic-1-1}}}\\\hline
\EncabezadoTabla{Paquete y resultado} & \EncabezadoTabla{Frontera, aceptación e insumos}\\\hline\endhead
\hline\endfoot
\hline\endlastfoot
\texttt{1.\allowbreak{}1.\allowbreak{}1.\allowbreak{}1}\newline{}\textbf{Acta de arranque, interlocutores y decisiones}\newline{}Ignacio Silva\newline{}1 instancia.\newline{}\textbf{Ventana:} 1. & \textbf{Contenido y frontera.} Responsabilidades, reuniones y habilitaciones; excluye ejecución de obra y construcción.\newline{}\textbf{Aceptación.} Interlocutores, autoridad, calendario de decisión y pendientes quedan registrados. Ignacio Silva responde por el paquete.\newline{}\textbf{Condiciones.} Predecesoras del catálogo: Inicio. Referencia: BA gobierno; SD1/SD2. Las habilitaciones técnicas se contrastan con arquitectura e inventario vigentes.\\\hline
\texttt{1.\allowbreak{}1.\allowbreak{}1.\allowbreak{}2}\newline{}\textbf{Línea base de alcance y trazabilidad}\newline{}Ignacio Silva\newline{}1 instancia.\newline{}\textbf{Ventana:} 1. & \textbf{Contenido y frontera.} Vinculación alcance-requisitos-entregables y control de cambios; excluye pruebas ejecutadas.\newline{}\textbf{Aceptación.} Se conservan IDs/enunciados y decisiones de cobertura; aprobación H1 se documenta sin convertir pendientes en conformidad. Ignacio Silva responde por el paquete.\newline{}\textbf{Condiciones.} Predecesoras del catálogo: 1.1.1.1. Referencia: E-25 H1; RT-19.01. Las habilitaciones técnicas se contrastan con arquitectura e inventario vigentes.\\\hline
\texttt{1.\allowbreak{}1.\allowbreak{}1.\allowbreak{}3}\newline{}\textbf{Definición de terminado y recepción}\newline{}Davor Serey\newline{}1 instancia.\newline{}\textbf{Ventana:} 2. & \textbf{Contenido y frontera.} Criterios código, pruebas, documentación, seguridad y despliegue por entregable.\newline{}\textbf{Aceptación.} Criterios medibles y protocolo T-17 tienen responsable y tratamiento de observaciones. Davor Serey responde por el paquete.\newline{}\textbf{Condiciones.} Predecesoras del catálogo: 1.1.1.2. Referencia: RT-20.07/08; BA18. Las habilitaciones técnicas se contrastan con arquitectura e inventario vigentes.\\\hline
\texttt{1.\allowbreak{}1.\allowbreak{}2.\allowbreak{}1}\newline{}\textbf{Registro de riesgos y respuestas presupuestadas}\newline{}Ignacio Silva\newline{}1 instancia.\newline{}\textbf{Ventana:} 2. & \textbf{Contenido y frontera.} Riesgos, disparadores y respuestas con cuenta de costo; excluye sumar otra vez acciones ya estimadas.\newline{}\textbf{Aceptación.} Cada respuesta tiene responsable, recurso, disparador, cuenta y efecto de calendario. Ignacio Silva responde por el paquete.\newline{}\textbf{Condiciones.} Predecesoras del catálogo: 1.1.1.1. Referencia: SD2; RT-19.06. Las habilitaciones técnicas se contrastan con arquitectura e inventario vigentes.\\\hline
\texttt{1.\allowbreak{}1.\allowbreak{}2.\allowbreak{}2}\newline{}\textbf{Línea base integrada de plazo, costo y recursos}\newline{}Ignacio Silva\newline{}1 instancia.\newline{}\textbf{Ventana:} 4. & \textbf{Contenido y frontera.} Consolidación EDT, red, cargas, calendario y presupuesto; excluye aceptar sobrecargas como disponibilidad.\newline{}\textbf{Aceptación.} Cobertura, recursos, holguras y reservas están demostrados; si fallan, se registra rechazo de línea base y plan de corrección. Ignacio Silva responde por el paquete.\newline{}\textbf{Condiciones.} Predecesoras del catálogo: H1,H2. Referencia: T-14/T-15; E-25; RT-19.07. Las habilitaciones técnicas se contrastan con arquitectura e inventario vigentes.\\\hline
\texttt{1.\allowbreak{}1.\allowbreak{}3.\allowbreak{}1}\newline{}\textbf{Plan inicial de reversibilidad}\newline{}Ignacio Silva\newline{}1 instancia.\newline{}\textbf{Ventana:} 3. & \textbf{Contenido y frontera.} Inventario, formatos, custodia, acompañamiento y sustitución; excluye ejecutar traslado a proveedor alternativo.\newline{}\textbf{Aceptación.} Plan se presenta dentro de 90 días de inicio e identifica trabajo portable/no portable y transferencia sin cobro adicional. Ignacio Silva responde por el paquete.\newline{}\textbf{Condiciones.} Predecesoras del catálogo: 1.1.1.1. Referencia: BA77; SD4 reversibilidad. Las habilitaciones técnicas se contrastan con arquitectura e inventario vigentes.\\\hline
\texttt{1.\allowbreak{}1.\allowbreak{}3.\allowbreak{}2}\newline{}\textbf{Expediente de custodia independiente de fuentes}\newline{}Ignacio Silva\newline{}1 instancia.\newline{}\textbf{Ventana:} 6. & \textbf{Contenido y frontera.} Preparación de contrato, fuentes y condiciones de liberación con tercero independiente; excluye atribuir custodia ya contratada.\newline{}\textbf{Aceptación.} Custodio y mecanismo tienen contrato, ubicación, acceso, actualización semestral y liberación por insolvencia/incumplimiento verificables. Ignacio Silva responde por el paquete.\newline{}\textbf{Condiciones.} Predecesoras del catálogo: 1.2.1.1. Referencia: BA propiedad de fuentes; SD3 entregables. Las habilitaciones técnicas se contrastan con arquitectura e inventario vigentes.\\\hline
\texttt{1.\allowbreak{}1.\allowbreak{}3.\allowbreak{}3.\allowbreak{}12}\newline{}\textbf{Actualización semestral de custodia de fuentes}\newline{}Ignacio Silva\newline{}1 instancia.\newline{}\textbf{Ventana:} 12. & \textbf{Contenido y frontera.} Depósito de fuentes y manifiesto ante custodio independiente; excluye reconstruir fuentes no versionadas.\newline{}\textbf{Aceptación.} Custodio confirma versión, integridad y condiciones de acceso/liberación y se conserva comprobante. Ignacio Silva responde por el paquete.\newline{}\textbf{Condiciones.} Predecesoras del catálogo: 1.1.3.2. Referencia: BA custodia independiente; actualización semestral. Las habilitaciones técnicas se contrastan con arquitectura e inventario vigentes.\\\hline
\texttt{1.\allowbreak{}1.\allowbreak{}3.\allowbreak{}3.\allowbreak{}18}\newline{}\textbf{Actualización semestral de custodia de fuentes}\newline{}Ignacio Silva\newline{}1 instancia.\newline{}\textbf{Ventana:} 18. & \textbf{Contenido y frontera.} Resultado individual del período o conjunto indicado. Excluye resultados de otras cuentas y horas de contraparte, obras o adquisiciones del cliente.\newline{}\textbf{Aceptación.} Custodio confirma versión, integridad y condiciones de acceso/liberación y se conserva comprobante. Ignacio Silva responde por el paquete.\newline{}\textbf{Insumos.} Versión/registro recibido, acceso autorizado, herramientas, participantes y ventana del período.\newline{}\textbf{Condiciones.} Predecesoras de habilitación: 1.1.3.2.\\\hline
\texttt{1.\allowbreak{}1.\allowbreak{}3.\allowbreak{}4}\newline{}\textbf{Plan de reversibilidad actualizado durante implementación}\newline{}Ignacio Silva\newline{}1 instancia.\newline{}\textbf{Ventana:} 15. & \textbf{Contenido y frontera.} Actualiza dependencias, inventario y estrategia de transferencia; excluye ejecutar traslado completo.\newline{}\textbf{Aceptación.} La actualización anual conserva versión, custodios, formatos y recursos requeridos. Ignacio Silva responde por el paquete.\newline{}\textbf{Condiciones.} Predecesoras del catálogo: 1.1.3.1. Referencia: BA77; SD4 reversibilidad. Las habilitaciones técnicas se contrastan con arquitectura e inventario vigentes.\\\hline
\texttt{1.\allowbreak{}1.\allowbreak{}4.\allowbreak{}1.\allowbreak{}1}\newline{}\textbf{Expediente quincenal de control}\newline{}Ignacio Silva\newline{}1 instancia.\newline{}\textbf{Ventana:} 1. & \textbf{Contenido y frontera.} Comités, acuerdos, riesgos, cambios, avance físico y costo, tablero y seguimiento; excluye ejecución de paquetes técnicos.\newline{}\textbf{Aceptación.} Se reportan valor ganado, SPI/CPI, acuerdos y desviaciones; actas en dos días hábiles y desvío mayor a 10\,\% comunicado en cinco días con recuperación. Ignacio Silva responde por el paquete.\newline{}\textbf{Condiciones.} Predecesoras del catálogo: 1.1.1.1. Referencia: RT-19.07--10. Las habilitaciones técnicas se contrastan con arquitectura e inventario vigentes.\\\hline
\texttt{1.\allowbreak{}1.\allowbreak{}4.\allowbreak{}\{1.\allowbreak{}2;\allowbreak{}\{2--20\}.\allowbreak{}\{1,\allowbreak{}2\}\}}\newline{}\textbf{Expediente quincenal de control}\newline{}Ignacio Silva\newline{}39 instancias.\newline{}\textbf{Ventana:} 1--20. & \textbf{Contenido y frontera.} Resultado individual del período o conjunto indicado. Excluye resultados de otras cuentas y horas de contraparte, obras o adquisiciones del cliente.\newline{}\textbf{Aceptación.} Se reportan valor ganado, SPI/CPI, acuerdos y desviaciones; actas en dos días hábiles y desvío mayor a 10\,\% comunicado en cinco días con recuperación. Ignacio Silva responde por el paquete.\newline{}\textbf{Insumos.} Versión/registro recibido, acceso autorizado, herramientas, participantes y ventana del período.\newline{}\textbf{Condiciones.} Predecesoras de habilitación: 1.1.1.1.\\\hline
\texttt{1.\allowbreak{}1.\allowbreak{}5.\allowbreak{}1}\newline{}\textbf{F29 · Expediente transversal de alcance y contratos}\newline{}Ignacio Silva\newline{}1 instancia.\newline{}\textbf{Ventana:} 2. & \textbf{Contenido y frontera.} Expediente de alcance y contratos con cláusulas aplicables, fuente/versiones, responsables, vigencias y evidencias de oferta o ejecución identificadas. No vuelve a construir los controles técnicos ni presenta una obligación futura como ejecutada.\newline{}\textbf{Aceptación.} Cláusulas, documentos, responsables y control de vigencias; se distingue evidencia de oferta de obligaciones futuras.\newline{}\textbf{Insumos.} Universo nominal/manifiesto por instancia, versión y política/protocolo aplicable recibidos, responsables, acceso autorizado y ventana admisible. Registro de instancia: código, población/activos, fecha, versión, evidencia y decisión; el conjunto no se acepta por cerrar solo su primera instancia.\newline{}\textbf{Condiciones.} Condiciones/predecesoras comunes del catálogo: 1.1.1.1. Referencia de detalle: FALTA REFERENCIA A SD7.\\\hline
\texttt{1.\allowbreak{}1.\allowbreak{}5.\allowbreak{}2}\newline{}\textbf{F30 · Expediente transversal de seguridad y privacidad}\newline{}Ignacio Silva\newline{}1 instancia.\newline{}\textbf{Ventana:} 3. & \textbf{Contenido y frontera.} Expediente de seguridad y privacidad con cláusulas aplicables, fuente/versiones, responsables, vigencias y evidencias de oferta o ejecución identificadas. No vuelve a construir los controles técnicos ni presenta una obligación futura como ejecutada.\newline{}\textbf{Aceptación.} Cláusulas, documentos, responsables y control de vigencias; se distingue evidencia de oferta de obligaciones futuras.\newline{}\textbf{Insumos.} Universo nominal/manifiesto por instancia, versión y política/protocolo aplicable recibidos, responsables, acceso autorizado y ventana admisible. Registro de instancia: código, población/activos, fecha, versión, evidencia y decisión; el conjunto no se acepta por cerrar solo su primera instancia.\newline{}\textbf{Condiciones.} Condiciones/predecesoras comunes del catálogo: 1.1.1.1. Referencia de detalle: FALTA REFERENCIA A SD7.\\\hline
\texttt{1.\allowbreak{}1.\allowbreak{}5.\allowbreak{}3}\newline{}\textbf{F31 · Expediente transversal de IA y trazabilidad}\newline{}Ignacio Silva\newline{}1 instancia.\newline{}\textbf{Ventana:} 3. & \textbf{Contenido y frontera.} Expediente de IA y trazabilidad con cláusulas aplicables, fuente/versiones, responsables, vigencias y evidencias de oferta o ejecución identificadas. No vuelve a construir los controles técnicos ni presenta una obligación futura como ejecutada.\newline{}\textbf{Aceptación.} Cláusulas, documentos, responsables y control de vigencias; se distingue evidencia de oferta de obligaciones futuras.\newline{}\textbf{Insumos.} Universo nominal/manifiesto por instancia, versión y política/protocolo aplicable recibidos, responsables, acceso autorizado y ventana admisible. Registro de instancia: código, población/activos, fecha, versión, evidencia y decisión; el conjunto no se acepta por cerrar solo su primera instancia.\newline{}\textbf{Condiciones.} Condiciones/predecesoras comunes del catálogo: 1.1.1.1. Referencia de detalle: FALTA REFERENCIA A SD7.\\\hline
\texttt{1.\allowbreak{}1.\allowbreak{}5.\allowbreak{}4}\newline{}\textbf{F32 · Expediente transversal de licencias y soporte}\newline{}Ignacio Silva\newline{}1 instancia.\newline{}\textbf{Ventana:} 3. & \textbf{Contenido y frontera.} Expediente de licencias y soporte con cláusulas aplicables, fuente/versiones, responsables, vigencias y evidencias de oferta o ejecución identificadas. No vuelve a construir los controles técnicos ni presenta una obligación futura como ejecutada.\newline{}\textbf{Aceptación.} Cláusulas, documentos, responsables y control de vigencias; se distingue evidencia de oferta de obligaciones futuras. RT-15.07--09: el expediente enlaza certificados vigentes y verificables del personal efectivamente asignado, perfiles/dedicaciones y control de reposición; evidencia de oferta en T-5/SD12 y vigencia contractual se distinguen. Se registran vencimientos y responsables, sin afirmar certificación por declaración. FALTA REFERENCIA A SD12.\newline{}\textbf{Insumos.} Universo nominal/manifiesto por instancia, versión y política/protocolo aplicable recibidos, responsables, acceso autorizado y ventana admisible. Registro de instancia: código, población/activos, fecha, versión, evidencia y decisión; el conjunto no se acepta por cerrar solo su primera instancia.\newline{}\textbf{Condiciones.} Condiciones/predecesoras comunes del catálogo: 1.1.1.1. Referencia de detalle: FALTA REFERENCIA A SD7.\\\hline
\texttt{1.\allowbreak{}1.\allowbreak{}5.\allowbreak{}5}\newline{}\textbf{F33 · Expediente transversal de custodia y propiedad}\newline{}Ignacio Silva\newline{}1 instancia.\newline{}\textbf{Ventana:} 2. & \textbf{Contenido y frontera.} Expediente de custodia y propiedad con cláusulas aplicables, fuente/versiones, responsables, vigencias y evidencias de oferta o ejecución identificadas. No vuelve a construir los controles técnicos ni presenta una obligación futura como ejecutada.\newline{}\textbf{Aceptación.} Cláusulas, documentos, responsables y control de vigencias; se distingue evidencia de oferta de obligaciones futuras.\newline{}\textbf{Insumos.} Universo nominal/manifiesto por instancia, versión y política/protocolo aplicable recibidos, responsables, acceso autorizado y ventana admisible. Registro de instancia: código, población/activos, fecha, versión, evidencia y decisión; el conjunto no se acepta por cerrar solo su primera instancia.\newline{}\textbf{Condiciones.} Condiciones/predecesoras comunes del catálogo: 1.1.1.1. Referencia de detalle: FALTA REFERENCIA A SD7.\\\hline
\texttt{1.\allowbreak{}1.\allowbreak{}5.\allowbreak{}6}\newline{}\textbf{F34 · Expediente transversal de garantías y oferta económica}\newline{}Ignacio Silva\newline{}1 instancia.\newline{}\textbf{Ventana:} 3. & \textbf{Contenido y frontera.} Expediente de garantías y oferta económica con cláusulas aplicables, fuente/versiones, responsables, vigencias y evidencias de oferta o ejecución identificadas. No vuelve a construir los controles técnicos ni presenta una obligación futura como ejecutada.\newline{}\textbf{Aceptación.} Cláusulas, documentos, responsables y control de vigencias; se distingue evidencia de oferta de obligaciones futuras.\newline{}\textbf{Insumos.} Universo nominal/manifiesto por instancia, versión y política/protocolo aplicable recibidos, responsables, acceso autorizado y ventana admisible. Registro de instancia: código, población/activos, fecha, versión, evidencia y decisión; el conjunto no se acepta por cerrar solo su primera instancia.\newline{}\textbf{Condiciones.} Condiciones/predecesoras comunes del catálogo: 1.1.1.1. Referencia de detalle: FALTA REFERENCIA A SD7.\\\hline
\end{longtable}\endgroup
\begingroup\setlength{\tabcolsep}{3pt}
\begin{longtable}{L{\dimexpr0.320000\textwidth-2\tabcolsep\relax}L{\dimexpr0.680000\textwidth-2\tabcolsep\relax}}
\caption{Diccionario: Plataforma híbrida}\label{tab:t14-dic-1-2}\\\hline
\EncabezadoTabla{Paquete y resultado} & \EncabezadoTabla{Frontera, aceptación e insumos}\\\hline\endfirsthead
\multicolumn{2}{l}{\textit{Continuación de la tabla \ref{tab:t14-dic-1-2}}}\\\hline
\EncabezadoTabla{Paquete y resultado} & \EncabezadoTabla{Frontera, aceptación e insumos}\\\hline\endhead
\hline\endfoot
\hline\endlastfoot
\texttt{1.\allowbreak{}2.\allowbreak{}1.\allowbreak{}1}\newline{}\textbf{Repositorio e infraestructura como código gobernados}\newline{}Fernando Guerrero\newline{}1 instancia.\newline{}\textbf{Ventana:} 3. & \textbf{Contenido y frontera.} Repositorio del cliente, cuenta raíz documentada, módulos OpenTofu, estado cifrado y bloqueo; excluye módulos de aplicación.\newline{}\textbf{Aceptación.} Titularidad, MFA, revisión independiente, estado recuperable y ausencia de secretos en Git quedan comprobados. Responsable: Fernando Guerrero (Arquitectura/Integración).\newline{}\textbf{Insumos.} cuentas y repositorio del cliente, infraestructura cloud/local, inventario T-11, herramientas de configuración y acceso autorizado.\newline{}\textbf{Condiciones.} Predecesoras del catálogo: H1. Referencia: SD4 gobierno de nube; RT-03.03. Las habilitaciones técnicas se contrastan con arquitectura e inventario vigentes.\\\hline
\texttt{1.\allowbreak{}2.\allowbreak{}1.\allowbreak{}2}\newline{}\textbf{Red cloud primaria reproducible}\newline{}Fernando Guerrero\newline{}1 instancia.\newline{}\textbf{Ventana:} 4. & \textbf{Contenido y frontera.} VPC, subredes privadas en dos zonas, rutas, endpoints y acceso de administración por código.\newline{}\textbf{Aceptación.} Se reproduce la red y se comprueba acceso permitido, bloqueo de exposición de datos y diferencias del plan IaC. Responsable: Fernando Guerrero (Arquitectura/Integración).\newline{}\textbf{Insumos.} cuentas y repositorio del cliente, infraestructura cloud/local, inventario T-11, herramientas de configuración y acceso autorizado.\newline{}\textbf{Condiciones.} Predecesoras del catálogo: 1.2.1.1. Referencia: SD4 DA-01/04; RT-03.03. Las habilitaciones técnicas se contrastan con arquitectura e inventario vigentes.\\\hline
\texttt{1.\allowbreak{}2.\allowbreak{}1.\allowbreak{}3}\newline{}\textbf{Servicios de persistencia y objetos habilitados}\newline{}Nicolás Torfan\newline{}1 instancia.\newline{}\textbf{Ventana:} 5. & \textbf{Contenido y frontera.} RDS PostgreSQL, almacén de proyecciones y S3 con políticas; excluye modelo funcional, migración y cargas de negocio.\newline{}\textbf{Aceptación.} Se identifican autoridades y recursos separados; conectividad TLS, políticas, respaldo y prohibición de acceso directo entre dominios se comprueban. Nicolás Torfan, Datos, responde por el resultado y presenta la evidencia; la contraparte mantiene la recepción contractual.\newline{}\textbf{Insumos.} cuentas y repositorio del cliente, infraestructura cloud/local, inventario T-11, herramientas de configuración y acceso autorizado.\newline{}\textbf{Condiciones.} Predecesoras del catálogo: 1.2.1.2,H2. Referencia: SD4 DA-08; SD5 persistencia. Las habilitaciones técnicas se contrastan con arquitectura e inventario vigentes. Diseño de persistencia recibido en 1.7.6.7/H2 antes de habilitar el recurso.\\\hline
\texttt{1.\allowbreak{}2.\allowbreak{}1.\allowbreak{}4}\newline{}\textbf{Plataforma de contenedores y exposición habilitada}\newline{}Fernando Guerrero\newline{}1 instancia.\newline{}\textbf{Ventana:} 5. & \textbf{Contenido y frontera.} ECS/Fargate, balanceo privado, API Gateway, DNS y protección de entrada por IaC; excluye servicios de negocio.\newline{}\textbf{Aceptación.} Imagen de comprobación se ejecuta por despliegue reproducible, sin publicar subredes de aplicación/datos ni claves estáticas. Responsable: Fernando Guerrero (Arquitectura/Integración).\newline{}\textbf{Insumos.} cuentas y repositorio del cliente, infraestructura cloud/local, inventario T-11, herramientas de configuración y acceso autorizado.\newline{}\textbf{Condiciones.} Predecesoras del catálogo: 1.2.1.2. Referencia: SD4 arquitectura física; RT-03.03. Las habilitaciones técnicas se contrastan con arquitectura e inventario vigentes.\\\hline
\texttt{1.\allowbreak{}2.\allowbreak{}1.\allowbreak{}5}\newline{}\textbf{Cadena de promoción y artefactos protegidos}\newline{}Fernando Guerrero\newline{}1 instancia.\newline{}\textbf{Ventana:} 5. & \textbf{Contenido y frontera.} Repositorio, registro OCI, firma, evidencias de dependencias y promoción; excluye construir módulos funcionales.\newline{}\textbf{Aceptación.} El mismo artefacto identificado se promueve sin recompilar por ambiente y una versión no autorizada queda bloqueada. Responsable: Fernando Guerrero (Arquitectura/Integración).\newline{}\textbf{Insumos.} cuentas y repositorio del cliente, infraestructura cloud/local, inventario T-11, herramientas de configuración y acceso autorizado.\newline{}\textbf{Condiciones.} Predecesoras del catálogo: 1.2.1.4. Referencia: SD4 ambientes; BT capítulos 3/19. Las habilitaciones técnicas se contrastan con arquitectura e inventario vigentes.\\\hline
\texttt{1.\allowbreak{}2.\allowbreak{}2.\allowbreak{}1}\newline{}\textbf{Ambiente DEV habilitado}\newline{}Fernando Guerrero\newline{}1 instancia.\newline{}\textbf{Ventana:} 5. & \textbf{Contenido y frontera.} Configuración, acceso y despliegue de comprobación DEV con datos sintéticos; excluye horas de construcción de software.\newline{}\textbf{Aceptación.} DEV se reproduce, usa sólo datos permitidos y mantiene aislamiento y horario versionado. Responsable: Fernando Guerrero (Arquitectura/Integración).\newline{}\textbf{Insumos.} cuentas y repositorio del cliente, infraestructura cloud/local, inventario T-11, herramientas de configuración y acceso autorizado.\newline{}\textbf{Condiciones.} Predecesoras del catálogo: 1.2.1.3,1.2.1.5. Referencia: RT-04.01; FALTA REFERENCIA A SD4. Las habilitaciones técnicas se contrastan con arquitectura e inventario vigentes.\\\hline
\texttt{1.\allowbreak{}2.\allowbreak{}2.\allowbreak{}2}\newline{}\textbf{Ambiente QA habilitado}\newline{}Davor Serey\newline{}1 instancia.\newline{}\textbf{Ventana:} 5. & \textbf{Contenido y frontera.} Configuración QA, datos autorizados y ejecución de comprobación; excluye regresión de todos los módulos.\newline{}\textbf{Aceptación.} QA queda aislado de producción, identifica versiones y sólo usa datos sintéticos o anonimizados aprobados. Davor Serey, Calidad, responde por el resultado y presenta la evidencia; la contraparte mantiene la recepción contractual.\newline{}\textbf{Insumos.} cuentas y repositorio del cliente, infraestructura cloud/local, inventario T-11, herramientas de configuración y acceso autorizado.\newline{}\textbf{Condiciones.} Predecesoras del catálogo: 1.2.1.3,1.2.1.5. Referencia: RT-04.01; FALTA REFERENCIA A SD4. Las habilitaciones técnicas se contrastan con arquitectura e inventario vigentes.\\\hline
\texttt{1.\allowbreak{}2.\allowbreak{}2.\allowbreak{}3}\newline{}\textbf{Ambiente PREPROD habilitado}\newline{}Fernando Guerrero\newline{}1 instancia.\newline{}\textbf{Ventana:} 6. & \textbf{Contenido y frontera.} Topología productiva a escala de prueba y restitución controlada; excluye ejecución de carga, UAT y ensayos de migración.\newline{}\textbf{Aceptación.} Se comprueban aislamiento, capacidad de restaurar configuración y control de copias excepcionales anonimizadas. Responsable: Fernando Guerrero (Arquitectura/Integración).\newline{}\textbf{Insumos.} cuentas y repositorio del cliente, infraestructura cloud/local, inventario T-11, herramientas de configuración y acceso autorizado.\newline{}\textbf{Condiciones.} Predecesoras del catálogo: 1.2.2.2. Referencia: RT-04.01; FALTA REFERENCIA A SD4. Las habilitaciones técnicas se contrastan con arquitectura e inventario vigentes.\\\hline
\texttt{1.\allowbreak{}2.\allowbreak{}2.\allowbreak{}4}\newline{}\textbf{Ambiente PROD habilitado}\newline{}José Lara Arce\newline{}1 instancia.\newline{}\textbf{Ventana:} 6. & \textbf{Contenido y frontera.} Configuración Multi-AZ, límites, permisos y comprobación técnica sin usuarios de negocio.\newline{}\textbf{Aceptación.} Producción tiene configuración identificada, aislamiento, monitoreo y acceso nominativo; habilitar ambiente no acredita paso a producción E1. José Lara Arce, Operación/SRE, responde por el resultado y presenta la evidencia; la contraparte mantiene la recepción contractual.\newline{}\textbf{Insumos.} cuentas y repositorio del cliente, infraestructura cloud/local, inventario T-11, herramientas de configuración y acceso autorizado.\newline{}\textbf{Condiciones.} Predecesoras del catálogo: 1.2.1.3,1.2.1.4. Referencia: RT-04.01; SD4 producción. Las habilitaciones técnicas se contrastan con arquitectura e inventario vigentes.\\\hline
\texttt{1.\allowbreak{}2.\allowbreak{}2.\allowbreak{}5}\newline{}\textbf{Ambiente DR habilitado}\newline{}José Lara Arce\newline{}1 instancia.\newline{}\textbf{Ventana:} 6. & \textbf{Contenido y frontera.} Warm standby México, réplica cifrada, imágenes, objetos y configuración mínima; excluye prueba integral de desastre.\newline{}\textbf{Aceptación.} Se comprueban salud, replicación, permisos y restitución del entorno; capacidad de respaldo no se declara por una réplica sin evidencia. José Lara Arce, Operación/SRE, responde por el resultado y presenta la evidencia; la contraparte mantiene la recepción contractual.\newline{}\textbf{Insumos.} cuentas y repositorio del cliente, infraestructura cloud/local, inventario T-11, herramientas de configuración y acceso autorizado.\newline{}\textbf{Condiciones.} Predecesoras del catálogo: 1.2.2.4,1.3.1.1. Referencia: RT-04.01; SD4 DA-17. Las habilitaciones técnicas se contrastan con arquitectura e inventario vigentes.\\\hline
\texttt{1.\allowbreak{}2.\allowbreak{}3.\allowbreak{}1}\newline{}\textbf{Expediente de habilitación del recinto y recepción de activos}\newline{}Fernando Guerrero\newline{}1 instancia.\newline{}\textbf{Ventana:} 4. & \textbf{Contenido y frontera.} Contraste de plano y condiciones SD4, compras/obras del cliente y T-11; excluye ejecutar obras o comprar equipamiento. Localiza las especificaciones del recinto RT-06 y del hardware RT-08, garantías, mantenimientos, certificaciones e inventario de ciclo de vida. Distingue habilitación externa, especificación/coordinación de Synaptix y configuración propia; cada obligación conserva su titular en T-11.\newline{}\textbf{Aceptación.} Cada habilitación tiene especificación, titular, fecha requerida y evidencia de recepción pendiente o conforme; ningún dato proyectado se atribuye al recinto existente. Responsable: Fernando Guerrero (Arquitectura/Integración). Matriz de habilitaciones identifica especificación/plano SD4, activo, criterio técnico, titular, fecha, evidencia y bloqueo de H3 cuando aplica. Dimensionamiento incluye carga, crecimiento, consumo y PUE; huella de carbono declara método, límites y fuente de intensidad de la región, sin dar una proyección por medición. Exigencias y responsabilidades no se reciben solo por incluir un enlace al plano.\newline{}\textbf{Insumos.} cuentas y repositorio del cliente, infraestructura cloud/local, inventario T-11, herramientas de configuración y acceso autorizado.\newline{}\textbf{Condiciones.} Predecesoras del catálogo: H2,EXT-RECINTO. Referencia: SD2 SUP-04; SD4 implantación del recinto; T-11. Las habilitaciones técnicas se contrastan con arquitectura e inventario vigentes.\\\hline
\texttt{1.\allowbreak{}2.\allowbreak{}3.\allowbreak{}2}\newline{}\textbf{Configuración de nodos locales y almacenamiento redundante}\newline{}Fernando Guerrero\newline{}1 instancia.\newline{}\textbf{Ventana:} 5. & \textbf{Contenido y frontera.} Configuración versionada de los dos nodos, almacenamiento y servicios locales; excluye construcción de reglas de dominio. Arquitectura/Integración integra y presenta el resultado de mes 5; ejecución técnica se asigna a recursos disponibles en T-15.\newline{}\textbf{Aceptación.} Inventario y firmware corresponden a activos recibidos; restauración de configuración, salud de discos y aislamiento comprobados con revisión independiente. Fernando Guerrero responde por la integración y evidencia del paquete. Operación/SRE recibe transferencia y comprueba preparación operativa desde su incorporación en mes 6, antes de H3; no se atribuye participación SRE en mes 5.\newline{}\textbf{Insumos.} Activos recibidos, inventario T-11, especificaciones vigentes, herramientas, repositorio y acceso autorizado; recursos de configuración y revisión disponibles en mes 5.\newline{}\textbf{Condiciones.} Predecesoras: 1.2.3.1,EXT-EQUIPOS. La entrega de configuración de mes 5 y su transferencia operativa de mes 6 son momentos distintos, sin duplicar construcción. H3 exige revisión/transferencia SRE completa. Conciliar responsable con T-15 y RACI; SD6 mantiene incorporación SRE desde mes 6. FALTA REFERENCIA A SD4; FALTA REFERENCIA A SD6.\\\hline
\texttt{1.\allowbreak{}2.\allowbreak{}3.\allowbreak{}3}\newline{}\textbf{Configuración de red local, segmentación y enlaces}\newline{}Fernando Guerrero\newline{}1 instancia.\newline{}\textbf{Ventana:} 5. & \textbf{Contenido y frontera.} Configuración versionada de cortafuegos, switches, VLAN, WAN y túneles; excluye obras y tendido físico del cliente.\newline{}\textbf{Aceptación.} Rutas autorizadas y segmentación se verifican; fallos de un enlace no abren acceso no autorizado ni exponen bases de datos. Responsable: Fernando Guerrero (Arquitectura/Integración). RT-03.24 C07: cobertura en cinco pantalanes y varadero, verificación sobre estructura flotante con marea en su rango completo y flexión del cableado; separación efectiva de redes de socios, administración, cámaras y operación; enlace principal y respaldo contratados y comprobados. Pruebas identifican puntos, condiciones, permisos, fallos y evidencia antes de H3; configuración no recibe por sí sola la obra externa.\newline{}\textbf{Insumos.} cuentas y repositorio del cliente, infraestructura cloud/local, inventario T-11, herramientas de configuración y acceso autorizado.\newline{}\textbf{Condiciones.} Predecesoras del catálogo: 1.2.3.1,EXT-RED. Referencia: SD4 DA-01 y enlaces WAN; SD2 habilitaciones. Las habilitaciones técnicas se contrastan con arquitectura e inventario vigentes.\\\hline
\texttt{1.\allowbreak{}2.\allowbreak{}3.\allowbreak{}4}\newline{}\textbf{Gestión y protección de terminales de referencia}\newline{}Bruno Toro\newline{}1 instancia.\newline{}\textbf{Ventana:} 6. & \textbf{Contenido y frontera.} Políticas UEM, cifrado, actualización y protección para una estación y un Android de referencia; enrolamiento del parque se estima por lotes.\newline{}\textbf{Aceptación.} Se verifican enrolamiento, bloqueo, pérdida, política, actualización y cobertura efectiva; la selección de una licencia no equivale a cobertura comprobada. Bruno Toro, Seguridad, responde por el resultado y presenta la evidencia; la contraparte mantiene la recepción contractual.\newline{}\textbf{Insumos.} cuentas y repositorio del cliente, infraestructura cloud/local, inventario T-11, herramientas de configuración y acceso autorizado.\newline{}\textbf{Condiciones.} Predecesoras del catálogo: 1.2.3.3,1.3.1.2,EXT-LICENCIAS. Referencia: FALTA REFERENCIA A SD4; T-11; BT capítulo 8. Las habilitaciones técnicas se contrastan con arquitectura e inventario vigentes.\\\hline
\texttt{1.\allowbreak{}2.\allowbreak{}3.\allowbreak{}5}\newline{}\textbf{Configuración de concentradores y broker de borde}\newline{}Fernando Guerrero\newline{}1 instancia.\newline{}\textbf{Ventana:} 6. & \textbf{Contenido y frontera.} Configuración de diez concentradores y EMQX, reloj, certificados y cola; excluye software de sincronización y puesta en servicio de medidores.\newline{}\textbf{Aceptación.} Cada concentrador tiene activo, identidad, UTC, muestreo por minuto, agregado cada 15 minutos y cola de 24 horas comprobados con señales de prueba. Responsable: Fernando Guerrero (Arquitectura/Integración).\newline{}\textbf{Insumos.} cuentas y repositorio del cliente, infraestructura cloud/local, inventario T-11, herramientas de configuración y acceso autorizado.\newline{}\textbf{Condiciones.} Predecesoras del catálogo: 1.2.3.3,1.3.1.1,EXT-TELEMETRIA. Referencia: SD4 DA-09; diez concentradores. Las habilitaciones técnicas se contrastan con arquitectura e inventario vigentes.\\\hline
\texttt{1.\allowbreak{}2.\allowbreak{}3.\allowbreak{}7.\allowbreak{}\{1--10\}}\newline{}\textbf{F03 · Lotes nominales de enrolamiento de terminales}\newline{}José Lara Arce\newline{}10 instancias.\newline{}\textbf{Ventana:} 6. & \textbf{Contenido y frontera.} Lote nominal de hasta cuatro terminales enrolados y comprobados contra la política recibida; inventario, propietario, licencia, cifrado, protección y evidencia por equipo. No reconstruye la política común de 1.2.3.4.\newline{}\textbf{Aceptación.} Equipo nominal inscrito, cifrado, pérdida, protección, licencia y revisión de política comprobados.\newline{}\textbf{Insumos.} Universo nominal/manifiesto por instancia, versión y política/protocolo aplicable recibidos, responsables, acceso autorizado y ventana admisible. Registro de instancia: código, población/activos, fecha, versión, evidencia y decisión; el conjunto no se acepta por cerrar solo su primera instancia.\newline{}\textbf{Condiciones.} Condiciones/predecesoras comunes del catálogo: 1.2.3.4. Referencia de detalle: FALTA REFERENCIA A SD4.\\\hline
\texttt{1.\allowbreak{}2.\allowbreak{}4.\allowbreak{}1}\newline{}\textbf{Observabilidad técnica y alertas habilitadas}\newline{}José Lara Arce\newline{}1 instancia.\newline{}\textbf{Ventana:} 6. & \textbf{Contenido y frontera.} OpenTelemetry, métricas, registros, trazas y alertas de infraestructura; excluye indicadores funcionales de cada módulo.\newline{}\textbf{Aceptación.} Las fallas simuladas generan alerta con activo, correlación y responsable; se comprueban cobertura y falta de agente. José Lara Arce, Operación/SRE, responde por el resultado y presenta la evidencia; la contraparte mantiene la recepción contractual. RT-14: nube y recinto correlacionan métricas/registros/trazas por transacción, con instrumentación OpenTelemetry; cliente accede a tableros y exportación. Alertas usan síntomas de negocio, agrupación y escalamiento. Retención de métricas/registros/trazas diferencia línea/archivo y remite su costo a la oferta; registros técnicos no contienen secretos ni datos sensibles.\newline{}\textbf{Insumos.} cuentas y repositorio del cliente, infraestructura cloud/local, inventario T-11, herramientas de configuración y acceso autorizado.\newline{}\textbf{Condiciones.} Predecesoras del catálogo: 1.2.2.3,1.2.2.4,1.2.3.2. Referencia: SD4 observabilidad; E-25 H3. Las habilitaciones técnicas se contrastan con arquitectura e inventario vigentes.\\\hline
\texttt{1.\allowbreak{}2.\allowbreak{}4.\allowbreak{}2}\newline{}\textbf{Procedimientos de respaldo y restauración técnica}\newline{}José Lara Arce\newline{}1 instancia.\newline{}\textbf{Ventana:} 6. & \textbf{Contenido y frontera.} Políticas de copia, retención, restauración de muestra y permisos; excluye recuperación integral y retorno de migración.\newline{}\textbf{Aceptación.} Una copia identificada se restaura con integridad, permisos y tiempos registrados; la existencia de backup no sustituye su restauración. José Lara Arce, Operación/SRE, responde por el resultado y presenta la evidencia; la contraparte mantiene la recepción contractual. RT-07.09--14: política 3-2-1-1-0, cifrado y claves independientes, inmutabilidad y restauración granular se comprueban con evidencia. Frecuencia y retención por dominio se enlazan a SD5; 1.12.3.2 ejecuta las comprobaciones mensuales posteriores, sin duplicar esta habilitación.\newline{}\textbf{Insumos.} cuentas y repositorio del cliente, infraestructura cloud/local, inventario T-11, herramientas de configuración y acceso autorizado.\newline{}\textbf{Condiciones.} Predecesoras del catálogo: 1.2.1.3,1.2.2.5. Referencia: SD4 continuidad; BT capítulo 15. Las habilitaciones técnicas se contrastan con arquitectura e inventario vigentes.\\\hline
\texttt{1.\allowbreak{}2.\allowbreak{}4.\allowbreak{}3}\newline{}\textbf{Expediente técnico de entrega de cinco ambientes}\newline{}Fernando Guerrero\newline{}1 instancia.\newline{}\textbf{Ventana:} 6. & \textbf{Contenido y frontera.} Consolidación de evidencia de infraestructura, observabilidad y habilitaciones; excluye repetir las comprobaciones anteriores.\newline{}\textbf{Aceptación.} DEV, QA, PREPROD, PROD y DR tienen evidencia y pendientes identificados; condiciones bloqueantes impiden recepción H3. Responsable: Fernando Guerrero (Arquitectura/Integración).\newline{}\textbf{Insumos.} cuentas y repositorio del cliente, infraestructura cloud/local, inventario T-11, herramientas de configuración y acceso autorizado.\newline{}\textbf{Condiciones.} Predecesoras del catálogo: 1.2.2.*,1.2.3.*,1.2.4.1,1.2.4.2. Referencia: RT-04.01; E-25 H3. Las habilitaciones técnicas se contrastan con arquitectura e inventario vigentes.\\\hline
\end{longtable}\endgroup
\begingroup\setlength{\tabcolsep}{3pt}
\begin{longtable}{L{\dimexpr0.320000\textwidth-2\tabcolsep\relax}L{\dimexpr0.680000\textwidth-2\tabcolsep\relax}}
\caption{Diccionario: Auditoría y seguridad}\label{tab:t14-dic-1-3}\\\hline
\EncabezadoTabla{Paquete y resultado} & \EncabezadoTabla{Frontera, aceptación e insumos}\\\hline\endfirsthead
\multicolumn{2}{l}{\textit{Continuación de la tabla \ref{tab:t14-dic-1-3}}}\\\hline
\EncabezadoTabla{Paquete y resultado} & \EncabezadoTabla{Frontera, aceptación e insumos}\\\hline\endhead
\hline\endfoot
\hline\endlastfoot
\texttt{1.\allowbreak{}3.\allowbreak{}1.\allowbreak{}1}\newline{}\textbf{Gestión de claves, secretos y certificados}\newline{}Bruno Toro\newline{}1 instancia.\newline{}\textbf{Ventana:} 5. & \textbf{Contenido y frontera.} KMS, almacén de secretos, certificados y procedimiento de rotación/revocación; excluye cifrado funcional del código de cada módulo.\newline{}\textbf{Aceptación.} Credenciales temporales, alcance mínimo, revocación, rotación y lectura autorizada quedan comprobados sin secretos en repositorio. Bruno Toro, Seguridad, responde por el resultado y presenta la evidencia; la contraparte mantiene la recepción contractual.\newline{}\textbf{Insumos.} identidades autorizadas, políticas, repositorio protegido y herramientas de seguridad.\newline{}\textbf{Condiciones.} Predecesoras del catálogo: 1.2.1.1. Referencia: SD4 seguridad y DA-16; BT capítulos 3/18. Las habilitaciones técnicas se contrastan con arquitectura e inventario vigentes. La protección funcional de campos sensibles se recibe en 1.3.4.4; la infraestructura de claves no sustituye esa prueba.\\\hline
\texttt{1.\allowbreak{}3.\allowbreak{}1.\allowbreak{}2}\newline{}\textbf{Identidad central y federación configuradas}\newline{}Bruno Toro\newline{}1 instancia.\newline{}\textbf{Ventana:} 5. & \textbf{Contenido y frontera.} Cognito, OIDC/SAML, MFA, identidades nominativas y recuperación DR; excluye autoridad local y permisos funcionales de los dominios. Incluye integración al directorio autorizado, SSO y cierre de sesión propagado, factores FIDO2 ofertados, autoservicio seguro de acceso externo y política de sesiones. Incluye acceso privilegiado temporal y acceso de emergencia con autorización, custodia y auditoría; no declara soportado un control por nombrar al proveedor.\newline{}\textbf{Aceptación.} Se rechazan cuentas compartidas y acceso sin factor exigido; recuperación DR no presupone replicación de contraseñas no soportada. Bruno Toro, Seguridad, responde por el resultado y presenta la evidencia; la contraparte mantiene la recepción contractual. RT-12: federación/directorio, SSO y logout se comprueban entre módulos; MFA se exige en administración, acceso privilegiado y acceso externo. Se verifican factor resistente a suplantación, expiración/inactividad, renovación, revocación inmediata y concurrencia según política recibida, credencial breve firmada y refresco rotatorio, sin identificador de sesión en URL. Elevación requiere autorización y registro de sesiones de mayor riesgo; acceso de emergencia se prueba con custodia y auditoría. Usuarios externos verifican identidad y recuperan acceso sin intervención informal.\newline{}\textbf{Insumos.} identidades autorizadas, políticas, repositorio protegido y herramientas de seguridad.\newline{}\textbf{Condiciones.} Predecesoras del catálogo: 1.3.1.1,EXT-IDENTIDAD. Referencia: SD4 DA-16 y región secundaria. Las habilitaciones técnicas se contrastan con arquitectura e inventario vigentes.\\\hline
\texttt{1.\allowbreak{}3.\allowbreak{}1.\allowbreak{}3}\newline{}\textbf{Matriz de permisos y separación de funciones}\newline{}Bruno Toro\newline{}1 instancia.\newline{}\textbf{Ventana:} 4. & \textbf{Contenido y frontera.} Roles, ámbitos y separación entre preparación/aprobación para los nueve dominios; excluye implementación en cada servicio.\newline{}\textbf{Aceptación.} Cada operación tiene actor, ámbito y denegación; el cliente valida las decisiones y no quedan permisos globales implícitos. Bruno Toro, Seguridad, responde por el resultado y presenta la evidencia; la contraparte mantiene la recepción contractual. RT-12.05: RBAC y atributos aplicables tienen matriz verificable, segregación y denegaciones. La configuración común se distingue de las autorizaciones que implementa cada dominio.\newline{}\textbf{Insumos.} identidades autorizadas, políticas, repositorio protegido y herramientas de seguridad.\newline{}\textbf{Condiciones.} Predecesoras del catálogo: H1. Referencia: SD3 nueve dominios; SD4 DA-16. Las habilitaciones técnicas se contrastan con arquitectura e inventario vigentes.\\\hline
\texttt{1.\allowbreak{}3.\allowbreak{}2.\allowbreak{}1}\newline{}\textbf{Autoridad local de acceso implementada}\newline{}Bruno Toro\newline{}1 instancia.\newline{}\textbf{Ventana:} 8. & \textbf{Contenido y frontera.} Componente local de permisos firmados y cifrados y vínculo a fuente laboral; excluye reglas propias de módulos y prueba integral offline.\newline{}\textbf{Aceptación.} Se comprueban identidad nominativa, ámbito y vencimiento máximo de 24 horas desde emisión, con denegación de administración global. Bruno Toro, Seguridad, responde por el resultado y presenta la evidencia; la contraparte mantiene la recepción contractual.\newline{}\textbf{Insumos.} identidades autorizadas, políticas, repositorio protegido y herramientas de seguridad.\newline{}\textbf{Condiciones.} Predecesoras del catálogo: 1.3.1.2,1.3.1.3,1.2.3.2. Referencia: SD4 DA-16; continuidad local. Las habilitaciones técnicas se contrastan con arquitectura e inventario vigentes.\\\hline
\texttt{1.\allowbreak{}3.\allowbreak{}2.\allowbreak{}2}\newline{}\textbf{Revocación y actualización de habilitación local}\newline{}Bruno Toro\newline{}1 instancia.\newline{}\textbf{Ventana:} 8. & \textbf{Contenido y frontera.} Propagación de cambios, bajas, permisos y fuente laboral actual; excluye provisión de datos laborales por el cliente.\newline{}\textbf{Aceptación.} Bajas y cambios se aplican con trazabilidad y fallo de fuente visible; no se mantiene acceso por caché indefinida. Bruno Toro, Seguridad, responde por el resultado y presenta la evidencia; la contraparte mantiene la recepción contractual. RT-12.10: altas/bajas automatizadas desde la fuente laboral autorizada y baja efectiva no superior a 24 horas desde desvinculación; emisión de un permiso local no reinicia ese plazo. Se comprueban bloqueo y actualización en canales/componentes alcanzados.\newline{}\textbf{Insumos.} identidades autorizadas, políticas, repositorio protegido y herramientas de seguridad.\newline{}\textbf{Condiciones.} Predecesoras del catálogo: 1.3.2.1,EXT-IDENTIDAD. Referencia: SD4 autoridad local; SD2 fuente laboral. Las habilitaciones técnicas se contrastan con arquitectura e inventario vigentes.\\\hline
\texttt{1.\allowbreak{}3.\allowbreak{}3.\allowbreak{}1}\newline{}\textbf{Registro de auditoría protegido}\newline{}Bruno Toro\newline{}1 instancia.\newline{}\textbf{Ventana:} 8. & \textbf{Contenido y frontera.} Eventos, correlación, cuenta separada y conservación protegida; excluye generación de hechos funcionales por módulos. Conserva auditoría de identidad, cambios de datos, consultas sensibles y exportaciones. Los dominios y adaptadores emiten los hechos; este resultado protege, correlaciona, consulta y exporta la evidencia, sin reconstruir su lógica funcional.\newline{}\textbf{Aceptación.} Actor, origen, instante y versión quedan preservados; alteración, acceso indebido y falta de envío se detectan. Bruno Toro, Seguridad, responde por el resultado y presenta la evidencia; la contraparte mantiene la recepción contractual. RT-12.09 y RT-16.06--09/30: eventos identifican actor, origen, fecha/hora/zona y operación; cambios conservan antes/después con acceso protegido. Toda consulta a salud y datos de menores, posiciones consentidas de embarcaciones y deuda de armadores queda registrada, además de modificaciones. Se comprueba consulta/exportación por persona, período, entidad y operación; ningún administrador puede alterar la evidencia. Se evita revelar valores sensibles o credenciales en registros técnicos. La retención procede de 1.7.6.5; emisión de eventos se comprueba en los componentes propietarios.\newline{}\textbf{Insumos.} identidades autorizadas, políticas, repositorio protegido y herramientas de seguridad.\newline{}\textbf{Condiciones.} Predecesoras del catálogo: 1.3.1.1,1.2.1.3. Referencia: SD4 auditoría; SD5 eventos. Las habilitaciones técnicas se contrastan con arquitectura e inventario vigentes.\\\hline
\texttt{1.\allowbreak{}3.\allowbreak{}3.\allowbreak{}2}\newline{}\textbf{Consentimiento y restricciones de finalidad implementados}\newline{}Bruno Toro\newline{}1 instancia.\newline{}\textbf{Ventana:} 9. & \textbf{Contenido y frontera.} Servicio de consentimiento y evento de cambios; excluye adhesión y reglas adicionales del pasaporte IN5.\newline{}\textbf{Aceptación.} Alta, revocación, finalidad y conservación tienen decisión trazada; la revocación impide nuevos usos dependientes sin borrar evidencia obligatoria. Bruno Toro, Seguridad, responde por el resultado y presenta la evidencia; la contraparte mantiene la recepción contractual.\newline{}\textbf{Insumos.} identidades autorizadas, políticas, repositorio protegido y herramientas de seguridad.\newline{}\textbf{Condiciones.} Predecesoras del catálogo: 1.3.1.3,1.3.3.1,H2. Referencia: SD5 Auditoría y seguridad; SD13 reutilización. Las habilitaciones técnicas se contrastan con arquitectura e inventario vigentes.\\\hline
\texttt{1.\allowbreak{}3.\allowbreak{}4.\allowbreak{}1}\newline{}\textbf{Modelo de amenazas y plan de seguridad aprobables}\newline{}Bruno Toro\newline{}1 instancia.\newline{}\textbf{Ventana:} 4. & \textbf{Contenido y frontera.} Amenazas por autoridad, canal, WAN, terminal y tercero, mitigación y comprobación prevista; excluye pruebas ofensivas integrales.\newline{}\textbf{Aceptación.} Cada amenaza relevante tiene control, cuenta de trabajo y método de verificación; se identifican pendientes antes de H2. Bruno Toro, Seguridad, responde por el resultado y presenta la evidencia; la contraparte mantiene la recepción contractual.\newline{}\textbf{Insumos.} identidades autorizadas, políticas, repositorio protegido y herramientas de seguridad.\newline{}\textbf{Condiciones.} Predecesoras del catálogo: H1. Referencia: E-25 H2; SD4 seguridad. Las habilitaciones técnicas se contrastan con arquitectura e inventario vigentes.\\\hline
\texttt{1.\allowbreak{}3.\allowbreak{}4.\allowbreak{}2}\newline{}\textbf{Detección, cobertura y respuesta técnica configuradas}\newline{}Bruno Toro\newline{}1 instancia.\newline{}\textbf{Ventana:} 6. & \textbf{Contenido y frontera.} Controles de endpoint, runtime y servicios administrados y procedimientos de contención; excluye operación continua SOC.\newline{}\textbf{Aceptación.} Se registra alcance efectivo por activo y fallo de agente; una alerta de prueba produce contención, evidencia y responsable identificados. Bruno Toro, Seguridad, responde por el resultado y presenta la evidencia; la contraparte mantiene la recepción contractual.\newline{}\textbf{Insumos.} identidades autorizadas, políticas, repositorio protegido y herramientas de seguridad.\newline{}\textbf{Condiciones.} Predecesoras del catálogo: 1.2.3.4,1.2.4.1,EXT-LICENCIAS. Referencia: FALTA REFERENCIA A SD4; GuardDuty y protección de terminales. Las habilitaciones técnicas se contrastan con arquitectura e inventario vigentes.\\\hline
\texttt{1.\allowbreak{}3.\allowbreak{}4.\allowbreak{}3}\newline{}\textbf{Expediente de privacidad y control de datos de prueba}\newline{}Bruno Toro\newline{}1 instancia.\newline{}\textbf{Ventana:} 6. & \textbf{Contenido y frontera.} Inventario de sensibilidad, reglas de anonimización y copias autorizadas; excluye anonimización masiva y migración.\newline{}\textbf{Aceptación.} Cada conjunto tiene finalidad, permiso y conservación; se impide uso productivo en DEV y uso de copia no aprobada en QA/PREPROD. Bruno Toro, Seguridad, responde por el resultado y presenta la evidencia; la contraparte mantiene la recepción contractual.\newline{}\textbf{Insumos.} identidades autorizadas, políticas, repositorio protegido y herramientas de seguridad.\newline{}\textbf{Condiciones.} Predecesoras del catálogo: H2,1.2.2.2. Referencia: SD4 ambientes; SD5 sensibilidad. Las habilitaciones técnicas se contrastan con arquitectura e inventario vigentes.\\\hline
\texttt{1.\allowbreak{}3.\allowbreak{}4.\allowbreak{}4}\newline{}\textbf{Protección funcional de categorías sensibles recibida}\newline{}Bruno Toro\newline{}1 instancia.\newline{}\textbf{Ventana:} 4--8,13--15. & \textbf{Contenido y frontera.} Clasificación de categorías y campos sensibles, reglas de seudonimización o cifrado de campo y adaptación de su escritura, consulta, exportación y auditoría. Usa claves de 1.3.1.1; no reconstruye KMS ni confunde cifrado de disco con protección del campo.\newline{}\textbf{Aceptación.} Inventario de campos y finalidad aprobado; personas autorizadas acceden según su rol, accesos no autorizados, exportaciones y registros técnicos no revelan valores sensibles. Pruebas comprueban escritura/lectura, revocación y tratamiento de copias conforme a RT-05.08. Seguridad responde; Desarrollo implementa adaptaciones y Datos valida categorías. Recepción E1 de salud/menores y demás categorías presentes; E2 incorpora únicamente categorías nuevas y conserva regresión E1.\newline{}\textbf{Insumos.} Para diseño: requisitos, modelo y responsables de datos. Para implementación: clasificación aprobada, 1.3.1.1, DEV/QA y componentes propietarios. Para recepción: versión y universo por etapa, casos autorizados y prohibidos.\newline{}\textbf{Condiciones.} Diseño antes de H2; configuración y adaptaciones antes de H4/H5 o H9/H10 según el componente. La aceptación identifica código, etapa, versión, campos y evidencia. FALTA REFERENCIA A SD5; FALTA REFERENCIA A SD9.\\\hline
\end{longtable}\endgroup
\begingroup\setlength{\tabcolsep}{3pt}
\begin{longtable}{L{\dimexpr0.320000\textwidth-2\tabcolsep\relax}L{\dimexpr0.680000\textwidth-2\tabcolsep\relax}}
\caption{Diccionario: Núcleo Operacional}\label{tab:t14-dic-1-4}\\\hline
\EncabezadoTabla{Paquete y resultado} & \EncabezadoTabla{Frontera, aceptación e insumos}\\\hline\endfirsthead
\multicolumn{2}{l}{\textit{Continuación de la tabla \ref{tab:t14-dic-1-4}}}\\\hline
\EncabezadoTabla{Paquete y resultado} & \EncabezadoTabla{Frontera, aceptación e insumos}\\\hline\endhead
\hline\endfoot
\hline\endlastfoot
\texttt{1.\allowbreak{}4.\allowbreak{}1.\allowbreak{}1.\allowbreak{}1}\newline{}\textbf{Capacidad funcional recibida: Movimientos de salida y regreso identificados}\newline{}Nikolai Navea\newline{}1 instancia.\newline{}\textbf{Ventana:} 3,8. & \textbf{Contenido y frontera.} Resultado único que reúne especificación, reglas, modelo propio, contratos, lógica, persistencia, vistas de tarea, pruebas unitarias y documentación de Movimientos de salida y regreso identificados. Reutiliza identidad, canales y adaptadores comunes; no reconstruye esas capacidades. Exclusiones: Otorgar zarpe de autoridad marítima y búsqueda o salvamento.\newline{}\textbf{Aceptación.} Regreso seguro exige confirmación del marinero u observación radial atribuida; declaración del armador no confirma ocupación física ni elimina discrepancias. La versión se reproduce desde fuentes controladas; casos positivos, rechazos, permisos, contratos y vistas quedan trazados. Funcional valida reglas; Datos, Seguridad y Calidad revisan sus ámbitos. Los ensayos integrales y la recepción contractual se verifican en 1.10 y T-17, sin duplicar las pruebas unitarias del componente.\newline{}\textbf{Insumos.} Para análisis: H1, requisitos, antecedentes y validadores autorizados. Antes de construir: decisiones de H2, modelos/contratos y DEV/QA recibidos. Antes de integrar: canales e interfaces comunes habilitados. El análisis no exige disponer anticipadamente de ambientes de construcción.\newline{}\textbf{Condiciones.} Etapa E1. Habilitaciones externas: 1.3.3.1, 1.7.1.3, 1.7.3.1, 1.7.3.3, H1, H2, H3. Actividades de análisis, diseño, construcción, verificación y documentación se desglosan una sola vez en T-15. Referencias: SD3 reparto y dominio, SD5 entidades/autoridad; RF-NAV-04/05; RN-01--03. Adaptaciones de ciclo de vida y campos sensibles se verifican con 1.7.6.5 y 1.3.4.4 cuando aplican. Los cinco actos sujetos a firma incorporan 1.7.6.6; no se duplica la capacidad compartida.\\\hline
\texttt{1.\allowbreak{}4.\allowbreak{}1.\allowbreak{}2.\allowbreak{}1}\newline{}\textbf{Capacidad funcional recibida: Planes voluntarios y discrepancias operacionales}\newline{}Nikolai Navea\newline{}1 instancia.\newline{}\textbf{Ventana:} 3,8. & \textbf{Contenido y frontera.} Resultado único que reúne especificación, reglas, modelo propio, contratos, lógica, persistencia, vistas de tarea, pruebas unitarias y documentación de Planes voluntarios y discrepancias operacionales. Reutiliza identidad, canales y adaptadores comunes; no reconstruye esas capacidades. Exclusiones: Notificación común y medición de ocupación por un sensor como autoridad única.\newline{}\textbf{Aceptación.} Plan conserva autor, vencimiento, cierre y discrepancia; cierre tardío no borra una actuación ni reabre un plan cerrado. La versión se reproduce desde fuentes controladas; casos positivos, rechazos, permisos, contratos y vistas quedan trazados. Funcional valida reglas; Datos, Seguridad y Calidad revisan sus ámbitos. Los ensayos integrales y la recepción contractual se verifican en 1.10 y T-17, sin duplicar las pruebas unitarias del componente.\newline{}\textbf{Insumos.} Para análisis: H1, requisitos, antecedentes y validadores autorizados. Antes de construir: decisiones de H2, modelos/contratos y DEV/QA recibidos. Antes de integrar: canales e interfaces comunes habilitados. El análisis no exige disponer anticipadamente de ambientes de construcción.\newline{}\textbf{Condiciones.} Etapa E1. Habilitaciones externas: 1.3.3.1, 1.7.1.3, 1.7.3.1, 1.7.3.3, H1, H2, H3. Actividades de análisis, diseño, construcción, verificación y documentación se desglosan una sola vez en T-15. Referencias: SD3 reparto y dominio, SD5 entidades/autoridad; RF-AUT-08; RF-NAV-06; RN-04/05. Adaptaciones de ciclo de vida y campos sensibles se verifican con 1.7.6.5 y 1.3.4.4 cuando aplican. Los cinco actos sujetos a firma incorporan 1.7.6.6; no se duplica la capacidad compartida.\\\hline
\texttt{1.\allowbreak{}4.\allowbreak{}1.\allowbreak{}3.\allowbreak{}1}\newline{}\textbf{Capacidad funcional recibida: Atención y escalamiento de planes vencidos}\newline{}Nikolai Navea\newline{}1 instancia.\newline{}\textbf{Ventana:} 3,8. & \textbf{Contenido y frontera.} Resultado único que reúne especificación, reglas, modelo propio, contratos, lógica, persistencia, vistas de tarea, pruebas unitarias y documentación de Atención y escalamiento de planes vencidos. Reutiliza identidad, canales y adaptadores comunes; no reconstruye esas capacidades. Exclusiones: Construcción del servicio común de notificaciones y ejecución de salvamento.\newline{}\textbf{Aceptación.} Vencimiento genera solicitud de alerta para guardia dentro de quince minutos, sin gracia; contacto y escalamiento conservan responsable y resultado. La versión se reproduce desde fuentes controladas; casos positivos, rechazos, permisos, contratos y vistas quedan trazados. Funcional valida reglas; Datos, Seguridad y Calidad revisan sus ámbitos. Los ensayos integrales y la recepción contractual se verifican en 1.10 y T-17, sin duplicar las pruebas unitarias del componente.\newline{}\textbf{Insumos.} Para análisis: H1, requisitos, antecedentes y validadores autorizados. Antes de construir: decisiones de H2, modelos/contratos y DEV/QA recibidos. Antes de integrar: canales e interfaces comunes habilitados. El análisis no exige disponer anticipadamente de ambientes de construcción.\newline{}\textbf{Condiciones.} Etapa E1. Habilitaciones externas: 1.3.3.1, 1.7.1.3, 1.7.3.1, 1.7.3.3, H1, H2, H3. Actividades de análisis, diseño, construcción, verificación y documentación se desglosan una sola vez en T-15. Referencias: SD3 reparto y dominio, SD5 entidades/autoridad; RF-NAV-06; DEC-21. Adaptaciones de ciclo de vida y campos sensibles se verifican con 1.7.6.5 y 1.3.4.4 cuando aplican. Los cinco actos sujetos a firma incorporan 1.7.6.6; no se duplica la capacidad compartida.\\\hline
\texttt{1.\allowbreak{}4.\allowbreak{}1.\allowbreak{}4.\allowbreak{}1}\newline{}\textbf{Capacidad funcional recibida: Autorización operacional de amarra y compatibilidad}\newline{}Nikolai Navea\newline{}1 instancia.\newline{}\textbf{Ventana:} 3,8. & \textbf{Contenido y frontera.} Resultado único que reúne especificación, reglas, modelo propio, contratos, lógica, persistencia, vistas de tarea, pruebas unitarias y documentación de Autorización operacional de amarra y compatibilidad. Reutiliza identidad, canales y adaptadores comunes; no reconstruye esas capacidades. Exclusiones: Gestión comercial del contrato y considerar reserva como ocupación comprobada.\newline{}\textbf{Aceptación.} Una autoridad local serializa puesto e intervalo y rechaza doble asignación; compatibilidad y última observación son explícitas. La versión se reproduce desde fuentes controladas; casos positivos, rechazos, permisos, contratos y vistas quedan trazados. Funcional valida reglas; Datos, Seguridad y Calidad revisan sus ámbitos. Los ensayos integrales y la recepción contractual se verifican en 1.10 y T-17, sin duplicar las pruebas unitarias del componente.\newline{}\textbf{Insumos.} Para análisis: H1, requisitos, antecedentes y validadores autorizados. Antes de construir: decisiones de H2, modelos/contratos y DEV/QA recibidos. Antes de integrar: canales e interfaces comunes habilitados. El análisis no exige disponer anticipadamente de ambientes de construcción.\newline{}\textbf{Condiciones.} Etapa E1. Habilitaciones externas: 1.3.3.1, 1.7.1.3, 1.7.3.1, 1.7.3.3, H1, H2, H3. Actividades de análisis, diseño, construcción, verificación y documentación se desglosan una sola vez en T-15. Referencias: SD3 reparto y dominio, SD5 entidades/autoridad; AMA; DEC-01/05; SD5 autoridad local. Adaptaciones de ciclo de vida y campos sensibles se verifican con 1.7.6.5 y 1.3.4.4 cuando aplican. Los cinco actos sujetos a firma incorporan 1.7.6.6; no se duplica la capacidad compartida.\\\hline
\texttt{1.\allowbreak{}4.\allowbreak{}2.\allowbreak{}1.\allowbreak{}1}\newline{}\textbf{Capacidad funcional recibida: Clases, nóminas y participación por intervalo}\newline{}Nikolai Navea\newline{}1 instancia.\newline{}\textbf{Ventana:} 3,8. & \textbf{Contenido y frontera.} Resultado único que reúne especificación, reglas, modelo propio, contratos, lógica, persistencia, vistas de tarea, pruebas unitarias y documentación de Clases, nóminas y participación por intervalo. Reutiliza identidad, canales y adaptadores comunes; no reconstruye esas capacidades. Exclusiones: Matrícula y autorizaciones migradas, plataforma de canal y prueba integral de carga.\newline{}\textbf{Aceptación.} Alumno, bote y monitor se vinculan por intervalo con cambios trazados; hasta 18 embarcaciones salen en 60 segundos en protocolo previsto. La versión se reproduce desde fuentes controladas; casos positivos, rechazos, permisos, contratos y vistas quedan trazados. Funcional valida reglas; Datos, Seguridad y Calidad revisan sus ámbitos. Los ensayos integrales y la recepción contractual se verifican en 1.10 y T-17, sin duplicar las pruebas unitarias del componente.\newline{}\textbf{Insumos.} Para análisis: H1, requisitos, antecedentes y validadores autorizados. Antes de construir: decisiones de H2, modelos/contratos y DEV/QA recibidos. Antes de integrar: canales e interfaces comunes habilitados. El análisis no exige disponer anticipadamente de ambientes de construcción.\newline{}\textbf{Condiciones.} Etapa E1. Habilitaciones externas: 1.3.3.1, 1.7.1.3, 1.7.3.1, 1.7.3.3, H1, H2, H3. Actividades de análisis, diseño, construcción, verificación y documentación se desglosan una sola vez en T-15. Referencias: SD3 reparto y dominio, SD5 entidades/autoridad; RF-ESC-04/06/07; RN-10--13. Adaptaciones de ciclo de vida y campos sensibles se verifican con 1.7.6.5 y 1.3.4.4 cuando aplican. Los cinco actos sujetos a firma incorporan 1.7.6.6; no se duplica la capacidad compartida.\\\hline
\texttt{1.\allowbreak{}4.\allowbreak{}2.\allowbreak{}2.\allowbreak{}1}\newline{}\textbf{Capacidad funcional recibida: Retiros coordinados y entrega única de alumno}\newline{}Nikolai Navea\newline{}1 instancia.\newline{}\textbf{Ventana:} 3,8. & \textbf{Contenido y frontera.} Resultado único que reúne especificación, reglas, modelo propio, contratos, lógica, persistencia, vistas de tarea, pruebas unitarias y documentación de Retiros coordinados y entrega única de alumno. Reutiliza identidad, canales y adaptadores comunes; no reconstruye esas capacidades. Exclusiones: Solicitud del apoderado tratada como entrega efectiva y reemplazar coordinación segura por radio.\newline{}\textbf{Aceptación.} Adulto e identidad/autorización se verifican; se confirma una sola entrega efectiva y se bloquea confirmación contradictoria sin estado comprobable. La versión se reproduce desde fuentes controladas; casos positivos, rechazos, permisos, contratos y vistas quedan trazados. Funcional valida reglas; Datos, Seguridad y Calidad revisan sus ámbitos. Los ensayos integrales y la recepción contractual se verifican en 1.10 y T-17, sin duplicar las pruebas unitarias del componente.\newline{}\textbf{Insumos.} Para análisis: H1, requisitos, antecedentes y validadores autorizados. Antes de construir: decisiones de H2, modelos/contratos y DEV/QA recibidos. Antes de integrar: canales e interfaces comunes habilitados. El análisis no exige disponer anticipadamente de ambientes de construcción.\newline{}\textbf{Condiciones.} Etapa E1. Habilitaciones externas: 1.3.3.1, 1.7.1.3, 1.7.3.1, 1.7.3.3, H1, H2, H3. Actividades de análisis, diseño, construcción, verificación y documentación se desglosan una sola vez en T-15. Referencias: SD3 reparto y dominio, SD5 entidades/autoridad; RF-ESC-08/09/10. Adaptaciones de ciclo de vida y campos sensibles se verifican con 1.7.6.5 y 1.3.4.4 cuando aplican. Los cinco actos sujetos a firma incorporan 1.7.6.6; no se duplica la capacidad compartida.\\\hline
\texttt{1.\allowbreak{}4.\allowbreak{}2.\allowbreak{}3.\allowbreak{}1}\newline{}\textbf{Capacidad funcional recibida: Retorno de participantes y vista mínima de salud}\newline{}Nikolai Navea\newline{}1 instancia.\newline{}\textbf{Ventana:} 3,8. & \textbf{Contenido y frontera.} Resultado único que reúne especificación, reglas, modelo propio, contratos, lógica, persistencia, vistas de tarea, pruebas unitarias y documentación de Retorno de participantes y vista mínima de salud. Reutiliza identidad, canales y adaptadores comunes; no reconstruye esas capacidades. Exclusiones: Difusión completa de ficha de salud y reconstrucción de datos ausentes.\newline{}\textbf{Aceptación.} Cada participante se confirma antes de cerrar clase; monitor ve sólo antecedentes mínimos de alumnos vinculados y se solicita aviso al apoderado. La versión se reproduce desde fuentes controladas; casos positivos, rechazos, permisos, contratos y vistas quedan trazados. Funcional valida reglas; Datos, Seguridad y Calidad revisan sus ámbitos. Los ensayos integrales y la recepción contractual se verifican en 1.10 y T-17, sin duplicar las pruebas unitarias del componente.\newline{}\textbf{Insumos.} Para análisis: H1, requisitos, antecedentes y validadores autorizados. Antes de construir: decisiones de H2, modelos/contratos y DEV/QA recibidos. Antes de integrar: canales e interfaces comunes habilitados. El análisis no exige disponer anticipadamente de ambientes de construcción.\newline{}\textbf{Condiciones.} Etapa E1. Habilitaciones externas: 1.3.3.1, 1.7.1.3, 1.7.3.1, 1.7.3.3, H1, H2, H3. Actividades de análisis, diseño, construcción, verificación y documentación se desglosan una sola vez en T-15. Referencias: SD3 reparto y dominio, SD5 entidades/autoridad; RF-ESC-14; SD5 Escuela náutica. Adaptaciones de ciclo de vida y campos sensibles se verifican con 1.7.6.5 y 1.3.4.4 cuando aplican. Los cinco actos sujetos a firma incorporan 1.7.6.6; no se duplica la capacidad compartida.\\\hline
\texttt{1.\allowbreak{}4.\allowbreak{}3.\allowbreak{}1.\allowbreak{}1}\newline{}\textbf{Capacidad funcional recibida: Inventario y condiciones de puestos y fondeos}\newline{}Nikolai Navea\newline{}1 instancia.\newline{}\textbf{Ventana:} 3,8. & \textbf{Contenido y frontera.} Resultado único que reúne especificación, reglas, modelo propio, contratos, lógica, persistencia, vistas de tarea, pruebas unitarias y documentación de Inventario y condiciones de puestos y fondeos. Reutiliza identidad, canales y adaptadores comunes; no reconstruye esas capacidades. Exclusiones: Instrumentación física adicional o piloto no seleccionado.\newline{}\textbf{Aceptación.} Puesto, boya y tren de fondeo tienen identidad, condición y última observación con antigüedad visible; una reserva no prueba presencia. La versión se reproduce desde fuentes controladas; casos positivos, rechazos, permisos, contratos y vistas quedan trazados. Funcional valida reglas; Datos, Seguridad y Calidad revisan sus ámbitos. Los ensayos integrales y la recepción contractual se verifican en 1.10 y T-17, sin duplicar las pruebas unitarias del componente.\newline{}\textbf{Insumos.} Para análisis: H1, requisitos, antecedentes y validadores autorizados. Antes de construir: decisiones de H2, modelos/contratos y DEV/QA recibidos. Antes de integrar: canales e interfaces comunes habilitados. El análisis no exige disponer anticipadamente de ambientes de construcción.\newline{}\textbf{Condiciones.} Etapa E1. Habilitaciones externas: 1.3.3.1, 1.7.1.3, 1.7.3.1, 1.7.3.3, H1, H2, H3. Actividades de análisis, diseño, construcción, verificación y documentación se desglosan una sola vez en T-15. Referencias: SD3 reparto y dominio, SD5 entidades/autoridad; BOY; AMA; SD5 Infraestructura. Adaptaciones de ciclo de vida y campos sensibles se verifican con 1.7.6.5 y 1.3.4.4 cuando aplican. Los cinco actos sujetos a firma incorporan 1.7.6.6; no se duplica la capacidad compartida.\\\hline
\texttt{1.\allowbreak{}4.\allowbreak{}3.\allowbreak{}2.\allowbreak{}1}\newline{}\textbf{Capacidad funcional recibida: Inspecciones de boyas y evidencia desconectada}\newline{}Nikolai Navea\newline{}1 instancia.\newline{}\textbf{Ventana:} 3,8. & \textbf{Contenido y frontera.} Resultado único que reúne especificación, reglas, modelo propio, contratos, lógica, persistencia, vistas de tarea, pruebas unitarias y documentación de Inspecciones de boyas y evidencia desconectada. Reutiliza identidad, canales y adaptadores comunes; no reconstruye esas capacidades. Exclusiones: Construcción de cola móvil compartida y suministro de energía/comunicaciones no seleccionadas.\newline{}\textbf{Aceptación.} Observación y fotos conservan activo, autor, tiempo y pendientes; inspección de ocho horas sin cobertura se prepara para comprobación integral. La versión se reproduce desde fuentes controladas; casos positivos, rechazos, permisos, contratos y vistas quedan trazados. Funcional valida reglas; Datos, Seguridad y Calidad revisan sus ámbitos. Los ensayos integrales y la recepción contractual se verifican en 1.10 y T-17, sin duplicar las pruebas unitarias del componente.\newline{}\textbf{Insumos.} Para análisis: H1, requisitos, antecedentes y validadores autorizados. Antes de construir: decisiones de H2, modelos/contratos y DEV/QA recibidos. Antes de integrar: canales e interfaces comunes habilitados. El análisis no exige disponer anticipadamente de ambientes de construcción.\newline{}\textbf{Condiciones.} Etapa E1. Habilitaciones externas: 1.3.3.1, 1.7.1.3, 1.7.3.1, 1.7.3.3, H1, H2, H3. Actividades de análisis, diseño, construcción, verificación y documentación se desglosan una sola vez en T-15. Referencias: SD3 reparto y dominio, SD5 entidades/autoridad; BOY; CA-13; SD3 cubierta. Adaptaciones de ciclo de vida y campos sensibles se verifican con 1.7.6.5 y 1.3.4.4 cuando aplican. Los cinco actos sujetos a firma incorporan 1.7.6.6; no se duplica la capacidad compartida.\\\hline
\texttt{1.\allowbreak{}4.\allowbreak{}3.\allowbreak{}3.\allowbreak{}1}\newline{}\textbf{Capacidad funcional recibida: Asociaciones temporales de medidores e instalaciones}\newline{}Nikolai Navea\newline{}1 instancia.\newline{}\textbf{Ventana:} 3,8. & \textbf{Contenido y frontera.} Resultado único que reúne especificación, reglas, modelo propio, contratos, lógica, persistencia, vistas de tarea, pruebas unitarias y documentación de Asociaciones temporales de medidores e instalaciones. Reutiliza identidad, canales y adaptadores comunes; no reconstruye esas capacidades. Exclusiones: Cálculo del consumo, adquisición e instalación física de medidores.\newline{}\textbf{Aceptación.} Un medidor queda asociado a servicio/amarra por período sin superposición; sustitución y época de contador conservan historial. La versión se reproduce desde fuentes controladas; casos positivos, rechazos, permisos, contratos y vistas quedan trazados. Funcional valida reglas; Datos, Seguridad y Calidad revisan sus ámbitos. Los ensayos integrales y la recepción contractual se verifican en 1.10 y T-17, sin duplicar las pruebas unitarias del componente.\newline{}\textbf{Insumos.} Para análisis: H1, requisitos, antecedentes y validadores autorizados. Antes de construir: decisiones de H2, modelos/contratos y DEV/QA recibidos. Antes de integrar: canales e interfaces comunes habilitados. El análisis no exige disponer anticipadamente de ambientes de construcción.\newline{}\textbf{Condiciones.} Etapa E1. Habilitaciones externas: 1.3.3.1, 1.7.1.3, 1.7.3.1, 1.7.3.3, H1, H2, H3. Actividades de análisis, diseño, construcción, verificación y documentación se desglosan una sola vez en T-15. Referencias: SD3 reparto y dominio, SD5 entidades/autoridad; CON; SD5 Infraestructura. Adaptaciones de ciclo de vida y campos sensibles se verifican con 1.7.6.5 y 1.3.4.4 cuando aplican. Los cinco actos sujetos a firma incorporan 1.7.6.6; no se duplica la capacidad compartida.\\\hline
\end{longtable}\endgroup
\begingroup\setlength{\tabcolsep}{3pt}
\begin{longtable}{L{\dimexpr0.320000\textwidth-2\tabcolsep\relax}L{\dimexpr0.680000\textwidth-2\tabcolsep\relax}}
\caption{Diccionario: Núcleo de Servicios}\label{tab:t14-dic-1-5}\\\hline
\EncabezadoTabla{Paquete y resultado} & \EncabezadoTabla{Frontera, aceptación e insumos}\\\hline\endfirsthead
\multicolumn{2}{l}{\textit{Continuación de la tabla \ref{tab:t14-dic-1-5}}}\\\hline
\EncabezadoTabla{Paquete y resultado} & \EncabezadoTabla{Frontera, aceptación e insumos}\\\hline\endhead
\hline\endfoot
\hline\endlastfoot
\texttt{1.\allowbreak{}5.\allowbreak{}1.\allowbreak{}1.\allowbreak{}1}\newline{}\textbf{Capacidad funcional recibida: Maestros de faena y habilitaciones de consulta inicial}\newline{}Nikolai Navea\newline{}1 instancia.\newline{}\textbf{Ventana:} 3,8. & \textbf{Contenido y frontera.} Resultado único que reúne especificación, reglas, modelo propio, contratos, lógica, persistencia, vistas de tarea, pruebas unitarias y documentación de Maestros de faena y habilitaciones de consulta inicial. Reutiliza identidad, canales y adaptadores comunes; no reconstruye esas capacidades. Exclusiones: Agenda y órdenes completas E2 y fuente laboral externa.\newline{}\textbf{Aceptación.} Activos, contratistas y vigencias necesarios para E1 tienen identidad y autoridad; una copia visible no autoriza una faena. La versión se reproduce desde fuentes controladas; casos positivos, rechazos, permisos, contratos y vistas quedan trazados. Funcional valida reglas; Datos, Seguridad y Calidad revisan sus ámbitos. Los ensayos integrales y la recepción contractual se verifican en 1.10 y T-17, sin duplicar las pruebas unitarias del componente.\newline{}\textbf{Insumos.} Para análisis: H1, requisitos, antecedentes y validadores autorizados. Antes de construir: decisiones de H2, modelos/contratos y DEV/QA recibidos. Antes de integrar: canales e interfaces comunes habilitados. El análisis no exige disponer anticipadamente de ambientes de construcción.\newline{}\textbf{Condiciones.} Etapa E1. Habilitaciones externas: 1.3.3.1, 1.7.1.3, 1.7.3.1, 1.7.3.3, H1, H2, H3. Actividades de análisis, diseño, construcción, verificación y documentación se desglosan una sola vez en T-15. Referencias: SD3 reparto y dominio, SD5 entidades/autoridad; VAR/CTR; reparto E1 SD3. Adaptaciones de ciclo de vida y campos sensibles se verifican con 1.7.6.5 y 1.3.4.4 cuando aplican. Los cinco actos sujetos a firma incorporan 1.7.6.6; no se duplica la capacidad compartida.\\\hline
\texttt{1.\allowbreak{}5.\allowbreak{}1.\allowbreak{}2.\allowbreak{}1}\newline{}\textbf{Capacidad funcional recibida: Agenda de varadero y posiciones compatibles}\newline{}Nikolai Navea\newline{}1 instancia.\newline{}\textbf{Ventana:} 13--15. & \textbf{Contenido y frontera.} Resultado único que reúne especificación, reglas, modelo propio, contratos, lógica, persistencia, vistas de tarea, pruebas unitarias y documentación de Agenda de varadero y posiciones compatibles. Reutiliza identidad, canales y adaptadores comunes; no reconstruye esas capacidades. Exclusiones: Inventario de puestos de Infraestructura y asignar capacidad adicional no prevista.\newline{}\textbf{Aceptación.} Agenda comprueba posición, recurso, intervalo y dependencias, y conserva autor de cambios y reprogramaciones. La versión se reproduce desde fuentes controladas; casos positivos, rechazos, permisos, contratos y vistas quedan trazados. Funcional valida reglas; Datos, Seguridad y Calidad revisan sus ámbitos. Los ensayos integrales y la recepción contractual se verifican en 1.10 y T-17, sin duplicar las pruebas unitarias del componente.\newline{}\textbf{Insumos.} Para análisis: H1, requisitos, antecedentes y validadores autorizados. Antes de construir: decisiones de H2, modelos/contratos y DEV/QA recibidos. Antes de integrar: canales e interfaces comunes habilitados. El análisis no exige disponer anticipadamente de ambientes de construcción.\newline{}\textbf{Condiciones.} Etapa E2. Habilitaciones externas: 1.3.3.1, 1.7.1.3, 1.7.3.1, 1.7.3.3, BASE-E1, H2, H8. Actividades de análisis, diseño, construcción, verificación y documentación se desglosan una sola vez en T-15. Referencias: SD3 reparto y dominio, SD5 entidades/autoridad; RF-VAR; 180 puestos SD3. Adaptaciones de ciclo de vida y campos sensibles se verifican con 1.7.6.5 y 1.3.4.4 cuando aplican. Los cinco actos sujetos a firma incorporan 1.7.6.6; no se duplica la capacidad compartida.\\\hline
\texttt{1.\allowbreak{}5.\allowbreak{}1.\allowbreak{}3.\allowbreak{}1}\newline{}\textbf{Capacidad funcional recibida: Plan de izaje, validación y suspensión de faena}\newline{}Nikolai Navea\newline{}1 instancia.\newline{}\textbf{Ventana:} 13--15. & \textbf{Contenido y frontera.} Resultado único que reúne especificación, reglas, modelo propio, contratos, lógica, persistencia, vistas de tarea, pruebas unitarias y documentación de Plan de izaje, validación y suspensión de faena. Reutiliza identidad, canales y adaptadores comunes; no reconstruye esas capacidades. Exclusiones: Enclavamiento físico del travelift y sustituir la decisión de un operador calificado.\newline{}\textbf{Aceptación.} Plan identifica eslingado, restricciones y validador calificado; aviso oficial suspende; no se exige interacción con carga suspendida. La versión se reproduce desde fuentes controladas; casos positivos, rechazos, permisos, contratos y vistas quedan trazados. Funcional valida reglas; Datos, Seguridad y Calidad revisan sus ámbitos. Los ensayos integrales y la recepción contractual se verifican en 1.10 y T-17, sin duplicar las pruebas unitarias del componente.\newline{}\textbf{Insumos.} Para análisis: H1, requisitos, antecedentes y validadores autorizados. Antes de construir: decisiones de H2, modelos/contratos y DEV/QA recibidos. Antes de integrar: canales e interfaces comunes habilitados. El análisis no exige disponer anticipadamente de ambientes de construcción.\newline{}\textbf{Condiciones.} Etapa E2. Habilitaciones externas: 1.3.3.1, 1.7.1.3, 1.7.3.1, 1.7.3.3, BASE-E1, H2, H8. Actividades de análisis, diseño, construcción, verificación y documentación se desglosan una sola vez en T-15. Referencias: SD3 reparto y dominio, SD5 entidades/autoridad; RF-VAR-04--08; DEC-12/13; RN-15/16. Adaptaciones de ciclo de vida y campos sensibles se verifican con 1.7.6.5 y 1.3.4.4 cuando aplican. Los cinco actos sujetos a firma incorporan 1.7.6.6; no se duplica la capacidad compartida.\\\hline
\texttt{1.\allowbreak{}5.\allowbreak{}1.\allowbreak{}4.\allowbreak{}1}\newline{}\textbf{Capacidad funcional recibida: Órdenes, ejecución y reprogramación de trabajos}\newline{}Nikolai Navea\newline{}1 instancia.\newline{}\textbf{Ventana:} 13--15. & \textbf{Contenido y frontera.} Resultado único que reúne especificación, reglas, modelo propio, contratos, lógica, persistencia, vistas de tarea, pruebas unitarias y documentación de Órdenes, ejecución y reprogramación de trabajos. Reutiliza identidad, canales y adaptadores comunes; no reconstruye esas capacidades. Exclusiones: Solicitudes adicionales de mantenimiento IN1 y mantenimiento preventivo IN2.\newline{}\textbf{Aceptación.} Orden conserva aceptación, actuación, evidencia y cambio de calendario; cada reprogramación comprueba posiciones y dependencias antes de notificar. La versión se reproduce desde fuentes controladas; casos positivos, rechazos, permisos, contratos y vistas quedan trazados. Funcional valida reglas; Datos, Seguridad y Calidad revisan sus ámbitos. Los ensayos integrales y la recepción contractual se verifican en 1.10 y T-17, sin duplicar las pruebas unitarias del componente.\newline{}\textbf{Insumos.} Para análisis: H1, requisitos, antecedentes y validadores autorizados. Antes de construir: decisiones de H2, modelos/contratos y DEV/QA recibidos. Antes de integrar: canales e interfaces comunes habilitados. El análisis no exige disponer anticipadamente de ambientes de construcción.\newline{}\textbf{Condiciones.} Etapa E2. Habilitaciones externas: 1.3.3.1, 1.7.1.3, 1.7.3.1, 1.7.3.3, BASE-E1, H2, H8. Actividades de análisis, diseño, construcción, verificación y documentación se desglosan una sola vez en T-15. Referencias: SD3 reparto y dominio, SD5 entidades/autoridad; SD5 Varadero y contratistas. Adaptaciones de ciclo de vida y campos sensibles se verifican con 1.7.6.5 y 1.3.4.4 cuando aplican. Los cinco actos sujetos a firma incorporan 1.7.6.6; no se duplica la capacidad compartida.\\\hline
\texttt{1.\allowbreak{}5.\allowbreak{}1.\allowbreak{}5.\allowbreak{}1}\newline{}\textbf{Capacidad funcional recibida: Acreditación, inducción y acceso a intervención}\newline{}Nikolai Navea\newline{}1 instancia.\newline{}\textbf{Ventana:} 13--15. & \textbf{Contenido y frontera.} Resultado único que reúne especificación, reglas, modelo propio, contratos, lógica, persistencia, vistas de tarea, pruebas unitarias y documentación de Acreditación, inducción y acceso a intervención. Reutiliza identidad, canales y adaptadores comunes; no reconstruye esas capacidades. Exclusiones: Administrar sistema laboral externo y automatizar un control físico sin contrato compatible.\newline{}\textbf{Aceptación.} Acceso exige identidad, vigencia, inducción e intervención autorizada; QR identifica sin conceder habilitación autónoma. La versión se reproduce desde fuentes controladas; casos positivos, rechazos, permisos, contratos y vistas quedan trazados. Funcional valida reglas; Datos, Seguridad y Calidad revisan sus ámbitos. Los ensayos integrales y la recepción contractual se verifican en 1.10 y T-17, sin duplicar las pruebas unitarias del componente.\newline{}\textbf{Insumos.} Para análisis: H1, requisitos, antecedentes y validadores autorizados. Antes de construir: decisiones de H2, modelos/contratos y DEV/QA recibidos. Antes de integrar: canales e interfaces comunes habilitados. El análisis no exige disponer anticipadamente de ambientes de construcción.\newline{}\textbf{Condiciones.} Etapa E2. Habilitaciones externas: 1.3.3.1, 1.7.1.3, 1.7.3.1, 1.7.3.3, BASE-E1, H2, H8. Actividades de análisis, diseño, construcción, verificación y documentación se desglosan una sola vez en T-15. Referencias: SD3 reparto y dominio, SD5 entidades/autoridad; RF-CTR; SD3 contratistas/portería. Adaptaciones de ciclo de vida y campos sensibles se verifican con 1.7.6.5 y 1.3.4.4 cuando aplican. Los cinco actos sujetos a firma incorporan 1.7.6.6; no se duplica la capacidad compartida.\\\hline
\texttt{1.\allowbreak{}5.\allowbreak{}2.\allowbreak{}1.\allowbreak{}1}\newline{}\textbf{Capacidad funcional recibida: Despacho identificado y captura segura de combustible}\newline{}Nikolai Navea\newline{}1 instancia.\newline{}\textbf{Ventana:} 3,8. & \textbf{Contenido y frontera.} Resultado único que reúne especificación, reglas, modelo propio, contratos, lógica, persistencia, vistas de tarea, pruebas unitarias y documentación de Despacho identificado y captura segura de combustible. Reutiliza identidad, canales y adaptadores comunes; no reconstruye esas capacidades. Exclusiones: Construcción de broker/sincronizador y certificación o sustitución del equipo físico.\newline{}\textbf{Aceptación.} Despacho conserva operador, nave, producto, instrumento, lecturas, unidad y evidencia en momento seguro; hecho local no depende de contabilidad. La versión se reproduce desde fuentes controladas; casos positivos, rechazos, permisos, contratos y vistas quedan trazados. Funcional valida reglas; Datos, Seguridad y Calidad revisan sus ámbitos. Los ensayos integrales y la recepción contractual se verifican en 1.10 y T-17, sin duplicar las pruebas unitarias del componente.\newline{}\textbf{Insumos.} Para análisis: H1, requisitos, antecedentes y validadores autorizados. Antes de construir: decisiones de H2, modelos/contratos y DEV/QA recibidos. Antes de integrar: canales e interfaces comunes habilitados. El análisis no exige disponer anticipadamente de ambientes de construcción.\newline{}\textbf{Condiciones.} Etapa E1. Habilitaciones externas: 1.3.3.1, 1.7.1.3, 1.7.3.1, 1.7.3.3, H1, H2, H3. Actividades de análisis, diseño, construcción, verificación y documentación se desglosan una sola vez en T-15. Referencias: SD3 reparto y dominio, SD5 entidades/autoridad; CON; C07 RT-03.10/RT-05.10. Adaptaciones de ciclo de vida y campos sensibles se verifican con 1.7.6.5 y 1.3.4.4 cuando aplican. Los cinco actos sujetos a firma incorporan 1.7.6.6; no se duplica la capacidad compartida.\\\hline
\texttt{1.\allowbreak{}5.\allowbreak{}2.\allowbreak{}2.\allowbreak{}1}\newline{}\textbf{Capacidad funcional recibida: Conciliación de tramos y rectificación de despachos}\newline{}Nikolai Navea\newline{}1 instancia.\newline{}\textbf{Ventana:} 3,8. & \textbf{Contenido y frontera.} Resultado único que reúne especificación, reglas, modelo propio, contratos, lógica, persistencia, vistas de tarea, pruebas unitarias y documentación de Conciliación de tramos y rectificación de despachos. Reutiliza identidad, canales y adaptadores comunes; no reconstruye esas capacidades. Exclusiones: Saldo contable autoritativo y convertir diferencia negativa en consumo válido.\newline{}\textbf{Aceptación.} Reinicio/escala o lectura inválida abre tramo/anomalía; sólo volumen verificable se consolida; rectificación conserva original. La versión se reproduce desde fuentes controladas; casos positivos, rechazos, permisos, contratos y vistas quedan trazados. Funcional valida reglas; Datos, Seguridad y Calidad revisan sus ámbitos. Los ensayos integrales y la recepción contractual se verifican en 1.10 y T-17, sin duplicar las pruebas unitarias del componente.\newline{}\textbf{Insumos.} Para análisis: H1, requisitos, antecedentes y validadores autorizados. Antes de construir: decisiones de H2, modelos/contratos y DEV/QA recibidos. Antes de integrar: canales e interfaces comunes habilitados. El análisis no exige disponer anticipadamente de ambientes de construcción.\newline{}\textbf{Condiciones.} Etapa E1. Habilitaciones externas: 1.3.3.1, 1.7.1.3, 1.7.3.1, 1.7.3.3, H1, H2, H3. Actividades de análisis, diseño, construcción, verificación y documentación se desglosan una sola vez en T-15. Referencias: SD3 reparto y dominio, SD5 entidades/autoridad; SD5 Combustible. Adaptaciones de ciclo de vida y campos sensibles se verifican con 1.7.6.5 y 1.3.4.4 cuando aplican. Los cinco actos sujetos a firma incorporan 1.7.6.6; no se duplica la capacidad compartida.\\\hline
\texttt{1.\allowbreak{}5.\allowbreak{}3.\allowbreak{}1.\allowbreak{}1}\newline{}\textbf{Capacidad funcional recibida: Intervalos y series conciliadas de agua y energía}\newline{}Nikolai Navea\newline{}1 instancia.\newline{}\textbf{Ventana:} 3,8. & \textbf{Contenido y frontera.} Resultado único que reúne especificación, reglas, modelo propio, contratos, lógica, persistencia, vistas de tarea, pruebas unitarias y documentación de Intervalos y series conciliadas de agua y energía. Reutiliza identidad, canales y adaptadores comunes; no reconstruye esas capacidades. Exclusiones: Captura técnica MQTT y cálculo comercial completo E2.\newline{}\textbf{Aceptación.} Intervalo mantiene lecturas originales, unidad, época, calidad y asociación temporal; no se suman totalizadores como volumen. La versión se reproduce desde fuentes controladas; casos positivos, rechazos, permisos, contratos y vistas quedan trazados. Funcional valida reglas; Datos, Seguridad y Calidad revisan sus ámbitos. Los ensayos integrales y la recepción contractual se verifican en 1.10 y T-17, sin duplicar las pruebas unitarias del componente.\newline{}\textbf{Insumos.} Para análisis: H1, requisitos, antecedentes y validadores autorizados. Antes de construir: decisiones de H2, modelos/contratos y DEV/QA recibidos. Antes de integrar: canales e interfaces comunes habilitados. El análisis no exige disponer anticipadamente de ambientes de construcción.\newline{}\textbf{Condiciones.} Etapa E1. Habilitaciones externas: 1.3.3.1, 1.7.1.3, 1.7.3.1, 1.7.3.3, H1, H2, H3. Actividades de análisis, diseño, construcción, verificación y documentación se desglosan una sola vez en T-15. Referencias: SD3 reparto y dominio, SD5 entidades/autoridad; CON/AMB; CA-18/19. Adaptaciones de ciclo de vida y campos sensibles se verifican con 1.7.6.5 y 1.3.4.4 cuando aplican. Los cinco actos sujetos a firma incorporan 1.7.6.6; no se duplica la capacidad compartida.\\\hline
\texttt{1.\allowbreak{}5.\allowbreak{}3.\allowbreak{}2.\allowbreak{}1}\newline{}\textbf{Capacidad funcional recibida: Consultas de consumo y evidencia ambiental inicial}\newline{}Nikolai Navea\newline{}1 instancia.\newline{}\textbf{Ventana:} 3,8. & \textbf{Contenido y frontera.} Resultado único que reúne especificación, reglas, modelo propio, contratos, lógica, persistencia, vistas de tarea, pruebas unitarias y documentación de Consultas de consumo y evidencia ambiental inicial. Reutiliza identidad, canales y adaptadores comunes; no reconstruye esas capacidades. Exclusiones: Pasaporte IN5 y presumir serie previa a 2029 sin calendario real.\newline{}\textbf{Aceptación.} Vista conserva período, cobertura, versión y fuente; evidencia incompleta impide afirmar cumplimiento ambiental o reducción. La versión se reproduce desde fuentes controladas; casos positivos, rechazos, permisos, contratos y vistas quedan trazados. Funcional valida reglas; Datos, Seguridad y Calidad revisan sus ámbitos. Los ensayos integrales y la recepción contractual se verifican en 1.10 y T-17, sin duplicar las pruebas unitarias del componente.\newline{}\textbf{Insumos.} Para análisis: H1, requisitos, antecedentes y validadores autorizados. Antes de construir: decisiones de H2, modelos/contratos y DEV/QA recibidos. Antes de integrar: canales e interfaces comunes habilitados. El análisis no exige disponer anticipadamente de ambientes de construcción.\newline{}\textbf{Condiciones.} Etapa E1. Habilitaciones externas: 1.3.3.1, 1.7.1.3, 1.7.3.1, 1.7.3.3, H1, H2, H3. Actividades de análisis, diseño, construcción, verificación y documentación se desglosan una sola vez en T-15. Referencias: SD3 reparto y dominio, SD5 entidades/autoridad; AMB; CA-19; SD3 historia ambiental. Adaptaciones de ciclo de vida y campos sensibles se verifican con 1.7.6.5 y 1.3.4.4 cuando aplican. Los cinco actos sujetos a firma incorporan 1.7.6.6; no se duplica la capacidad compartida.\\\hline
\texttt{1.\allowbreak{}5.\allowbreak{}3.\allowbreak{}3.\allowbreak{}1}\newline{}\textbf{Capacidad funcional recibida: Recepción de residuos y discrepancias de origen}\newline{}Nikolai Navea\newline{}1 instancia.\newline{}\textbf{Ventana:} 13--15. & \textbf{Contenido y frontera.} Resultado único que reúne especificación, reglas, modelo propio, contratos, lógica, persistencia, vistas de tarea, pruebas unitarias y documentación de Recepción de residuos y discrepancias de origen. Reutiliza identidad, canales y adaptadores comunes; no reconstruye esas capacidades. Exclusiones: Investigación física ejecutada por software y emitir certificado de gestor.\newline{}\textbf{Aceptación.} Registro conserva origen, cantidad, unidad y evidencia; origen desconocido permanece observado con investigación identificada. La versión se reproduce desde fuentes controladas; casos positivos, rechazos, permisos, contratos y vistas quedan trazados. Funcional valida reglas; Datos, Seguridad y Calidad revisan sus ámbitos. Los ensayos integrales y la recepción contractual se verifican en 1.10 y T-17, sin duplicar las pruebas unitarias del componente.\newline{}\textbf{Insumos.} Para análisis: H1, requisitos, antecedentes y validadores autorizados. Antes de construir: decisiones de H2, modelos/contratos y DEV/QA recibidos. Antes de integrar: canales e interfaces comunes habilitados. El análisis no exige disponer anticipadamente de ambientes de construcción.\newline{}\textbf{Condiciones.} Etapa E2. Habilitaciones externas: 1.3.3.1, 1.7.1.3, 1.7.3.1, 1.7.3.3, BASE-E1, H2, H8. Actividades de análisis, diseño, construcción, verificación y documentación se desglosan una sola vez en T-15. Referencias: SD3 reparto y dominio, SD5 entidades/autoridad; AMB; RN-19--21; SD5 residuos. Adaptaciones de ciclo de vida y campos sensibles se verifican con 1.7.6.5 y 1.3.4.4 cuando aplican. Los cinco actos sujetos a firma incorporan 1.7.6.6; no se duplica la capacidad compartida.\\\hline
\texttt{1.\allowbreak{}5.\allowbreak{}3.\allowbreak{}4.\allowbreak{}1}\newline{}\textbf{Capacidad funcional recibida: Entregas RESPEL parciales y certificados vinculados}\newline{}Nikolai Navea\newline{}1 instancia.\newline{}\textbf{Ventana:} 13--15. & \textbf{Contenido y frontera.} Resultado único que reúne especificación, reglas, modelo propio, contratos, lógica, persistencia, vistas de tarea, pruebas unitarias y documentación de Entregas RESPEL parciales y certificados vinculados. Reutiliza identidad, canales y adaptadores comunes; no reconstruye esas capacidades. Exclusiones: Otorgar certificación externa y borrar retiro real por discrepancia de origen.\newline{}\textbf{Aceptación.} Detalle permite dividir/agrupar residuos sin exceder disponibilidad; retiro y certificado se vinculan a entrega efectiva y diferencias quedan trazadas. La versión se reproduce desde fuentes controladas; casos positivos, rechazos, permisos, contratos y vistas quedan trazados. Funcional valida reglas; Datos, Seguridad y Calidad revisan sus ámbitos. Los ensayos integrales y la recepción contractual se verifican en 1.10 y T-17, sin duplicar las pruebas unitarias del componente.\newline{}\textbf{Insumos.} Para análisis: H1, requisitos, antecedentes y validadores autorizados. Antes de construir: decisiones de H2, modelos/contratos y DEV/QA recibidos. Antes de integrar: canales e interfaces comunes habilitados. El análisis no exige disponer anticipadamente de ambientes de construcción.\newline{}\textbf{Condiciones.} Etapa E2. Habilitaciones externas: 1.3.3.1, 1.7.1.3, 1.7.3.1, 1.7.3.3, BASE-E1, H2, H8. Actividades de análisis, diseño, construcción, verificación y documentación se desglosan una sola vez en T-15. Referencias: SD3 reparto y dominio, SD5 entidades/autoridad; AMB; SD5 Entrega a gestor. Adaptaciones de ciclo de vida y campos sensibles se verifican con 1.7.6.5 y 1.3.4.4 cuando aplican. Los cinco actos sujetos a firma incorporan 1.7.6.6; no se duplica la capacidad compartida.\\\hline
\texttt{1.\allowbreak{}5.\allowbreak{}3.\allowbreak{}5.\allowbreak{}1}\newline{}\textbf{Capacidad funcional recibida: Indicadores ambientales e imputación comercial de consumo}\newline{}Nikolai Navea\newline{}1 instancia.\newline{}\textbf{Ventana:} 13--15. & \textbf{Contenido y frontera.} Resultado único que reúne especificación, reglas, modelo propio, contratos, lógica, persistencia, vistas de tarea, pruebas unitarias y documentación de Indicadores ambientales e imputación comercial de consumo. Reutiliza identidad, canales y adaptadores comunes; no reconstruye esas capacidades. Exclusiones: Cobro tributario, incentivo IN4 y beneficio ambiental voluntario IN5.\newline{}\textbf{Aceptación.} Indicador tiene período, regla y registros de soporte; sólo intervalos conciliados producen hecho de consumo para Finanzas sin duplicados. La versión se reproduce desde fuentes controladas; casos positivos, rechazos, permisos, contratos y vistas quedan trazados. Funcional valida reglas; Datos, Seguridad y Calidad revisan sus ámbitos. Los ensayos integrales y la recepción contractual se verifican en 1.10 y T-17, sin duplicar las pruebas unitarias del componente.\newline{}\textbf{Insumos.} Para análisis: H1, requisitos, antecedentes y validadores autorizados. Antes de construir: decisiones de H2, modelos/contratos y DEV/QA recibidos. Antes de integrar: canales e interfaces comunes habilitados. El análisis no exige disponer anticipadamente de ambientes de construcción.\newline{}\textbf{Condiciones.} Etapa E2. Habilitaciones externas: 1.3.3.1, 1.7.1.3, 1.7.3.1, 1.7.3.3, BASE-E1, H2, H8. Actividades de análisis, diseño, construcción, verificación y documentación se desglosan una sola vez en T-15. Referencias: SD3 reparto y dominio, SD5 entidades/autoridad; CA-18/19; SD5 MediciónConsolidada. Adaptaciones de ciclo de vida y campos sensibles se verifican con 1.7.6.5 y 1.3.4.4 cuando aplican. Los cinco actos sujetos a firma incorporan 1.7.6.6; no se duplica la capacidad compartida.\\\hline
\end{longtable}\endgroup
\begingroup\setlength{\tabcolsep}{3pt}
\begin{longtable}{L{\dimexpr0.320000\textwidth-2\tabcolsep\relax}L{\dimexpr0.680000\textwidth-2\tabcolsep\relax}}
\caption{Diccionario: Núcleo Administrativo}\label{tab:t14-dic-1-6}\\\hline
\EncabezadoTabla{Paquete y resultado} & \EncabezadoTabla{Frontera, aceptación e insumos}\\\hline\endfirsthead
\multicolumn{2}{l}{\textit{Continuación de la tabla \ref{tab:t14-dic-1-6}}}\\\hline
\EncabezadoTabla{Paquete y resultado} & \EncabezadoTabla{Frontera, aceptación e insumos}\\\hline\endhead
\hline\endfoot
\hline\endlastfoot
\texttt{1.\allowbreak{}6.\allowbreak{}1.\allowbreak{}1.\allowbreak{}1}\newline{}\textbf{Capacidad funcional recibida: Maestros de personas, organizaciones y embarcaciones}\newline{}Nikolai Navea\newline{}1 instancia.\newline{}\textbf{Ventana:} 3,8. & \textbf{Contenido y frontera.} Resultado único que reúne especificación, reglas, modelo propio, contratos, lógica, persistencia, vistas de tarea, pruebas unitarias y documentación de Maestros de personas, organizaciones y embarcaciones. Reutiliza identidad, canales y adaptadores comunes; no reconstruye esas capacidades. Exclusiones: Extracción/saneamiento de legado y autenticación técnica Cognito.\newline{}\textbf{Aceptación.} Identidad estable, vínculos y vigencias se consultan por contrato; persona y embarcación no se duplican por cambio de nombre o canal. La versión se reproduce desde fuentes controladas; casos positivos, rechazos, permisos, contratos y vistas quedan trazados. Funcional valida reglas; Datos, Seguridad y Calidad revisan sus ámbitos. Los ensayos integrales y la recepción contractual se verifican en 1.10 y T-17, sin duplicar las pruebas unitarias del componente.\newline{}\textbf{Insumos.} Para análisis: H1, requisitos, antecedentes y validadores autorizados. Antes de construir: decisiones de H2, modelos/contratos y DEV/QA recibidos. Antes de integrar: canales e interfaces comunes habilitados. El análisis no exige disponer anticipadamente de ambientes de construcción.\newline{}\textbf{Condiciones.} Etapa E1. Habilitaciones externas: 1.3.3.1, 1.7.1.3, 1.7.3.1, 1.7.3.3, H1, H2, H3. Actividades de análisis, diseño, construcción, verificación y documentación se desglosan una sola vez en T-15. Referencias: SD3 reparto y dominio, SD5 entidades/autoridad; ADM/AUT; SD5 Clientes y contratos. Adaptaciones de ciclo de vida y campos sensibles se verifican con 1.7.6.5 y 1.3.4.4 cuando aplican. Los cinco actos sujetos a firma incorporan 1.7.6.6; no se duplica la capacidad compartida.\\\hline
\texttt{1.\allowbreak{}6.\allowbreak{}1.\allowbreak{}2.\allowbreak{}1}\newline{}\textbf{Capacidad funcional recibida: Contratos y documentos vigentes de consulta E1}\newline{}Nikolai Navea\newline{}1 instancia.\newline{}\textbf{Ventana:} 3,8. & \textbf{Contenido y frontera.} Resultado único que reúne especificación, reglas, modelo propio, contratos, lógica, persistencia, vistas de tarea, pruebas unitarias y documentación de Contratos y documentos vigentes de consulta E1. Reutiliza identidad, canales y adaptadores comunes; no reconstruye esas capacidades. Exclusiones: OCR IN3 y revisión documental íntegra de migración.\newline{}\textbf{Aceptación.} Contrato, parte, documento y vigencia tienen autoridad y versión; acceso respeta ámbito y un documento pendiente no se declara vigente. La versión se reproduce desde fuentes controladas; casos positivos, rechazos, permisos, contratos y vistas quedan trazados. Funcional valida reglas; Datos, Seguridad y Calidad revisan sus ámbitos. Los ensayos integrales y la recepción contractual se verifican en 1.10 y T-17, sin duplicar las pruebas unitarias del componente.\newline{}\textbf{Insumos.} Para análisis: H1, requisitos, antecedentes y validadores autorizados. Antes de construir: decisiones de H2, modelos/contratos y DEV/QA recibidos. Antes de integrar: canales e interfaces comunes habilitados. El análisis no exige disponer anticipadamente de ambientes de construcción.\newline{}\textbf{Condiciones.} Etapa E1. Habilitaciones externas: 1.3.3.1, 1.7.1.3, 1.7.3.1, 1.7.3.3, H1, H2, H3. Actividades de análisis, diseño, construcción, verificación y documentación se desglosan una sola vez en T-15. Referencias: SD3 reparto y dominio, SD5 entidades/autoridad; SD3 reparto E1; SD5 contratos. Adaptaciones de ciclo de vida y campos sensibles se verifican con 1.7.6.5 y 1.3.4.4 cuando aplican. Los cinco actos sujetos a firma incorporan 1.7.6.6; no se duplica la capacidad compartida.\\\hline
\texttt{1.\allowbreak{}6.\allowbreak{}1.\allowbreak{}3.\allowbreak{}1}\newline{}\textbf{Capacidad funcional recibida: Atención a visitas y asignación transitoria inicial}\newline{}Nikolai Navea\newline{}1 instancia.\newline{}\textbf{Ventana:} 3,8. & \textbf{Contenido y frontera.} Resultado único que reúne especificación, reglas, modelo propio, contratos, lógica, persistencia, vistas de tarea, pruebas unitarias y documentación de Atención a visitas y asignación transitoria inicial. Reutiliza identidad, canales y adaptadores comunes; no reconstruye esas capacidades. Exclusiones: Emisión tributaria y doble autoridad sobre ocupación física.\newline{}\textbf{Aceptación.} Solicitud conserva identidad, período y estado; confirmación requiere compatibilidad y autoridad de puesto, con información español/inglés prevista. La versión se reproduce desde fuentes controladas; casos positivos, rechazos, permisos, contratos y vistas quedan trazados. Funcional valida reglas; Datos, Seguridad y Calidad revisan sus ámbitos. Los ensayos integrales y la recepción contractual se verifican en 1.10 y T-17, sin duplicar las pruebas unitarias del componente.\newline{}\textbf{Insumos.} Para análisis: H1, requisitos, antecedentes y validadores autorizados. Antes de construir: decisiones de H2, modelos/contratos y DEV/QA recibidos. Antes de integrar: canales e interfaces comunes habilitados. El análisis no exige disponer anticipadamente de ambientes de construcción.\newline{}\textbf{Condiciones.} Etapa E1. Habilitaciones externas: 1.3.3.1, 1.7.1.3, 1.7.3.1, 1.7.3.3, H1, H2, H3. Actividades de análisis, diseño, construcción, verificación y documentación se desglosan una sola vez en T-15. Referencias: SD3 reparto y dominio, SD5 entidades/autoridad; AMA/AUT; SD3 atención náutica. Adaptaciones de ciclo de vida y campos sensibles se verifican con 1.7.6.5 y 1.3.4.4 cuando aplican. Los cinco actos sujetos a firma incorporan 1.7.6.6; no se duplica la capacidad compartida.\\\hline
\texttt{1.\allowbreak{}6.\allowbreak{}1.\allowbreak{}4.\allowbreak{}1}\newline{}\textbf{Capacidad funcional recibida: Espera y gestión contractual completa}\newline{}Nikolai Navea\newline{}1 instancia.\newline{}\textbf{Ventana:} 13--15. & \textbf{Contenido y frontera.} Resultado único que reúne especificación, reglas, modelo propio, contratos, lógica, persistencia, vistas de tarea, pruebas unitarias y documentación de Espera y gestión contractual completa. Reutiliza identidad, canales y adaptadores comunes; no reconstruye esas capacidades. Exclusiones: Decidir prioridades comerciales por Synaptix y rehacer compatibilidad operacional E1.\newline{}\textbf{Aceptación.} Orden de espera, prioridad y excepción quedan autorizados; contrato y asignación se concilian sin cambiar hechos de ocupación de otro dominio. La versión se reproduce desde fuentes controladas; casos positivos, rechazos, permisos, contratos y vistas quedan trazados. Funcional valida reglas; Datos, Seguridad y Calidad revisan sus ámbitos. Los ensayos integrales y la recepción contractual se verifican en 1.10 y T-17, sin duplicar las pruebas unitarias del componente.\newline{}\textbf{Insumos.} Para análisis: H1, requisitos, antecedentes y validadores autorizados. Antes de construir: decisiones de H2, modelos/contratos y DEV/QA recibidos. Antes de integrar: canales e interfaces comunes habilitados. El análisis no exige disponer anticipadamente de ambientes de construcción.\newline{}\textbf{Condiciones.} Etapa E2. Habilitaciones externas: 1.3.3.1, 1.7.1.3, 1.7.3.1, 1.7.3.3, BASE-E1, H2, H8. Actividades de análisis, diseño, construcción, verificación y documentación se desglosan una sola vez en T-15. Referencias: SD3 reparto y dominio, SD5 entidades/autoridad; AMA/ADM; SD3 reparto E2. Adaptaciones de ciclo de vida y campos sensibles se verifican con 1.7.6.5 y 1.3.4.4 cuando aplican. Los cinco actos sujetos a firma incorporan 1.7.6.6; no se duplica la capacidad compartida.\\\hline
\texttt{1.\allowbreak{}6.\allowbreak{}2.\allowbreak{}1.\allowbreak{}1}\newline{}\textbf{Capacidad funcional recibida: Hechos facturables, tarifas y lotes de envío}\newline{}Nikolai Navea\newline{}1 instancia.\newline{}\textbf{Ventana:} 13--15. & \textbf{Contenido y frontera.} Resultado único que reúne especificación, reglas, modelo propio, contratos, lógica, persistencia, vistas de tarea, pruebas unitarias y documentación de Hechos facturables, tarifas y lotes de envío. Reutiliza identidad, canales y adaptadores comunes; no reconstruye esas capacidades. Exclusiones: Adaptador contable técnico y emisión de documentos tributarios.\newline{}\textbf{Aceptación.} Cada hecho conserva servicio, cantidad, unidad, moneda, tarifa aplicable y responsable de pago; lote no duplica hechos remitidos. La versión se reproduce desde fuentes controladas; casos positivos, rechazos, permisos, contratos y vistas quedan trazados. Funcional valida reglas; Datos, Seguridad y Calidad revisan sus ámbitos. Los ensayos integrales y la recepción contractual se verifican en 1.10 y T-17, sin duplicar las pruebas unitarias del componente.\newline{}\textbf{Insumos.} Para análisis: H1, requisitos, antecedentes y validadores autorizados. Antes de construir: decisiones de H2, modelos/contratos y DEV/QA recibidos. Antes de integrar: canales e interfaces comunes habilitados. El análisis no exige disponer anticipadamente de ambientes de construcción.\newline{}\textbf{Condiciones.} Etapa E2. Habilitaciones externas: 1.3.3.1, 1.7.1.3, 1.7.3.1, 1.7.3.3, BASE-E1, H2, H8. Actividades de análisis, diseño, construcción, verificación y documentación se desglosan una sola vez en T-15. Referencias: SD3 reparto y dominio, SD5 entidades/autoridad; FAC; SD5 Finanzas y cobranza. Adaptaciones de ciclo de vida y campos sensibles se verifican con 1.7.6.5 y 1.3.4.4 cuando aplican. Los cinco actos sujetos a firma incorporan 1.7.6.6; no se duplica la capacidad compartida.\\\hline
\texttt{1.\allowbreak{}6.\allowbreak{}2.\allowbreak{}2.\allowbreak{}1}\newline{}\textbf{Capacidad funcional recibida: Respuesta contable, pagos y saldos informados}\newline{}Nikolai Navea\newline{}1 instancia.\newline{}\textbf{Ventana:} 13--15. & \textbf{Contenido y frontera.} Resultado único que reúne especificación, reglas, modelo propio, contratos, lógica, persistencia, vistas de tarea, pruebas unitarias y documentación de Respuesta contable, pagos y saldos informados. Reutiliza identidad, canales y adaptadores comunes; no reconstruye esas capacidades. Exclusiones: Construcción de webhooks/pasarela y modificar saldo externo sin autorización.\newline{}\textbf{Aceptación.} Aceptación/rechazo y saldo conservan fuente y fecha; un saldo local atrasado se identifica y no sustituye autoridad contable. La versión se reproduce desde fuentes controladas; casos positivos, rechazos, permisos, contratos y vistas quedan trazados. Funcional valida reglas; Datos, Seguridad y Calidad revisan sus ámbitos. Los ensayos integrales y la recepción contractual se verifican en 1.10 y T-17, sin duplicar las pruebas unitarias del componente.\newline{}\textbf{Insumos.} Para análisis: H1, requisitos, antecedentes y validadores autorizados. Antes de construir: decisiones de H2, modelos/contratos y DEV/QA recibidos. Antes de integrar: canales e interfaces comunes habilitados. El análisis no exige disponer anticipadamente de ambientes de construcción.\newline{}\textbf{Condiciones.} Etapa E2. Habilitaciones externas: 1.3.3.1, 1.7.1.3, 1.7.3.1, 1.7.3.3, BASE-E1, H2, H8. Actividades de análisis, diseño, construcción, verificación y documentación se desglosan una sola vez en T-15. Referencias: SD3 reparto y dominio, SD5 entidades/autoridad; FAC; SD5 autoridad externa. Adaptaciones de ciclo de vida y campos sensibles se verifican con 1.7.6.5 y 1.3.4.4 cuando aplican. Los cinco actos sujetos a firma incorporan 1.7.6.6; no se duplica la capacidad compartida.\\\hline
\texttt{1.\allowbreak{}6.\allowbreak{}2.\allowbreak{}3.\allowbreak{}1}\newline{}\textbf{Capacidad funcional recibida: Cobranza y medidas con decisión competente}\newline{}Nikolai Navea\newline{}1 instancia.\newline{}\textbf{Ventana:} 13--15. & \textbf{Contenido y frontera.} Resultado único que reúne especificación, reglas, modelo propio, contratos, lógica, persistencia, vistas de tarea, pruebas unitarias y documentación de Cobranza y medidas con decisión competente. Reutiliza identidad, canales y adaptadores comunes; no reconstruye esas capacidades. Exclusiones: Adoptar decisiones jurídicas o físicas por software.\newline{}\textbf{Aceptación.} Actuación registra propuesta, autorización y ejecución; tramo de deuda no activa restricción física ni desposesión automática. La versión se reproduce desde fuentes controladas; casos positivos, rechazos, permisos, contratos y vistas quedan trazados. Funcional valida reglas; Datos, Seguridad y Calidad revisan sus ámbitos. Los ensayos integrales y la recepción contractual se verifican en 1.10 y T-17, sin duplicar las pruebas unitarias del componente.\newline{}\textbf{Insumos.} Para análisis: H1, requisitos, antecedentes y validadores autorizados. Antes de construir: decisiones de H2, modelos/contratos y DEV/QA recibidos. Antes de integrar: canales e interfaces comunes habilitados. El análisis no exige disponer anticipadamente de ambientes de construcción.\newline{}\textbf{Condiciones.} Etapa E2. Habilitaciones externas: 1.3.3.1, 1.7.1.3, 1.7.3.1, 1.7.3.3, BASE-E1, H2, H8. Actividades de análisis, diseño, construcción, verificación y documentación se desglosan una sola vez en T-15. Referencias: SD3 reparto y dominio, SD5 entidades/autoridad; FAC; DEC-17/18. Adaptaciones de ciclo de vida y campos sensibles se verifican con 1.7.6.5 y 1.3.4.4 cuando aplican. Los cinco actos sujetos a firma incorporan 1.7.6.6; no se duplica la capacidad compartida.\\\hline
\texttt{1.\allowbreak{}6.\allowbreak{}2.\allowbreak{}4.\allowbreak{}1}\newline{}\textbf{Capacidad funcional recibida: Expedientes de custodia y exportación íntegra}\newline{}Nikolai Navea\newline{}1 instancia.\newline{}\textbf{Ventana:} 13--15. & \textbf{Contenido y frontera.} Resultado único que reúne especificación, reglas, modelo propio, contratos, lógica, persistencia, vistas de tarea, pruebas unitarias y documentación de Expedientes de custodia y exportación íntegra. Reutiliza identidad, canales y adaptadores comunes; no reconstruye esas capacidades. Exclusiones: Inventar antecedentes faltantes, repetir digitalización/revisión de migración y resolver litigios.\newline{}\textbf{Aceptación.} Índice conserva folio, versión, documento, origen e integridad; exportación mantiene relaciones y señala faltantes para los 14 expedientes. La versión se reproduce desde fuentes controladas; casos positivos, rechazos, permisos, contratos y vistas quedan trazados. Funcional valida reglas; Datos, Seguridad y Calidad revisan sus ámbitos. Los ensayos integrales y la recepción contractual se verifican en 1.10 y T-17, sin duplicar las pruebas unitarias del componente.\newline{}\textbf{Insumos.} Para análisis: H1, requisitos, antecedentes y validadores autorizados. Antes de construir: decisiones de H2, modelos/contratos y DEV/QA recibidos. Antes de integrar: canales e interfaces comunes habilitados. El análisis no exige disponer anticipadamente de ambientes de construcción.\newline{}\textbf{Condiciones.} Etapa E2. Habilitaciones externas: 1.3.3.1, 1.7.1.3, 1.7.3.1, 1.7.3.3, BASE-E1, H2, H8. Actividades de análisis, diseño, construcción, verificación y documentación se desglosan una sola vez en T-15. Referencias: SD3 reparto y dominio, SD5 entidades/autoridad; FAC; C07 RT-05.15; 14 abandonadas. Adaptaciones de ciclo de vida y campos sensibles se verifican con 1.7.6.5 y 1.3.4.4 cuando aplican. Los cinco actos sujetos a firma incorporan 1.7.6.6; no se duplica la capacidad compartida.\\\hline
\end{longtable}\endgroup
\begingroup\setlength{\tabcolsep}{3pt}
\begin{longtable}{L{\dimexpr0.320000\textwidth-2\tabcolsep\relax}L{\dimexpr0.680000\textwidth-2\tabcolsep\relax}}
\caption{Diccionario: Canales, integración y datos compartidos}\label{tab:t14-dic-1-7}\\\hline
\EncabezadoTabla{Paquete y resultado} & \EncabezadoTabla{Frontera, aceptación e insumos}\\\hline\endfirsthead
\multicolumn{2}{l}{\textit{Continuación de la tabla \ref{tab:t14-dic-1-7}}}\\\hline
\EncabezadoTabla{Paquete y resultado} & \EncabezadoTabla{Frontera, aceptación e insumos}\\\hline\endhead
\hline\endfoot
\hline\endlastfoot
\texttt{1.\allowbreak{}7.\allowbreak{}1.\allowbreak{}1}\newline{}\textbf{Contratos de interfaces y eventos comunes}\newline{}Fernando Guerrero\newline{}1 instancia.\newline{}\textbf{Ventana:} 4. & \textbf{Contenido y frontera.} OpenAPI, sobre de eventos, autoridades, idempotencia, correlación y errores; excluye especificar reglas funcionales completas. Incluye registro de formatos y canales sectoriales: zarpe/avisos marítimos, residuos peligrosos, certificación ambiental de marinas y mensajería de baja capacidad. Identifica fuente, versión, justificación, disponibilidad efectiva y consumidor; no presupone una API oficial disponible.\newline{}\textbf{Aceptación.} Contratos tienen versiones y ejemplos de rechazo y reintento y no conceden acceso directo a tablas de otro dominio. Responsable: Fernando Guerrero (Arquitectura/Integración). Cada interfaz sectorial tiene contrato y medio verificados o una dependencia abierta con responsable y fecha; la recepción de diseño no acredita integración operativa. RT-05.23; FALTA REFERENCIA A SD5.\newline{}\textbf{Insumos.} Requisitos, decisiones de arquitectura y datos, contratos externos disponibles, responsables y herramientas de especificación. Los contratos versionados son resultado de esta ficha; DEV/QA se exigen al implementar/verificar sus consumidores, no para especificarlos en mes 4.\newline{}\textbf{Condiciones.} Predecesora de elaboración: H1, requisitos vigentes y decisiones de arquitectura/datos disponibles. Los contratos de diseño se integran con 1.7.8.1 y se aprueban conjuntamente en H2, sin requerir H2 recibido para producir el diseño que lo integra. H2 habilita la construcción definitiva de sus consumidores; no convierte una dependencia externa abierta en integración operativa. Referencia: SD4 DA-05/06/07; SD5 contratos. FALTA REFERENCIA A SD4; FALTA REFERENCIA A SD5.\\\hline
\texttt{1.\allowbreak{}7.\allowbreak{}1.\allowbreak{}2}\newline{}\textbf{Infraestructura de mensajería y DLQ configurada}\newline{}Fernando Guerrero\newline{}1 instancia.\newline{}\textbf{Ventana:} 6. & \textbf{Contenido y frontera.} EventBridge, SQS/SNS, rutas y dead-letter queues por IaC; excluye código productor/consumidor de cada módulo.\newline{}\textbf{Aceptación.} Permisos, entrega, reintento, DLQ y monitoreo se comprueban con mensajes de prueba sin duplicar hechos de negocio. Responsable: Fernando Guerrero (Arquitectura/Integración).\newline{}\textbf{Insumos.} DEV/QA, contratos versionados, sistemas de prueba, terminales y credenciales autorizadas.\newline{}\textbf{Condiciones.} Predecesoras del catálogo: 1.2.1.2,1.3.1.1,1.7.1.1. Referencia: SD4 DA-07. Las habilitaciones técnicas se contrastan con arquitectura e inventario vigentes.\\\hline
\texttt{1.\allowbreak{}7.\allowbreak{}1.\allowbreak{}3}\newline{}\textbf{Biblioteca de outbox, idempotencia y correlación}\newline{}Fernando Guerrero\newline{}1 instancia.\newline{}\textbf{Ventana:} 8. & \textbf{Contenido y frontera.} Componente compartido para emisión y recepción verificables; excluye implementar handlers de cada dominio.\newline{}\textbf{Aceptación.} Duplicados, reintentos y fallo entre escritura/emisión conservan un efecto y recuperación trazable. Responsable: Fernando Guerrero (Arquitectura/Integración).\newline{}\textbf{Insumos.} DEV/QA, contratos versionados, sistemas de prueba, terminales y credenciales autorizadas.\newline{}\textbf{Condiciones.} Predecesoras del catálogo: 1.7.1.1,1.7.1.2. Referencia: SD4 DA-05/06; SD5 autoridad. Las habilitaciones técnicas se contrastan con arquitectura e inventario vigentes.\\\hline
\texttt{1.\allowbreak{}7.\allowbreak{}2.\allowbreak{}1}\newline{}\textbf{Sincronizador de lotes y confirmaciones}\newline{}Fernando Guerrero\newline{}1 instancia.\newline{}\textbf{Ventana:} 8. & \textbf{Contenido y frontera.} Agente y coordinador HTTPS/mTLS con secuencia y ACK; excluye reglas de conciliación particulares y retorno de migración.\newline{}\textbf{Aceptación.} La cola sólo se elimina tras confirmación; interrupción y reintento preservan orden e integridad y muestran pendientes. Responsable: Fernando Guerrero (Arquitectura/Integración).\newline{}\textbf{Insumos.} DEV/QA, contratos versionados, sistemas de prueba, terminales y credenciales autorizadas.\newline{}\textbf{Condiciones.} Predecesoras del catálogo: 1.7.1.3,1.2.3.2,1.3.1.1. Referencia: SD4 coordinador de sincronización; SD5. Las habilitaciones técnicas se contrastan con arquitectura e inventario vigentes.\\\hline
\texttt{1.\allowbreak{}7.\allowbreak{}2.\allowbreak{}2}\newline{}\textbf{Conciliación de autoridad y conflictos implementada}\newline{}Nicolás Torfan\newline{}1 instancia.\newline{}\textbf{Ventana:} 8. & \textbf{Contenido y frontera.} Reglas comunes de autoridad por agregado, conflictos y revisión; excluye decidir hechos monetarios o cambios de negocio por el cliente.\newline{}\textbf{Aceptación.} Conflictos se conservan y escalan sin última escritura silenciosa; decisión y replay tienen evidencia. Nicolás Torfan, Datos, responde por el resultado y presenta la evidencia; la contraparte mantiene la recepción contractual.\newline{}\textbf{Insumos.} DEV/QA, contratos versionados, sistemas de prueba, terminales y credenciales autorizadas.\newline{}\textbf{Condiciones.} Predecesoras del catálogo: 1.7.2.1,H2. Referencia: SD4 DA-05; SD5 sincronización. Las habilitaciones técnicas se contrastan con arquitectura e inventario vigentes.\\\hline
\texttt{1.\allowbreak{}7.\allowbreak{}2.\allowbreak{}3}\newline{}\textbf{Ingesta y consolidación de telemetría implementadas}\newline{}Fernando Guerrero\newline{}1 instancia.\newline{}\textbf{Ventana:} 8. & \textbf{Contenido y frontera.} Consumo local MQTT, lotes cloud y series con calidad; excluye facturación, medición física e IN5.\newline{}\textbf{Aceptación.} Se validan secuencia, UTC, unidad, calidad y origen; no se suman totalizadores ni se habilita consumo directo del broker desde cloud. Responsable: Fernando Guerrero (Arquitectura/Integración).\newline{}\textbf{Insumos.} DEV/QA, contratos versionados, sistemas de prueba, terminales y credenciales autorizadas.\newline{}\textbf{Condiciones.} Predecesoras del catálogo: 1.2.3.5,1.7.1.3,H2. Referencia: SD4 DA-09; SD5 MediciónConsolidada. Las habilitaciones técnicas se contrastan con arquitectura e inventario vigentes.\\\hline
\texttt{1.\allowbreak{}7.\allowbreak{}2.\allowbreak{}5.\allowbreak{}\{1--16\}}\newline{}\textbf{F04 · Recepción por lote de 40 puntos de medición}\newline{}Fernando Guerrero\newline{}16 instancias.\newline{}\textbf{Ventana:} 5--6. & \textbf{Contenido y frontera.} Lote de 40 puntos de medición identificado por servicio/amarra y manifiesto. Mes 5: identificación e inspección preparatorias sobre activos disponibles; mes 6: configuración, comprobación y recepción después de habilitar concentradores y broker. Obras e instalación eléctrica corresponden al cliente. Los 16 lotes tienen conjuntos disjuntos; la inspección preparatoria no recibe el lote.\newline{}\textbf{Aceptación.} Por instancia: 40 puntos del manifiesto recibidos, con servicio/amarra, calibración, error, flexión, tiempo, corte y lectura trazados. Cierre agregado: 640 puntos únicos al recibir los 16 lotes; la instalación eléctrica corresponde al cliente.\newline{}\textbf{Insumos.} Universo nominal/manifiesto por instancia, versión y política/protocolo aplicable recibidos, responsables, acceso autorizado y ventana admisible. Registro de instancia: código, población/activos, fecha, versión, evidencia y decisión; el conjunto no se acepta por cerrar solo su primera instancia.\newline{}\textbf{Condiciones.} Para preparar: manifiesto, activos y acceso autorizados; no requiere declarar recibido el broker. Para configurar y recibir: 1.2.3.5 y EXT-TELEMETRIA conformes. La recepción de cada lote ocurre después de habilitar su concentrador/broker en mes 6 y antes de H3; T-15 ordena configuración, pruebas y revisión con recursos y duración efectivos. Si falla una condición, H3 queda bloqueado; no se anticipa aceptación ni se desplaza automáticamente el plazo contractual. FALTA REFERENCIA A SD4.\\\hline
\texttt{1.\allowbreak{}7.\allowbreak{}3.\allowbreak{}1}\newline{}\textbf{Base Android de operación en terreno}\newline{}Nikolai Navea\newline{}1 instancia.\newline{}\textbf{Ventana:} 8. & \textbf{Contenido y frontera.} Navegación común, sesión y almacenamiento cifrado durable; excluye pantallas y reglas de cada dominio.\newline{}\textbf{Aceptación.} La sesión limita rol y ámbito; reinicio conserva pendientes protegidos y el canal no accede directo a la base. Nikolai Navea, Desarrollo, responde por el resultado y presenta la evidencia; la contraparte mantiene la recepción contractual.\newline{}\textbf{Insumos.} DEV/QA, contratos versionados, sistemas de prueba, terminales y credenciales autorizadas.\newline{}\textbf{Condiciones.} Predecesoras del catálogo: 1.3.1.2,1.3.2.1,1.7.1.1. Referencia: SD4 DA-19; SD3 canal móvil. Las habilitaciones técnicas se contrastan con arquitectura e inventario vigentes.\\\hline
\texttt{1.\allowbreak{}7.\allowbreak{}3.\allowbreak{}2}\newline{}\textbf{Cola móvil, evidencia y sincronización de canal}\newline{}Nikolai Navea\newline{}1 instancia.\newline{}\textbf{Ventana:} 8. & \textbf{Contenido y frontera.} Cola local, fotos, revisión de pendientes y adaptación a sincronizador; excluye registrar inspecciones funcionales.\newline{}\textbf{Aceptación.} Desconexión, reinicio y ACK tardío no pierden fotos/hechos ni crean duplicados; permisos vencidos bloquean nuevas acciones. Nikolai Navea, Desarrollo, responde por el resultado y presenta la evidencia; la contraparte mantiene la recepción contractual.\newline{}\textbf{Insumos.} DEV/QA, contratos versionados, sistemas de prueba, terminales y credenciales autorizadas.\newline{}\textbf{Condiciones.} Predecesoras del catálogo: 1.7.3.1,1.7.2.1. Referencia: SD4 Android; SD5 evidencia offline. Las habilitaciones técnicas se contrastan con arquitectura e inventario vigentes.\\\hline
\texttt{1.\allowbreak{}7.\allowbreak{}3.\allowbreak{}3}\newline{}\textbf{Base de portal y consola accesibles}\newline{}Nikolai Navea\newline{}1 instancia.\newline{}\textbf{Ventana:} 8. & \textbf{Contenido y frontera.} Estructura React/PWA, sesión, navegación, errores y componentes comunes; excluye vistas funcionales de módulos e innovaciones. Incluye interfaz común de administración, bandeja unificada y búsqueda, con adaptaciones de parámetros, flujos y documentos de cada dominio propietario. La interfaz no sustituye sus reglas ni sus controles de datos; consume diseño validado de 1.7.7.1 y delta 1.7.7.2.\newline{}\textbf{Aceptación.} Se comprueban teclado, foco, mensajes, aislamiento de sesión e idiomas previstos; la base no declara WCAG integral de toda aplicación. Nikolai Navea, Desarrollo, responde por el resultado y presenta la evidencia; la contraparte mantiene la recepción contractual. RT-13/16: se verifica diseño responsivo y sistema de diseño recibido, teclado/foco, navegador soportado, mensajes de usuario y adaptación de terreno. Administración permite gestionar usuarios/roles/unidades y parámetros autorizados; cambios de impacto operacional requieren segundo perfil, justificación y auditoría. Catálogo distingue parámetros de cambios que exigen desarrollo; bandeja, búsqueda y autoatención respetan permisos. Flujos y documentos se comprueban con la versión funcional propietaria; la base de interfaz sola no acredita toda la aplicación.\newline{}\textbf{Insumos.} DEV/QA, contratos versionados, sistemas de prueba, terminales y credenciales autorizadas.\newline{}\textbf{Condiciones.} Predecesoras del catálogo: 1.3.1.2,1.7.1.1. Referencia: SD4 DA-19; BT accesibilidad. Las habilitaciones técnicas se contrastan con arquitectura e inventario vigentes.\\\hline
\texttt{1.\allowbreak{}7.\allowbreak{}4.\allowbreak{}1}\newline{}\textbf{Servicio común de notificaciones y preferencias}\newline{}Fernando Guerrero\newline{}1 instancia.\newline{}\textbf{Ventana:} 8--9. & \textbf{Contenido y frontera.} Canales, preferencias, evidencias de envío y reintento; excluye decidir destinatarios por las reglas de cada dominio. Mes 8: construcción y comprobaciones técnicas del motor común; mes 9: integración y recepción con consentimiento/finalidad recibidos. Las comprobaciones técnicas no acreditan por sí solas la recepción del servicio integrado. Incluye adaptaciones comunes de canales sectoriales seleccionados en 1.7.1.1; reglas y contenidos específicos permanecen en el dominio propietario. Interfaces externas no disponibles bloquean la recepción del intercambio afectado.\newline{}\textbf{Aceptación.} Un fallo conserva estado y responsable; mensajes repetidos no duplican efectos y cada envío respeta autorización y finalidad. Responsable: Fernando Guerrero (Arquitectura/Integración). Para RT-05.23 se verifica intercambio efectivo, formato, permisos, error y evidencia con el destinatario o canal disponible; un emulador solo demuestra una prueba técnica provisional. RT-16.20--25: canales ofertados, proveedor y contrato identificados; plantillas versionadas y preferencias respetan avisos obligatorios y baja cuando corresponde. Se comprueban envío asíncrono, reintentos, deduplicación y estados de entrega/apertura/error disponibles, sin atribuir apertura a canales que no la evidencian. Volumen y tratamiento del costo variable remiten a la Oferta Económica; la ficha no fija una tarifa nueva.\newline{}\textbf{Insumos.} DEV/QA, contratos versionados, sistemas de prueba, terminales y credenciales autorizadas.\newline{}\textbf{Condiciones.} Para construir/comprobar el motor: 1.7.1.2, contratos aplicables y DEV/QA recibidos, con datos y destinatarios de prueba autorizados. Para recibir el servicio integrado en mes 9: 1.3.3.2 recibido, preferencias/finalidad y contratos/canales efectivos comprobados; no se habilitan envíos reales que requieran consentimiento antes de esa condición. Se reconoce una sola construcción; T-15 distingue trabajo y evidencia de meses 8 y 9, sin duplicar esfuerzo ni valor ganado. FALTA REFERENCIA A SD3; FALTA REFERENCIA A SD4.\\\hline
\texttt{1.\allowbreak{}7.\allowbreak{}4.\allowbreak{}2}\newline{}\textbf{Adaptador con sistema laboral y fuente de identidad}\newline{}Fernando Guerrero\newline{}1 instancia.\newline{}\textbf{Ventana:} 8. & \textbf{Contenido y frontera.} Intercambio documentado de altas/bajas y habilitación; excluye administrar el sistema laboral externo.\newline{}\textbf{Aceptación.} Cambios y fallos quedan visibles, se conserva persona-identidad y se rechaza tratar una respuesta vencida como habilitación actual. Responsable: Fernando Guerrero (Arquitectura/Integración). El intercambio permite creación y baja automatizadas de RT-12.10, con origen/fecha de desvinculación y plazo máximo de 24 horas; un cambio recibido en el adaptador no acredita la revocación en los componentes consumidores.\newline{}\textbf{Insumos.} DEV/QA, contratos versionados, sistemas de prueba, terminales y credenciales autorizadas.\newline{}\textbf{Condiciones.} Predecesoras del catálogo: 1.3.2.2,1.7.1.1,EXT-IDENTIDAD. Referencia: SD4 autoridad local; SD3 frontera. Las habilitaciones técnicas se contrastan con arquitectura e inventario vigentes.\\\hline
\texttt{1.\allowbreak{}7.\allowbreak{}4.\allowbreak{}3}\newline{}\textbf{Adaptador contable y conciliación técnica}\newline{}Fernando Guerrero\newline{}1 instancia.\newline{}\textbf{Ventana:} 14. & \textbf{Contenido y frontera.} Intercambio de hechos/saldos con autoridad contable externa; excluye reglas de cobranza, emisión tributaria y conciliación de negocio.\newline{}\textbf{Aceptación.} Acuse, rechazo, reintentos e identificadores se conservan; no se emite un documento tributario desde un canal no autorizado. Responsable: Fernando Guerrero (Arquitectura/Integración).\newline{}\textbf{Insumos.} DEV/QA, contratos versionados, sistemas de prueba, terminales y credenciales autorizadas.\newline{}\textbf{Condiciones.} Predecesoras del catálogo: 1.7.1.1,BASE-ADM,EXT-CONTABLE. Referencia: SD3 Finanzas y cobranza; SD4 integración. Las habilitaciones técnicas se contrastan con arquitectura e inventario vigentes.\\\hline
\texttt{1.\allowbreak{}7.\allowbreak{}4.\allowbreak{}4}\newline{}\textbf{Adaptador de pagos y webhooks protegidos}\newline{}Fernando Guerrero\newline{}1 instancia.\newline{}\textbf{Ventana:} 14. & \textbf{Contenido y frontera.} Webhook, autenticidad, idempotencia y estados técnicos; excluye reglas monetarias y reversión de cobros por cliente.\newline{}\textbf{Aceptación.} Eventos repetidos o fuera de orden no crean otro pago; rechazo, conciliación pendiente y origen quedan trazados. Responsable: Fernando Guerrero (Arquitectura/Integración).\newline{}\textbf{Insumos.} DEV/QA, contratos versionados, sistemas de prueba, terminales y credenciales autorizadas.\newline{}\textbf{Condiciones.} Predecesoras del catálogo: 1.7.1.1,BASE-ADM,EXT-PAGOS. Referencia: SD3 pagos; SD4 integraciones. Las habilitaciones técnicas se contrastan con arquitectura e inventario vigentes.\\\hline
\texttt{1.\allowbreak{}7.\allowbreak{}6.\allowbreak{}1}\newline{}\textbf{ANAL-5: autoservicio analítico y modelo semántico recibidos}\newline{}Fernando Guerrero\newline{}1 instancia.\newline{}\textbf{Ventana:} 4--12,13--18. & \textbf{Contenido y frontera.} Capa común de medidas, tableros, filtros, profundización y diseño de informes del cliente. Lectura analítica separada de la transaccional; los dominios aportan únicamente medidas, reglas y proyecciones propias. No incluye envío automático ni duplicación de reportes ya definidos.\newline{}\textbf{Aceptación.} Usuario autorizado construye y exporta un informe nuevo sin intervención técnica; filtros y profundización llevan a evidencia de origen. Modelo semántico, diccionario y permisos versionados; se comprueba separación y protección de la operación bajo carga. RT-05.05 y RT-05.25--27; método y umbrales del SD5 recibido, con protocolo específico en T-13/T-17. Recepción E1: manifiesto de dominios, contratos, funciones/datos/informes incluidos, versión y pruebas cerradas para esa etapa; recepción E2: delta nominal de dominios y funciones, nueva versión, pruebas del delta y regresión de E1. Cada acta identifica código y etapa. El criterio de E1 no requiere dominios reservados a E2. T-15 distribuye una sola vez actividades/HH y valor ganado entre incrementos; la recepción E2 no reconoce nuevamente la construcción común recibida en E1.\newline{}\textbf{Insumos.} Modelos y reglas de SD5, servicios propietarios, identidades, almacén de lectura y contratos habilitados.\newline{}\textbf{Condiciones.} Base para los dominios E1 antes de H4/H5; extensión a los dominios E2 antes de H9/H10. Referencia: SD5, 5.2.3.1; FALTA REFERENCIA A SD4; FALTA REFERENCIA A SD9.\\\hline
\texttt{1.\allowbreak{}7.\allowbreak{}6.\allowbreak{}2}\newline{}\textbf{MASS-5: carga y descarga masiva ordinaria recibidas}\newline{}Fernando Guerrero\newline{}1 instancia.\newline{}\textbf{Ventana:} 4--12,13--18. & \textbf{Contenido y frontera.} Servicio permanente de la consola para iniciar, previsualizar, confirmar, cancelar y reanudar intercambio ordinario en CSV UTF-8 y JSON Lines. Invoca comandos del dominio propietario; no modifica directamente sus tablas ni cambia la autoridad contable. Se distingue de la migración inicial de 1.8.\newline{}\textbf{Aceptación.} Validación previa y diagnóstico por fila/grupo; procesamiento parcial expresamente autorizado; idempotencia y reanudación sin duplicados. Descarga íntegra con manifiesto y permisos. Se exige al menos 10.000 registros por minuto con validación activa y mezcla contractual de SD5; se prueban concurrencia, cancelación, reinicio y resultado desconocido. No se utiliza como recepción el perfil preliminar retirado por SD5. RT-05.22. Recepción E1: manifiesto de dominios, contratos, funciones/datos/informes incluidos, versión y pruebas cerradas para esa etapa; recepción E2: delta nominal de dominios y funciones, nueva versión, pruebas del delta y regresión de E1. Cada acta identifica código y etapa. El criterio de E1 no requiere dominios reservados a E2. T-15 distribuye una sola vez actividades/HH y valor ganado entre incrementos; la recepción E2 no reconoce nuevamente la construcción común recibida en E1.\newline{}\textbf{Insumos.} Esquemas, reglas, API de comando, auditoría, correlación y datos representativos autorizados; interfaz contable efectiva para los casos que la requieren.\newline{}\textbf{Condiciones.} Base E1 y ampliación E2 antes de las respectivas certificaciones; no acredita disponibilidad de terceros mediante un emulador. Referencia: SD5, 5.2.3.2; FALTA REFERENCIA A SD9.\\\hline
\texttt{1.\allowbreak{}7.\allowbreak{}6.\allowbreak{}3}\newline{}\textbf{ENV-5: registro, calendario y entrega automática de informes recibidos}\newline{}Fernando Guerrero\newline{}1 instancia.\newline{}\textbf{Ventana:} 4--12,13--18. & \textbf{Contenido y frontera.} REG-INF-5 y programación/envío de todos los informes de la solución, incluidos los diseños del cliente en ANAL-5. Reutiliza exportadores y permisos, sin reconstruir cálculos por dominio.\newline{}\textbf{Aceptación.} Se concilia el inventario completo de informes de cada etapa con REG-INF-5: cada informe dispone de exportador abierto, metadatos, permisos, filtros y programación habilitada. Todo diseño nuevo de ANAL-5 se registra y habilita automáticamente según sus permisos; se prueba esa regla de incorporación. Un informe nuevo y muestras justificadas por dominio comprueban calendario, cambios de hora, fin de mes, revocación, reinicio, reintento y carga; no sustituyen la comprobación de incorporación completa. Manifiesto, destinatario, versión, huellas y recibos trazables; entregado se distingue de abierto. RT-05.28. Recepciones E1 y E2 tienen inventario/versiones y evidencia propios; E2 comprueba delta y regresión sin volver a aceptar ni imputar el motor común.\newline{}\textbf{Insumos.} ANAL-5, REG-INF-5, exportadores, identidad, transportes y auditoría habilitados; dominios incorporados por etapa.\newline{}\textbf{Condiciones.} Para E1 depende de la versión E1 recibida de 1.7.6.1 y exportadores de sus dominios antes de H5; para E2 depende de su incremento E2 y dominios adicionales antes de H10. No depende del cierre completo E2 para recibir E1. FALTA REFERENCIA A SD5; FALTA REFERENCIA A SD9.\\\hline
\texttt{1.\allowbreak{}7.\allowbreak{}6.\allowbreak{}4}\newline{}\textbf{Actualidad de proyecciones y consultas por dominio verificada}\newline{}Fernando Guerrero\newline{}1 instancia.\newline{}\textbf{Ventana:} 4--12,13--18. & \textbf{Contenido y frontera.} Ruta común de actualización, correlación y medición extremo a extremo desde ocurrencia hasta vista. Cada dominio mantiene su autoridad y proyección; no es una segunda implementación de sus reglas ni otra ingesta de migración.\newline{}\textbf{Aceptación.} Ocupación no supera 60 s, consumo se actualiza al menos cada hora, estado de cuenta y escuela reflejan cambios en tiempo real conforme al caso, y ambiente consolida diariamente. Se registran captura, enlace, cola, aplicación, vista e incertidumbre de reloj; una copia atrasada se identifica. Interfaz contable y coordinación escolar efectivas probadas; no se declara cerrado el requisito mientras esas dependencias no estén disponibles. RT-05.29; criterio y presupuesto del SD5 recibido. Recepción E1: manifiesto de dominios, contratos, funciones/datos/informes incluidos, versión y pruebas cerradas para esa etapa; recepción E2: delta nominal de dominios y funciones, nueva versión, pruebas del delta y regresión de E1. Cada acta identifica código y etapa. El criterio de E1 no requiere dominios reservados a E2. T-15 distribuye una sola vez actividades/HH y valor ganado entre incrementos; la recepción E2 no reconoce nuevamente la construcción común recibida en E1.\newline{}\textbf{Insumos.} Eventos y fuentes de cada dominio, reloj controlado, proyecciones, canales y contrato contable efectivo; escenarios de conectividad autorizados.\newline{}\textbf{Condiciones.} Certificación por etapa según dominios habilitados; no se rebaja la obligación por pérdida de comunicación. Referencia: SD5, 5.2.3.4; FALTA REFERENCIA A SD4; FALTA REFERENCIA A SD9.\\\hline
\texttt{1.\allowbreak{}7.\allowbreak{}6.\allowbreak{}5}\newline{}\textbf{Ciclo de vida de datos por dominio recibido}\newline{}Nicolás Torfan\newline{}1 instancia.\newline{}\textbf{Ventana:} 4--8,13--15. & \textbf{Contenido y frontera.} Política versionada y ejecución de conservación, archivo y eliminación segura de datos operacionales, históricos y auditoría por dominio. Incluye excepciones autorizadas, relaciones, copias y evidencia; no repite la migración ni el respaldo técnico.\newline{}\textbf{Aceptación.} RT-05.07 y complemento C07 de RT-05.10: escuela de menores hasta los 18 años más cinco años, mínimo diez; salida/regreso y zarpe cinco años; contratos, custodia y estado de cuenta seis; residuos peligrosos y derrames seis; consumos cinco; fondeo vida útil más cinco; combustible cinco; videovigilancia seis meses. Política identifica evento inicial, categoría, autoridad, conservación y eliminación; pruebas de límites y excepciones impiden eliminación prematura y producen evidencia de eliminación autorizada. Auditoría conserva el período aplicable conforme a RT-16.10. E1 recibe sus categorías y E2 las restantes; no se declara recibida una categoría no implementada.\newline{}\textbf{Insumos.} Modelo y fuentes, decisiones de finalidad/conservación, validadores autorizados; servicios propietarios, claves, repositorios y casos de prueba. La adquisición y operación de cámaras conserva la frontera contractual del cliente y su interfaz se verifica.\newline{}\textbf{Condiciones.} Política de diseño antes de H2; ejecución por categoría antes de su certificación H5/H10. T-13/T-17 detallan pruebas, sin convertir un backup o una exportación final en eliminación ordinaria. FALTA REFERENCIA A SD5; FALTA REFERENCIA A SD9.\\\hline
\texttt{1.\allowbreak{}7.\allowbreak{}6.\allowbreak{}6}\newline{}\textbf{Firma, verificación y evidencia documental compartidas recibidas}\newline{}Fernando Guerrero\newline{}1 instancia.\newline{}\textbf{Ventana:} 4--8,13--15. & \textbf{Contenido y frontera.} Capacidad compartida para vincular documento y versión con identidad, voluntad, modalidad de firma admitida y evidencia verificable. Adaptaciones de contrato/renovación de amarra, autorización del apoderado, recepción/entrega de custodia, acta de izaje y entrega RESPEL. No decide la modalidad jurídica ni reconstruye las reglas de cada proceso.\newline{}\textbf{Aceptación.} Para cada tipo de acto se identifica modalidad admisible y fuente aplicable con validación competente; documento, firmante, versión, integridad y evidencia de verificación quedan vinculados. Pruebas positivas, rechazo por alteración o identidad no válida y conservación/exportación del expediente. No se recibe el acto por una casilla de interfaz ni por una firma de artefacto de software. E1 recibe autorizaciones escolares y actos incluidos efectivamente en esa etapa; E2 recibe las nuevas clases de acto y conserva regresión E1.\newline{}\textbf{Insumos.} Identidades, documentos y autoridades de cada dominio; decisión competente sobre modalidad; contratos y medios de firma efectivamente disponibles. Para diseño no se requieren DEV/QA; implementación y prueba sí los requieren.\newline{}\textbf{Condiciones.} Complemento obligatorio C07 RT-16.14. Diseño antes de H2 y adaptación antes de H5/H10 según etapa; el proveedor o mecanismo concreto se verifica antes de comprometer su recepción. FALTA REFERENCIA A SD4; FALTA REFERENCIA A SD5; FALTA REFERENCIA A SD9.\\\hline
\texttt{1.\allowbreak{}7.\allowbreak{}6.\allowbreak{}7}\newline{}\textbf{Decisiones de persistencia por dominio recibidas}\newline{}Fernando Guerrero\newline{}1 instancia.\newline{}\textbf{Ventana:} 4. & \textbf{Contenido y frontera.} Registro de decisiones del paradigma y motor por dominio, garantías transaccionales, autoridad de datos y elección de consistencia/disponibilidad frente a partición. No habilita recursos RDS/S3 ni sustituye modelo y diccionario de SD5.\newline{}\textbf{Aceptación.} Cada dominio identifica decisión, fundamento, restricciones, escenario de partición y efecto sobre operación/consulta. Datos y Arquitectura verifican correspondencia con modelos y contratos; recepción de H2 conserva decisiones y versiones. RT-05.02.\newline{}\textbf{Insumos.} Requisitos, modelo, escenarios operacionales, arquitectura y responsables de datos; herramientas de diseño. No exige ambientes de construcción.\newline{}\textbf{Condiciones.} Precede 1.2.1.3 y construcción de los componentes; decisión recibida en H2. Se referencia una única definición vigente del diseño. FALTA REFERENCIA A SD4; FALTA REFERENCIA A SD5.\\\hline
\texttt{1.\allowbreak{}7.\allowbreak{}7.\allowbreak{}1}\newline{}\textbf{Diseño UX E1 investigado y validado}\newline{}Nikolai Navea\newline{}1 instancia.\newline{}\textbf{Ventana:} 2--3. & \textbf{Contenido y frontera.} Investigación con personas usuarias reales de los perfiles y tareas del alcance E1; prototipo de investigación y diseño de interacción anterior a la construcción definitiva. Integra el sistema de diseño, mapa de navegación, flujos críticos y matriz de navegadores. Excluye construir las vistas productivas, impartir formación y sustituir la verificación WCAG final.\newline{}\textbf{Aceptación.} Universo de perfiles/tareas e indicadores de tiempo, pasos, errores y aprendizaje identificados y aprobados; investigación y pruebas con usuarios reales documentadas. Hallazgos se vinculan a cambios de diseño o decisiones justificadas; los bloqueantes se resuelven antes de construir el componente afectado. Se comprueban máximo tres interacciones para funciones principales, objetivos táctiles de al menos 44 por 44 píxeles, uso de terreno y alternativas accesibles. Diseño reutilizable documentado, con no más de cinco colores principales y dos familias tipográficas. No se atribuyen entrevistas, participantes o resultados aún no obtenidos. Desarrollo responde, Calidad revisa la evidencia y la contraparte decide la recepción.\newline{}\textbf{Insumos.} Alcance E1 y perfiles de SD3/T-12, tareas reales y condiciones de uso del caso; disponibilidad acordada con usuarios, herramientas de prototipado y responsables de diseño/validación. Si falta acceso a usuarios, se registra bloqueo; una revisión interna no reemplaza la investigación exigida.\newline{}\textbf{Condiciones.} Predecesoras: H1. La versión de diseño precede construcción definitiva de las capacidades E1 y se integra en H2; T-15 separa investigación/diseño de construcción. RT-13.03/04/09 y criterios de interfaz aplicables. El prototipo académico de Informe 3 es un antecedente distinto: solo se reutiliza si se identifican versión, finalidad y evidencia. FALTA REFERENCIA A SD3; FALTA REFERENCIA A SD4; FALTA REFERENCIA A SD9.\\\hline
\texttt{1.\allowbreak{}7.\allowbreak{}7.\allowbreak{}2}\newline{}\textbf{Incremento de diseño UX E2 validado}\newline{}Nikolai Navea\newline{}1 instancia.\newline{}\textbf{Ventana:} 13--14. & \textbf{Contenido y frontera.} Investigación y validación del cambio de perfiles, tareas y flujos que introduce E2. Reutiliza el diseño E1 recibido; documenta el delta y sus efectos, sin repetir ni imputar nuevamente la investigación conservada.\newline{}\textbf{Aceptación.} Manifiesto E2 identifica tareas/perfiles nuevos o modificados, indicadores aplicables, prototipo del delta, participantes efectivamente consultados, hallazgos y decisiones. Se validan con usuarios reales los cambios antes de su construcción definitiva y se comprueba regresión de los flujos E1 afectados. Bloqueantes cerrados antes de construir el componente; el diseño de E1 no acredita automáticamente el incremento. Desarrollo responde, Calidad revisa y la contraparte recibe.\newline{}\textbf{Insumos.} 1.7.7.1 recibido, alcance E2 de SD3/T-12, hallazgos de E1, usuarios y ventanas acordadas. Los indicadores conservados mantienen su definición; cambios se justifican y versionan.\newline{}\textbf{Condiciones.} Recepción de diseño detallado en H8 antes de construir definitivamente el delta; las actividades de mes 13 anteriores a esa recepción se limitan a análisis/prototipado. Precede vistas y adaptaciones E2 de 1.4--1.7; no exige una segunda construcción del sistema de diseño común. FALTA REFERENCIA A SD3; FALTA REFERENCIA A SD4; FALTA REFERENCIA A SD9.\\\hline
\texttt{1.\allowbreak{}7.\allowbreak{}8.\allowbreak{}1}\newline{}\textbf{Línea base de arquitectura integrada y decisiones vigentes}\newline{}Fernando Guerrero\newline{}1 instancia.\newline{}\textbf{Ventana:} 3--4,13--14. & \textbf{Contenido y frontera.} Integración y control de las vistas lógica, procesos, despliegue, datos y seguridad de SD4, con límites de contexto, interfaces y decisiones ADR. Referencia las especificaciones propietarias de plataforma, persistencia, seguridad y contratos; no vuelve a producirlas ni construye infraestructura. Mantiene versión inicial y delta E2 de una única arquitectura.\newline{}\textbf{Aceptación.} Cada vista identifica componentes, relaciones, responsables y correspondencia con el alcance de su etapa; ADR fechados conservan alternativas y criterio. Se declaran puntos únicos de falla remanentes, escalamiento y comportamiento degradado, con controles y pendientes explícitos. Se comprueba coherencia entre vistas y con SD3/SD5. La contraparte recibe la línea base inicial en H2 y el delta E2 en H8; el delta no vuelve a reconocer el diseño conservado.\newline{}\textbf{Insumos.} H1, requisitos vigentes, diseño de 1.7.1.1, decisiones de persistencia 1.7.6.7, modelo de amenazas 1.3.4.1 y especificaciones de plataforma/recinto de SD4. Los diseños se integran antes de recibir H2/H8; no se exige infraestructura operativa, obra concluida ni expediente de habilitación 1.2.3.1 recibido como insumo del diseño.\newline{}\textbf{Condiciones.} RT-02.01--04/11/13; vistas según ISO/IEC/IEEE 42010. Precede construcción definitiva e hitos de habilitación; cambios posteriores actualizan la arquitectura y sus ADR mediante el control de cambios existente en 1.1 y 1.12, sin crear una segunda cuenta de diseño. FALTA REFERENCIA A SD4; FALTA REFERENCIA A SD5; FALTA REFERENCIA A SD6.\\\hline
\end{longtable}\endgroup
\begingroup\setlength{\tabcolsep}{3pt}
\begin{longtable}{L{\dimexpr0.320000\textwidth-2\tabcolsep\relax}L{\dimexpr0.680000\textwidth-2\tabcolsep\relax}}
\caption{Diccionario: Migración e historia}\label{tab:t14-dic-1-8}\\\hline
\EncabezadoTabla{Paquete y resultado} & \EncabezadoTabla{Frontera, aceptación e insumos}\\\hline\endfirsthead
\multicolumn{2}{l}{\textit{Continuación de la tabla \ref{tab:t14-dic-1-8}}}\\\hline
\EncabezadoTabla{Paquete y resultado} & \EncabezadoTabla{Frontera, aceptación e insumos}\\\hline\endhead
\hline\endfoot
\hline\endlastfoot
\texttt{1.\allowbreak{}8.\allowbreak{}1.\allowbreak{}1}\newline{}\textbf{Catálogo de fuentes y accesos}\newline{}Nicolás Torfan\newline{}1 instancia.\newline{}\textbf{Ventana:} 1. & \textbf{Contenido y frontera.} Inventario de fuentes, ubicación, formatos, custodios, accesos y restricciones. Excluye el esfuerzo de los ensayos completos y del acompañamiento posterior.\newline{}\textbf{Aceptación.} Los seis conjuntos están identificados y se delimita la historia de doce años; faltantes y permisos quedan registrados sin suponer acceso concedido.\newline{}\textbf{Insumos.} Fuentes autorizadas, copia protegida, SD5, herramientas y validadores de negocio.\newline{}\textbf{Condiciones.} Diagnóstico precede reglas; reglas preceden saneamiento/carga; herramientas y datos conciliados preceden ensayos. El alcance operacional/histórico se fija por conjunto.\\\hline
\texttt{1.\allowbreak{}8.\allowbreak{}1.\allowbreak{}2}\newline{}\textbf{Extracción diagnóstica reproducible}\newline{}Nicolás Torfan\newline{}1 instancia.\newline{}\textbf{Ventana:} 2. & \textbf{Contenido y frontera.} Procedimiento de extracción diagnóstica de solo lectura, controles y registro de ejecución. Excluye el esfuerzo de los ensayos completos y del acompañamiento posterior.\newline{}\textbf{Aceptación.} La extracción se reproduce, conserva origen y hash/manifiesto, protege información sensible y no altera el legado.\newline{}\textbf{Insumos.} Fuentes autorizadas, copia protegida, SD5, herramientas y validadores de negocio.\newline{}\textbf{Condiciones.} Diagnóstico precede reglas; reglas preceden saneamiento/carga; herramientas y datos conciliados preceden ensayos. El alcance operacional/histórico se fija por conjunto.\\\hline
\texttt{1.\allowbreak{}8.\allowbreak{}1.\allowbreak{}3}\newline{}\textbf{Informe de calidad del origen}\newline{}Nicolás Torfan\newline{}1 instancia.\newline{}\textbf{Ventana:} 2. & \textbf{Contenido y frontera.} Perfilado de completitud, unicidad, validez, integridad y consistencia. Excluye el esfuerzo de los ensayos completos y del acompañamiento posterior.\newline{}\textbf{Aceptación.} Se muestran universo, reglas, resultados y excepciones reproducibles; no se atribuye saneamiento ejecutado al informe.\newline{}\textbf{Insumos.} Fuentes autorizadas, copia protegida, SD5, herramientas y validadores de negocio.\newline{}\textbf{Condiciones.} Diagnóstico precede reglas; reglas preceden saneamiento/carga; herramientas y datos conciliados preceden ensayos. El alcance operacional/histórico se fija por conjunto.\\\hline
\texttt{1.\allowbreak{}8.\allowbreak{}1.\allowbreak{}4}\newline{}\textbf{Matriz de cobertura y decisiones}\newline{}Nicolás Torfan\newline{}1 instancia.\newline{}\textbf{Ventana:} 3. & \textbf{Contenido y frontera.} Correspondencia de conjuntos con obligaciones y decisiones de tratamiento. Excluye el esfuerzo de los ensayos completos y del acompañamiento posterior.\newline{}\textbf{Aceptación.} Cada conjunto tiene destino operativo/histórico, reglas, responsable y decisión; los pendientes materiales impiden cerrar las reglas afectadas.\newline{}\textbf{Insumos.} Fuentes autorizadas, copia protegida, SD5, herramientas y validadores de negocio.\newline{}\textbf{Condiciones.} Diagnóstico precede reglas; reglas preceden saneamiento/carga; herramientas y datos conciliados preceden ensayos. El alcance operacional/histórico se fija por conjunto.\\\hline
\texttt{1.\allowbreak{}8.\allowbreak{}1.\allowbreak{}5}\newline{}\textbf{Extracción y primer cierre paralelo del legado reproducidos}\newline{}Nicolás Torfan\newline{}1 instancia.\newline{}\textbf{Ventana:} 6. & \textbf{Contenido y frontera.} Medida preventiva de Synaptix ante la dependencia de la persona experta: extracción reproducible, procedimiento disponible y primer simulacro de cierre mensual paralelo. Incluye reproducción por la persona sustituta; no sustituye ensayos integrales ni el corte posterior.\newline{}\textbf{Aceptación.} Persona sustituta reproduce extracción y cierre sobre copia autorizada; manifiesto, versión, saldos y diferencias quedan conciliados y documentados. Datos/Funcional revisan con la persona experta; SRE participa desde su incorporación. El acta identifica pendientes y riesgo residual sin declarar eliminada toda dependencia.\newline{}\textbf{Insumos.} Catálogo de fuentes, extracción diagnóstica, acceso/copia autorizados, persona experta y sustituta disponibles; reglas de saldos y seguridad. No exige que las herramientas finales de migración de meses 7--8 estén recibidas.\newline{}\textbf{Condiciones.} Cierre en mes 6; disponibilidad de contraparte y suplencia coordinadas en T-15/T-16. Referencia: SD3, 3.1.2.1; SD6, 6.1.1. No es un plazo impuesto por las bases.\\\hline
\texttt{1.\allowbreak{}8.\allowbreak{}2.\allowbreak{}1}\newline{}\textbf{Reglas: Socios y armadores}\newline{}Nicolás Torfan\newline{}1 instancia.\newline{}\textbf{Ventana:} 3. & \textbf{Contenido y frontera.} Correspondencia de origen/destino, claves, relaciones y reglas de Socios y armadores. Excluye el esfuerzo de los ensayos completos y del acompañamiento posterior.\newline{}\textbf{Aceptación.} Campos y vínculos obligatorios cubiertos; ejemplos válidos, ausentes e inválidos, autoridad y decisión funcional registrados antes de implementar la transformación.\newline{}\textbf{Insumos.} Fuentes autorizadas, copia protegida, SD5, herramientas y validadores de negocio.\newline{}\textbf{Condiciones.} Diagnóstico precede reglas; reglas preceden saneamiento/carga; herramientas y datos conciliados preceden ensayos. El alcance operacional/histórico se fija por conjunto.\\\hline
\texttt{1.\allowbreak{}8.\allowbreak{}2.\allowbreak{}2}\newline{}\textbf{Reglas: Contratos vigentes y documentos}\newline{}Nicolás Torfan\newline{}1 instancia.\newline{}\textbf{Ventana:} 3. & \textbf{Contenido y frontera.} Correspondencia de origen/destino, claves, relaciones y reglas de Contratos vigentes y documentos. Excluye el esfuerzo de los ensayos completos y del acompañamiento posterior.\newline{}\textbf{Aceptación.} Campos y vínculos obligatorios cubiertos; ejemplos válidos, ausentes e inválidos, autoridad y decisión funcional registrados antes de implementar la transformación.\newline{}\textbf{Insumos.} Fuentes autorizadas, copia protegida, SD5, herramientas y validadores de negocio.\newline{}\textbf{Condiciones.} Diagnóstico precede reglas; reglas preceden saneamiento/carga; herramientas y datos conciliados preceden ensayos. El alcance operacional/histórico se fija por conjunto.\\\hline
\texttt{1.\allowbreak{}8.\allowbreak{}2.\allowbreak{}3}\newline{}\textbf{Reglas: Seis años de pagos y saldos}\newline{}Nicolás Torfan\newline{}1 instancia.\newline{}\textbf{Ventana:} 3. & \textbf{Contenido y frontera.} Correspondencia de origen/destino, claves, relaciones y reglas de Seis años de pagos y saldos. Excluye el esfuerzo de los ensayos completos y del acompañamiento posterior.\newline{}\textbf{Aceptación.} Campos y vínculos obligatorios cubiertos; ejemplos válidos, ausentes e inválidos, autoridad y decisión funcional registrados antes de implementar la transformación.\newline{}\textbf{Insumos.} Fuentes autorizadas, copia protegida, SD5, herramientas y validadores de negocio.\newline{}\textbf{Condiciones.} Diagnóstico precede reglas; reglas preceden saneamiento/carga; herramientas y datos conciliados preceden ensayos. El alcance operacional/histórico se fija por conjunto.\\\hline
\texttt{1.\allowbreak{}8.\allowbreak{}2.\allowbreak{}4}\newline{}\textbf{Reglas: Embarcaciones e invernadas}\newline{}Nicolás Torfan\newline{}1 instancia.\newline{}\textbf{Ventana:} 3. & \textbf{Contenido y frontera.} Correspondencia de origen/destino, claves, relaciones y reglas de Embarcaciones e invernadas. Excluye el esfuerzo de los ensayos completos y del acompañamiento posterior.\newline{}\textbf{Aceptación.} Campos y vínculos obligatorios cubiertos; ejemplos válidos, ausentes e inválidos, autoridad y decisión funcional registrados antes de implementar la transformación.\newline{}\textbf{Insumos.} Fuentes autorizadas, copia protegida, SD5, herramientas y validadores de negocio.\newline{}\textbf{Condiciones.} Diagnóstico precede reglas; reglas preceden saneamiento/carga; herramientas y datos conciliados preceden ensayos. El alcance operacional/histórico se fija por conjunto.\\\hline
\texttt{1.\allowbreak{}8.\allowbreak{}2.\allowbreak{}5}\newline{}\textbf{Reglas: 14 expedientes de abandono}\newline{}Nicolás Torfan\newline{}1 instancia.\newline{}\textbf{Ventana:} 3. & \textbf{Contenido y frontera.} Correspondencia de origen/destino, claves, relaciones y reglas de 14 expedientes de abandono. Excluye el esfuerzo de los ensayos completos y del acompañamiento posterior.\newline{}\textbf{Aceptación.} Campos y vínculos obligatorios cubiertos; ejemplos válidos, ausentes e inválidos, autoridad y decisión funcional registrados antes de implementar la transformación.\newline{}\textbf{Insumos.} Fuentes autorizadas, copia protegida, SD5, herramientas y validadores de negocio.\newline{}\textbf{Condiciones.} Diagnóstico precede reglas; reglas preceden saneamiento/carga; herramientas y datos conciliados preceden ensayos. El alcance operacional/histórico se fija por conjunto.\\\hline
\texttt{1.\allowbreak{}8.\allowbreak{}2.\allowbreak{}6}\newline{}\textbf{Reglas: Tres años de matrícula escolar}\newline{}Nicolás Torfan\newline{}1 instancia.\newline{}\textbf{Ventana:} 3. & \textbf{Contenido y frontera.} Correspondencia de origen/destino, claves, relaciones y reglas de Tres años de matrícula escolar. Excluye el esfuerzo de los ensayos completos y del acompañamiento posterior.\newline{}\textbf{Aceptación.} Campos y vínculos obligatorios cubiertos; ejemplos válidos, ausentes e inválidos, autoridad y decisión funcional registrados antes de implementar la transformación.\newline{}\textbf{Insumos.} Fuentes autorizadas, copia protegida, SD5, herramientas y validadores de negocio.\newline{}\textbf{Condiciones.} Diagnóstico precede reglas; reglas preceden saneamiento/carga; herramientas y datos conciliados preceden ensayos. El alcance operacional/histórico se fija por conjunto.\\\hline
\texttt{1.\allowbreak{}8.\allowbreak{}3.\allowbreak{}1}\newline{}\textbf{Saneamiento automatizado: Socios y armadores}\newline{}Nicolás Torfan\newline{}1 instancia.\newline{}\textbf{Ventana:} 7. & \textbf{Contenido y frontera.} Procedimiento automatizado y excepciones de saneamiento de Socios y armadores. Excluye el esfuerzo de los ensayos completos y del acompañamiento posterior.\newline{}\textbf{Aceptación.} Reglas versionadas, copia de origen conservada, ejecución reproducible y reporte de cambios/errores; no se inventa información para completar faltantes.\newline{}\textbf{Insumos.} Fuentes autorizadas, copia protegida, SD5, herramientas y validadores de negocio.\newline{}\textbf{Condiciones.} Diagnóstico precede reglas; reglas preceden saneamiento/carga; herramientas y datos conciliados preceden ensayos. El alcance operacional/histórico se fija por conjunto.\\\hline
\texttt{1.\allowbreak{}8.\allowbreak{}3.\allowbreak{}2}\newline{}\textbf{Saneamiento automatizado: Contratos vigentes y documentos}\newline{}Nicolás Torfan\newline{}1 instancia.\newline{}\textbf{Ventana:} 7. & \textbf{Contenido y frontera.} Procedimiento automatizado y excepciones de saneamiento de Contratos vigentes y documentos. Excluye el esfuerzo de los ensayos completos y del acompañamiento posterior.\newline{}\textbf{Aceptación.} Reglas versionadas, copia de origen conservada, ejecución reproducible y reporte de cambios/errores; no se inventa información para completar faltantes.\newline{}\textbf{Insumos.} Fuentes autorizadas, copia protegida, SD5, herramientas y validadores de negocio.\newline{}\textbf{Condiciones.} Diagnóstico precede reglas; reglas preceden saneamiento/carga; herramientas y datos conciliados preceden ensayos. El alcance operacional/histórico se fija por conjunto.\\\hline
\texttt{1.\allowbreak{}8.\allowbreak{}3.\allowbreak{}3}\newline{}\textbf{Saneamiento automatizado: Seis años de pagos y saldos}\newline{}Nicolás Torfan\newline{}1 instancia.\newline{}\textbf{Ventana:} 7. & \textbf{Contenido y frontera.} Procedimiento automatizado y excepciones de saneamiento de Seis años de pagos y saldos. Excluye el esfuerzo de los ensayos completos y del acompañamiento posterior.\newline{}\textbf{Aceptación.} Reglas versionadas, copia de origen conservada, ejecución reproducible y reporte de cambios/errores; no se inventa información para completar faltantes.\newline{}\textbf{Insumos.} Fuentes autorizadas, copia protegida, SD5, herramientas y validadores de negocio.\newline{}\textbf{Condiciones.} Diagnóstico precede reglas; reglas preceden saneamiento/carga; herramientas y datos conciliados preceden ensayos. El alcance operacional/histórico se fija por conjunto.\\\hline
\texttt{1.\allowbreak{}8.\allowbreak{}3.\allowbreak{}4}\newline{}\textbf{Saneamiento automatizado: Embarcaciones e invernadas}\newline{}Nicolás Torfan\newline{}1 instancia.\newline{}\textbf{Ventana:} 7. & \textbf{Contenido y frontera.} Procedimiento automatizado y excepciones de saneamiento de Embarcaciones e invernadas. Excluye el esfuerzo de los ensayos completos y del acompañamiento posterior.\newline{}\textbf{Aceptación.} Reglas versionadas, copia de origen conservada, ejecución reproducible y reporte de cambios/errores; no se inventa información para completar faltantes.\newline{}\textbf{Insumos.} Fuentes autorizadas, copia protegida, SD5, herramientas y validadores de negocio.\newline{}\textbf{Condiciones.} Diagnóstico precede reglas; reglas preceden saneamiento/carga; herramientas y datos conciliados preceden ensayos. El alcance operacional/histórico se fija por conjunto.\\\hline
\texttt{1.\allowbreak{}8.\allowbreak{}3.\allowbreak{}5}\newline{}\textbf{Saneamiento automatizado: 14 expedientes de abandono}\newline{}Nicolás Torfan\newline{}1 instancia.\newline{}\textbf{Ventana:} 7. & \textbf{Contenido y frontera.} Procedimiento automatizado y excepciones de saneamiento de 14 expedientes de abandono. Excluye el esfuerzo de los ensayos completos y del acompañamiento posterior.\newline{}\textbf{Aceptación.} Reglas versionadas, copia de origen conservada, ejecución reproducible y reporte de cambios/errores; no se inventa información para completar faltantes.\newline{}\textbf{Insumos.} Fuentes autorizadas, copia protegida, SD5, herramientas y validadores de negocio.\newline{}\textbf{Condiciones.} Diagnóstico precede reglas; reglas preceden saneamiento/carga; herramientas y datos conciliados preceden ensayos. El alcance operacional/histórico se fija por conjunto.\\\hline
\texttt{1.\allowbreak{}8.\allowbreak{}3.\allowbreak{}6}\newline{}\textbf{Saneamiento automatizado: Tres años de matrícula escolar}\newline{}Nicolás Torfan\newline{}1 instancia.\newline{}\textbf{Ventana:} 7. & \textbf{Contenido y frontera.} Procedimiento automatizado y excepciones de saneamiento de Tres años de matrícula escolar. Excluye el esfuerzo de los ensayos completos y del acompañamiento posterior.\newline{}\textbf{Aceptación.} Reglas versionadas, copia de origen conservada, ejecución reproducible y reporte de cambios/errores; no se inventa información para completar faltantes.\newline{}\textbf{Insumos.} Fuentes autorizadas, copia protegida, SD5, herramientas y validadores de negocio.\newline{}\textbf{Condiciones.} Diagnóstico precede reglas; reglas preceden saneamiento/carga; herramientas y datos conciliados preceden ensayos. El alcance operacional/histórico se fija por conjunto.\\\hline
\texttt{1.\allowbreak{}8.\allowbreak{}4.\allowbreak{}1}\newline{}\textbf{Carga y conciliación: Socios y armadores}\newline{}Nicolás Torfan\newline{}1 instancia.\newline{}\textbf{Ventana:} 8. & \textbf{Contenido y frontera.} Carga y conciliación de Socios y armadores. Excluye el esfuerzo de los ensayos completos y del acompañamiento posterior.\newline{}\textbf{Aceptación.} Conteos, claves, relaciones, importes y documentos según el conjunto reconciliados; diferencias explicadas, reejecución idempotente y sin efectos duplicados.\newline{}\textbf{Insumos.} Fuentes autorizadas, copia protegida, SD5, herramientas y validadores de negocio.\newline{}\textbf{Condiciones.} Diagnóstico precede reglas; reglas preceden saneamiento/carga; herramientas y datos conciliados preceden ensayos. El alcance operacional/histórico se fija por conjunto.\\\hline
\texttt{1.\allowbreak{}8.\allowbreak{}4.\allowbreak{}2}\newline{}\textbf{Carga y conciliación: Contratos vigentes y documentos}\newline{}Nicolás Torfan\newline{}1 instancia.\newline{}\textbf{Ventana:} 8. & \textbf{Contenido y frontera.} Carga y conciliación de Contratos vigentes y documentos. Excluye el esfuerzo de los ensayos completos y del acompañamiento posterior.\newline{}\textbf{Aceptación.} Conteos, claves, relaciones, importes y documentos según el conjunto reconciliados; diferencias explicadas, reejecución idempotente y sin efectos duplicados.\newline{}\textbf{Insumos.} Fuentes autorizadas, copia protegida, SD5, herramientas y validadores de negocio.\newline{}\textbf{Condiciones.} Diagnóstico precede reglas; reglas preceden saneamiento/carga; herramientas y datos conciliados preceden ensayos. El alcance operacional/histórico se fija por conjunto.\\\hline
\texttt{1.\allowbreak{}8.\allowbreak{}4.\allowbreak{}3}\newline{}\textbf{Carga y conciliación: Seis años de pagos y saldos}\newline{}Nicolás Torfan\newline{}1 instancia.\newline{}\textbf{Ventana:} 8. & \textbf{Contenido y frontera.} Carga y conciliación de Seis años de pagos y saldos. Excluye el esfuerzo de los ensayos completos y del acompañamiento posterior.\newline{}\textbf{Aceptación.} Conteos, claves, relaciones, importes y documentos según el conjunto reconciliados; diferencias explicadas, reejecución idempotente y sin efectos duplicados.\newline{}\textbf{Insumos.} Fuentes autorizadas, copia protegida, SD5, herramientas y validadores de negocio.\newline{}\textbf{Condiciones.} Diagnóstico precede reglas; reglas preceden saneamiento/carga; herramientas y datos conciliados preceden ensayos. El alcance operacional/histórico se fija por conjunto.\\\hline
\texttt{1.\allowbreak{}8.\allowbreak{}4.\allowbreak{}4}\newline{}\textbf{Carga y conciliación: Embarcaciones e invernadas}\newline{}Nicolás Torfan\newline{}1 instancia.\newline{}\textbf{Ventana:} 8. & \textbf{Contenido y frontera.} Carga y conciliación de Embarcaciones e invernadas. Excluye el esfuerzo de los ensayos completos y del acompañamiento posterior.\newline{}\textbf{Aceptación.} Conteos, claves, relaciones, importes y documentos según el conjunto reconciliados; diferencias explicadas, reejecución idempotente y sin efectos duplicados.\newline{}\textbf{Insumos.} Fuentes autorizadas, copia protegida, SD5, herramientas y validadores de negocio.\newline{}\textbf{Condiciones.} Diagnóstico precede reglas; reglas preceden saneamiento/carga; herramientas y datos conciliados preceden ensayos. El alcance operacional/histórico se fija por conjunto.\\\hline
\texttt{1.\allowbreak{}8.\allowbreak{}4.\allowbreak{}5}\newline{}\textbf{Carga y conciliación: 14 expedientes de abandono}\newline{}Nicolás Torfan\newline{}1 instancia.\newline{}\textbf{Ventana:} 8. & \textbf{Contenido y frontera.} Carga y conciliación de 14 expedientes de abandono. Excluye el esfuerzo de los ensayos completos y del acompañamiento posterior.\newline{}\textbf{Aceptación.} Conteos, claves, relaciones, importes y documentos según el conjunto reconciliados; diferencias explicadas, reejecución idempotente y sin efectos duplicados.\newline{}\textbf{Insumos.} Fuentes autorizadas, copia protegida, SD5, herramientas y validadores de negocio.\newline{}\textbf{Condiciones.} Diagnóstico precede reglas; reglas preceden saneamiento/carga; herramientas y datos conciliados preceden ensayos. El alcance operacional/histórico se fija por conjunto.\\\hline
\texttt{1.\allowbreak{}8.\allowbreak{}4.\allowbreak{}6}\newline{}\textbf{Carga y conciliación: Tres años de matrícula escolar}\newline{}Nicolás Torfan\newline{}1 instancia.\newline{}\textbf{Ventana:} 8. & \textbf{Contenido y frontera.} Carga y conciliación de Tres años de matrícula escolar. Excluye el esfuerzo de los ensayos completos y del acompañamiento posterior.\newline{}\textbf{Aceptación.} Conteos, claves, relaciones, importes y documentos según el conjunto reconciliados; diferencias explicadas, reejecución idempotente y sin efectos duplicados.\newline{}\textbf{Insumos.} Fuentes autorizadas, copia protegida, SD5, herramientas y validadores de negocio.\newline{}\textbf{Condiciones.} Diagnóstico precede reglas; reglas preceden saneamiento/carga; herramientas y datos conciliados preceden ensayos. El alcance operacional/histórico se fija por conjunto.\\\hline
\texttt{1.\allowbreak{}8.\allowbreak{}5.\allowbreak{}1}\newline{}\textbf{Repositorio e índice histórico}\newline{}Nicolás Torfan\newline{}1 instancia.\newline{}\textbf{Ventana:} 8. & \textbf{Contenido y frontera.} Repositorio e índice histórico independiente del software de 2014. Excluye el esfuerzo de los ensayos completos y del acompañamiento posterior.\newline{}\textbf{Aceptación.} Objetos, metadatos, integridad, origen y retención identificados; navegación de muestra y reconstrucción comprobadas.\newline{}\textbf{Insumos.} Fuentes autorizadas, copia protegida, SD5, herramientas y validadores de negocio.\newline{}\textbf{Condiciones.} Diagnóstico precede reglas; reglas preceden saneamiento/carga; herramientas y datos conciliados preceden ensayos. El alcance operacional/histórico se fija por conjunto.\\\hline
\texttt{1.\allowbreak{}8.\allowbreak{}5.\allowbreak{}2}\newline{}\textbf{Consulta histórica con permisos}\newline{}Nicolás Torfan\newline{}1 instancia.\newline{}\textbf{Ventana:} 8. & \textbf{Contenido y frontera.} Consulta histórica con control de acceso. Excluye el esfuerzo de los ensayos completos y del acompañamiento posterior.\newline{}\textbf{Aceptación.} Consultas y denegaciones por perfil verificadas; enlaces y resultados trazables sin exigir ejecutar el sistema legado.\newline{}\textbf{Insumos.} Fuentes autorizadas, copia protegida, SD5, herramientas y validadores de negocio.\newline{}\textbf{Condiciones.} Diagnóstico precede reglas; reglas preceden saneamiento/carga; herramientas y datos conciliados preceden ensayos. El alcance operacional/histórico se fija por conjunto.\\\hline
\texttt{1.\allowbreak{}8.\allowbreak{}5.\allowbreak{}3}\newline{}\textbf{Exportación y comprobación de integridad}\newline{}Nicolás Torfan\newline{}1 instancia.\newline{}\textbf{Ventana:} 8. & \textbf{Contenido y frontera.} Exportación e integridad del archivo. Excluye el esfuerzo de los ensayos completos y del acompañamiento posterior.\newline{}\textbf{Aceptación.} Formato abierto, manifestación de objetos, hashes y muestras de restauración/exportación comprobados.\newline{}\textbf{Insumos.} Fuentes autorizadas, copia protegida, SD5, herramientas y validadores de negocio.\newline{}\textbf{Condiciones.} Diagnóstico precede reglas; reglas preceden saneamiento/carga; herramientas y datos conciliados preceden ensayos. El alcance operacional/histórico se fija por conjunto.\\\hline
\texttt{1.\allowbreak{}8.\allowbreak{}6.\allowbreak{}1}\newline{}\textbf{Detección y carga idempotente de delta}\newline{}Nicolás Torfan\newline{}1 instancia.\newline{}\textbf{Ventana:} 8. & \textbf{Contenido y frontera.} Detección y carga de delta entre origen y destino. Excluye el esfuerzo de los ensayos completos y del acompañamiento posterior.\newline{}\textbf{Aceptación.} Registros nuevos, modificados y eliminados tratados según reglas; reintentos no duplican hechos y diferencias se registran.\newline{}\textbf{Insumos.} Fuentes autorizadas, copia protegida, SD5, herramientas y validadores de negocio.\newline{}\textbf{Condiciones.} Diagnóstico precede reglas; reglas preceden saneamiento/carga; herramientas y datos conciliados preceden ensayos. El alcance operacional/histórico se fija por conjunto.\\\hline
\texttt{1.\allowbreak{}8.\allowbreak{}6.\allowbreak{}2}\newline{}\textbf{Herramientas de retorno previo}\newline{}Nicolás Torfan\newline{}1 instancia.\newline{}\textbf{Ventana:} 8. & \textbf{Contenido y frontera.} Herramientas de retorno antes de la activación. Excluye el esfuerzo de los ensayos completos y del acompañamiento posterior.\newline{}\textbf{Aceptación.} Copia y procedimientos restablecen el estado previo sin pérdida; el ensayo integral acredita los tiempos efectivos.\newline{}\textbf{Insumos.} Fuentes autorizadas, copia protegida, SD5, herramientas y validadores de negocio.\newline{}\textbf{Condiciones.} Diagnóstico precede reglas; reglas preceden saneamiento/carga; herramientas y datos conciliados preceden ensayos. El alcance operacional/histórico se fija por conjunto.\\\hline
\texttt{1.\allowbreak{}8.\allowbreak{}6.\allowbreak{}3}\newline{}\textbf{Herramientas de retorno posterior y diario}\newline{}Nicolás Torfan\newline{}1 instancia.\newline{}\textbf{Ventana:} 8. & \textbf{Contenido y frontera.} Retorno posterior y diario de cambios. Excluye el esfuerzo de los ensayos completos y del acompañamiento posterior.\newline{}\textbf{Aceptación.} Delta posterior conservado y reconciliado; restitución no pierde datos ni repite cobros, comunicaciones o decisiones externas.\newline{}\textbf{Insumos.} Fuentes autorizadas, copia protegida, SD5, herramientas y validadores de negocio.\newline{}\textbf{Condiciones.} Diagnóstico precede reglas; reglas preceden saneamiento/carga; herramientas y datos conciliados preceden ensayos. El alcance operacional/histórico se fija por conjunto.\\\hline
\texttt{1.\allowbreak{}8.\allowbreak{}7.\allowbreak{}1.\allowbreak{}\{1--14\}}\newline{}\textbf{F01 · Expedientes individuales de abandono}\newline{}Nicolás Torfan\newline{}14 instancias.\newline{}\textbf{Ventana:} 7. & \textbf{Contenido y frontera.} Expediente individual de abandono con folio, fuente, documentos, huellas, vínculos y decisiones pendientes. Se registra cada una de las 14 embarcaciones y se diferencia evidencia conservada de resolución jurídica.\newline{}\textbf{Aceptación.} Folio, origen, integridad y vínculos recibidos; faltantes y decisiones jurídicas identificados.\newline{}\textbf{Insumos.} Universo nominal/manifiesto por instancia, versión y política/protocolo aplicable recibidos, responsables, acceso autorizado y ventana admisible. Registro de instancia: código, población/activos, fecha, versión, evidencia y decisión; el conjunto no se acepta por cerrar solo su primera instancia.\newline{}\textbf{Condiciones.} Condiciones/predecesoras comunes del catálogo: 1.8.1.2,1.8.2.5. Referencia de detalle: FALTA REFERENCIA A SD5.\\\hline
\texttt{1.\allowbreak{}8.\allowbreak{}7.\allowbreak{}2.\allowbreak{}\{1--56\}}\newline{}\textbf{F02 · Lotes nominales de autorizaciones escolares}\newline{}Nicolás Torfan\newline{}56 instancias.\newline{}\textbf{Ventana:} 7. & \textbf{Contenido y frontera.} Lote nominal de hasta 32 autorizaciones escolares: lectura humana, menor/apoderado, vigencia, consentimiento y rechazo de faltantes. Cada autorización se vincula al alumno y al registro de origen; el manifiesto evita incluir la misma autorización en dos lotes.\newline{}\textbf{Aceptación.} Lectura humana, vigencia, menor/apoderado y evidencia; autorización ausente bloquea participación.\newline{}\textbf{Insumos.} Universo nominal/manifiesto por instancia, versión y política/protocolo aplicable recibidos, responsables, acceso autorizado y ventana admisible. Registro de instancia: código, población/activos, fecha, versión, evidencia y decisión; el conjunto no se acepta por cerrar solo su primera instancia.\newline{}\textbf{Condiciones.} Condiciones/predecesoras comunes del catálogo: 1.8.2.6. Referencia de detalle: FALTA REFERENCIA A SD5.\\\hline
\texttt{1.\allowbreak{}8.\allowbreak{}8.\allowbreak{}\{1--4\}}\newline{}\textbf{F05 · Adaptaciones identificadas del ensayo E2}\newline{}Nicolás Torfan\newline{}4 instancias.\newline{}\textbf{Ventana:} 14. & \textbf{Contenido y frontera.} Adaptación individual de esquema, reglas, universo y manifiesto al ensayo E2. Se identifica qué cambia respecto de E1; no ejecuta ni imputa nuevamente el ensayo completo de 1.10.\newline{}\textbf{Aceptación.} Manifiesto E2 identifica esquemas/contratos, transformaciones y universos incluidos, sin usar vistas como modelo.\newline{}\textbf{Insumos.} Universo nominal/manifiesto por instancia, versión y política/protocolo aplicable recibidos, responsables, acceso autorizado y ventana admisible. Registro de instancia: código, población/activos, fecha, versión, evidencia y decisión; el conjunto no se acepta por cerrar solo su primera instancia.\newline{}\textbf{Condiciones.} Condiciones/predecesoras comunes del catálogo: H8. Referencia de detalle: FALTA REFERENCIA A SD5.\\\hline
\end{longtable}\endgroup
\begingroup\setlength{\tabcolsep}{3pt}
\begin{longtable}{L{\dimexpr0.320000\textwidth-2\tabcolsep\relax}L{\dimexpr0.680000\textwidth-2\tabcolsep\relax}}
\caption{Diccionario: Innovaciones}\label{tab:t14-dic-1-9}\\\hline
\EncabezadoTabla{Paquete y resultado} & \EncabezadoTabla{Frontera, aceptación e insumos}\\\hline\endfirsthead
\multicolumn{2}{l}{\textit{Continuación de la tabla \ref{tab:t14-dic-1-9}}}\\\hline
\EncabezadoTabla{Paquete y resultado} & \EncabezadoTabla{Frontera, aceptación e insumos}\\\hline\endhead
\hline\endfoot
\hline\endlastfoot
\texttt{1.\allowbreak{}9.\allowbreak{}1.\allowbreak{}1.\allowbreak{}1}\newline{}\textbf{Flujo, estados y matriz de roles revisados}\newline{}Fernando Guerrero\newline{}1 instancia.\newline{}\textbf{Ventana:} 14--15. & \textbf{Contenido y frontera.} Se entregará flujo, estados y matriz de roles revisados, correspondiente a IN1.1.1. Excluye los resultados de las fichas predecesoras, la construcción de capacidades base y el servicio recurrente.\newline{}\textbf{Aceptación.} Cada transición tiene actor, precondición, rechazo y autoridad identificados; se valida el aislamiento entre armadores y contratistas. Fernando Guerrero, Arquitectura/Integración, responde por el resultado; los demás perfiles ejecutan sus aportes y el cliente valida las decisiones de su ámbito.\newline{}\textbf{Insumos.} repositorio de evidencia, registros autorizados, herramientas de análisis y validadores del proceso.\newline{}\textbf{Condiciones.} Predecesoras del catálogo: BASE-CTR. Referencia: FALTA REFERENCIA A SD13. Las habilitaciones técnicas se contrastan con arquitectura e inventario vigentes.\\\hline
\texttt{1.\allowbreak{}9.\allowbreak{}1.\allowbreak{}1.\allowbreak{}2}\newline{}\textbf{Especificación de las interfaces}\newline{}Fernando Guerrero\newline{}1 instancia.\newline{}\textbf{Ventana:} 14--15. & \textbf{Contenido y frontera.} Se entregará especificación de las interfaces, correspondiente a IN1.1.2. Excluye los resultados de las fichas predecesoras, la construcción de capacidades base y el servicio recurrente.\newline{}\textbf{Aceptación.} Contratos OpenAPI y eventos versionados identifican errores, idempotencia, correlación y pruebas de compatibilidad. Fernando Guerrero, Arquitectura/Integración, responde por el resultado; los demás perfiles ejecutan sus aportes y el cliente valida las decisiones de su ámbito.\newline{}\textbf{Insumos.} DEV/QA, contratos de interfaz, repositorio y herramientas de automatización.\newline{}\textbf{Condiciones.} Predecesoras del catálogo: 1.9.1.1.1. Referencia: FALTA REFERENCIA A SD13. Las habilitaciones técnicas se contrastan con arquitectura e inventario vigentes.\\\hline
\texttt{1.\allowbreak{}9.\allowbreak{}1.\allowbreak{}1.\allowbreak{}3}\newline{}\textbf{Protocolo de medición y aceptación}\newline{}Davor Serey\newline{}1 instancia.\newline{}\textbf{Ventana:} 14--15. & \textbf{Contenido y frontera.} Se entregará protocolo de medición y aceptación, correspondiente a IN1.1.3. Excluye los resultados de las fichas predecesoras, la construcción de capacidades base y el servicio recurrente.\newline{}\textbf{Aceptación.} Quedan fijados cohorte, registro de primera respuesta, tratamiento de casos abiertos y criterio de aceptación antes de medir. Davor Serey, Calidad, responde por el resultado; los demás perfiles ejecutan sus aportes y el cliente valida las decisiones de su ámbito.\newline{}\textbf{Insumos.} repositorio de evidencia, registros autorizados, herramientas de análisis y validadores del proceso.\newline{}\textbf{Condiciones.} Predecesoras del catálogo: 1.9.1.1.1. Referencia: FALTA REFERENCIA A SD13. Las habilitaciones técnicas se contrastan con arquitectura e inventario vigentes.\\\hline
\texttt{1.\allowbreak{}9.\allowbreak{}1.\allowbreak{}2.\allowbreak{}1}\newline{}\textbf{Servicio de solicitudes y transiciones}\newline{}Nikolai Navea\newline{}1 instancia.\newline{}\textbf{Ventana:} 16--17. & \textbf{Contenido y frontera.} Se entregará servicio de solicitudes y transiciones, correspondiente a IN1.2.1. Comprende implementación, pruebas unitarias y documentación del resultado. Excluye los resultados de las fichas predecesoras, la construcción de capacidades base y el servicio recurrente.\newline{}\textbf{Aceptación.} Casos válidos e inválidos y escrituras repetidas conservan una sola transición autorizada y su evidencia. Nikolai Navea, Desarrollo, responde por el resultado; los demás perfiles ejecutan sus aportes y el cliente valida las decisiones de su ámbito.\newline{}\textbf{Insumos.} DEV/QA, contratos de interfaz, repositorio y herramientas de automatización.\newline{}\textbf{Condiciones.} Predecesoras del catálogo: 1.9.1.1.2. Referencia: FALTA REFERENCIA A SD13. Las habilitaciones técnicas se contrastan con arquitectura e inventario vigentes.\\\hline
\texttt{1.\allowbreak{}9.\allowbreak{}1.\allowbreak{}2.\allowbreak{}2}\newline{}\textbf{Formularios y vistas en AUT}\newline{}Nikolai Navea\newline{}1 instancia.\newline{}\textbf{Ventana:} 16--17. & \textbf{Contenido y frontera.} Se entregará formularios y vistas en aut, correspondiente a IN1.2.2. Comprende implementación, pruebas unitarias y documentación del resultado. Excluye los resultados de las fichas predecesoras, la construcción de capacidades base y el servicio recurrente.\newline{}\textbf{Aceptación.} Armador y contratista sólo ven y modifican su ámbito; las vistas muestran estados y errores y admiten navegación accesible. Nikolai Navea, Desarrollo, responde por el resultado; los demás perfiles ejecutan sus aportes y el cliente valida las decisiones de su ámbito.\newline{}\textbf{Insumos.} DEV/QA, contratos de interfaz, repositorio y herramientas de automatización.\newline{}\textbf{Condiciones.} Predecesoras del catálogo: 1.9.1.1.2. Referencia: FALTA REFERENCIA A SD13. Las habilitaciones técnicas se contrastan con arquitectura e inventario vigentes.\\\hline
\texttt{1.\allowbreak{}9.\allowbreak{}1.\allowbreak{}2.\allowbreak{}3}\newline{}\textbf{Historial y registro asistido auditable}\newline{}Nikolai Navea\newline{}1 instancia.\newline{}\textbf{Ventana:} 16--17. & \textbf{Contenido y frontera.} Se entregará historial y registro asistido auditable, correspondiente a IN1.2.3. Comprende implementación, pruebas unitarias y documentación del resultado. Excluye los resultados de las fichas predecesoras, la construcción de capacidades base y el servicio recurrente.\newline{}\textbf{Aceptación.} Cada cambio registra actor, fecha, estado anterior y nuevo; el ingreso asistido identifica quién lo realizó. Nikolai Navea, Desarrollo, responde por el resultado; los demás perfiles ejecutan sus aportes y el cliente valida las decisiones de su ámbito.\newline{}\textbf{Insumos.} DEV/QA, contratos de interfaz, repositorio y herramientas de automatización.\newline{}\textbf{Condiciones.} Predecesoras del catálogo: 1.9.1.2.1. Referencia: FALTA REFERENCIA A SD13. Las habilitaciones técnicas se contrastan con arquitectura e inventario vigentes.\\\hline
\texttt{1.\allowbreak{}9.\allowbreak{}1.\allowbreak{}3.\allowbreak{}1}\newline{}\textbf{Conexiones con CTR y VAR verificadas}\newline{}Fernando Guerrero\newline{}1 instancia.\newline{}\textbf{Ventana:} 16--17. & \textbf{Contenido y frontera.} Se entregará conexiones con ctr y var verificadas, correspondiente a IN1.3.1. Comprende implementación, pruebas unitarias y documentación del resultado. Excluye los resultados de las fichas predecesoras, la construcción de capacidades base y el servicio recurrente.\newline{}\textbf{Aceptación.} Solicitudes y faenas se vinculan por contrato sin escribir tablas de otro dominio y toleran reintentos. Fernando Guerrero, Arquitectura/Integración, responde por el resultado; los demás perfiles ejecutan sus aportes y el cliente valida las decisiones de su ámbito.\newline{}\textbf{Insumos.} DEV/QA, contratos de interfaz, repositorio y herramientas de automatización.\newline{}\textbf{Condiciones.} Predecesoras del catálogo: 1.9.1.2.1. Referencia: FALTA REFERENCIA A SD13. Las habilitaciones técnicas se contrastan con arquitectura e inventario vigentes.\\\hline
\texttt{1.\allowbreak{}9.\allowbreak{}1.\allowbreak{}3.\allowbreak{}2}\newline{}\textbf{Identidad y vínculo de armador integrados}\newline{}Bruno Toro\newline{}1 instancia.\newline{}\textbf{Ventana:} 16--17. & \textbf{Contenido y frontera.} Se entregará identidad y vínculo de armador integrados, correspondiente a IN1.3.2. Comprende implementación, pruebas unitarias y documentación del resultado. Excluye los resultados de las fichas predecesoras, la construcción de capacidades base y el servicio recurrente.\newline{}\textbf{Aceptación.} Los casos de suplantación y acceso cruzado se rechazan y los vínculos autorizados se comprueban. Bruno Toro, Seguridad, responde por el resultado; los demás perfiles ejecutan sus aportes y el cliente valida las decisiones de su ámbito.\newline{}\textbf{Insumos.} DEV/QA, contratos de interfaz, repositorio y herramientas de automatización.\newline{}\textbf{Condiciones.} Predecesoras del catálogo: 1.9.1.2.2. Referencia: FALTA REFERENCIA A SD13. Las habilitaciones técnicas se contrastan con arquitectura e inventario vigentes.\\\hline
\texttt{1.\allowbreak{}9.\allowbreak{}1.\allowbreak{}3.\allowbreak{}3}\newline{}\textbf{Notificaciones y reintentos integrados}\newline{}Nikolai Navea\newline{}1 instancia.\newline{}\textbf{Ventana:} 16--17. & \textbf{Contenido y frontera.} Se entregará notificaciones y reintentos integrados, correspondiente a IN1.3.3. Comprende implementación, pruebas unitarias y documentación del resultado. Excluye los resultados de las fichas predecesoras, la construcción de capacidades base y el servicio recurrente.\newline{}\textbf{Aceptación.} Avisos fallidos quedan visibles y los reintentos no duplican efectos ni conceden una habilitación. Nikolai Navea, Desarrollo, responde por el resultado; los demás perfiles ejecutan sus aportes y el cliente valida las decisiones de su ámbito.\newline{}\textbf{Insumos.} DEV/QA, contratos de interfaz, repositorio y herramientas de automatización.\newline{}\textbf{Condiciones.} Predecesoras del catálogo: 1.9.1.3.1. Referencia: FALTA REFERENCIA A SD13. Las habilitaciones técnicas se contrastan con arquitectura e inventario vigentes.\\\hline
\texttt{1.\allowbreak{}9.\allowbreak{}1.\allowbreak{}4.\allowbreak{}1}\newline{}\textbf{Informe de pruebas funcionales}\newline{}Davor Serey\newline{}1 instancia.\newline{}\textbf{Ventana:} 18. & \textbf{Contenido y frontera.} Se entregará informe de pruebas funcionales, correspondiente a IN1.4.1. Comprende preparación, ejecución de comprobaciones y expediente. Excluye los resultados de las fichas predecesoras, la construcción de capacidades base y el servicio recurrente.\newline{}\textbf{Aceptación.} Se comprueban estados, duplicados, concurrencia, habilitación y flujos completos; cada defecto queda trazado. Davor Serey, Calidad, responde por el resultado; los demás perfiles ejecutan sus aportes y el cliente valida las decisiones de su ámbito.\newline{}\textbf{Insumos.} PREPROD, versión identificada, datos autorizados y herramientas de pruebas.\newline{}\textbf{Condiciones.} Predecesoras del catálogo: 1.9.1.2.3,1.9.1.3.1,1.9.1.3.2,1.9.1.3.3. Referencia: FALTA REFERENCIA A SD13. Las habilitaciones técnicas se contrastan con arquitectura e inventario vigentes.\\\hline
\texttt{1.\allowbreak{}9.\allowbreak{}1.\allowbreak{}4.\allowbreak{}2}\newline{}\textbf{Informe de seguridad y recuperación}\newline{}Bruno Toro\newline{}1 instancia.\newline{}\textbf{Ventana:} 18. & \textbf{Contenido y frontera.} Se entregará informe de seguridad y recuperación, correspondiente a IN1.4.2. Comprende preparación, ejecución de comprobaciones y expediente. Excluye los resultados de las fichas predecesoras, la construcción de capacidades base y el servicio recurrente.\newline{}\textbf{Aceptación.} No quedan vulnerabilidades críticas o altas abiertas; aislamiento y recuperación se prueban sobre versión identificada. Bruno Toro, Seguridad, responde por el resultado; los demás perfiles ejecutan sus aportes y el cliente valida las decisiones de su ámbito.\newline{}\textbf{Insumos.} PREPROD, versión identificada, datos autorizados y herramientas de pruebas.\newline{}\textbf{Condiciones.} Predecesoras del catálogo: 1.9.1.4.1. Referencia: FALTA REFERENCIA A SD13. Las habilitaciones técnicas se contrastan con arquitectura e inventario vigentes.\\\hline
\texttt{1.\allowbreak{}9.\allowbreak{}1.\allowbreak{}5.\allowbreak{}1}\newline{}\textbf{Guía operativa y capacitación registrada}\newline{}José Lara Arce\newline{}1 instancia.\newline{}\textbf{Ventana:} 19--20. & \textbf{Contenido y frontera.} Se entregará guía operativa y capacitación registrada, correspondiente a IN1.5.1. Comprende registro, validación y emisión del resultado. Excluye los resultados de las fichas predecesoras, la construcción de capacidades base y el servicio recurrente.\newline{}\textbf{Aceptación.} Los participantes reproducen los flujos; quedan material, asistencia, evaluación y soporte identificados. José Lara Arce, Implantación, responde por el resultado; los demás perfiles ejecutan sus aportes y el cliente valida las decisiones de su ámbito.\newline{}\textbf{Insumos.} repositorio de evidencia, registros autorizados, herramientas de análisis y validadores del proceso.\newline{}\textbf{Condiciones.} Inicio posterior al cierre de 1.9.1.5.3 y anterior a 1.9.1.5.4. El intervalo efectivo se calcula en T-15/T-18 con disponibilidad de la cohorte; no interviene durante la medición de línea base. Referencia: SD3, 3.1.2.2; FALTA REFERENCIA A SD13.\\\hline
\texttt{1.\allowbreak{}9.\allowbreak{}1.\allowbreak{}5.\allowbreak{}2}\newline{}\textbf{Registro base de la primera quincena}\newline{}Davor Serey\newline{}1 instancia.\newline{}\textbf{Ventana:} 19. & \textbf{Contenido y frontera.} Se entregará registro base de la primera quincena, correspondiente a IN1.5.2. Comprende registro, validación y emisión del resultado. Excluye los resultados de las fichas predecesoras, la construcción de capacidades base y el servicio recurrente.\newline{}\textbf{Aceptación.} Las solicitudes de la base conservan apertura, respuesta y casos sin cerrar; el flujo adicional sigue desactivado para la cohorte. Davor Serey, Calidad, responde por el resultado; los demás perfiles ejecutan sus aportes y el cliente valida las decisiones de su ámbito.\newline{}\textbf{Insumos.} repositorio de evidencia, registros autorizados, herramientas de análisis y validadores del proceso.\newline{}\textbf{Condiciones.} Predecesoras del catálogo: 1.9.1.1.3,1.9.1.4.2. Referencia: FALTA REFERENCIA A SD13. Las habilitaciones técnicas se contrastan con arquitectura e inventario vigentes.\\\hline
\texttt{1.\allowbreak{}9.\allowbreak{}1.\allowbreak{}5.\allowbreak{}3}\newline{}\textbf{Registro base de la segunda quincena}\newline{}Davor Serey\newline{}1 instancia.\newline{}\textbf{Ventana:} 19. & \textbf{Contenido y frontera.} Se entregará registro base de la segunda quincena, correspondiente a IN1.5.3. Comprende registro, validación y emisión del resultado. Excluye los resultados de las fichas predecesoras, la construcción de capacidades base y el servicio recurrente.\newline{}\textbf{Aceptación.} El segundo registro usa el mismo método y cohorte, identifica exclusiones y cierra el período base. Davor Serey, Calidad, responde por el resultado; los demás perfiles ejecutan sus aportes y el cliente valida las decisiones de su ámbito.\newline{}\textbf{Insumos.} repositorio de evidencia, registros autorizados, herramientas de análisis y validadores del proceso.\newline{}\textbf{Condiciones.} Predecesoras del catálogo: 1.9.1.5.2. Referencia: FALTA REFERENCIA A SD13. Las habilitaciones técnicas se contrastan con arquitectura e inventario vigentes.\\\hline
\texttt{1.\allowbreak{}9.\allowbreak{}1.\allowbreak{}5.\allowbreak{}4}\newline{}\textbf{Registro de piloto de la primera quincena}\newline{}José Lara Arce\newline{}1 instancia.\newline{}\textbf{Ventana:} 20. & \textbf{Contenido y frontera.} Se entregará registro de piloto de la primera quincena, correspondiente a IN1.5.4. Comprende registro, validación y emisión del resultado. Excluye los resultados de las fichas predecesoras, la construcción de capacidades base y el servicio recurrente.\newline{}\textbf{Aceptación.} Se conserva el universo de solicitudes y la primera respuesta de cada una, incluidos rechazos y casos abiertos. José Lara Arce, Implantación, responde por el resultado; los demás perfiles ejecutan sus aportes y el cliente valida las decisiones de su ámbito.\newline{}\textbf{Insumos.} repositorio de evidencia, registros autorizados, herramientas de análisis y validadores del proceso.\newline{}\textbf{Condiciones.} Predecesoras del catálogo: 1.9.1.5.1. Referencia: FALTA REFERENCIA A SD13. Las habilitaciones técnicas se contrastan con arquitectura e inventario vigentes. Piloto posterior a línea base completa y capacitación 1.9.1.5.1 recibida.\\\hline
\texttt{1.\allowbreak{}9.\allowbreak{}1.\allowbreak{}5.\allowbreak{}5}\newline{}\textbf{Registro de piloto de la segunda quincena}\newline{}José Lara Arce\newline{}1 instancia.\newline{}\textbf{Ventana:} 20. & \textbf{Contenido y frontera.} Se entregará registro de piloto de la segunda quincena, correspondiente a IN1.5.5. Comprende registro, validación y emisión del resultado. Excluye los resultados de las fichas predecesoras, la construcción de capacidades base y el servicio recurrente.\newline{}\textbf{Aceptación.} El registro conserva método comparable, incidencias, cobertura y pérdidas de observación explicadas. José Lara Arce, Implantación, responde por el resultado; los demás perfiles ejecutan sus aportes y el cliente valida las decisiones de su ámbito.\newline{}\textbf{Insumos.} repositorio de evidencia, registros autorizados, herramientas de análisis y validadores del proceso.\newline{}\textbf{Condiciones.} Predecesoras del catálogo: 1.9.1.5.4. Referencia: FALTA REFERENCIA A SD13. Las habilitaciones técnicas se contrastan con arquitectura e inventario vigentes. Piloto posterior a línea base completa y capacitación 1.9.1.5.1 recibida.\\\hline
\texttt{1.\allowbreak{}9.\allowbreak{}1.\allowbreak{}5.\allowbreak{}6}\newline{}\textbf{Informe comparativo y acta de evaluación}\newline{}Davor Serey\newline{}1 instancia.\newline{}\textbf{Ventana:} 20. & \textbf{Contenido y frontera.} Se entregará informe comparativo y acta de evaluación, correspondiente a IN1.5.6. Comprende registro, validación y emisión del resultado. Excluye los resultados de las fichas predecesoras, la construcción de capacidades base y el servicio recurrente.\newline{}\textbf{Aceptación.} La mediana y la trazabilidad se comparan con la base sin inventar beneficio; la decisión y observaciones quedan registradas. Davor Serey, Calidad, responde por el resultado; los demás perfiles ejecutan sus aportes y el cliente valida las decisiones de su ámbito.\newline{}\textbf{Insumos.} repositorio de evidencia, registros autorizados, herramientas de análisis y validadores del proceso.\newline{}\textbf{Condiciones.} Predecesoras del catálogo: 1.9.1.5.5. Referencia: FALTA REFERENCIA A SD13. Las habilitaciones técnicas se contrastan con arquitectura e inventario vigentes.\\\hline
\texttt{1.\allowbreak{}9.\allowbreak{}1.\allowbreak{}6.\allowbreak{}1}\newline{}\textbf{Servicio desplegado y transferido a operación}\newline{}José Lara Arce\newline{}1 instancia.\newline{}\textbf{Ventana:} 21. & \textbf{Contenido y frontera.} Se entregará servicio desplegado y transferido a operación, correspondiente a IN1.6.1. Comprende registro, validación y emisión del resultado. Excluye los resultados de las fichas predecesoras, la construcción de capacidades base y el servicio recurrente.\newline{}\textbf{Aceptación.} Versión, permisos, procedimiento, soporte y decisión de habilitación quedan recibidos; no se presenta un piloto como aceptación. José Lara Arce, Operación/SRE, responde por el resultado; los demás perfiles ejecutan sus aportes y el cliente valida las decisiones de su ámbito.\newline{}\textbf{Insumos.} repositorio de evidencia, registros autorizados, herramientas de análisis y validadores del proceso.\newline{}\textbf{Condiciones.} Predecesoras del catálogo: 1.9.1.5.6. Referencia: FALTA REFERENCIA A SD13. Las habilitaciones técnicas se contrastan con arquitectura e inventario vigentes.\\\hline
\texttt{1.\allowbreak{}9.\allowbreak{}2.\allowbreak{}1.\allowbreak{}1}\newline{}\textbf{Catálogo de mantenimiento y reglas de planificación definidos}\newline{}Ulysses Barreto\newline{}1 instancia.\newline{}\textbf{Ventana:} 14--15. & \textbf{Contenido y frontera.} Se entregará catálogo de mantenimiento y reglas de planificación definidos, correspondiente a IN2.1.1. Excluye los resultados de las fichas predecesoras, la construcción de capacidades base y el servicio recurrente.\newline{}\textbf{Aceptación.} Cada tarea tiene condición, frecuencia, antecedente, responsable y límite de recomendación; no se diagnostica una embarcación sin evidencia. Ulysses Barreto, Funcional, responde por el resultado; los demás perfiles ejecutan sus aportes y el cliente valida las decisiones de su ámbito.\newline{}\textbf{Insumos.} repositorio de evidencia, registros autorizados, herramientas de análisis y validadores del proceso.\newline{}\textbf{Condiciones.} Predecesoras del catálogo: BASE-VAR. Referencia: FALTA REFERENCIA A SD13. Las habilitaciones técnicas se contrastan con arquitectura e inventario vigentes.\\\hline
\texttt{1.\allowbreak{}9.\allowbreak{}2.\allowbreak{}1.\allowbreak{}2}\newline{}\textbf{Ficha, responsabilidades e interfaces especificadas}\newline{}Fernando Guerrero\newline{}1 instancia.\newline{}\textbf{Ventana:} 14--15. & \textbf{Contenido y frontera.} Se entregará ficha, responsabilidades e interfaces especificadas, correspondiente a IN2.1.2. Excluye los resultados de las fichas predecesoras, la construcción de capacidades base y el servicio recurrente.\newline{}\textbf{Aceptación.} Se distinguen recomendación, aceptación del armador y ejecución por contratista; cada interfaz y permiso tiene autoridad. Fernando Guerrero, Arquitectura/Integración, responde por el resultado; los demás perfiles ejecutan sus aportes y el cliente valida las decisiones de su ámbito.\newline{}\textbf{Insumos.} repositorio de evidencia, registros autorizados, herramientas de análisis y validadores del proceso.\newline{}\textbf{Condiciones.} Predecesoras del catálogo: 1.9.2.1.1,1.9.1.1.2. Referencia: FALTA REFERENCIA A SD13. Las habilitaciones técnicas se contrastan con arquitectura e inventario vigentes.\\\hline
\texttt{1.\allowbreak{}9.\allowbreak{}2.\allowbreak{}2.\allowbreak{}1}\newline{}\textbf{Ficha de ingreso de antecedentes disponible en AUT}\newline{}Nikolai Navea\newline{}1 instancia.\newline{}\textbf{Ventana:} 16--17. & \textbf{Contenido y frontera.} Se entregará ficha de ingreso de antecedentes disponible en aut, correspondiente a IN2.2.1. Comprende implementación, pruebas unitarias y documentación del resultado. Excluye los resultados de las fichas predecesoras, la construcción de capacidades base y el servicio recurrente.\newline{}\textbf{Aceptación.} Campos ausentes se identifican y no se convierten en datos técnicos confirmados; acceso restringido al armador autorizado. Nikolai Navea, Desarrollo, responde por el resultado; los demás perfiles ejecutan sus aportes y el cliente valida las decisiones de su ámbito.\newline{}\textbf{Insumos.} DEV/QA, contratos de interfaz, repositorio y herramientas de automatización.\newline{}\textbf{Condiciones.} Predecesoras del catálogo: 1.9.2.1.2. Referencia: FALTA REFERENCIA A SD13. Las habilitaciones técnicas se contrastan con arquitectura e inventario vigentes.\\\hline
\texttt{1.\allowbreak{}9.\allowbreak{}2.\allowbreak{}2.\allowbreak{}2}\newline{}\textbf{Catálogo configurable y aceptación de propuestas implementados}\newline{}Nikolai Navea\newline{}1 instancia.\newline{}\textbf{Ventana:} 16--17. & \textbf{Contenido y frontera.} Se entregará catálogo configurable y aceptación de propuestas implementados, correspondiente a IN2.2.2. Comprende implementación, pruebas unitarias y documentación del resultado. Excluye los resultados de las fichas predecesoras, la construcción de capacidades base y el servicio recurrente.\newline{}\textbf{Aceptación.} Cambios de catálogo se versionan y ninguna propuesta genera una orden sin aceptación explícita. Nikolai Navea, Desarrollo, responde por el resultado; los demás perfiles ejecutan sus aportes y el cliente valida las decisiones de su ámbito.\newline{}\textbf{Insumos.} DEV/QA, contratos de interfaz, repositorio y herramientas de automatización.\newline{}\textbf{Condiciones.} Predecesoras del catálogo: 1.9.2.2.1. Referencia: FALTA REFERENCIA A SD13. Las habilitaciones técnicas se contrastan con arquitectura e inventario vigentes.\\\hline
\texttt{1.\allowbreak{}9.\allowbreak{}2.\allowbreak{}3.\allowbreak{}1}\newline{}\textbf{Calendario integrado con posiciones, botaduras y habilitación}\newline{}Fernando Guerrero\newline{}1 instancia.\newline{}\textbf{Ventana:} 17--18. & \textbf{Contenido y frontera.} Se entregará calendario integrado con posiciones, botaduras y habilitación, correspondiente a IN2.3.1. Comprende implementación, pruebas unitarias y documentación del resultado. Excluye los resultados de las fichas predecesoras, la construcción de capacidades base y el servicio recurrente.\newline{}\textbf{Aceptación.} Se detectan incompatibilidades de posición, fecha y habilitación y las decisiones conservan su origen. Fernando Guerrero, Arquitectura/Integración, responde por el resultado; los demás perfiles ejecutan sus aportes y el cliente valida las decisiones de su ámbito.\newline{}\textbf{Insumos.} DEV/QA, contratos de interfaz, repositorio y herramientas de automatización.\newline{}\textbf{Condiciones.} Predecesoras del catálogo: 1.9.2.2.2. Referencia: FALTA REFERENCIA A SD13. Las habilitaciones técnicas se contrastan con arquitectura e inventario vigentes.\\\hline
\texttt{1.\allowbreak{}9.\allowbreak{}2.\allowbreak{}3.\allowbreak{}2}\newline{}\textbf{Trabajos aceptados vinculados a solicitudes IN1}\newline{}Nikolai Navea\newline{}1 instancia.\newline{}\textbf{Ventana:} 17--18. & \textbf{Contenido y frontera.} Se entregará trabajos aceptados vinculados a solicitudes in1, correspondiente a IN2.3.2. Comprende implementación, pruebas unitarias y documentación del resultado. Excluye los resultados de las fichas predecesoras, la construcción de capacidades base y el servicio recurrente.\newline{}\textbf{Aceptación.} Cada trabajo aceptado crea o vincula una solicitud una sola vez y conserva actor, consentimiento y correlación. Nikolai Navea, Desarrollo, responde por el resultado; los demás perfiles ejecutan sus aportes y el cliente valida las decisiones de su ámbito.\newline{}\textbf{Insumos.} DEV/QA, contratos de interfaz, repositorio y herramientas de automatización.\newline{}\textbf{Condiciones.} Predecesoras del catálogo: 1.9.2.3.1,1.9.1.3.1. Referencia: FALTA REFERENCIA A SD13. Las habilitaciones técnicas se contrastan con arquitectura e inventario vigentes.\\\hline
\texttt{1.\allowbreak{}9.\allowbreak{}2.\allowbreak{}3.\allowbreak{}3}\newline{}\textbf{Permisos y notificaciones configurados}\newline{}Bruno Toro\newline{}1 instancia.\newline{}\textbf{Ventana:} 17--18. & \textbf{Contenido y frontera.} Se entregará permisos y notificaciones configurados, correspondiente a IN2.3.3. Comprende implementación, pruebas unitarias y documentación del resultado. Excluye los resultados de las fichas predecesoras, la construcción de capacidades base y el servicio recurrente.\newline{}\textbf{Aceptación.} El contratista ve sólo su encargo y los avisos no equivalen a autorización ni duplican solicitudes. Bruno Toro, Seguridad, responde por el resultado; los demás perfiles ejecutan sus aportes y el cliente valida las decisiones de su ámbito.\newline{}\textbf{Insumos.} DEV/QA, contratos de interfaz, repositorio y herramientas de automatización.\newline{}\textbf{Condiciones.} Predecesoras del catálogo: 1.9.2.3.2. Referencia: FALTA REFERENCIA A SD13. Las habilitaciones técnicas se contrastan con arquitectura e inventario vigentes.\\\hline
\texttt{1.\allowbreak{}9.\allowbreak{}2.\allowbreak{}4.\allowbreak{}1}\newline{}\textbf{Evidencias de pruebas funcionales y de integración}\newline{}Davor Serey\newline{}1 instancia.\newline{}\textbf{Ventana:} 19--20. & \textbf{Contenido y frontera.} Se entregará evidencias de pruebas funcionales y de integración, correspondiente a IN2.4.1. Comprende preparación, ejecución de comprobaciones y expediente. Excluye los resultados de las fichas predecesoras, la construcción de capacidades base y el servicio recurrente.\newline{}\textbf{Aceptación.} Se prueban propuestas, rechazo, aceptación, cambio de fecha, contratista no habilitado y reintentos. Davor Serey, Calidad, responde por el resultado; los demás perfiles ejecutan sus aportes y el cliente valida las decisiones de su ámbito.\newline{}\textbf{Insumos.} PREPROD, versión identificada, datos autorizados y herramientas de pruebas.\newline{}\textbf{Condiciones.} Predecesoras del catálogo: 1.9.2.3.3. Referencia: FALTA REFERENCIA A SD13. Las habilitaciones técnicas se contrastan con arquitectura e inventario vigentes.\\\hline
\texttt{1.\allowbreak{}9.\allowbreak{}2.\allowbreak{}4.\allowbreak{}2}\newline{}\textbf{Seguridad, recuperación y aceptación técnica}\newline{}Bruno Toro\newline{}1 instancia.\newline{}\textbf{Ventana:} 19--20. & \textbf{Contenido y frontera.} Se entregará seguridad, recuperación y aceptación técnica, correspondiente a IN2.4.2. Comprende preparación, ejecución de comprobaciones y expediente. Excluye los resultados de las fichas predecesoras, la construcción de capacidades base y el servicio recurrente.\newline{}\textbf{Aceptación.} Se comprueban aislamiento, continuidad y retorno; no quedan defectos críticos o altos abiertos y se registra decisión técnica. Bruno Toro, Seguridad, responde por el resultado; los demás perfiles ejecutan sus aportes y el cliente valida las decisiones de su ámbito.\newline{}\textbf{Insumos.} PREPROD, versión identificada, datos autorizados y herramientas de pruebas.\newline{}\textbf{Condiciones.} Predecesoras del catálogo: 1.9.2.4.1. Referencia: FALTA REFERENCIA A SD13. Las habilitaciones técnicas se contrastan con arquitectura e inventario vigentes.\\\hline
\texttt{1.\allowbreak{}9.\allowbreak{}2.\allowbreak{}5.\allowbreak{}1}\newline{}\textbf{Participantes inscritos, fichas y capacitación registrada}\newline{}José Lara Arce\newline{}1 instancia.\newline{}\textbf{Ventana:} 30. & \textbf{Contenido y frontera.} Se entregará participantes inscritos, fichas y capacitación registrada, correspondiente a IN2.5.1. Comprende registro, validación y emisión del resultado. Excluye los resultados de las fichas predecesoras, la construcción de capacidades base y el servicio recurrente.\newline{}\textbf{Aceptación.} Participación autorizada, fichas completas o ausencias identificadas y evaluación de capacitación por participante. José Lara Arce, Implantación, responde por el resultado; los demás perfiles ejecutan sus aportes y el cliente valida las decisiones de su ámbito.\newline{}\textbf{Insumos.} repositorio de evidencia, registros autorizados, herramientas de análisis y validadores del proceso.\newline{}\textbf{Condiciones.} Predecesoras del catálogo: 1.9.2.4.2,1.9.1.6.1. Referencia: FALTA REFERENCIA A SD13. Las habilitaciones técnicas se contrastan con arquitectura e inventario vigentes.\\\hline
\texttt{1.\allowbreak{}9.\allowbreak{}2.\allowbreak{}5.\allowbreak{}2}\newline{}\textbf{Planes aceptados y calendario confirmado}\newline{}Ulysses Barreto\newline{}1 instancia.\newline{}\textbf{Ventana:} 30. & \textbf{Contenido y frontera.} Se entregará planes aceptados y calendario confirmado, correspondiente a IN2.5.2. Comprende registro, validación y emisión del resultado. Excluye los resultados de las fichas predecesoras, la construcción de capacidades base y el servicio recurrente.\newline{}\textbf{Aceptación.} Cada trabajo tiene aceptación, contratista habilitado, fecha y vínculo de solicitud; el cliente confirma ventanas. Ulysses Barreto, Funcional, responde por el resultado; los demás perfiles ejecutan sus aportes y el cliente valida las decisiones de su ámbito.\newline{}\textbf{Insumos.} repositorio de evidencia, registros autorizados, herramientas de análisis y validadores del proceso.\newline{}\textbf{Condiciones.} Predecesoras del catálogo: 1.9.2.5.1. Referencia: FALTA REFERENCIA A SD13. Las habilitaciones técnicas se contrastan con arquitectura e inventario vigentes.\\\hline
\texttt{1.\allowbreak{}9.\allowbreak{}2.\allowbreak{}5.\allowbreak{}3.\allowbreak{}0.\allowbreak{}1}\newline{}\textbf{Informe quincenal del piloto invernal}\newline{}José Lara Arce\newline{}1 instancia.\newline{}\textbf{Ventana:} 31. & \textbf{Contenido y frontera.} Se entregará informe quincenal del piloto invernal, correspondiente a IN2.5.3.0.1. Comprende registro, validación y emisión del resultado. Excluye los resultados de las fichas predecesoras, la construcción de capacidades base y el servicio recurrente.\newline{}\textbf{Aceptación.} El informe conserva planes, solicitudes, trabajos aceptados y pendientes, incidencias y cambios de calendario sin atribuir trabajos no realizados. José Lara Arce, Implantación, responde por el resultado; los demás perfiles ejecutan sus aportes y el cliente valida las decisiones de su ámbito.\newline{}\textbf{Insumos.} repositorio de evidencia, registros autorizados, herramientas de análisis y validadores del proceso.\newline{}\textbf{Condiciones.} Predecesoras del catálogo: 1.9.2.5.2. Referencia: FALTA REFERENCIA A SD13. Las habilitaciones técnicas se contrastan con arquitectura e inventario vigentes.\\\hline
\texttt{1.\allowbreak{}9.\allowbreak{}2.\allowbreak{}5.\allowbreak{}3.\allowbreak{}0.\allowbreak{}2}\newline{}\textbf{Informe quincenal del piloto invernal}\newline{}José Lara Arce\newline{}1 instancia.\newline{}\textbf{Ventana:} 31. & \textbf{Contenido y frontera.} Se entregará informe quincenal del piloto invernal, correspondiente a IN2.5.3.0.2. Comprende registro, validación y emisión del resultado. Excluye los resultados de las fichas predecesoras, la construcción de capacidades base y el servicio recurrente.\newline{}\textbf{Aceptación.} El informe conserva planes, solicitudes, trabajos aceptados y pendientes, incidencias y cambios de calendario sin atribuir trabajos no realizados. José Lara Arce, Implantación, responde por el resultado; los demás perfiles ejecutan sus aportes y el cliente valida las decisiones de su ámbito.\newline{}\textbf{Insumos.} repositorio de evidencia, registros autorizados, herramientas de análisis y validadores del proceso.\newline{}\textbf{Condiciones.} Predecesoras del catálogo: 1.9.2.5.2. Referencia: FALTA REFERENCIA A SD13. Las habilitaciones técnicas se contrastan con arquitectura e inventario vigentes.\\\hline
\texttt{1.\allowbreak{}9.\allowbreak{}2.\allowbreak{}5.\allowbreak{}3.\allowbreak{}1.\allowbreak{}1}\newline{}\textbf{Informe quincenal del piloto invernal}\newline{}José Lara Arce\newline{}1 instancia.\newline{}\textbf{Ventana:} 32. & \textbf{Contenido y frontera.} Se entregará informe quincenal del piloto invernal, correspondiente a IN2.5.3.1.1. Excluye los resultados de las fichas predecesoras, la construcción de capacidades base y el servicio recurrente.\newline{}\textbf{Aceptación.} El informe conserva planes, solicitudes, trabajos aceptados y pendientes, incidencias y cambios de calendario sin atribuir trabajos no realizados. José Lara Arce, Implantación, responde por el resultado; los demás perfiles ejecutan sus aportes y el cliente valida las decisiones de su ámbito.\newline{}\textbf{Insumos.} repositorio de evidencia, registros autorizados, herramientas de análisis y validadores del proceso.\newline{}\textbf{Condiciones.} Predecesoras del catálogo: 1.9.2.5.2. Referencia: FALTA REFERENCIA A SD13. Las habilitaciones técnicas se contrastan con arquitectura e inventario vigentes.\\\hline
\texttt{1.\allowbreak{}9.\allowbreak{}2.\allowbreak{}5.\allowbreak{}3.\allowbreak{}1.\allowbreak{}2}\newline{}\textbf{Informe quincenal del piloto invernal}\newline{}José Lara Arce\newline{}1 instancia.\newline{}\textbf{Ventana:} 32. & \textbf{Contenido y frontera.} Se entregará informe quincenal del piloto invernal, correspondiente a IN2.5.3.1.2. Excluye los resultados de las fichas predecesoras, la construcción de capacidades base y el servicio recurrente.\newline{}\textbf{Aceptación.} El informe conserva planes, solicitudes, trabajos aceptados y pendientes, incidencias y cambios de calendario sin atribuir trabajos no realizados. José Lara Arce, Implantación, responde por el resultado; los demás perfiles ejecutan sus aportes y el cliente valida las decisiones de su ámbito.\newline{}\textbf{Insumos.} repositorio de evidencia, registros autorizados, herramientas de análisis y validadores del proceso.\newline{}\textbf{Condiciones.} Predecesoras del catálogo: 1.9.2.5.2. Referencia: FALTA REFERENCIA A SD13. Las habilitaciones técnicas se contrastan con arquitectura e inventario vigentes.\\\hline
\texttt{1.\allowbreak{}9.\allowbreak{}2.\allowbreak{}5.\allowbreak{}3.\allowbreak{}2.\allowbreak{}1}\newline{}\textbf{Informe quincenal del piloto invernal}\newline{}José Lara Arce\newline{}1 instancia.\newline{}\textbf{Ventana:} 33. & \textbf{Contenido y frontera.} Se entregará informe quincenal del piloto invernal, correspondiente a IN2.5.3.2.1. Comprende registro, validación y emisión del resultado. Excluye los resultados de las fichas predecesoras, la construcción de capacidades base y el servicio recurrente.\newline{}\textbf{Aceptación.} El informe conserva planes, solicitudes, trabajos aceptados y pendientes, incidencias y cambios de calendario sin atribuir trabajos no realizados. José Lara Arce, Implantación, responde por el resultado; los demás perfiles ejecutan sus aportes y el cliente valida las decisiones de su ámbito.\newline{}\textbf{Insumos.} repositorio de evidencia, registros autorizados, herramientas de análisis y validadores del proceso.\newline{}\textbf{Condiciones.} Predecesoras del catálogo: 1.9.2.5.2. Referencia: FALTA REFERENCIA A SD13. Las habilitaciones técnicas se contrastan con arquitectura e inventario vigentes.\\\hline
\texttt{1.\allowbreak{}9.\allowbreak{}2.\allowbreak{}5.\allowbreak{}3.\allowbreak{}2.\allowbreak{}2}\newline{}\textbf{Informe quincenal del piloto invernal}\newline{}José Lara Arce\newline{}1 instancia.\newline{}\textbf{Ventana:} 33. & \textbf{Contenido y frontera.} Se entregará informe quincenal del piloto invernal, correspondiente a IN2.5.3.2.2. Comprende registro, validación y emisión del resultado. Excluye los resultados de las fichas predecesoras, la construcción de capacidades base y el servicio recurrente.\newline{}\textbf{Aceptación.} El informe conserva planes, solicitudes, trabajos aceptados y pendientes, incidencias y cambios de calendario sin atribuir trabajos no realizados. José Lara Arce, Implantación, responde por el resultado; los demás perfiles ejecutan sus aportes y el cliente valida las decisiones de su ámbito.\newline{}\textbf{Insumos.} repositorio de evidencia, registros autorizados, herramientas de análisis y validadores del proceso.\newline{}\textbf{Condiciones.} Predecesoras del catálogo: 1.9.2.5.2. Referencia: FALTA REFERENCIA A SD13. Las habilitaciones técnicas se contrastan con arquitectura e inventario vigentes.\\\hline
\texttt{1.\allowbreak{}9.\allowbreak{}2.\allowbreak{}5.\allowbreak{}4.\allowbreak{}\{1--5\}}\newline{}\textbf{F26 · Revisiones previas a botadura por lote nominal}\newline{}José Lara Arce\newline{}5 instancias.\newline{}\textbf{Ventana:} 33--35. & \textbf{Contenido y frontera.} Revisión previa de un lote nominal de hasta cuatro embarcaciones del piloto IN2, individualizado por manifiesto. Se contrasta plan, tareas, responsables y botadura; no se interactúa durante izaje ni se añade una segunda planificación de varadero.\newline{}\textbf{Aceptación.} Plan, responsabilidades y trabajos revisados al menos tres semanas antes de botadura, sin interacción durante izaje.\newline{}\textbf{Insumos.} Universo nominal/manifiesto por instancia, versión y política/protocolo aplicable recibidos, responsables, acceso autorizado y ventana admisible. Registro de instancia: código, población/activos, fecha, versión, evidencia y decisión; el conjunto no se acepta por cerrar solo su primera instancia.\newline{}\textbf{Condiciones.} Condiciones/predecesoras comunes del catálogo: 1.9.2.5.2. Referencia de detalle: FALTA REFERENCIA A SD13.\\\hline
\texttt{1.\allowbreak{}9.\allowbreak{}2.\allowbreak{}5.\allowbreak{}5.\allowbreak{}3.\allowbreak{}1}\newline{}\textbf{Registro quincenal de resultados de botaduras}\newline{}Davor Serey\newline{}1 instancia.\newline{}\textbf{Ventana:} 34. & \textbf{Contenido y frontera.} Se entregará registro quincenal de resultados de botaduras, correspondiente a IN2.5.5.3.1. Comprende registro, validación y emisión del resultado. Excluye los resultados de las fichas predecesoras, la construcción de capacidades base y el servicio recurrente.\newline{}\textbf{Aceptación.} Cada botadura observada se relaciona con el plan y su revisión previa; se registran resultados, ausencias y reprogramaciones. Davor Serey, Calidad, responde por el resultado; los demás perfiles ejecutan sus aportes y el cliente valida las decisiones de su ámbito.\newline{}\textbf{Insumos.} repositorio de evidencia, registros autorizados, herramientas de análisis y validadores del proceso.\newline{}\textbf{Condiciones.} Predecesoras del catálogo: 1.9.2.5.4.*,1.9.2.5.2. Referencia: FALTA REFERENCIA A SD13. Las habilitaciones técnicas se contrastan con arquitectura e inventario vigentes.\\\hline
\texttt{1.\allowbreak{}9.\allowbreak{}2.\allowbreak{}5.\allowbreak{}5.\allowbreak{}3.\allowbreak{}2}\newline{}\textbf{Registro quincenal de resultados de botaduras}\newline{}Davor Serey\newline{}1 instancia.\newline{}\textbf{Ventana:} 34. & \textbf{Contenido y frontera.} Se entregará registro quincenal de resultados de botaduras, correspondiente a IN2.5.5.3.2. Comprende registro, validación y emisión del resultado. Excluye los resultados de las fichas predecesoras, la construcción de capacidades base y el servicio recurrente.\newline{}\textbf{Aceptación.} Cada botadura observada se relaciona con el plan y su revisión previa; se registran resultados, ausencias y reprogramaciones. Davor Serey, Calidad, responde por el resultado; los demás perfiles ejecutan sus aportes y el cliente valida las decisiones de su ámbito.\newline{}\textbf{Insumos.} repositorio de evidencia, registros autorizados, herramientas de análisis y validadores del proceso.\newline{}\textbf{Condiciones.} Predecesoras del catálogo: 1.9.2.5.4.*,1.9.2.5.2. Referencia: FALTA REFERENCIA A SD13. Las habilitaciones técnicas se contrastan con arquitectura e inventario vigentes.\\\hline
\texttt{1.\allowbreak{}9.\allowbreak{}2.\allowbreak{}5.\allowbreak{}5.\allowbreak{}4.\allowbreak{}1}\newline{}\textbf{Registro quincenal de resultados de botaduras}\newline{}Davor Serey\newline{}1 instancia.\newline{}\textbf{Ventana:} 35. & \textbf{Contenido y frontera.} Se entregará registro quincenal de resultados de botaduras, correspondiente a IN2.5.5.4.1. Comprende preparación, ejecución de comprobaciones y expediente. Excluye los resultados de las fichas predecesoras, la construcción de capacidades base y el servicio recurrente.\newline{}\textbf{Aceptación.} Cada botadura observada se relaciona con el plan y su revisión previa; se registran resultados, ausencias y reprogramaciones. Davor Serey, Calidad, responde por el resultado; los demás perfiles ejecutan sus aportes y el cliente valida las decisiones de su ámbito.\newline{}\textbf{Insumos.} PREPROD, versión identificada, datos autorizados y herramientas de pruebas.\newline{}\textbf{Condiciones.} Predecesoras del catálogo: 1.9.2.5.4.*,1.9.2.5.2. Referencia: FALTA REFERENCIA A SD13. Las habilitaciones técnicas se contrastan con arquitectura e inventario vigentes.\\\hline
\texttt{1.\allowbreak{}9.\allowbreak{}2.\allowbreak{}5.\allowbreak{}5.\allowbreak{}4.\allowbreak{}2}\newline{}\textbf{Registro quincenal de resultados de botaduras}\newline{}Davor Serey\newline{}1 instancia.\newline{}\textbf{Ventana:} 35. & \textbf{Contenido y frontera.} Se entregará registro quincenal de resultados de botaduras, correspondiente a IN2.5.5.4.2. Comprende preparación, ejecución de comprobaciones y expediente. Excluye los resultados de las fichas predecesoras, la construcción de capacidades base y el servicio recurrente.\newline{}\textbf{Aceptación.} Cada botadura observada se relaciona con el plan y su revisión previa; se registran resultados, ausencias y reprogramaciones. Davor Serey, Calidad, responde por el resultado; los demás perfiles ejecutan sus aportes y el cliente valida las decisiones de su ámbito.\newline{}\textbf{Insumos.} PREPROD, versión identificada, datos autorizados y herramientas de pruebas.\newline{}\textbf{Condiciones.} Predecesoras del catálogo: 1.9.2.5.4.*,1.9.2.5.2. Referencia: FALTA REFERENCIA A SD13. Las habilitaciones técnicas se contrastan con arquitectura e inventario vigentes.\\\hline
\texttt{1.\allowbreak{}9.\allowbreak{}2.\allowbreak{}5.\allowbreak{}5.\allowbreak{}5.\allowbreak{}1}\newline{}\textbf{Registro quincenal de resultados de botaduras}\newline{}Davor Serey\newline{}1 instancia.\newline{}\textbf{Ventana:} 36. & \textbf{Contenido y frontera.} Se entregará registro quincenal de resultados de botaduras, correspondiente a IN2.5.5.5.1. Comprende registro, validación y emisión del resultado. Excluye los resultados de las fichas predecesoras, la construcción de capacidades base y el servicio recurrente.\newline{}\textbf{Aceptación.} Cada botadura observada se relaciona con el plan y su revisión previa; se registran resultados, ausencias y reprogramaciones. Davor Serey, Calidad, responde por el resultado; los demás perfiles ejecutan sus aportes y el cliente valida las decisiones de su ámbito.\newline{}\textbf{Insumos.} repositorio de evidencia, registros autorizados, herramientas de análisis y validadores del proceso.\newline{}\textbf{Condiciones.} Predecesoras del catálogo: 1.9.2.5.4.*,1.9.2.5.2. Referencia: FALTA REFERENCIA A SD13. Las habilitaciones técnicas se contrastan con arquitectura e inventario vigentes.\\\hline
\texttt{1.\allowbreak{}9.\allowbreak{}2.\allowbreak{}5.\allowbreak{}5.\allowbreak{}5.\allowbreak{}2}\newline{}\textbf{Registro quincenal de resultados de botaduras}\newline{}Davor Serey\newline{}1 instancia.\newline{}\textbf{Ventana:} 36. & \textbf{Contenido y frontera.} Se entregará registro quincenal de resultados de botaduras, correspondiente a IN2.5.5.5.2. Comprende registro, validación y emisión del resultado. Excluye los resultados de las fichas predecesoras, la construcción de capacidades base y el servicio recurrente.\newline{}\textbf{Aceptación.} Cada botadura observada se relaciona con el plan y su revisión previa; se registran resultados, ausencias y reprogramaciones. Davor Serey, Calidad, responde por el resultado; los demás perfiles ejecutan sus aportes y el cliente valida las decisiones de su ámbito.\newline{}\textbf{Insumos.} repositorio de evidencia, registros autorizados, herramientas de análisis y validadores del proceso.\newline{}\textbf{Condiciones.} Predecesoras del catálogo: 1.9.2.5.4.*,1.9.2.5.2. Referencia: FALTA REFERENCIA A SD13. Las habilitaciones técnicas se contrastan con arquitectura e inventario vigentes.\\\hline
\texttt{1.\allowbreak{}9.\allowbreak{}2.\allowbreak{}5.\allowbreak{}6}\newline{}\textbf{Informe comparativo de resultados y acta de evaluación}\newline{}Davor Serey\newline{}1 instancia.\newline{}\textbf{Ventana:} 37. & \textbf{Contenido y frontera.} Se entregará informe comparativo de resultados y acta de evaluación, correspondiente a IN2.5.6. Comprende registro, validación y emisión del resultado. Excluye los resultados de las fichas predecesoras, la construcción de capacidades base y el servicio recurrente.\newline{}\textbf{Aceptación.} Se comparan planes, cumplimiento y botaduras con casos pendientes incluidos; acta identifica límites y decisión. Davor Serey, Calidad, responde por el resultado; los demás perfiles ejecutan sus aportes y el cliente valida las decisiones de su ámbito.\newline{}\textbf{Insumos.} repositorio de evidencia, registros autorizados, herramientas de análisis y validadores del proceso.\newline{}\textbf{Condiciones.} Predecesoras del catálogo: 1.9.2.5.3.*,1.9.2.5.4.*,1.9.2.5.5.*. Referencia: FALTA REFERENCIA A SD13. Las habilitaciones técnicas se contrastan con arquitectura e inventario vigentes.\\\hline
\texttt{1.\allowbreak{}9.\allowbreak{}2.\allowbreak{}6.\allowbreak{}1}\newline{}\textbf{Procedimiento, documentación y transferencia a operación}\newline{}José Lara Arce\newline{}1 instancia.\newline{}\textbf{Ventana:} 37. & \textbf{Contenido y frontera.} Se entregará procedimiento, documentación y transferencia a operación, correspondiente a IN2.6.1. Comprende registro, validación y emisión del resultado. Excluye los resultados de las fichas predecesoras, la construcción de capacidades base y el servicio recurrente.\newline{}\textbf{Aceptación.} Quedan versión del procedimiento, catálogo, responsables de actualización y decisión de transferencia. José Lara Arce, Operación/SRE, responde por el resultado; los demás perfiles ejecutan sus aportes y el cliente valida las decisiones de su ámbito.\newline{}\textbf{Insumos.} repositorio de evidencia, registros autorizados, herramientas de análisis y validadores del proceso.\newline{}\textbf{Condiciones.} Predecesoras del catálogo: 1.9.2.5.6. Referencia: FALTA REFERENCIA A SD13. Las habilitaciones técnicas se contrastan con arquitectura e inventario vigentes.\\\hline
\texttt{1.\allowbreak{}9.\allowbreak{}3.\allowbreak{}1.\allowbreak{}1}\newline{}\textbf{Inventario de formatos, volumen y campos requeridos}\newline{}Nicolás Torfan\newline{}1 instancia.\newline{}\textbf{Ventana:} 13. & \textbf{Contenido y frontera.} Se entregará inventario de formatos, volumen y campos requeridos, correspondiente a IN3.1.1. Excluye los resultados de las fichas predecesoras, la construcción de capacidades base y el servicio recurrente.\newline{}\textbf{Aceptación.} Se identifican formatos, campos obligatorios, volumen observado y casos ilegibles sin asumir una muestra representativa. Nicolás Torfan, Datos, responde por el resultado; los demás perfiles ejecutan sus aportes y el cliente valida las decisiones de su ámbito.\newline{}\textbf{Insumos.} repositorio de evidencia, registros autorizados, herramientas de análisis y validadores del proceso.\newline{}\textbf{Condiciones.} Predecesoras del catálogo: BASE-NAV. Referencia: FALTA REFERENCIA A SD13. Las habilitaciones técnicas se contrastan con arquitectura e inventario vigentes.\\\hline
\texttt{1.\allowbreak{}9.\allowbreak{}3.\allowbreak{}1.\allowbreak{}2}\newline{}\textbf{Interfaces y controles de IA especificados}\newline{}Fernando Guerrero\newline{}1 instancia.\newline{}\textbf{Ventana:} 14. & \textbf{Contenido y frontera.} Se entregará interfaces y controles de ia especificados, correspondiente a IN3.1.2. Excluye los resultados de las fichas predecesoras, la construcción de capacidades base y el servicio recurrente.\newline{}\textbf{Aceptación.} Se define procesamiento, revisión humana, finalidades, errores, retención y frontera con SD4 y requisitos de IA. Fernando Guerrero, Arquitectura/Integración, responde por el resultado; los demás perfiles ejecutan sus aportes y el cliente valida las decisiones de su ámbito.\newline{}\textbf{Insumos.} repositorio de evidencia, registros autorizados, herramientas de análisis y validadores del proceso.\newline{}\textbf{Condiciones.} Predecesoras del catálogo: 1.9.3.1.1. Referencia: FALTA REFERENCIA A SD13. Las habilitaciones técnicas se contrastan con arquitectura e inventario vigentes.\\\hline
\texttt{1.\allowbreak{}9.\allowbreak{}3.\allowbreak{}2.\allowbreak{}1}\newline{}\textbf{Adaptador del extractor y procesamiento asíncrono}\newline{}Nikolai Navea\newline{}1 instancia.\newline{}\textbf{Ventana:} 14. & \textbf{Contenido y frontera.} Se entregará adaptador del extractor y procesamiento asíncrono, correspondiente a IN3.2.1. Comprende implementación, pruebas unitarias y documentación del resultado. Excluye los resultados de las fichas predecesoras, la construcción de capacidades base y el servicio recurrente.\newline{}\textbf{Aceptación.} Cada trabajo identifica versión, documento y respuesta; resultados tardíos y reintentos no alteran datos aceptados. Nikolai Navea, Desarrollo, responde por el resultado; los demás perfiles ejecutan sus aportes y el cliente valida las decisiones de su ámbito.\newline{}\textbf{Insumos.} DEV/QA, contratos de interfaz, repositorio y herramientas de automatización.\newline{}\textbf{Condiciones.} Predecesoras del catálogo: 1.9.3.1.2. Referencia: FALTA REFERENCIA A SD13. Las habilitaciones técnicas se contrastan con arquitectura e inventario vigentes.\\\hline
\texttt{1.\allowbreak{}9.\allowbreak{}3.\allowbreak{}2.\allowbreak{}2}\newline{}\textbf{Reglas de campos, discrepancias y datos pendientes}\newline{}Nicolás Torfan\newline{}1 instancia.\newline{}\textbf{Ventana:} 14. & \textbf{Contenido y frontera.} Se entregará reglas de campos, discrepancias y datos pendientes, correspondiente a IN3.2.2. Comprende implementación, pruebas unitarias y documentación del resultado. Excluye los resultados de las fichas predecesoras, la construcción de capacidades base y el servicio recurrente.\newline{}\textbf{Aceptación.} Campos faltantes o contradictorios quedan pendientes de revisión; no se completan como hechos del documento. Nicolás Torfan, Datos, responde por el resultado; los demás perfiles ejecutan sus aportes y el cliente valida las decisiones de su ámbito.\newline{}\textbf{Insumos.} DEV/QA, contratos de interfaz, repositorio y herramientas de automatización.\newline{}\textbf{Condiciones.} Predecesoras del catálogo: 1.9.3.1.2. Referencia: FALTA REFERENCIA A SD13. Las habilitaciones técnicas se contrastan con arquitectura e inventario vigentes.\\\hline
\texttt{1.\allowbreak{}9.\allowbreak{}3.\allowbreak{}3.\allowbreak{}1}\newline{}\textbf{Formulario de revisión y aprobación humana}\newline{}Nikolai Navea\newline{}1 instancia.\newline{}\textbf{Ventana:} 15. & \textbf{Contenido y frontera.} Se entregará formulario de revisión y aprobación humana, correspondiente a IN3.3.1. Comprende implementación, pruebas unitarias y documentación del resultado. Excluye los resultados de las fichas predecesoras, la construcción de capacidades base y el servicio recurrente.\newline{}\textbf{Aceptación.} El revisor compara fuente y campos y confirma o corrige antes de registrar; cada decisión queda atribuida. Nikolai Navea, Desarrollo, responde por el resultado; los demás perfiles ejecutan sus aportes y el cliente valida las decisiones de su ámbito.\newline{}\textbf{Insumos.} DEV/QA, contratos de interfaz, repositorio y herramientas de automatización.\newline{}\textbf{Condiciones.} Predecesoras del catálogo: 1.9.3.2.1,1.9.3.2.2. Referencia: FALTA REFERENCIA A SD13. Las habilitaciones técnicas se contrastan con arquitectura e inventario vigentes.\\\hline
\texttt{1.\allowbreak{}9.\allowbreak{}3.\allowbreak{}3.\allowbreak{}2}\newline{}\textbf{Integración con expediente, NAV y auditoría}\newline{}Fernando Guerrero\newline{}1 instancia.\newline{}\textbf{Ventana:} 15. & \textbf{Contenido y frontera.} Se entregará integración con expediente, nav y auditoría, correspondiente a IN3.3.2. Comprende implementación, pruebas unitarias y documentación del resultado. Excluye los resultados de las fichas predecesoras, la construcción de capacidades base y el servicio recurrente.\newline{}\textbf{Aceptación.} El registro aprobado conserva documento, actor y versión, y accede a NAV por su contrato autorizado. Fernando Guerrero, Arquitectura/Integración, responde por el resultado; los demás perfiles ejecutan sus aportes y el cliente valida las decisiones de su ámbito.\newline{}\textbf{Insumos.} DEV/QA, contratos de interfaz, repositorio y herramientas de automatización.\newline{}\textbf{Condiciones.} Predecesoras del catálogo: 1.9.3.3.1. Referencia: FALTA REFERENCIA A SD13. Las habilitaciones técnicas se contrastan con arquitectura e inventario vigentes.\\\hline
\texttt{1.\allowbreak{}9.\allowbreak{}3.\allowbreak{}3.\allowbreak{}3}\newline{}\textbf{Ingreso manual, reintentos y resultados tardíos}\newline{}Nikolai Navea\newline{}1 instancia.\newline{}\textbf{Ventana:} 15. & \textbf{Contenido y frontera.} Se entregará ingreso manual, reintentos y resultados tardíos, correspondiente a IN3.3.3. Comprende implementación, pruebas unitarias y documentación del resultado. Excluye los resultados de las fichas predecesoras, la construcción de capacidades base y el servicio recurrente.\newline{}\textbf{Aceptación.} El ingreso manual mantiene los mismos controles y un resultado tardío no sobrescribe una decisión humana. Nikolai Navea, Desarrollo, responde por el resultado; los demás perfiles ejecutan sus aportes y el cliente valida las decisiones de su ámbito.\newline{}\textbf{Insumos.} DEV/QA, contratos de interfaz, repositorio y herramientas de automatización.\newline{}\textbf{Condiciones.} Predecesoras del catálogo: 1.9.3.3.1. Referencia: FALTA REFERENCIA A SD13. Las habilitaciones técnicas se contrastan con arquitectura e inventario vigentes.\\\hline
\texttt{1.\allowbreak{}9.\allowbreak{}3.\allowbreak{}4.\allowbreak{}1}\newline{}\textbf{Pruebas de extracción, integración y continuidad}\newline{}Davor Serey\newline{}1 instancia.\newline{}\textbf{Ventana:} 16. & \textbf{Contenido y frontera.} Se entregará pruebas de extracción, integración y continuidad, correspondiente a IN3.4.1. Comprende preparación, ejecución de comprobaciones y expediente. Excluye los resultados de las fichas predecesoras, la construcción de capacidades base y el servicio recurrente.\newline{}\textbf{Aceptación.} Se comprueban legibilidad, discrepancias, duplicados, fallos del servicio, ingreso manual y recuperación. Davor Serey, Calidad, responde por el resultado; los demás perfiles ejecutan sus aportes y el cliente valida las decisiones de su ámbito.\newline{}\textbf{Insumos.} PREPROD, versión identificada, datos autorizados y herramientas de pruebas.\newline{}\textbf{Condiciones.} Predecesoras del catálogo: 1.9.3.3.2,1.9.3.3.3. Referencia: FALTA REFERENCIA A SD13. Las habilitaciones técnicas se contrastan con arquitectura e inventario vigentes.\\\hline
\texttt{1.\allowbreak{}9.\allowbreak{}3.\allowbreak{}4.\allowbreak{}2}\newline{}\textbf{Matriz y evidencias de controles de IA}\newline{}Bruno Toro\newline{}1 instancia.\newline{}\textbf{Ventana:} 16. & \textbf{Contenido y frontera.} Se entregará matriz y evidencias de controles de ia, correspondiente a IN3.4.2. Comprende preparación, ejecución de comprobaciones y expediente. Excluye los resultados de las fichas predecesoras, la construcción de capacidades base y el servicio recurrente.\newline{}\textbf{Aceptación.} Cada control aplicable se vincula con diseño y evidencia; se verifican aislamiento, revisión y trazabilidad sin declarar cumplimiento por usar IA. Bruno Toro, Seguridad, responde por el resultado; los demás perfiles ejecutan sus aportes y el cliente valida las decisiones de su ámbito.\newline{}\textbf{Insumos.} PREPROD, versión identificada, datos autorizados y herramientas de pruebas.\newline{}\textbf{Condiciones.} Predecesoras del catálogo: 1.9.3.4.1. Referencia: FALTA REFERENCIA A SD13. Las habilitaciones técnicas se contrastan con arquitectura e inventario vigentes.\\\hline
\texttt{1.\allowbreak{}9.\allowbreak{}3.\allowbreak{}5.\allowbreak{}1}\newline{}\textbf{Muestra autorizada, protocolo y capacitación}\newline{}Nicolás Torfan\newline{}1 instancia.\newline{}\textbf{Ventana:} 16. & \textbf{Contenido y frontera.} Se entregará muestra autorizada, protocolo y capacitación, correspondiente a IN3.5.1. Comprende registro, validación y emisión del resultado. Excluye los resultados de las fichas predecesoras, la construcción de capacidades base y el servicio recurrente.\newline{}\textbf{Aceptación.} Muestra, autorizaciones, método manual/asistido y evaluación de usuarios están documentados antes de medir. Nicolás Torfan, Datos, responde por el resultado; los demás perfiles ejecutan sus aportes y el cliente valida las decisiones de su ámbito.\newline{}\textbf{Insumos.} repositorio de evidencia, registros autorizados, herramientas de análisis y validadores del proceso.\newline{}\textbf{Condiciones.} Predecesoras del catálogo: 1.9.3.4.2. Referencia: FALTA REFERENCIA A SD13. Las habilitaciones técnicas se contrastan con arquitectura e inventario vigentes.\\\hline
\texttt{1.\allowbreak{}9.\allowbreak{}3.\allowbreak{}5.\allowbreak{}2.\allowbreak{}19.\allowbreak{}1}\newline{}\textbf{Registro quincenal de ejecución manual y asistida}\newline{}José Lara Arce\newline{}1 instancia.\newline{}\textbf{Ventana:} 18. & \textbf{Contenido y frontera.} Se entregará registro quincenal de ejecución manual y asistida, correspondiente a IN3.5.2.19.1. Comprende registro, validación y emisión del resultado. Excluye los resultados de las fichas predecesoras, la construcción de capacidades base y el servicio recurrente.\newline{}\textbf{Aceptación.} El registro conserva documentos autorizados, tiempos, errores, correcciones y costo de procesamiento; incluye casos no completados. José Lara Arce, Implantación, responde por el resultado; los demás perfiles ejecutan sus aportes y el cliente valida las decisiones de su ámbito.\newline{}\textbf{Insumos.} repositorio de evidencia, registros autorizados, herramientas de análisis y validadores del proceso.\newline{}\textbf{Condiciones.} Predecesoras del catálogo: 1.9.3.5.1. Referencia: FALTA REFERENCIA A SD13. Las habilitaciones técnicas se contrastan con arquitectura e inventario vigentes.\\\hline
\texttt{1.\allowbreak{}9.\allowbreak{}3.\allowbreak{}5.\allowbreak{}2.\allowbreak{}19.\allowbreak{}2}\newline{}\textbf{Registro quincenal de ejecución manual y asistida}\newline{}José Lara Arce\newline{}1 instancia.\newline{}\textbf{Ventana:} 18. & \textbf{Contenido y frontera.} Se entregará registro quincenal de ejecución manual y asistida, correspondiente a IN3.5.2.19.2. Comprende registro, validación y emisión del resultado. Excluye los resultados de las fichas predecesoras, la construcción de capacidades base y el servicio recurrente.\newline{}\textbf{Aceptación.} El registro conserva documentos autorizados, tiempos, errores, correcciones y costo de procesamiento; incluye casos no completados. José Lara Arce, Implantación, responde por el resultado; los demás perfiles ejecutan sus aportes y el cliente valida las decisiones de su ámbito.\newline{}\textbf{Insumos.} repositorio de evidencia, registros autorizados, herramientas de análisis y validadores del proceso.\newline{}\textbf{Condiciones.} Predecesoras del catálogo: 1.9.3.5.2.19.1. Referencia: FALTA REFERENCIA A SD13. Las habilitaciones técnicas se contrastan con arquitectura e inventario vigentes.\\\hline
\texttt{1.\allowbreak{}9.\allowbreak{}3.\allowbreak{}5.\allowbreak{}2.\allowbreak{}20.\allowbreak{}1}\newline{}\textbf{Registro quincenal de ejecución manual y asistida}\newline{}José Lara Arce\newline{}1 instancia.\newline{}\textbf{Ventana:} 19. & \textbf{Contenido y frontera.} Se entregará registro quincenal de ejecución manual y asistida, correspondiente a IN3.5.2.20.1. Comprende registro, validación y emisión del resultado. Excluye los resultados de las fichas predecesoras, la construcción de capacidades base y el servicio recurrente.\newline{}\textbf{Aceptación.} El registro conserva documentos autorizados, tiempos, errores, correcciones y costo de procesamiento; incluye casos no completados. José Lara Arce, Implantación, responde por el resultado; los demás perfiles ejecutan sus aportes y el cliente valida las decisiones de su ámbito.\newline{}\textbf{Insumos.} repositorio de evidencia, registros autorizados, herramientas de análisis y validadores del proceso.\newline{}\textbf{Condiciones.} Predecesoras del catálogo: 1.9.3.5.2.19.2. Referencia: FALTA REFERENCIA A SD13. Las habilitaciones técnicas se contrastan con arquitectura e inventario vigentes.\\\hline
\texttt{1.\allowbreak{}9.\allowbreak{}3.\allowbreak{}5.\allowbreak{}2.\allowbreak{}20.\allowbreak{}2}\newline{}\textbf{Registro quincenal de ejecución manual y asistida}\newline{}José Lara Arce\newline{}1 instancia.\newline{}\textbf{Ventana:} 19. & \textbf{Contenido y frontera.} Se entregará registro quincenal de ejecución manual y asistida, correspondiente a IN3.5.2.20.2. Comprende registro, validación y emisión del resultado. Excluye los resultados de las fichas predecesoras, la construcción de capacidades base y el servicio recurrente.\newline{}\textbf{Aceptación.} El registro conserva documentos autorizados, tiempos, errores, correcciones y costo de procesamiento; incluye casos no completados. José Lara Arce, Implantación, responde por el resultado; los demás perfiles ejecutan sus aportes y el cliente valida las decisiones de su ámbito.\newline{}\textbf{Insumos.} repositorio de evidencia, registros autorizados, herramientas de análisis y validadores del proceso.\newline{}\textbf{Condiciones.} Predecesoras del catálogo: 1.9.3.5.2.20.1. Referencia: FALTA REFERENCIA A SD13. Las habilitaciones técnicas se contrastan con arquitectura e inventario vigentes.\\\hline
\texttt{1.\allowbreak{}9.\allowbreak{}3.\allowbreak{}5.\allowbreak{}3}\newline{}\textbf{Informe de tiempos, errores, costos y beneficio}\newline{}Davor Serey\newline{}1 instancia.\newline{}\textbf{Ventana:} 19. & \textbf{Contenido y frontera.} Se entregará informe de tiempos, errores, costos y beneficio, correspondiente a IN3.5.3. Comprende registro, validación y emisión del resultado. Excluye los resultados de las fichas predecesoras, la construcción de capacidades base y el servicio recurrente.\newline{}\textbf{Aceptación.} La comparación conserva errores, rechazos, revisión humana y costos; evidencia insuficiente no se transforma en mejora demostrada. Davor Serey, Calidad, responde por el resultado; los demás perfiles ejecutan sus aportes y el cliente valida las decisiones de su ámbito.\newline{}\textbf{Insumos.} repositorio de evidencia, registros autorizados, herramientas de análisis y validadores del proceso.\newline{}\textbf{Condiciones.} Predecesoras del catálogo: 1.9.3.5.2.*. Referencia: FALTA REFERENCIA A SD13. Las habilitaciones técnicas se contrastan con arquitectura e inventario vigentes.\\\hline
\texttt{1.\allowbreak{}9.\allowbreak{}3.\allowbreak{}6.\allowbreak{}1}\newline{}\textbf{Configuración, manuales y transferencia}\newline{}José Lara Arce\newline{}1 instancia.\newline{}\textbf{Ventana:} 21. & \textbf{Contenido y frontera.} Se entregará configuración, manuales y transferencia, correspondiente a IN3.6.1. Comprende registro, validación y emisión del resultado. Excluye los resultados de las fichas predecesoras, la construcción de capacidades base y el servicio recurrente.\newline{}\textbf{Aceptación.} Se reciben configuración, límites, ingreso manual, monitoreo de consumo y procedimiento de incidentes. José Lara Arce, Operación/SRE, responde por el resultado; los demás perfiles ejecutan sus aportes y el cliente valida las decisiones de su ámbito.\newline{}\textbf{Insumos.} repositorio de evidencia, registros autorizados, herramientas de análisis y validadores del proceso.\newline{}\textbf{Condiciones.} Predecesoras del catálogo: 1.9.3.5.3. Referencia: FALTA REFERENCIA A SD13. Las habilitaciones técnicas se contrastan con arquitectura e inventario vigentes.\\\hline
\texttt{1.\allowbreak{}9.\allowbreak{}4.\allowbreak{}1.\allowbreak{}1}\newline{}\textbf{Ficha de medición y atribución del resultado}\newline{}Ulysses Barreto\newline{}1 instancia.\newline{}\textbf{Ventana:} 2. & \textbf{Contenido y frontera.} Se entregará ficha de medición y atribución del resultado, correspondiente a IN4.1.1. Excluye los resultados de las fichas predecesoras, la construcción de capacidades base y el servicio recurrente.\newline{}\textbf{Aceptación.} Se define tiempo activo, tareas comparables, calidad y exclusiones; las esperas y ahorros de otras innovaciones no se duplican. Ulysses Barreto, Funcional, responde por el resultado; los demás perfiles ejecutan sus aportes y el cliente valida las decisiones de su ámbito.\newline{}\textbf{Insumos.} repositorio de evidencia, registros autorizados, herramientas de análisis y validadores del proceso.\newline{}\textbf{Condiciones.} Predecesoras del catálogo: BASE-VAR. Referencia: FALTA REFERENCIA A SD13. Las habilitaciones técnicas se contrastan con arquitectura e inventario vigentes.\\\hline
\texttt{1.\allowbreak{}9.\allowbreak{}4.\allowbreak{}1.\allowbreak{}2}\newline{}\textbf{Regla comercial y escenarios de incentivo}\newline{}Ignacio Silva\newline{}1 instancia.\newline{}\textbf{Ventana:} 4. & \textbf{Contenido y frontera.} Se entregará regla comercial y escenarios de incentivo, correspondiente a IN4.1.2. Excluye los resultados de las fichas predecesoras, la construcción de capacidades base y el servicio recurrente.\newline{}\textbf{Aceptación.} La fórmula, tope, rechazo y factor cero están documentados; no modifica las obligaciones ni los hitos E-25. Ignacio Silva, Proyecto, responde por el resultado; los demás perfiles ejecutan sus aportes y el cliente valida las decisiones de su ámbito.\newline{}\textbf{Insumos.} repositorio de evidencia, registros autorizados, herramientas de análisis y validadores del proceso.\newline{}\textbf{Condiciones.} Predecesoras del catálogo: 1.9.4.1.1. Referencia: FALTA REFERENCIA A SD13. Las habilitaciones técnicas se contrastan con arquitectura e inventario vigentes.\\\hline
\texttt{1.\allowbreak{}9.\allowbreak{}4.\allowbreak{}2.\allowbreak{}1}\newline{}\textbf{Interfaces y expediente de medición especificados}\newline{}Fernando Guerrero\newline{}1 instancia.\newline{}\textbf{Ventana:} 14--15. & \textbf{Contenido y frontera.} Se entregará interfaces y expediente de medición especificados, correspondiente a IN4.2.1. Comprende registro, validación y emisión del resultado. Excluye los resultados de las fichas predecesoras, la construcción de capacidades base y el servicio recurrente.\newline{}\textbf{Aceptación.} Se especifican captura de intervalos, solicitudes, autoridad, versiones y separación entre preparación y aprobación. Fernando Guerrero, Arquitectura/Integración, responde por el resultado; los demás perfiles ejecutan sus aportes y el cliente valida las decisiones de su ámbito.\newline{}\textbf{Insumos.} repositorio de evidencia, registros autorizados, herramientas de análisis y validadores del proceso.\newline{}\textbf{Condiciones.} Predecesoras del catálogo: 1.9.4.1.2,1.9.1.1.1. La especificación final y su recepción son posteriores a la recepción del flujo, estados y matriz de roles de IN1. Compartir la ventana 14--15 no autoriza recibir esta especificación con ese diseño pendiente. Referencia: FALTA REFERENCIA A SD13. Las habilitaciones técnicas se contrastan con arquitectura e inventario vigentes.\\\hline
\texttt{1.\allowbreak{}9.\allowbreak{}4.\allowbreak{}2.\allowbreak{}2}\newline{}\textbf{Captura de intervalos y validaciones implementadas}\newline{}Nikolai Navea\newline{}1 instancia.\newline{}\textbf{Ventana:} 16. & \textbf{Contenido y frontera.} Se entregará captura de intervalos y validaciones implementadas, correspondiente a IN4.2.2. Comprende implementación, pruebas unitarias y documentación del resultado. Excluye los resultados de las fichas predecesoras, la construcción de capacidades base y el servicio recurrente.\newline{}\textbf{Aceptación.} Cada intervalo tiene solicitud, persona, inicio, término y origen; se detectan superposiciones y datos incompletos. Nikolai Navea, Desarrollo, responde por el resultado; los demás perfiles ejecutan sus aportes y el cliente valida las decisiones de su ámbito.\newline{}\textbf{Insumos.} DEV/QA, contratos de interfaz, repositorio y herramientas de automatización.\newline{}\textbf{Condiciones.} Predecesoras del catálogo: 1.9.4.2.1. Referencia: FALTA REFERENCIA A SD13. Las habilitaciones técnicas se contrastan con arquitectura e inventario vigentes.\\\hline
\texttt{1.\allowbreak{}9.\allowbreak{}4.\allowbreak{}2.\allowbreak{}3}\newline{}\textbf{Cálculo versionado y expediente de conciliación}\newline{}Nikolai Navea\newline{}1 instancia.\newline{}\textbf{Ventana:} 17. & \textbf{Contenido y frontera.} Se entregará cálculo versionado y expediente de conciliación, correspondiente a IN4.2.3. Comprende implementación, pruebas unitarias y documentación del resultado. Excluye los resultados de las fichas predecesoras, la construcción de capacidades base y el servicio recurrente.\newline{}\textbf{Aceptación.} El cálculo es reproducible, conserva sus entradas y versiones y no emite cobro automático. Nikolai Navea, Desarrollo, responde por el resultado; los demás perfiles ejecutan sus aportes y el cliente valida las decisiones de su ámbito.\newline{}\textbf{Insumos.} DEV/QA, contratos de interfaz, repositorio y herramientas de automatización.\newline{}\textbf{Condiciones.} Predecesoras del catálogo: 1.9.4.2.2. Referencia: FALTA REFERENCIA A SD13. Las habilitaciones técnicas se contrastan con arquitectura e inventario vigentes.\\\hline
\texttt{1.\allowbreak{}9.\allowbreak{}4.\allowbreak{}2.\allowbreak{}4}\newline{}\textbf{Evidencia de seguridad y cálculo del incentivo}\newline{}Davor Serey\newline{}1 instancia.\newline{}\textbf{Ventana:} 17. & \textbf{Contenido y frontera.} Se entregará evidencia de seguridad y cálculo del incentivo, correspondiente a IN4.2.4. Comprende registro, validación y emisión del resultado. Excluye los resultados de las fichas predecesoras, la construcción de capacidades base y el servicio recurrente.\newline{}\textbf{Aceptación.} Los casos sin evidencia, base cero, calidad fallida y muestra insuficiente producen factor cero; no hay cambios silenciosos. Davor Serey, Calidad, responde por el resultado; los demás perfiles ejecutan sus aportes y el cliente valida las decisiones de su ámbito.\newline{}\textbf{Insumos.} repositorio de evidencia, registros autorizados, herramientas de análisis y validadores del proceso.\newline{}\textbf{Condiciones.} Predecesoras del catálogo: 1.9.4.2.3. Las pruebas se ejecutan sobre la versión del cálculo terminada y comprobada; compartir el mes 17 no permite anticiparlas a esa condición. T-15 debe asignar duraciones y recursos que permitan completar cálculo, pruebas y correcciones antes de iniciar la medición de línea base del mes 18. Referencia: FALTA REFERENCIA A SD13. Las habilitaciones técnicas se contrastan con arquitectura e inventario vigentes.\\\hline
\texttt{1.\allowbreak{}9.\allowbreak{}4.\allowbreak{}3.\allowbreak{}1.\allowbreak{}1}\newline{}\textbf{Registro base del tiempo activo}\newline{}Davor Serey\newline{}1 instancia.\newline{}\textbf{Ventana:} 18. & \textbf{Contenido y frontera.} Se entregará registro base del tiempo activo, correspondiente a IN4.3.1.1. Excluye los resultados de las fichas predecesoras, la construcción de capacidades base y el servicio recurrente.\newline{}\textbf{Aceptación.} Cada intervalo identifica persona, solicitud y tarea; esperas y tiempos superpuestos quedan excluidos sin duplicar la medición de IN1. Davor Serey, Calidad, responde por el resultado; los demás perfiles ejecutan sus aportes y el cliente valida las decisiones de su ámbito.\newline{}\textbf{Insumos.} repositorio de evidencia, registros autorizados, herramientas de análisis y validadores del proceso.\newline{}\textbf{Condiciones.} Predecesoras del catálogo: 1.9.4.2.4. Referencia: FALTA REFERENCIA A SD13. Las habilitaciones técnicas se contrastan con arquitectura e inventario vigentes.\\\hline
\texttt{1.\allowbreak{}9.\allowbreak{}4.\allowbreak{}3.\allowbreak{}1.\allowbreak{}2}\newline{}\textbf{Registro base del tiempo activo}\newline{}Davor Serey\newline{}1 instancia.\newline{}\textbf{Ventana:} 18. & \textbf{Contenido y frontera.} Se entregará registro base del tiempo activo, correspondiente a IN4.3.1.2. Excluye los resultados de las fichas predecesoras, la construcción de capacidades base y el servicio recurrente.\newline{}\textbf{Aceptación.} Cada intervalo identifica persona, solicitud y tarea; esperas y tiempos superpuestos quedan excluidos sin duplicar la medición de IN1. Davor Serey, Calidad, responde por el resultado; los demás perfiles ejecutan sus aportes y el cliente valida las decisiones de su ámbito.\newline{}\textbf{Insumos.} repositorio de evidencia, registros autorizados, herramientas de análisis y validadores del proceso.\newline{}\textbf{Condiciones.} Predecesoras del catálogo: 1.9.4.3.1.1. Referencia: FALTA REFERENCIA A SD13. Las habilitaciones técnicas se contrastan con arquitectura e inventario vigentes.\\\hline
\texttt{1.\allowbreak{}9.\allowbreak{}4.\allowbreak{}3.\allowbreak{}2.\allowbreak{}1}\newline{}\textbf{Registro de evaluación del tiempo activo}\newline{}Davor Serey\newline{}1 instancia.\newline{}\textbf{Ventana:} 19. & \textbf{Contenido y frontera.} Se entregará registro de evaluación del tiempo activo, correspondiente a IN4.3.2.1. Comprende registro, validación y emisión del resultado. Excluye los resultados de las fichas predecesoras, la construcción de capacidades base y el servicio recurrente.\newline{}\textbf{Aceptación.} Cada intervalo identifica persona, solicitud y tarea; esperas y tiempos superpuestos quedan excluidos sin duplicar la medición de IN1. Davor Serey, Calidad, responde por el resultado; los demás perfiles ejecutan sus aportes y el cliente valida las decisiones de su ámbito.\newline{}\textbf{Insumos.} repositorio de evidencia, registros autorizados, herramientas de análisis y validadores del proceso.\newline{}\textbf{Condiciones.} Predecesoras del catálogo: 1.9.4.3.1.2. Referencia: FALTA REFERENCIA A SD13. Las habilitaciones técnicas se contrastan con arquitectura e inventario vigentes.\\\hline
\texttt{1.\allowbreak{}9.\allowbreak{}4.\allowbreak{}3.\allowbreak{}2.\allowbreak{}2}\newline{}\textbf{Registro de evaluación del tiempo activo}\newline{}Davor Serey\newline{}1 instancia.\newline{}\textbf{Ventana:} 19. & \textbf{Contenido y frontera.} Se entregará registro de evaluación del tiempo activo, correspondiente a IN4.3.2.2. Comprende registro, validación y emisión del resultado. Excluye los resultados de las fichas predecesoras, la construcción de capacidades base y el servicio recurrente.\newline{}\textbf{Aceptación.} Cada intervalo identifica persona, solicitud y tarea; esperas y tiempos superpuestos quedan excluidos sin duplicar la medición de IN1. Davor Serey, Calidad, responde por el resultado; los demás perfiles ejecutan sus aportes y el cliente valida las decisiones de su ámbito.\newline{}\textbf{Insumos.} repositorio de evidencia, registros autorizados, herramientas de análisis y validadores del proceso.\newline{}\textbf{Condiciones.} Predecesoras del catálogo: 1.9.4.3.2.1. Referencia: FALTA REFERENCIA A SD13. Las habilitaciones técnicas se contrastan con arquitectura e inventario vigentes.\\\hline
\texttt{1.\allowbreak{}9.\allowbreak{}4.\allowbreak{}3.\allowbreak{}3}\newline{}\textbf{Evaluación del piloto sin efectos de pago}\newline{}Davor Serey\newline{}1 instancia.\newline{}\textbf{Ventana:} 19. & \textbf{Contenido y frontera.} Se entregará evaluación del piloto sin efectos de pago, correspondiente a IN4.3.3. Comprende registro, validación y emisión del resultado. Excluye los resultados de las fichas predecesoras, la construcción de capacidades base y el servicio recurrente.\newline{}\textbf{Aceptación.} Se compara tiempo activo y calidad con la regla aprobada y se conserva decisión; el piloto no genera pagos variables. Davor Serey, Calidad, responde por el resultado; los demás perfiles ejecutan sus aportes y el cliente valida las decisiones de su ámbito.\newline{}\textbf{Insumos.} repositorio de evidencia, registros autorizados, herramientas de análisis y validadores del proceso.\newline{}\textbf{Condiciones.} Predecesoras del catálogo: 1.9.4.3.1.*,1.9.4.3.2.*. Referencia: FALTA REFERENCIA A SD13. Las habilitaciones técnicas se contrastan con arquitectura e inventario vigentes.\\\hline
\texttt{1.\allowbreak{}9.\allowbreak{}4.\allowbreak{}4.\allowbreak{}1}\newline{}\textbf{Procedimiento de conciliación contractual habilitado}\newline{}Ignacio Silva\newline{}1 instancia.\newline{}\textbf{Ventana:} 21. & \textbf{Contenido y frontera.} Se entregará procedimiento de conciliación contractual habilitado, correspondiente a IN4.4.1. Comprende registro, validación y emisión del resultado. Excluye los resultados de las fichas predecesoras, la construcción de capacidades base y el servicio recurrente.\newline{}\textbf{Aceptación.} Cliente y Synaptix disponen de procedimiento, evidencia, reclamación y decisión registrada antes de habilitar el incentivo. Ignacio Silva, Proyecto, responde por el resultado; los demás perfiles ejecutan sus aportes y el cliente valida las decisiones de su ámbito.\newline{}\textbf{Insumos.} repositorio de evidencia, registros autorizados, herramientas de análisis y validadores del proceso.\newline{}\textbf{Condiciones.} Predecesoras del catálogo: 1.9.4.3.3,1.9.1.6.1. Referencia: FALTA REFERENCIA A SD13. Las habilitaciones técnicas se contrastan con arquitectura e inventario vigentes.\\\hline
\texttt{1.\allowbreak{}9.\allowbreak{}5.\allowbreak{}1.\allowbreak{}1}\newline{}\textbf{Plan ambiental y protocolo de comparación}\newline{}Ulysses Barreto\newline{}1 instancia.\newline{}\textbf{Ventana:} 13. & \textbf{Contenido y frontera.} Se entregará plan ambiental y protocolo de comparación, correspondiente a IN5.1.1. Excluye los resultados de las fichas predecesoras, la construcción de capacidades base y el servicio recurrente.\newline{}\textbf{Aceptación.} Se fijan indicador de consumo por permanencia, ventanas de 28 días, criterios de comparabilidad y evidencia de acciones. Ulysses Barreto, Funcional, responde por el resultado; los demás perfiles ejecutan sus aportes y el cliente valida las decisiones de su ámbito.\newline{}\textbf{Insumos.} repositorio de evidencia, registros autorizados, herramientas de análisis y validadores del proceso.\newline{}\textbf{Condiciones.} Predecesoras del catálogo: BASE-AMB. Referencia: FALTA REFERENCIA A SD13. Las habilitaciones técnicas se contrastan con arquitectura e inventario vigentes.\\\hline
\texttt{1.\allowbreak{}9.\allowbreak{}5.\allowbreak{}1.\allowbreak{}2}\newline{}\textbf{Consentimiento, modelo y contratos definidos}\newline{}Fernando Guerrero\newline{}1 instancia.\newline{}\textbf{Ventana:} 14. & \textbf{Contenido y frontera.} Se entregará consentimiento, modelo y contratos definidos, correspondiente a IN5.1.2. Excluye los resultados de las fichas predecesoras, la construcción de capacidades base y el servicio recurrente.\newline{}\textbf{Aceptación.} Adhesión, plan, acción y evaluación tienen estados, conservación y permisos; se especifica revocación y autoridad de lecturas. Fernando Guerrero, Arquitectura/Integración, responde por el resultado; los demás perfiles ejecutan sus aportes y el cliente valida las decisiones de su ámbito.\newline{}\textbf{Insumos.} repositorio de evidencia, registros autorizados, herramientas de análisis y validadores del proceso.\newline{}\textbf{Condiciones.} Predecesoras del catálogo: 1.9.5.1.1. Referencia: FALTA REFERENCIA A SD13. Las habilitaciones técnicas se contrastan con arquitectura e inventario vigentes.\\\hline
\texttt{1.\allowbreak{}9.\allowbreak{}5.\allowbreak{}2.\allowbreak{}1}\newline{}\textbf{Historial validado y cálculo por permanencia}\newline{}Nikolai Navea\newline{}1 instancia.\newline{}\textbf{Ventana:} 14. & \textbf{Contenido y frontera.} Se entregará historial validado y cálculo por permanencia, correspondiente a IN5.2.1. Comprende implementación, pruebas unitarias y documentación del resultado. Excluye los resultados de las fichas predecesoras, la construcción de capacidades base y el servicio recurrente.\newline{}\textbf{Aceptación.} El historial usa intervalos conciliados, conserva originales y no suma lecturas acumulativas ni períodos repetidos. Nikolai Navea, Desarrollo, responde por el resultado; los demás perfiles ejecutan sus aportes y el cliente valida las decisiones de su ámbito.\newline{}\textbf{Insumos.} DEV/QA, contratos de interfaz, repositorio y herramientas de automatización.\newline{}\textbf{Condiciones.} Predecesoras del catálogo: 1.9.5.1.2. Referencia: FALTA REFERENCIA A SD13. Las habilitaciones técnicas se contrastan con arquitectura e inventario vigentes.\\\hline
\texttt{1.\allowbreak{}9.\allowbreak{}5.\allowbreak{}2.\allowbreak{}2}\newline{}\textbf{Adhesión, acciones y portal ambiental}\newline{}Nikolai Navea\newline{}1 instancia.\newline{}\textbf{Ventana:} 15. & \textbf{Contenido y frontera.} Se entregará adhesión, acciones y portal ambiental, correspondiente a IN5.2.2. Comprende implementación, pruebas unitarias y documentación del resultado. Excluye los resultados de las fichas predecesoras, la construcción de capacidades base y el servicio recurrente.\newline{}\textbf{Aceptación.} El armador acepta el plan y consulta sólo sus datos; cada acción tiene responsable, fechas y evidencia. Nikolai Navea, Desarrollo, responde por el resultado; los demás perfiles ejecutan sus aportes y el cliente valida las decisiones de su ámbito.\newline{}\textbf{Insumos.} DEV/QA, contratos de interfaz, repositorio y herramientas de automatización.\newline{}\textbf{Condiciones.} Predecesoras del catálogo: 1.9.5.1.2. Referencia: FALTA REFERENCIA A SD13. Las habilitaciones técnicas se contrastan con arquitectura e inventario vigentes.\\\hline
\texttt{1.\allowbreak{}9.\allowbreak{}5.\allowbreak{}2.\allowbreak{}3}\newline{}\textbf{Integración, consentimiento y actualización de cortes}\newline{}Fernando Guerrero\newline{}1 instancia.\newline{}\textbf{Ventana:} 15. & \textbf{Contenido y frontera.} Se entregará integración, consentimiento y actualización de cortes, correspondiente a IN5.2.3. Comprende implementación, pruebas unitarias y documentación del resultado. Excluye los resultados de las fichas predecesoras, la construcción de capacidades base y el servicio recurrente.\newline{}\textbf{Aceptación.} Revocación y correcciones propagan límites de uso; la vista identifica cobertura y último corte validado. Fernando Guerrero, Arquitectura/Integración, responde por el resultado; los demás perfiles ejecutan sus aportes y el cliente valida las decisiones de su ámbito.\newline{}\textbf{Insumos.} DEV/QA, contratos de interfaz, repositorio y herramientas de automatización.\newline{}\textbf{Condiciones.} Predecesoras del catálogo: 1.9.5.2.1,1.9.5.2.2. Referencia: FALTA REFERENCIA A SD13. Las habilitaciones técnicas se contrastan con arquitectura e inventario vigentes.\\\hline
\texttt{1.\allowbreak{}9.\allowbreak{}5.\allowbreak{}2.\allowbreak{}4}\newline{}\textbf{Pruebas de medición, privacidad y accesibilidad}\newline{}Davor Serey\newline{}1 instancia.\newline{}\textbf{Ventana:} 16. & \textbf{Contenido y frontera.} Se entregará pruebas de medición, privacidad y accesibilidad, correspondiente a IN5.2.4. Comprende registro, validación y emisión del resultado. Excluye los resultados de las fichas predecesoras, la construcción de capacidades base y el servicio recurrente.\newline{}\textbf{Aceptación.} Se prueban acceso entre flotas, exportaciones, revocación, diferencias de lectura, permanencia y vistas accesibles. Davor Serey, Calidad, responde por el resultado; los demás perfiles ejecutan sus aportes y el cliente valida las decisiones de su ámbito.\newline{}\textbf{Insumos.} repositorio de evidencia, registros autorizados, herramientas de análisis y validadores del proceso.\newline{}\textbf{Condiciones.} Predecesoras del catálogo: 1.9.5.2.3. Referencia: FALTA REFERENCIA A SD13. Las habilitaciones técnicas se contrastan con arquitectura e inventario vigentes.\\\hline
\texttt{1.\allowbreak{}9.\allowbreak{}5.\allowbreak{}3.\allowbreak{}1}\newline{}\textbf{Participantes y acciones preparados para el piloto}\newline{}José Lara Arce\newline{}1 instancia.\newline{}\textbf{Ventana:} 16. & \textbf{Contenido y frontera.} Se entregará participantes y acciones preparados para el piloto, correspondiente a IN5.3.1. Comprende registro, validación y emisión del resultado. Excluye los resultados de las fichas predecesoras, la construcción de capacidades base y el servicio recurrente.\newline{}\textbf{Aceptación.} Hay adhesión voluntaria, medición validada, permanencias conciliadas y usuarios preparados; no se presume participación efectiva. José Lara Arce, Implantación, responde por el resultado; los demás perfiles ejecutan sus aportes y el cliente valida las decisiones de su ámbito.\newline{}\textbf{Insumos.} repositorio de evidencia, registros autorizados, herramientas de análisis y validadores del proceso.\newline{}\textbf{Condiciones.} Predecesoras del catálogo: 1.9.5.2.4. Referencia: FALTA REFERENCIA A SD13. Las habilitaciones técnicas se contrastan con arquitectura e inventario vigentes.\\\hline
\texttt{1.\allowbreak{}9.\allowbreak{}5.\allowbreak{}3.\allowbreak{}2.\allowbreak{}1}\newline{}\textbf{Registro base del consumo normalizado}\newline{}Davor Serey\newline{}1 instancia.\newline{}\textbf{Ventana:} 18. & \textbf{Contenido y frontera.} Se entregará registro base del consumo normalizado, correspondiente a IN5.3.2.1. Comprende registro, validación y emisión del resultado. Excluye los resultados de las fichas predecesoras, la construcción de capacidades base y el servicio recurrente.\newline{}\textbf{Aceptación.} Se conservan intervalos de consumo conciliados y permanencia de la cohorte; no se afirma reducción con datos faltantes o acumulativos duplicados. Davor Serey, Calidad, responde por el resultado; los demás perfiles ejecutan sus aportes y el cliente valida las decisiones de su ámbito.\newline{}\textbf{Insumos.} repositorio de evidencia, registros autorizados, herramientas de análisis y validadores del proceso.\newline{}\textbf{Condiciones.} Predecesoras del catálogo: 1.9.5.3.1. Referencia: FALTA REFERENCIA A SD13. Las habilitaciones técnicas se contrastan con arquitectura e inventario vigentes.\\\hline
\texttt{1.\allowbreak{}9.\allowbreak{}5.\allowbreak{}3.\allowbreak{}2.\allowbreak{}2}\newline{}\textbf{Registro base del consumo normalizado}\newline{}Davor Serey\newline{}1 instancia.\newline{}\textbf{Ventana:} 18. & \textbf{Contenido y frontera.} Se entregará registro base del consumo normalizado, correspondiente a IN5.3.2.2. Comprende registro, validación y emisión del resultado. Excluye los resultados de las fichas predecesoras, la construcción de capacidades base y el servicio recurrente.\newline{}\textbf{Aceptación.} Se conservan intervalos de consumo conciliados y permanencia de la cohorte; no se afirma reducción con datos faltantes o acumulativos duplicados. Davor Serey, Calidad, responde por el resultado; los demás perfiles ejecutan sus aportes y el cliente valida las decisiones de su ámbito.\newline{}\textbf{Insumos.} repositorio de evidencia, registros autorizados, herramientas de análisis y validadores del proceso.\newline{}\textbf{Condiciones.} Predecesoras del catálogo: 1.9.5.3.2.1. Referencia: FALTA REFERENCIA A SD13. Las habilitaciones técnicas se contrastan con arquitectura e inventario vigentes.\\\hline
\texttt{1.\allowbreak{}9.\allowbreak{}5.\allowbreak{}3.\allowbreak{}3.\allowbreak{}1}\newline{}\textbf{Registro de evaluación del consumo normalizado}\newline{}Davor Serey\newline{}1 instancia.\newline{}\textbf{Ventana:} 19. & \textbf{Contenido y frontera.} Se entregará registro de evaluación del consumo normalizado, correspondiente a IN5.3.3.1. Comprende registro, validación y emisión del resultado. Excluye los resultados de las fichas predecesoras, la construcción de capacidades base y el servicio recurrente.\newline{}\textbf{Aceptación.} Se conservan intervalos de consumo conciliados y permanencia de la cohorte; no se afirma reducción con datos faltantes o acumulativos duplicados. Davor Serey, Calidad, responde por el resultado; los demás perfiles ejecutan sus aportes y el cliente valida las decisiones de su ámbito.\newline{}\textbf{Insumos.} repositorio de evidencia, registros autorizados, herramientas de análisis y validadores del proceso.\newline{}\textbf{Condiciones.} Predecesoras del catálogo: 1.9.5.3.2.2. Referencia: FALTA REFERENCIA A SD13. Las habilitaciones técnicas se contrastan con arquitectura e inventario vigentes.\\\hline
\texttt{1.\allowbreak{}9.\allowbreak{}5.\allowbreak{}3.\allowbreak{}3.\allowbreak{}2}\newline{}\textbf{Registro de evaluación del consumo normalizado}\newline{}Davor Serey\newline{}1 instancia.\newline{}\textbf{Ventana:} 19. & \textbf{Contenido y frontera.} Se entregará registro de evaluación del consumo normalizado, correspondiente a IN5.3.3.2. Comprende registro, validación y emisión del resultado. Excluye los resultados de las fichas predecesoras, la construcción de capacidades base y el servicio recurrente.\newline{}\textbf{Aceptación.} Se conservan intervalos de consumo conciliados y permanencia de la cohorte; no se afirma reducción con datos faltantes o acumulativos duplicados. Davor Serey, Calidad, responde por el resultado; los demás perfiles ejecutan sus aportes y el cliente valida las decisiones de su ámbito.\newline{}\textbf{Insumos.} repositorio de evidencia, registros autorizados, herramientas de análisis y validadores del proceso.\newline{}\textbf{Condiciones.} Predecesoras del catálogo: 1.9.5.3.3.1. Referencia: FALTA REFERENCIA A SD13. Las habilitaciones técnicas se contrastan con arquitectura e inventario vigentes.\\\hline
\texttt{1.\allowbreak{}9.\allowbreak{}5.\allowbreak{}3.\allowbreak{}4}\newline{}\textbf{Evaluación ambiental y acta del piloto}\newline{}Davor Serey\newline{}1 instancia.\newline{}\textbf{Ventana:} 19. & \textbf{Contenido y frontera.} Se entregará evaluación ambiental y acta del piloto, correspondiente a IN5.3.4. Comprende registro, validación y emisión del resultado. Excluye los resultados de las fichas predecesoras, la construcción de capacidades base y el servicio recurrente.\newline{}\textbf{Aceptación.} Se comparan ventanas equivalentes y acciones evidenciadas; faltantes o muestra insuficiente impiden afirmar reducción. Davor Serey, Calidad, responde por el resultado; los demás perfiles ejecutan sus aportes y el cliente valida las decisiones de su ámbito.\newline{}\textbf{Insumos.} repositorio de evidencia, registros autorizados, herramientas de análisis y validadores del proceso.\newline{}\textbf{Condiciones.} Predecesoras del catálogo: 1.9.5.3.2.*,1.9.5.3.3.*. Referencia: FALTA REFERENCIA A SD13. Las habilitaciones técnicas se contrastan con arquitectura e inventario vigentes.\\\hline
\texttt{1.\allowbreak{}9.\allowbreak{}5.\allowbreak{}4.\allowbreak{}1}\newline{}\textbf{Pasaporte y seguimiento transferidos a operación}\newline{}José Lara Arce\newline{}1 instancia.\newline{}\textbf{Ventana:} 21. & \textbf{Contenido y frontera.} Se entregará pasaporte y seguimiento transferidos a operación, correspondiente a IN5.4.1. Comprende registro, validación y emisión del resultado. Excluye los resultados de las fichas predecesoras, la construcción de capacidades base y el servicio recurrente.\newline{}\textbf{Aceptación.} Se reciben configuración, consentimiento, monitoreo de calidad de lecturas y responsables del informe previo a temporada 2029. José Lara Arce, Operación/SRE, responde por el resultado; los demás perfiles ejecutan sus aportes y el cliente valida las decisiones de su ámbito.\newline{}\textbf{Insumos.} repositorio de evidencia, registros autorizados, herramientas de análisis y validadores del proceso.\newline{}\textbf{Condiciones.} Predecesoras del catálogo: 1.9.5.3.4. Referencia: FALTA REFERENCIA A SD13. Las habilitaciones técnicas se contrastan con arquitectura e inventario vigentes.\\\hline
\end{longtable}\endgroup
\begingroup\setlength{\tabcolsep}{3pt}
\begin{longtable}{L{\dimexpr0.320000\textwidth-2\tabcolsep\relax}L{\dimexpr0.680000\textwidth-2\tabcolsep\relax}}
\caption{Diccionario: Verificación y aceptación}\label{tab:t14-dic-1-10}\\\hline
\EncabezadoTabla{Paquete y resultado} & \EncabezadoTabla{Frontera, aceptación e insumos}\\\hline\endfirsthead
\multicolumn{2}{l}{\textit{Continuación de la tabla \ref{tab:t14-dic-1-10}}}\\\hline
\EncabezadoTabla{Paquete y resultado} & \EncabezadoTabla{Frontera, aceptación e insumos}\\\hline\endhead
\hline\endfoot
\hline\endlastfoot
\texttt{1.\allowbreak{}10.\allowbreak{}1.\allowbreak{}1.\allowbreak{}\{1,\allowbreak{}2\}.\allowbreak{}1}\newline{}\textbf{Ensayo completo E1, ensayos 1 y 2: PREPROD, origen y protocolo preparados}\newline{}Davor Serey\newline{}2 instancias.\newline{}\textbf{Ventana:} Antes de C del grupo E1: ensayo 1 en C menos 3 semanas; ensayo 2 en C menos 1 semana. & \textbf{Contenido y frontera.} Resultado parcial del ensayo integral; excluye construir nuevamente las herramientas de migración y repetir la misma ejecución en T-18.\newline{}\textbf{Aceptación.} PREPROD aislado, copia protegida, fuentes, versiones, accesos y protocolo de los seis conjuntos identificados; ninguna preparación sustituye el ensayo completo.\newline{}\textbf{Insumos.} Versión/registro recibido, acceso autorizado, herramientas, participantes y ventana del período.\newline{}\textbf{Condiciones.} Cada instancia identifica etapa, ensayo, grupo, C y versión. C es el cambio de autoridad del grupo, no necesariamente el paso oficial a producción. Dos ensayos completos aprobados antes del primer C; T-18 fija cortes y actualización/repetición si cambian sus condiciones. Revisión de habilitación C menos 2 días hábiles. Referencia: SD5, tabla 5.13; FALTA REFERENCIA A SD9.\\\hline
\texttt{1.\allowbreak{}10.\allowbreak{}1.\allowbreak{}1.\allowbreak{}\{1,\allowbreak{}2\}.\allowbreak{}2}\newline{}\textbf{Ensayo completo E1, ensayos 1 y 2: Ejecución, manifiestos y tiempos}\newline{}Davor Serey\newline{}2 instancias.\newline{}\textbf{Ventana:} Antes de C del grupo E1: ensayo 1 en C menos 3 semanas; ensayo 2 en C menos 1 semana. & \textbf{Contenido y frontera.} Resultado parcial del ensayo integral; excluye construir nuevamente las herramientas de migración y repetir la misma ejecución en T-18.\newline{}\textbf{Aceptación.} Extracción, transformación y carga recorridas de extremo a extremo; manifiestos, errores y tiempos medidos sobre las versiones y el delta de la etapa.\newline{}\textbf{Insumos.} Versión/registro recibido, acceso autorizado, herramientas, participantes y ventana del período.\newline{}\textbf{Condiciones.} Cada instancia identifica etapa, ensayo, grupo, C y versión. C es el cambio de autoridad del grupo, no necesariamente el paso oficial a producción. Dos ensayos completos aprobados antes del primer C; T-18 fija cortes y actualización/repetición si cambian sus condiciones. Revisión de habilitación C menos 2 días hábiles. Referencia: SD5, tabla 5.13; FALTA REFERENCIA A SD9.\\\hline
\texttt{1.\allowbreak{}10.\allowbreak{}1.\allowbreak{}1.\allowbreak{}\{1,\allowbreak{}2\}.\allowbreak{}3}\newline{}\textbf{Ensayo completo E1, ensayos 1 y 2: Conciliación y revisión funcional}\newline{}Davor Serey\newline{}2 instancias.\newline{}\textbf{Ventana:} Antes de C del grupo E1: ensayo 1 en C menos 3 semanas; ensayo 2 en C menos 1 semana. & \textbf{Contenido y frontera.} Resultado parcial del ensayo integral; excluye construir nuevamente las herramientas de migración y repetir la misma ejecución en T-18.\newline{}\textbf{Aceptación.} Conteos, importes, claves, vínculos, expedientes y autorizaciones contrastados; diferencias explicadas y decisiones funcionales registradas.\newline{}\textbf{Insumos.} Versión/registro recibido, acceso autorizado, herramientas, participantes y ventana del período.\newline{}\textbf{Condiciones.} Cada instancia identifica etapa, ensayo, grupo, C y versión. C es el cambio de autoridad del grupo, no necesariamente el paso oficial a producción. Dos ensayos completos aprobados antes del primer C; T-18 fija cortes y actualización/repetición si cambian sus condiciones. Revisión de habilitación C menos 2 días hábiles. Referencia: SD5, tabla 5.13; FALTA REFERENCIA A SD9.\\\hline
\texttt{1.\allowbreak{}10.\allowbreak{}1.\allowbreak{}1.\allowbreak{}\{1,\allowbreak{}2\}.\allowbreak{}4}\newline{}\textbf{Ensayo completo E1, ensayos 1 y 2: Retorno previo y posterior comprobados}\newline{}Davor Serey\newline{}2 instancias.\newline{}\textbf{Ventana:} Antes de C del grupo E1: ensayo 1 en C menos 3 semanas; ensayo 2 en C menos 1 semana. & \textbf{Contenido y frontera.} Resultado parcial del ensayo integral; excluye construir nuevamente las herramientas de migración y repetir la misma ejecución en T-18.\newline{}\textbf{Aceptación.} Retorno antes y después de activar comprobado sin pérdida ni doble efecto; punto de recuperación, delta y tiempos registrados.\newline{}\textbf{Insumos.} Versión/registro recibido, acceso autorizado, herramientas, participantes y ventana del período.\newline{}\textbf{Condiciones.} Cada instancia identifica etapa, ensayo, grupo, C y versión. C es el cambio de autoridad del grupo, no necesariamente el paso oficial a producción. Dos ensayos completos aprobados antes del primer C; T-18 fija cortes y actualización/repetición si cambian sus condiciones. Revisión de habilitación C menos 2 días hábiles. Referencia: SD5, tabla 5.13; FALTA REFERENCIA A SD9.\\\hline
\texttt{1.\allowbreak{}10.\allowbreak{}1.\allowbreak{}1.\allowbreak{}\{1,\allowbreak{}2\}.\allowbreak{}5}\newline{}\textbf{Ensayo completo E1, ensayos 1 y 2: Expediente presentado a recepción}\newline{}Davor Serey\newline{}2 instancias.\newline{}\textbf{Ventana:} Antes de C del grupo E1: ensayo 1 en C menos 3 semanas; ensayo 2 en C menos 1 semana. & \textbf{Contenido y frontera.} Resultado parcial del ensayo integral; excluye construir nuevamente las herramientas de migración y repetir la misma ejecución en T-18.\newline{}\textbf{Aceptación.} Protocolo, ejecución, conciliación, retorno, observaciones y decisión enlazados al ensayo; el segundo incorpora suplentes técnico y funcional.\newline{}\textbf{Insumos.} Versión/registro recibido, acceso autorizado, herramientas, participantes y ventana del período.\newline{}\textbf{Condiciones.} Cada instancia identifica etapa, ensayo, grupo, C y versión. C es el cambio de autoridad del grupo, no necesariamente el paso oficial a producción. Dos ensayos completos aprobados antes del primer C; T-18 fija cortes y actualización/repetición si cambian sus condiciones. Revisión de habilitación C menos 2 días hábiles. Referencia: SD5, tabla 5.13; FALTA REFERENCIA A SD9.\\\hline
\texttt{1.\allowbreak{}10.\allowbreak{}1.\allowbreak{}2.\allowbreak{}\{1,\allowbreak{}2\}.\allowbreak{}1}\newline{}\textbf{Ensayo completo E2, ensayos 1 y 2: PREPROD, origen y protocolo preparados}\newline{}Davor Serey\newline{}2 instancias.\newline{}\textbf{Ventana:} Antes de C del grupo E2: ensayo 1 en C menos 3 semanas; ensayo 2 en C menos 1 semana. & \textbf{Contenido y frontera.} Resultado parcial del ensayo integral; excluye construir nuevamente las herramientas de migración y repetir la misma ejecución en T-18.\newline{}\textbf{Aceptación.} PREPROD aislado, copia protegida, fuentes, versiones, accesos y protocolo de los seis conjuntos identificados; ninguna preparación sustituye el ensayo completo.\newline{}\textbf{Insumos.} Versión/registro recibido, acceso autorizado, herramientas, participantes y ventana del período.\newline{}\textbf{Condiciones.} Cada instancia identifica etapa, ensayo, grupo, C y versión. C es el cambio de autoridad del grupo, no necesariamente el paso oficial a producción. Dos ensayos completos aprobados antes del primer C; T-18 fija cortes y actualización/repetición si cambian sus condiciones. Revisión de habilitación C menos 2 días hábiles. Referencia: SD5, tabla 5.13; FALTA REFERENCIA A SD9.\\\hline
\texttt{1.\allowbreak{}10.\allowbreak{}1.\allowbreak{}2.\allowbreak{}\{1,\allowbreak{}2\}.\allowbreak{}2}\newline{}\textbf{Ensayo completo E2, ensayos 1 y 2: Ejecución, manifiestos y tiempos}\newline{}Davor Serey\newline{}2 instancias.\newline{}\textbf{Ventana:} Antes de C del grupo E2: ensayo 1 en C menos 3 semanas; ensayo 2 en C menos 1 semana. & \textbf{Contenido y frontera.} Resultado parcial del ensayo integral; excluye construir nuevamente las herramientas de migración y repetir la misma ejecución en T-18.\newline{}\textbf{Aceptación.} Extracción, transformación y carga recorridas de extremo a extremo; manifiestos, errores y tiempos medidos sobre las versiones y el delta de la etapa.\newline{}\textbf{Insumos.} Versión/registro recibido, acceso autorizado, herramientas, participantes y ventana del período.\newline{}\textbf{Condiciones.} Cada instancia identifica etapa, ensayo, grupo, C y versión. C es el cambio de autoridad del grupo, no necesariamente el paso oficial a producción. Dos ensayos completos aprobados antes del primer C; T-18 fija cortes y actualización/repetición si cambian sus condiciones. Revisión de habilitación C menos 2 días hábiles. Referencia: SD5, tabla 5.13; FALTA REFERENCIA A SD9.\\\hline
\texttt{1.\allowbreak{}10.\allowbreak{}1.\allowbreak{}2.\allowbreak{}\{1,\allowbreak{}2\}.\allowbreak{}3}\newline{}\textbf{Ensayo completo E2, ensayos 1 y 2: Conciliación y revisión funcional}\newline{}Davor Serey\newline{}2 instancias.\newline{}\textbf{Ventana:} Antes de C del grupo E2: ensayo 1 en C menos 3 semanas; ensayo 2 en C menos 1 semana. & \textbf{Contenido y frontera.} Resultado parcial del ensayo integral; excluye construir nuevamente las herramientas de migración y repetir la misma ejecución en T-18.\newline{}\textbf{Aceptación.} Conteos, importes, claves, vínculos, expedientes y autorizaciones contrastados; diferencias explicadas y decisiones funcionales registradas.\newline{}\textbf{Insumos.} Versión/registro recibido, acceso autorizado, herramientas, participantes y ventana del período.\newline{}\textbf{Condiciones.} Cada instancia identifica etapa, ensayo, grupo, C y versión. C es el cambio de autoridad del grupo, no necesariamente el paso oficial a producción. Dos ensayos completos aprobados antes del primer C; T-18 fija cortes y actualización/repetición si cambian sus condiciones. Revisión de habilitación C menos 2 días hábiles. Referencia: SD5, tabla 5.13; FALTA REFERENCIA A SD9.\\\hline
\texttt{1.\allowbreak{}10.\allowbreak{}1.\allowbreak{}2.\allowbreak{}\{1,\allowbreak{}2\}.\allowbreak{}4}\newline{}\textbf{Ensayo completo E2, ensayos 1 y 2: Retorno previo y posterior comprobados}\newline{}Davor Serey\newline{}2 instancias.\newline{}\textbf{Ventana:} Antes de C del grupo E2: ensayo 1 en C menos 3 semanas; ensayo 2 en C menos 1 semana. & \textbf{Contenido y frontera.} Resultado parcial del ensayo integral; excluye construir nuevamente las herramientas de migración y repetir la misma ejecución en T-18.\newline{}\textbf{Aceptación.} Retorno antes y después de activar comprobado sin pérdida ni doble efecto; punto de recuperación, delta y tiempos registrados.\newline{}\textbf{Insumos.} Versión/registro recibido, acceso autorizado, herramientas, participantes y ventana del período.\newline{}\textbf{Condiciones.} Cada instancia identifica etapa, ensayo, grupo, C y versión. C es el cambio de autoridad del grupo, no necesariamente el paso oficial a producción. Dos ensayos completos aprobados antes del primer C; T-18 fija cortes y actualización/repetición si cambian sus condiciones. Revisión de habilitación C menos 2 días hábiles. Referencia: SD5, tabla 5.13; FALTA REFERENCIA A SD9.\\\hline
\texttt{1.\allowbreak{}10.\allowbreak{}1.\allowbreak{}2.\allowbreak{}\{1,\allowbreak{}2\}.\allowbreak{}5}\newline{}\textbf{Ensayo completo E2, ensayos 1 y 2: Expediente presentado a recepción}\newline{}Davor Serey\newline{}2 instancias.\newline{}\textbf{Ventana:} Antes de C del grupo E2: ensayo 1 en C menos 3 semanas; ensayo 2 en C menos 1 semana. & \textbf{Contenido y frontera.} Resultado parcial del ensayo integral; excluye construir nuevamente las herramientas de migración y repetir la misma ejecución en T-18.\newline{}\textbf{Aceptación.} Protocolo, ejecución, conciliación, retorno, observaciones y decisión enlazados al ensayo; el segundo incorpora suplentes técnico y funcional.\newline{}\textbf{Insumos.} Versión/registro recibido, acceso autorizado, herramientas, participantes y ventana del período.\newline{}\textbf{Condiciones.} Cada instancia identifica etapa, ensayo, grupo, C y versión. C es el cambio de autoridad del grupo, no necesariamente el paso oficial a producción. Dos ensayos completos aprobados antes del primer C; T-18 fija cortes y actualización/repetición si cambian sus condiciones. Revisión de habilitación C menos 2 días hábiles. Referencia: SD5, tabla 5.13; FALTA REFERENCIA A SD9.\\\hline
\texttt{1.\allowbreak{}10.\allowbreak{}2.\allowbreak{}1.\allowbreak{}1}\newline{}\textbf{Protocolo integral y datos de prueba E1}\newline{}Davor Serey\newline{}1 instancia.\newline{}\textbf{Ventana:} 9. & \textbf{Contenido y frontera.} Casos, requisitos, datos, versiones y criterios; excluye unidades ya incluidas en construcción.\newline{}\textbf{Aceptación.} Cada caso y criterio CA tiene escenario, umbral, autoridad y evidencia prevista. Davor Serey responde por el paquete.\newline{}\textbf{Condiciones.} Predecesoras del catálogo: H2. Referencia: BT20.1; SD3 CA; SD4/SD5. Las habilitaciones técnicas se contrastan con arquitectura e inventario vigentes.\\\hline
\texttt{1.\allowbreak{}10.\allowbreak{}2.\allowbreak{}1.\allowbreak{}2}\newline{}\textbf{Puerta de QA, cobertura y contratos E1}\newline{}Davor Serey\newline{}1 instancia.\newline{}\textbf{Ventana:} 10. & \textbf{Contenido y frontera.} Consolidación de cobertura y flujos QA; excluye construir unidades y contratos otra vez.\newline{}\textbf{Aceptación.} Lógica de negocio alcanza 70\,\% con puerta bloqueante y sin pruebas fallidas; integraciones de la etapa ejecutadas sin error. Davor Serey responde por el paquete.\newline{}\textbf{Condiciones.} Predecesoras del catálogo: VERSION-E1. Referencia: BT20.1; SD3 CA; SD4/SD5. Las habilitaciones técnicas se contrastan con arquitectura e inventario vigentes.\\\hline
\texttt{1.\allowbreak{}10.\allowbreak{}2.\allowbreak{}1.\allowbreak{}3}\newline{}\textbf{Regresión de sistema y convivencia E1}\newline{}Davor Serey\newline{}1 instancia.\newline{}\textbf{Ventana:} 11. & \textbf{Contenido y frontera.} Batería completa por versión y protección de E1; excluye regresión específica ya ejecutada de innovación.\newline{}\textbf{Aceptación.} No hay regresión abierta y los flujos/canales conservan permisos, hechos y autoridad. Davor Serey responde por el paquete.\newline{}\textbf{Condiciones.} Predecesoras del catálogo: VERSION-E1. Referencia: BT20.1; SD3 CA; SD4/SD5. Las habilitaciones técnicas se contrastan con arquitectura e inventario vigentes.\\\hline
\texttt{1.\allowbreak{}10.\allowbreak{}2.\allowbreak{}1.\allowbreak{}4}\newline{}\textbf{Carga y estrés con mezcla operacional E1}\newline{}Davor Serey\newline{}1 instancia.\newline{}\textbf{Ventana:} 11. & \textbf{Contenido y frontera.} Preparación y ejecución de carga/estrés, curvas y recuperación; excluye fijar límites por promedios simples.\newline{}\textbf{Aceptación.} Umbrales se cumplen a 1.5 veces peak declarado y se registra saturación y degradación sin pérdida silenciosa. Davor Serey responde por el paquete.\newline{}\textbf{Condiciones.} Predecesoras del catálogo: VERSION-E1. Referencia: BT20.1; SD3 CA; SD4/SD5. Las habilitaciones técnicas se contrastan con arquitectura e inventario vigentes.\\\hline
\texttt{1.\allowbreak{}10.\allowbreak{}2.\allowbreak{}1.\allowbreak{}5}\newline{}\textbf{Continuidad local y terminales desconectados E1}\newline{}Davor Serey\newline{}1 instancia.\newline{}\textbf{Ventana:} 11. & \textbf{Contenido y frontera.} Prueba conjunta de autoridad y reconciliación; excluye medir sólo la cola vacía.\newline{}\textbf{Aceptación.} Recinto opera 24 horas, terminales 12 e inspección 8 con mezcla y fallos previstos, sin pérdida ni duplicados. Davor Serey responde por el paquete.\newline{}\textbf{Condiciones.} Predecesoras del catálogo: VERSION-E1. Referencia: BT20.1; SD3 CA; SD4/SD5. Las habilitaciones técnicas se contrastan con arquitectura e inventario vigentes.\\\hline
\texttt{1.\allowbreak{}10.\allowbreak{}2.\allowbreak{}1.\allowbreak{}6}\newline{}\textbf{Resiliencia de componentes y recuperación E1}\newline{}José Lara Arce\newline{}1 instancia.\newline{}\textbf{Ventana:} 11. & \textbf{Contenido y frontera.} Falla controlada de enlace, nodo y dependencia, con registro y recuperación.\newline{}\textbf{Aceptación.} Se conserva operación/degradación controlada y recuperación, con tiempos, transacciones y brechas documentados. José Lara Arce responde por el paquete.\newline{}\textbf{Condiciones.} Predecesoras del catálogo: VERSION-E1. Referencia: BT20.1; SD3 CA; SD4/SD5. Las habilitaciones técnicas se contrastan con arquitectura e inventario vigentes.\\\hline
\texttt{1.\allowbreak{}10.\allowbreak{}2.\allowbreak{}1.\allowbreak{}7}\newline{}\textbf{Conmutación real de desastre y retorno E1}\newline{}José Lara Arce\newline{}1 instancia.\newline{}\textbf{Ventana:} 11. & \textbf{Contenido y frontera.} Conmutación regional real, verificación de datos/identidad y retorno; excluye usar existencia de réplica como evidencia.\newline{}\textbf{Aceptación.} RTO/RPO comprometidos se alcanzan y evidencia identifica último dato, tiempos, permisos y conciliación. José Lara Arce responde por el paquete.\newline{}\textbf{Condiciones.} Predecesoras del catálogo: VERSION-E1. Referencia: BT20.1; SD3 CA; SD4/SD5. Las habilitaciones técnicas se contrastan con arquitectura e inventario vigentes.\\\hline
\texttt{1.\allowbreak{}10.\allowbreak{}2.\allowbreak{}1.\allowbreak{}8}\newline{}\textbf{Prueba ofensiva y cierre de hallazgos E1}\newline{}Bruno Toro\newline{}1 instancia.\newline{}\textbf{Ventana:} 11. & \textbf{Contenido y frontera.} Ejecución ofensiva de superficie completa y comprobación de resolución; correcciones de código se imputan a construcción.\newline{}\textbf{Aceptación.} No quedan hallazgos críticos ni altos abiertos y cada resultado corresponde a versión recibida. Bruno Toro responde por el paquete.\newline{}\textbf{Condiciones.} Predecesoras del catálogo: VERSION-E1. Referencia: BT20.1; SD3 CA; SD4/SD5. Las habilitaciones técnicas se contrastan con arquitectura e inventario vigentes.\\\hline
\texttt{1.\allowbreak{}10.\allowbreak{}2.\allowbreak{}1.\allowbreak{}9}\newline{}\textbf{Verificación de accesibilidad WCAG 2.2 AA E1}\newline{}Davor Serey\newline{}1 instancia.\newline{}\textbf{Ventana:} 11. & \textbf{Contenido y frontera.} Auditoría automática y manual de tareas/canales; excluye acreditar AA sólo por biblioteca usada.\newline{}\textbf{Aceptación.} Criterios aplicables AA están verificados, barreras corregidas y evidencia de teclado, foco, mensajes y ayudas registrada. Davor Serey responde por el paquete.\newline{}\textbf{Condiciones.} Predecesoras del catálogo: VERSION-E1. Referencia: BT20.1; SD3 CA; SD4/SD5. Las habilitaciones técnicas se contrastan con arquitectura e inventario vigentes.\\\hline
\texttt{1.\allowbreak{}10.\allowbreak{}2.\allowbreak{}1.\allowbreak{}10}\newline{}\textbf{Aceptación de usuario y expediente T-17 E1}\newline{}Davor Serey\newline{}1 instancia.\newline{}\textbf{Ventana:} 12. & \textbf{Contenido y frontera.} Ejecución conjunta de escenarios CA y expediente, sin repetir ensayos de migración 1.10.1.\newline{}\textbf{Aceptación.} Contraparte decide y firma aceptación; observaciones resueltas, versiones y pruebas trazadas antes de certificación. Davor Serey responde por el paquete.\newline{}\textbf{Condiciones.} Predecesoras del catálogo: 1.10.2.1.1,1.10.2.1.2,1.10.2.1.3,1.10.2.1.4,1.10.2.1.5,1.10.2.1.6,1.10.2.1.7,1.10.2.1.8,1.10.2.1.9,ENS-VENT-1-2. Referencia: BT20.1; SD3 CA; SD4/SD5. Las habilitaciones técnicas se contrastan con arquitectura e inventario vigentes.\\\hline
\texttt{1.\allowbreak{}10.\allowbreak{}2.\allowbreak{}2.\allowbreak{}1}\newline{}\textbf{Protocolo integral y datos de prueba E2}\newline{}Davor Serey\newline{}1 instancia.\newline{}\textbf{Ventana:} 15. & \textbf{Contenido y frontera.} Casos, requisitos, datos, versiones y criterios; excluye unidades ya incluidas en construcción.\newline{}\textbf{Aceptación.} Cada caso y criterio CA tiene escenario, umbral, autoridad y evidencia prevista. Davor Serey responde por el paquete.\newline{}\textbf{Condiciones.} Predecesoras del catálogo: H8. Referencia: BT20.1; SD3 CA; SD4/SD5. Las habilitaciones técnicas se contrastan con arquitectura e inventario vigentes.\\\hline
\texttt{1.\allowbreak{}10.\allowbreak{}2.\allowbreak{}2.\allowbreak{}2}\newline{}\textbf{Puerta de QA, cobertura y contratos E2}\newline{}Davor Serey\newline{}1 instancia.\newline{}\textbf{Ventana:} 17. & \textbf{Contenido y frontera.} Consolidación de cobertura y flujos QA; excluye construir unidades y contratos otra vez.\newline{}\textbf{Aceptación.} Lógica de negocio alcanza 70\,\% con puerta bloqueante y sin pruebas fallidas; integraciones de la etapa ejecutadas sin error. Davor Serey responde por el paquete.\newline{}\textbf{Condiciones.} Predecesoras del catálogo: VERSION-E2. Referencia: BT20.1; SD3 CA; SD4/SD5. Las habilitaciones técnicas se contrastan con arquitectura e inventario vigentes.\\\hline
\texttt{1.\allowbreak{}10.\allowbreak{}2.\allowbreak{}2.\allowbreak{}3}\newline{}\textbf{Regresión de sistema y convivencia E2}\newline{}Davor Serey\newline{}1 instancia.\newline{}\textbf{Ventana:} 17. & \textbf{Contenido y frontera.} Batería completa por versión y protección de E1; excluye regresión específica ya ejecutada de innovación.\newline{}\textbf{Aceptación.} No hay regresión abierta y los flujos/canales conservan permisos, hechos y autoridad. Davor Serey responde por el paquete.\newline{}\textbf{Condiciones.} Predecesoras del catálogo: VERSION-E2. Referencia: BT20.1; SD3 CA; SD4/SD5. Las habilitaciones técnicas se contrastan con arquitectura e inventario vigentes.\\\hline
\texttt{1.\allowbreak{}10.\allowbreak{}2.\allowbreak{}2.\allowbreak{}4}\newline{}\textbf{Carga y estrés con mezcla operacional E2}\newline{}Davor Serey\newline{}1 instancia.\newline{}\textbf{Ventana:} 17. & \textbf{Contenido y frontera.} Preparación y ejecución de carga/estrés, curvas y recuperación; excluye fijar límites por promedios simples.\newline{}\textbf{Aceptación.} Umbrales se cumplen a 1.5 veces peak declarado y se registra saturación y degradación sin pérdida silenciosa. Davor Serey responde por el paquete.\newline{}\textbf{Condiciones.} Predecesoras del catálogo: VERSION-E2. Referencia: BT20.1; SD3 CA; SD4/SD5. Las habilitaciones técnicas se contrastan con arquitectura e inventario vigentes.\\\hline
\texttt{1.\allowbreak{}10.\allowbreak{}2.\allowbreak{}2.\allowbreak{}5}\newline{}\textbf{Continuidad local y terminales desconectados E2}\newline{}Davor Serey\newline{}1 instancia.\newline{}\textbf{Ventana:} 17. & \textbf{Contenido y frontera.} Prueba conjunta de autoridad y reconciliación; excluye medir sólo la cola vacía.\newline{}\textbf{Aceptación.} Recinto opera 24 horas, terminales 12 e inspección 8 con mezcla y fallos previstos, sin pérdida ni duplicados. Davor Serey responde por el paquete.\newline{}\textbf{Condiciones.} Predecesoras del catálogo: VERSION-E2. Referencia: BT20.1; SD3 CA; SD4/SD5. Las habilitaciones técnicas se contrastan con arquitectura e inventario vigentes.\\\hline
\texttt{1.\allowbreak{}10.\allowbreak{}2.\allowbreak{}2.\allowbreak{}6}\newline{}\textbf{Resiliencia de componentes y recuperación E2}\newline{}José Lara Arce\newline{}1 instancia.\newline{}\textbf{Ventana:} 17. & \textbf{Contenido y frontera.} Falla controlada de enlace, nodo y dependencia, con registro y recuperación.\newline{}\textbf{Aceptación.} Se conserva operación/degradación controlada y recuperación, con tiempos, transacciones y brechas documentados. José Lara Arce responde por el paquete.\newline{}\textbf{Condiciones.} Predecesoras del catálogo: VERSION-E2. Referencia: BT20.1; SD3 CA; SD4/SD5. Las habilitaciones técnicas se contrastan con arquitectura e inventario vigentes.\\\hline
\texttt{1.\allowbreak{}10.\allowbreak{}2.\allowbreak{}2.\allowbreak{}7}\newline{}\textbf{Conmutación real de desastre y retorno E2}\newline{}José Lara Arce\newline{}1 instancia.\newline{}\textbf{Ventana:} 17. & \textbf{Contenido y frontera.} Conmutación regional real, verificación de datos/identidad y retorno; excluye usar existencia de réplica como evidencia.\newline{}\textbf{Aceptación.} RTO/RPO comprometidos se alcanzan y evidencia identifica último dato, tiempos, permisos y conciliación. José Lara Arce responde por el paquete.\newline{}\textbf{Condiciones.} Predecesoras del catálogo: VERSION-E2. Referencia: BT20.1; SD3 CA; SD4/SD5. Las habilitaciones técnicas se contrastan con arquitectura e inventario vigentes.\\\hline
\texttt{1.\allowbreak{}10.\allowbreak{}2.\allowbreak{}2.\allowbreak{}8}\newline{}\textbf{Prueba ofensiva y cierre de hallazgos E2}\newline{}Bruno Toro\newline{}1 instancia.\newline{}\textbf{Ventana:} 17. & \textbf{Contenido y frontera.} Ejecución ofensiva de superficie completa y comprobación de resolución; correcciones de código se imputan a construcción.\newline{}\textbf{Aceptación.} No quedan hallazgos críticos ni altos abiertos y cada resultado corresponde a versión recibida. Bruno Toro responde por el paquete.\newline{}\textbf{Condiciones.} Predecesoras del catálogo: VERSION-E2. Referencia: BT20.1; SD3 CA; SD4/SD5. Las habilitaciones técnicas se contrastan con arquitectura e inventario vigentes.\\\hline
\texttt{1.\allowbreak{}10.\allowbreak{}2.\allowbreak{}2.\allowbreak{}9}\newline{}\textbf{Verificación de accesibilidad WCAG 2.2 AA E2}\newline{}Davor Serey\newline{}1 instancia.\newline{}\textbf{Ventana:} 17. & \textbf{Contenido y frontera.} Auditoría automática y manual de tareas/canales; excluye acreditar AA sólo por biblioteca usada.\newline{}\textbf{Aceptación.} Criterios aplicables AA están verificados, barreras corregidas y evidencia de teclado, foco, mensajes y ayudas registrada. Davor Serey responde por el paquete.\newline{}\textbf{Condiciones.} Predecesoras del catálogo: VERSION-E2. Referencia: BT20.1; SD3 CA; SD4/SD5. Las habilitaciones técnicas se contrastan con arquitectura e inventario vigentes.\\\hline
\texttt{1.\allowbreak{}10.\allowbreak{}2.\allowbreak{}2.\allowbreak{}10}\newline{}\textbf{Aceptación de usuario y expediente T-17 E2}\newline{}Davor Serey\newline{}1 instancia.\newline{}\textbf{Ventana:} 18. & \textbf{Contenido y frontera.} Ejecución conjunta de escenarios CA y expediente, sin repetir ensayos de migración 1.10.1.\newline{}\textbf{Aceptación.} Contraparte decide y firma aceptación; observaciones resueltas, versiones y pruebas trazadas antes de certificación. Davor Serey responde por el paquete.\newline{}\textbf{Condiciones.} Predecesoras del catálogo: 1.10.2.2.1,1.10.2.2.2,1.10.2.2.3,1.10.2.2.4,1.10.2.2.5,1.10.2.2.6,1.10.2.2.7,1.10.2.2.8,1.10.2.2.9,ENSAYOS-E2. Referencia: BT20.1; SD3 CA; SD4/SD5. Las habilitaciones técnicas se contrastan con arquitectura e inventario vigentes.\\\hline
\texttt{1.\allowbreak{}10.\allowbreak{}3.\allowbreak{}1.\allowbreak{}\{1--4\}}\newline{}\textbf{F27 · Bloques de subsanación y reensayo E1}\newline{}Nikolai Navea\newline{}4 instancias.\newline{}\textbf{Ventana:} Desde el hallazgo del ensayo de la etapa hasta su reensayo aprobado, antes de autorizar C. & \textbf{Contenido y frontera.} Resultado por bloque autorizado de subsanación/reensayo E1: hallazgo identificado, corrección, regresión y evidencia de cierre. Se activa por un hallazgo del ensayo 1 o 2, sin esperar ambos ensayos aprobados para corregir el primero. Reserva no activada no es trabajo ejecutado; corrección y prueba se imputan una vez, sin duplicar construcción ni una ejecución sin cambios.\newline{}\textbf{Aceptación.} Hallazgos del bloque cerrados con cambio/versiones trazables y regresión/reensayo aprobados sobre la versión que se pretende activar. El bloque tiene decisión propia; recibirlo no recibe los otros bloques ni autoriza por sí solo el corte. Trabajo no ejecutado permanece como reserva.\newline{}\textbf{Insumos.} Hallazgo y evidencia identificados del ensayo 1 o 2 de 1.10.1.1, con etapa, grupo, fecha C, versión, responsable y bloque activado; herramientas/datos autorizados y recursos efectivos para corregir y repetir las comprobaciones afectadas.\newline{}\textbf{Condiciones.} Predecesora de activación: hallazgo identificado de una instancia del ensayo 1 o 2 de 1.10.1.1, con su resultado de preparación, ejecución, conciliación, retorno o expediente individualizado; no se requiere la aprobación de ambos ensayos para activar una corrección. Para hallazgos del primer ensayo, corrección y precarga se planifican desde su registro, normalmente en C menos 2 semanas, antes del segundo ensayo en C menos 1 semana. Un hallazgo del segundo ensayo bloquea C hasta corrección y reensayo aprobados; T-18 revalida C y vigencia de evidencia, y T-15 comprueba duración/recursos sin desplazar automáticamente las fases contractuales. FALTA REFERENCIA A SD5; FALTA REFERENCIA A SD9.\\\hline
\texttt{1.\allowbreak{}10.\allowbreak{}3.\allowbreak{}2.\allowbreak{}\{1--4\}}\newline{}\textbf{F28 · Bloques de subsanación y reensayo E2}\newline{}Nikolai Navea\newline{}4 instancias.\newline{}\textbf{Ventana:} Desde el hallazgo del ensayo de la etapa hasta su reensayo aprobado, antes de autorizar C. & \textbf{Contenido y frontera.} Resultado por bloque autorizado de subsanación/reensayo E2: hallazgo identificado, corrección, regresión y evidencia de cierre. Se activa por un hallazgo del ensayo 1 o 2, sin esperar ambos ensayos aprobados para corregir el primero. Reserva no activada no es trabajo ejecutado; corrección y prueba se imputan una vez, sin duplicar construcción ni una ejecución sin cambios.\newline{}\textbf{Aceptación.} Hallazgos del bloque cerrados con cambio/versiones trazables y regresión/reensayo aprobados sobre la versión que se pretende activar. El bloque tiene decisión propia; recibirlo no recibe los otros bloques ni autoriza por sí solo el corte. Trabajo no ejecutado permanece como reserva.\newline{}\textbf{Insumos.} Hallazgo y evidencia identificados del ensayo 1 o 2 de 1.10.1.2, con etapa, grupo, fecha C, versión, responsable y bloque activado; herramientas/datos autorizados y recursos efectivos para corregir y repetir las comprobaciones afectadas.\newline{}\textbf{Condiciones.} Predecesora de activación: hallazgo identificado de una instancia del ensayo 1 o 2 de 1.10.1.2, con su resultado de preparación, ejecución, conciliación, retorno o expediente individualizado; no se requiere la aprobación de ambos ensayos para activar una corrección. Para hallazgos del primer ensayo, corrección y precarga se planifican desde su registro, normalmente en C menos 2 semanas, antes del segundo ensayo en C menos 1 semana. Un hallazgo del segundo ensayo bloquea C hasta corrección y reensayo aprobados; T-18 revalida C y vigencia de evidencia, y T-15 comprueba duración/recursos sin desplazar automáticamente las fases contractuales. FALTA REFERENCIA A SD5; FALTA REFERENCIA A SD9.\\\hline
\end{longtable}\endgroup
\begingroup\setlength{\tabcolsep}{3pt}
\begin{longtable}{L{\dimexpr0.320000\textwidth-2\tabcolsep\relax}L{\dimexpr0.680000\textwidth-2\tabcolsep\relax}}
\caption{Diccionario: Formación e implantación}\label{tab:t14-dic-1-11}\\\hline
\EncabezadoTabla{Paquete y resultado} & \EncabezadoTabla{Frontera, aceptación e insumos}\\\hline\endfirsthead
\multicolumn{2}{l}{\textit{Continuación de la tabla \ref{tab:t14-dic-1-11}}}\\\hline
\EncabezadoTabla{Paquete y resultado} & \EncabezadoTabla{Frontera, aceptación e insumos}\\\hline\endhead
\hline\endfoot
\hline\endlastfoot
\texttt{1.\allowbreak{}11.\allowbreak{}1.\allowbreak{}1.\allowbreak{}1}\newline{}\textbf{Grupo 1 (maestros e inventario): lista de habilitación y formación}\newline{}José Lara Arce\newline{}1 instancia.\newline{}\textbf{Ventana:} 12. & \textbf{Contenido y frontera.} Habilitación específica del grupo de datos; excluye la capacitación general y la recepción oficial anticipada de E1/E2.\newline{}\textbf{Aceptación.} Datos conciliados, permisos, usuarios preparados y ventana autorizada; formación específica del cambio de datos identificada.\newline{}\textbf{Insumos.} Versión/registro recibido, acceso autorizado, herramientas, participantes y ventana del período.\newline{}\textbf{Condiciones.} Predecesoras de habilitación: herramientas, datos y decisiones de la etapa recibidos.\\\hline
\texttt{1.\allowbreak{}11.\allowbreak{}1.\allowbreak{}1.\allowbreak{}2}\newline{}\textbf{Grupo 1 (maestros e inventario): activación y decisión documentadas}\newline{}José Lara Arce\newline{}1 instancia.\newline{}\textbf{Ventana:} 13. & \textbf{Contenido y frontera.} Habilitación específica del grupo de datos; excluye la capacitación general y la recepción oficial anticipada de E1/E2.\newline{}\textbf{Aceptación.} Versión, tiempos, copia recuperable y decisión del autorizador registrados; activación supervisada no adelanta la autoridad oficial.\newline{}\textbf{Insumos.} Versión/registro recibido, acceso autorizado, herramientas, participantes y ventana del período.\newline{}\textbf{Condiciones.} Predecesoras de habilitación: herramientas, datos y decisiones de la etapa recibidos.\\\hline
\texttt{1.\allowbreak{}11.\allowbreak{}1.\allowbreak{}1.\allowbreak{}3}\newline{}\textbf{Grupo 1 (maestros e inventario): conciliación inicial}\newline{}José Lara Arce\newline{}1 instancia.\newline{}\textbf{Ventana:} 13. & \textbf{Contenido y frontera.} Habilitación específica del grupo de datos; excluye la capacitación general y la recepción oficial anticipada de E1/E2.\newline{}\textbf{Aceptación.} Autoridad de escritura, datos y diferencias comprobados; no hay doble digitación ni efectos externos duplicados.\newline{}\textbf{Insumos.} Versión/registro recibido, acceso autorizado, herramientas, participantes y ventana del período.\newline{}\textbf{Condiciones.} Predecesoras de habilitación: herramientas, datos y decisiones de la etapa recibidos.\\\hline
\texttt{1.\allowbreak{}11.\allowbreak{}1.\allowbreak{}2.\allowbreak{}1}\newline{}\textbf{Grupo 2 (antecedentes operacionales): lista de habilitación y formación}\newline{}José Lara Arce\newline{}1 instancia.\newline{}\textbf{Ventana:} 12. & \textbf{Contenido y frontera.} Habilitación específica del grupo de datos; excluye la capacitación general y la recepción oficial anticipada de E1/E2.\newline{}\textbf{Aceptación.} Datos conciliados, permisos, usuarios preparados y ventana autorizada; formación específica del cambio de datos identificada.\newline{}\textbf{Insumos.} Versión/registro recibido, acceso autorizado, herramientas, participantes y ventana del período.\newline{}\textbf{Condiciones.} Predecesoras de habilitación: herramientas, datos y decisiones de la etapa recibidos.\\\hline
\texttt{1.\allowbreak{}11.\allowbreak{}1.\allowbreak{}2.\allowbreak{}2}\newline{}\textbf{Grupo 2 (antecedentes operacionales): activación y decisión documentadas}\newline{}José Lara Arce\newline{}1 instancia.\newline{}\textbf{Ventana:} 13. & \textbf{Contenido y frontera.} Habilitación específica del grupo de datos; excluye la capacitación general y la recepción oficial anticipada de E1/E2.\newline{}\textbf{Aceptación.} Versión, tiempos, copia recuperable y decisión del autorizador registrados; activación supervisada no adelanta la autoridad oficial.\newline{}\textbf{Insumos.} Versión/registro recibido, acceso autorizado, herramientas, participantes y ventana del período.\newline{}\textbf{Condiciones.} Predecesoras de habilitación: herramientas, datos y decisiones de la etapa recibidos.\\\hline
\texttt{1.\allowbreak{}11.\allowbreak{}1.\allowbreak{}2.\allowbreak{}3}\newline{}\textbf{Grupo 2 (antecedentes operacionales): conciliación inicial}\newline{}José Lara Arce\newline{}1 instancia.\newline{}\textbf{Ventana:} 13. & \textbf{Contenido y frontera.} Habilitación específica del grupo de datos; excluye la capacitación general y la recepción oficial anticipada de E1/E2.\newline{}\textbf{Aceptación.} Autoridad de escritura, datos y diferencias comprobados; no hay doble digitación ni efectos externos duplicados.\newline{}\textbf{Insumos.} Versión/registro recibido, acceso autorizado, herramientas, participantes y ventana del período.\newline{}\textbf{Condiciones.} Predecesoras de habilitación: herramientas, datos y decisiones de la etapa recibidos.\\\hline
\texttt{1.\allowbreak{}11.\allowbreak{}1.\allowbreak{}3.\allowbreak{}1}\newline{}\textbf{Grupo 3 (contratos, saldos y sustitución administrativa): lista de habilitación y formación}\newline{}José Lara Arce\newline{}1 instancia.\newline{}\textbf{Ventana:} 18. & \textbf{Contenido y frontera.} Habilitación específica del grupo de datos; excluye la capacitación general y la recepción oficial anticipada de E1/E2.\newline{}\textbf{Aceptación.} Datos conciliados, permisos, usuarios preparados y ventana autorizada; formación específica del cambio de datos identificada.\newline{}\textbf{Insumos.} Versión/registro recibido, acceso autorizado, herramientas, participantes y ventana del período.\newline{}\textbf{Condiciones.} Predecesoras de habilitación: herramientas, datos y decisiones de la etapa recibidos.\\\hline
\texttt{1.\allowbreak{}11.\allowbreak{}1.\allowbreak{}3.\allowbreak{}2}\newline{}\textbf{Grupo 3 (contratos, saldos y sustitución administrativa): activación y decisión documentadas}\newline{}José Lara Arce\newline{}1 instancia.\newline{}\textbf{Ventana:} 19. & \textbf{Contenido y frontera.} Habilitación específica del grupo de datos; excluye la capacitación general y la recepción oficial anticipada de E1/E2.\newline{}\textbf{Aceptación.} Versión, tiempos, copia recuperable y decisión del autorizador registrados; activación supervisada no adelanta la autoridad oficial.\newline{}\textbf{Insumos.} Versión/registro recibido, acceso autorizado, herramientas, participantes y ventana del período.\newline{}\textbf{Condiciones.} Predecesoras de habilitación: herramientas, datos y decisiones de la etapa recibidos.\\\hline
\texttt{1.\allowbreak{}11.\allowbreak{}1.\allowbreak{}3.\allowbreak{}3}\newline{}\textbf{Grupo 3 (contratos, saldos y sustitución administrativa): conciliación inicial}\newline{}José Lara Arce\newline{}1 instancia.\newline{}\textbf{Ventana:} 19. & \textbf{Contenido y frontera.} Habilitación específica del grupo de datos; excluye la capacitación general y la recepción oficial anticipada de E1/E2.\newline{}\textbf{Aceptación.} Autoridad de escritura, datos y diferencias comprobados; no hay doble digitación ni efectos externos duplicados.\newline{}\textbf{Insumos.} Versión/registro recibido, acceso autorizado, herramientas, participantes y ventana del período.\newline{}\textbf{Condiciones.} Predecesoras de habilitación: herramientas, datos y decisiones de la etapa recibidos.\\\hline
\texttt{1.\allowbreak{}11.\allowbreak{}1.\allowbreak{}EST.\allowbreak{}1.\allowbreak{}\{1,\allowbreak{}2\}}\newline{}\textbf{Estabilización de datos E1}\newline{}José Lara Arce\newline{}2 instancias.\newline{}\textbf{Ventana:} 16. & \textbf{Contenido y frontera.} Seguimiento técnico de datos después del paso oficial; excluye los dos puestos de acompañamiento general y la cobertura NOC/SOC/mesa.\newline{}\textbf{Aceptación.} Incidentes, correcciones, estabilidad diaria y pendientes transferidos al servicio; las fechas reales provienen de la ventana posterior al paso oficial.\newline{}\textbf{Insumos.} Versión/registro recibido, acceso autorizado, herramientas, participantes y ventana del período.\newline{}\textbf{Condiciones.} Predecesoras de habilitación: herramientas, datos y decisiones de la etapa recibidos.\\\hline
\texttt{1.\allowbreak{}11.\allowbreak{}1.\allowbreak{}EST.\allowbreak{}2.\allowbreak{}\{1,\allowbreak{}2\}}\newline{}\textbf{Estabilización de datos E2}\newline{}José Lara Arce\newline{}2 instancias.\newline{}\textbf{Ventana:} 21. & \textbf{Contenido y frontera.} Seguimiento técnico de datos después del paso oficial; excluye los dos puestos de acompañamiento general y la cobertura NOC/SOC/mesa.\newline{}\textbf{Aceptación.} Incidentes, correcciones, estabilidad diaria y pendientes transferidos al servicio; las fechas reales provienen de la ventana posterior al paso oficial.\newline{}\textbf{Insumos.} Versión/registro recibido, acceso autorizado, herramientas, participantes y ventana del período.\newline{}\textbf{Condiciones.} Predecesoras de habilitación: herramientas, datos y decisiones de la etapa recibidos.\\\hline
\texttt{1.\allowbreak{}11.\allowbreak{}1.\allowbreak{}MB.\allowbreak{}1.\allowbreak{}\{13--15\}.\allowbreak{}\{1,\allowbreak{}2\}}\newline{}\textbf{Expediente de convivencia E1}\newline{}José Lara Arce\newline{}6 instancias.\newline{}\textbf{Ventana:} 13--15. & \textbf{Contenido y frontera.} Revisión y evidencia de convivencia; excluye guardia continua, acompañamiento presencial de olas y conciliación inicial del grupo.\newline{}\textbf{Aceptación.} Registros diarios con incidencias, denominadores, disponibilidad, respuesta y conciliación; cuatro semanas finales reales incorporadas al cierre, sin reemplazar el acta.\newline{}\textbf{Insumos.} Versión/registro recibido, acceso autorizado, herramientas, participantes y ventana del período.\newline{}\textbf{Condiciones.} Predecesoras de habilitación: herramientas, datos y decisiones de la etapa recibidos.\\\hline
\texttt{1.\allowbreak{}11.\allowbreak{}1.\allowbreak{}MB.\allowbreak{}2.\allowbreak{}\{19,\allowbreak{}20\}.\allowbreak{}\{1,\allowbreak{}2\}}\newline{}\textbf{Expediente de convivencia E2}\newline{}José Lara Arce\newline{}4 instancias.\newline{}\textbf{Ventana:} 19,20. & \textbf{Contenido y frontera.} Revisión y evidencia de convivencia; excluye guardia continua, acompañamiento presencial de olas y conciliación inicial del grupo.\newline{}\textbf{Aceptación.} Registros diarios con incidencias, denominadores, disponibilidad, respuesta y conciliación; cuatro semanas finales reales incorporadas al cierre, sin reemplazar el acta.\newline{}\textbf{Insumos.} Versión/registro recibido, acceso autorizado, herramientas, participantes y ventana del período.\newline{}\textbf{Condiciones.} Predecesoras de habilitación: herramientas, datos y decisiones de la etapa recibidos.\\\hline
\texttt{1.\allowbreak{}11.\allowbreak{}2.\allowbreak{}1.\allowbreak{}1}\newline{}\textbf{Materiales y ejercicios perfil 1 E1}\newline{}José Lara Arce\newline{}1 instancia.\newline{}\textbf{Ventana:} 11. & \textbf{Contenido y frontera.} Manual, guía, FAQ, guion/video y ejercicios editables de su perfil; excluye impartir todas las sesiones y materiales incrementales de innovación.\newline{}\textbf{Aceptación.} Cliente recibe fuentes en español, ejemplos, contingencia y ejercicios; autoformación registra avance y material se vincula a versión. José Lara Arce responde por el paquete.\newline{}\textbf{Condiciones.} Predecesoras del catálogo: VERSION-E1. Referencia: RT-22.01--03. Las habilitaciones técnicas se contrastan con arquitectura e inventario vigentes.\\\hline
\texttt{1.\allowbreak{}11.\allowbreak{}2.\allowbreak{}1.\allowbreak{}2}\newline{}\textbf{Materiales y ejercicios perfil 2 E1}\newline{}José Lara Arce\newline{}1 instancia.\newline{}\textbf{Ventana:} 11. & \textbf{Contenido y frontera.} Manual, guía, FAQ, guion/video y ejercicios editables de su perfil; excluye impartir todas las sesiones y materiales incrementales de innovación.\newline{}\textbf{Aceptación.} Cliente recibe fuentes en español, ejemplos, contingencia y ejercicios; autoformación registra avance y material se vincula a versión. José Lara Arce responde por el paquete.\newline{}\textbf{Condiciones.} Predecesoras del catálogo: VERSION-E1. Referencia: RT-22.01--03. Las habilitaciones técnicas se contrastan con arquitectura e inventario vigentes.\\\hline
\texttt{1.\allowbreak{}11.\allowbreak{}2.\allowbreak{}1.\allowbreak{}3}\newline{}\textbf{Materiales y ejercicios perfil 3 E1}\newline{}José Lara Arce\newline{}1 instancia.\newline{}\textbf{Ventana:} 11. & \textbf{Contenido y frontera.} Manual, guía, FAQ, guion/video y ejercicios editables de su perfil; excluye impartir todas las sesiones y materiales incrementales de innovación.\newline{}\textbf{Aceptación.} Cliente recibe fuentes en español, ejemplos, contingencia y ejercicios; autoformación registra avance y material se vincula a versión. José Lara Arce responde por el paquete.\newline{}\textbf{Condiciones.} Predecesoras del catálogo: VERSION-E1. Referencia: RT-22.01--03. Las habilitaciones técnicas se contrastan con arquitectura e inventario vigentes.\\\hline
\texttt{1.\allowbreak{}11.\allowbreak{}2.\allowbreak{}1.\allowbreak{}4}\newline{}\textbf{Materiales y ejercicios perfil 4 E1}\newline{}José Lara Arce\newline{}1 instancia.\newline{}\textbf{Ventana:} 11. & \textbf{Contenido y frontera.} Manual, guía, FAQ, guion/video y ejercicios editables de su perfil; excluye impartir todas las sesiones y materiales incrementales de innovación.\newline{}\textbf{Aceptación.} Cliente recibe fuentes en español, ejemplos, contingencia y ejercicios; autoformación registra avance y material se vincula a versión. José Lara Arce responde por el paquete.\newline{}\textbf{Condiciones.} Predecesoras del catálogo: VERSION-E1. Referencia: RT-22.01--03. Las habilitaciones técnicas se contrastan con arquitectura e inventario vigentes.\\\hline
\texttt{1.\allowbreak{}11.\allowbreak{}2.\allowbreak{}1.\allowbreak{}5}\newline{}\textbf{Materiales y ejercicios perfil 5 E1}\newline{}José Lara Arce\newline{}1 instancia.\newline{}\textbf{Ventana:} 11. & \textbf{Contenido y frontera.} Manual, guía, FAQ, guion/video y ejercicios editables de su perfil; excluye impartir todas las sesiones y materiales incrementales de innovación.\newline{}\textbf{Aceptación.} Cliente recibe fuentes en español, ejemplos, contingencia y ejercicios; autoformación registra avance y material se vincula a versión. José Lara Arce responde por el paquete.\newline{}\textbf{Condiciones.} Predecesoras del catálogo: VERSION-E1. Referencia: RT-22.01--03. Las habilitaciones técnicas se contrastan con arquitectura e inventario vigentes. Base de conocimiento de soporte editable/versionada, fuentes, responsable, acceso, contenidos y métricas de uso/utilidad identificados; actualización hasta transferencia, luego en 1.11.8.1 y 1.12.14.1. El delta E2 no reconstruye el contenido E1 conservado.\\\hline
\texttt{1.\allowbreak{}11.\allowbreak{}2.\allowbreak{}2.\allowbreak{}1}\newline{}\textbf{Materiales y ejercicios perfil 1 E2}\newline{}José Lara Arce\newline{}1 instancia.\newline{}\textbf{Ventana:} 17. & \textbf{Contenido y frontera.} Manual, guía, FAQ, guion/video y ejercicios editables de su perfil; excluye impartir todas las sesiones y materiales incrementales de innovación.\newline{}\textbf{Aceptación.} Cliente recibe fuentes en español, ejemplos, contingencia y ejercicios; autoformación registra avance y material se vincula a versión. José Lara Arce responde por el paquete.\newline{}\textbf{Condiciones.} Predecesoras del catálogo: VERSION-E2. Referencia: RT-22.01--03. Las habilitaciones técnicas se contrastan con arquitectura e inventario vigentes.\\\hline
\texttt{1.\allowbreak{}11.\allowbreak{}2.\allowbreak{}2.\allowbreak{}2}\newline{}\textbf{Materiales y ejercicios perfil 2 E2}\newline{}José Lara Arce\newline{}1 instancia.\newline{}\textbf{Ventana:} 17. & \textbf{Contenido y frontera.} Manual, guía, FAQ, guion/video y ejercicios editables de su perfil; excluye impartir todas las sesiones y materiales incrementales de innovación.\newline{}\textbf{Aceptación.} Cliente recibe fuentes en español, ejemplos, contingencia y ejercicios; autoformación registra avance y material se vincula a versión. José Lara Arce responde por el paquete.\newline{}\textbf{Condiciones.} Predecesoras del catálogo: VERSION-E2. Referencia: RT-22.01--03. Las habilitaciones técnicas se contrastan con arquitectura e inventario vigentes.\\\hline
\texttt{1.\allowbreak{}11.\allowbreak{}2.\allowbreak{}2.\allowbreak{}3}\newline{}\textbf{Materiales y ejercicios perfil 3 E2}\newline{}José Lara Arce\newline{}1 instancia.\newline{}\textbf{Ventana:} 17. & \textbf{Contenido y frontera.} Manual, guía, FAQ, guion/video y ejercicios editables de su perfil; excluye impartir todas las sesiones y materiales incrementales de innovación.\newline{}\textbf{Aceptación.} Cliente recibe fuentes en español, ejemplos, contingencia y ejercicios; autoformación registra avance y material se vincula a versión. José Lara Arce responde por el paquete.\newline{}\textbf{Condiciones.} Predecesoras del catálogo: VERSION-E2. Referencia: RT-22.01--03. Las habilitaciones técnicas se contrastan con arquitectura e inventario vigentes.\\\hline
\texttt{1.\allowbreak{}11.\allowbreak{}2.\allowbreak{}2.\allowbreak{}4}\newline{}\textbf{Materiales y ejercicios perfil 4 E2}\newline{}José Lara Arce\newline{}1 instancia.\newline{}\textbf{Ventana:} 17. & \textbf{Contenido y frontera.} Manual, guía, FAQ, guion/video y ejercicios editables de su perfil; excluye impartir todas las sesiones y materiales incrementales de innovación.\newline{}\textbf{Aceptación.} Cliente recibe fuentes en español, ejemplos, contingencia y ejercicios; autoformación registra avance y material se vincula a versión. José Lara Arce responde por el paquete.\newline{}\textbf{Condiciones.} Predecesoras del catálogo: VERSION-E2. Referencia: RT-22.01--03. Las habilitaciones técnicas se contrastan con arquitectura e inventario vigentes.\\\hline
\texttt{1.\allowbreak{}11.\allowbreak{}2.\allowbreak{}2.\allowbreak{}5}\newline{}\textbf{Materiales y ejercicios perfil 5 E2}\newline{}José Lara Arce\newline{}1 instancia.\newline{}\textbf{Ventana:} 17. & \textbf{Contenido y frontera.} Manual, guía, FAQ, guion/video y ejercicios editables de su perfil; excluye impartir todas las sesiones y materiales incrementales de innovación.\newline{}\textbf{Aceptación.} Cliente recibe fuentes en español, ejemplos, contingencia y ejercicios; autoformación registra avance y material se vincula a versión. José Lara Arce responde por el paquete.\newline{}\textbf{Condiciones.} Predecesoras del catálogo: VERSION-E2. Referencia: RT-22.01--03. Las habilitaciones técnicas se contrastan con arquitectura e inventario vigentes. Base de conocimiento de soporte editable/versionada, fuentes, responsable, acceso, contenidos y métricas de uso/utilidad identificados; actualización hasta transferencia, luego en 1.11.8.1 y 1.12.14.1. El delta E2 no reconstruye el contenido E1 conservado.\\\hline
\texttt{1.\allowbreak{}11.\allowbreak{}3.\allowbreak{}1.\allowbreak{}1}\newline{}\textbf{Plan de formación y turnos E1}\newline{}José Lara Arce\newline{}1 instancia.\newline{}\textbf{Ventana:} 11. & \textbf{Contenido y frontera.} Manifiesto de formación E1: población por los cinco perfiles, turnos, sitios, grupos compatibles, códigos activados, capacidad, sesiones y calendario. Actualiza el dimensionamiento preliminar de H1 o H8 y enlaza el plan detallado de T-18; excluye duplicar su desarrollo o contar participación del cliente como personal de Synaptix.\newline{}\textbf{Aceptación.} Por perfil/etapa constan fuente y versión de población, particiones compatibles, cantidad calculada, códigos activados y lista de personas a cubrir; sin doble cuenta de asistentes dentro del mismo perfil y objetivo. Toda persona objetivo pertenece a un grupo admisible o tiene una recuperación/reprogramación identificada. José Lara Arce verifica y presenta el manifiesto; la contraparte valida población y calendario. No se cierra la cantidad ni la estimación de esa formación sin ese control. Congelamiento, turnos y capacitación estacional se respetan.\newline{}\textbf{Condiciones.} La población preliminar se solicita al constituir H1, se actualiza para E2 en H8 y se concilia antes de fijar cantidades/HH de formación en T-15. El plan de ejecución se recibe antes de activar cohortes; el dato faltante queda bloqueante y no se reemplaza por los sufijos disponibles. Predecesoras de ejecución: materiales y versión de la etapa recibidos. FALTA REFERENCIA A SD3; FALTA REFERENCIA A SD10.\\\hline
\texttt{1.\allowbreak{}11.\allowbreak{}3.\allowbreak{}1.\allowbreak{}2}\newline{}\textbf{Expediente de evaluación y certificación E1}\newline{}José Lara Arce\newline{}1 instancia.\newline{}\textbf{Ventana:} 15. & \textbf{Contenido y frontera.} Resultados, brechas y decisión para administradores/equipo técnico y soporte; excluye impartir sesiones de recuperación no dimensionadas.\newline{}\textbf{Aceptación.} Roles administradores/técnicos designados son evaluados y certificados antes del cierre; ausencia de área TI no elimina los responsables que recibirán transferencia. José Lara Arce responde por el paquete.\newline{}\textbf{Condiciones.} Predecesoras del catálogo: FORMACION-E1. Referencia: RT-22.05; BA17.3. Las habilitaciones técnicas se contrastan con arquitectura e inventario vigentes.\\\hline
\texttt{1.\allowbreak{}11.\allowbreak{}3.\allowbreak{}2.\allowbreak{}1}\newline{}\textbf{Plan de formación y turnos E2}\newline{}José Lara Arce\newline{}1 instancia.\newline{}\textbf{Ventana:} 17. & \textbf{Contenido y frontera.} Manifiesto de formación E2: población por los cinco perfiles, turnos, sitios, grupos compatibles, códigos activados, capacidad, sesiones y calendario. Actualiza el dimensionamiento preliminar de H1 o H8 y enlaza el plan detallado de T-18; excluye duplicar su desarrollo o contar participación del cliente como personal de Synaptix.\newline{}\textbf{Aceptación.} Por perfil/etapa constan fuente y versión de población, particiones compatibles, cantidad calculada, códigos activados y lista de personas a cubrir; sin doble cuenta de asistentes dentro del mismo perfil y objetivo. Toda persona objetivo pertenece a un grupo admisible o tiene una recuperación/reprogramación identificada. José Lara Arce verifica y presenta el manifiesto; la contraparte valida población y calendario. No se cierra la cantidad ni la estimación de esa formación sin ese control. Congelamiento, turnos y capacitación estacional se respetan.\newline{}\textbf{Condiciones.} La población preliminar se solicita al constituir H1, se actualiza para E2 en H8 y se concilia antes de fijar cantidades/HH de formación en T-15. El plan de ejecución se recibe antes de activar cohortes; el dato faltante queda bloqueante y no se reemplaza por los sufijos disponibles. Predecesoras de ejecución: materiales y versión de la etapa recibidos. FALTA REFERENCIA A SD3; FALTA REFERENCIA A SD10.\\\hline
\texttt{1.\allowbreak{}11.\allowbreak{}3.\allowbreak{}2.\allowbreak{}2}\newline{}\textbf{Expediente de evaluación y certificación E2}\newline{}José Lara Arce\newline{}1 instancia.\newline{}\textbf{Ventana:} 20. & \textbf{Contenido y frontera.} Resultados, brechas y decisión para administradores/equipo técnico y soporte; excluye impartir sesiones de recuperación no dimensionadas.\newline{}\textbf{Aceptación.} Roles administradores/técnicos designados son evaluados y certificados antes del cierre; ausencia de área TI no elimina los responsables que recibirán transferencia. José Lara Arce responde por el paquete.\newline{}\textbf{Condiciones.} Predecesoras del catálogo: FORMACION-E2. Referencia: RT-22.05; BA17.3. Las habilitaciones técnicas se contrastan con arquitectura e inventario vigentes.\\\hline
\texttt{1.\allowbreak{}11.\allowbreak{}4.\allowbreak{}1}\newline{}\textbf{Acompañamiento de implantación E1 recibido}\newline{}José Lara Arce\newline{}1 paquete por etapa\newline{}\textbf{Ventana:} 13--15. & \textbf{Contenido y frontera.} Acompañamiento presencial de las olas activadas de E1, consolidado en un único resultado de etapa. Comprende atención en puesto, preparación, relevo, seguimiento y cierre por ola/sitio. Los registros diarios son evidencia del paquete, no entregables terminales adicionales. Excluye turnos de mesa/NOC/SOC y estabilización posterior al paso oficial.\newline{}\textbf{Aceptación.} Manifiesto de olas recibido identifica objetivo, capacidades, grupo de datos, población, sitio, zona, fecha relativa de activación y criterio de avance. Por cada ola se conserva presencia, preparación, relevos, incidencias y decisión de cierre. Se verifica curva decreciente de presencia: diseño de 28 días por ola, 8 horas presenciales diarias los primeros 14 y 4 horas presenciales más 4 de preparación/seguimiento los siguientes 14. Pantalanes y varadero son los dos frentes de referencia; cada ola identifica los frentes efectivamente afectados y justifica la cobertura de otros sitios. No se acepta una ola por el solo registro de asistencia ni la etapa por cerrar una sola ola. Cambios de cobertura pasan por control de cambios y no reducen exigencias contractuales.\newline{}\textbf{Insumos.} Formación y versión E1 recibidas; manifiesto de olas de T-18, ventanas autorizadas, capacidad de ejecución/relevo y accesos. Los dos grupos de datos E1 y el grupo E2 no equivalen automáticamente a una misma cantidad de olas.\newline{}\textbf{Condiciones.} FORMACION-E1 precede al acompañamiento de las personas de su ola. T-18 debe identificar y justificar cantidad, orden y correspondencia de olas con 1.11.1 antes de cerrar esfuerzo en T-15; cuatro sufijos históricos no fijan cuatro ejecuciones. Las olas pueden solaparse solo con capacidad comprobada y deben caber en la ventana de implantación; no se presumen simultáneas ni consecutivas. RT-20.02/05; FALTA REFERENCIA A SD7.\\\hline
\texttt{1.\allowbreak{}11.\allowbreak{}4.\allowbreak{}2}\newline{}\textbf{Acompañamiento de implantación E2 recibido}\newline{}José Lara Arce\newline{}1 paquete por etapa\newline{}\textbf{Ventana:} 19--20. & \textbf{Contenido y frontera.} Acompañamiento presencial de las olas activadas de E2, consolidado en un único resultado de etapa. Comprende atención en puesto, preparación, relevo, seguimiento y cierre por ola/sitio. Los registros diarios son evidencia del paquete, no entregables terminales adicionales. Excluye turnos de mesa/NOC/SOC y estabilización posterior al paso oficial.\newline{}\textbf{Aceptación.} Manifiesto de olas recibido identifica objetivo, capacidades, grupo de datos, población, sitio, zona, fecha relativa de activación y criterio de avance. Por cada ola se conserva presencia, preparación, relevos, incidencias y decisión de cierre. Se verifica curva decreciente de presencia: diseño de 28 días por ola, 8 horas presenciales diarias los primeros 14 y 4 horas presenciales más 4 de preparación/seguimiento los siguientes 14. Pantalanes y varadero son los dos frentes de referencia; cada ola identifica los frentes efectivamente afectados y justifica la cobertura de otros sitios. No se acepta una ola por el solo registro de asistencia ni la etapa por cerrar una sola ola. Cambios de cobertura pasan por control de cambios y no reducen exigencias contractuales.\newline{}\textbf{Insumos.} Formación y versión E2 recibidas; manifiesto de olas de T-18, ventanas autorizadas, capacidad de ejecución/relevo y accesos. Los dos grupos de datos E1 y el grupo E2 no equivalen automáticamente a una misma cantidad de olas.\newline{}\textbf{Condiciones.} FORMACION-E2 precede al acompañamiento de las personas de su ola. T-18 debe identificar y justificar cantidad, orden y correspondencia de olas con 1.11.1 antes de cerrar esfuerzo en T-15; cuatro sufijos históricos no fijan cuatro ejecuciones. Las olas pueden solaparse solo con capacidad comprobada y deben caber en la ventana de implantación; no se presumen simultáneas ni consecutivas. RT-20.02/05; FALTA REFERENCIA A SD7.\\\hline
\texttt{1.\allowbreak{}11.\allowbreak{}5.\allowbreak{}1.\allowbreak{}1}\newline{}\textbf{Expediente de cierre E1 preparado}\newline{}Ignacio Silva\newline{}1 instancia.\newline{}\textbf{Ventana:} 15. & \textbf{Contenido y frontera.} Expediente de E1 preparado para consolidar la marcha blanca: estructura de evidencia, responsables, indicadores, conciliación, formación y observaciones. Integra referencias a pruebas y registros propietarios, sin repetir su ejecución ni anticipar el acta final.\newline{}\textbf{Aceptación.} Estructura, fuentes, versiones, responsables y pendientes completos; la evidencia todavía en producción tiene fecha requerida y dueño. La preparación es recibible como resultado documental y no requiere que ya concluyan las cuatro últimas semanas ni la certificación final. No declara cerrada la marcha blanca ni contiene una aceptación contractual anticipada.\newline{}\textbf{Insumos.} MB-INI-E1, protocolo de recepción, indicadores y registros disponibles de la etapa; plan de certificación y revisión de observaciones.\newline{}\textbf{Condiciones.} Predecesora de preparación: MB-INI-E1. Certificación 1.11.3.1.2 y evidencia final son condiciones del cierre en 1.11.5.1.2, no predecesoras que deban existir antes de preparar este expediente. BA17.3; RT-20.04; FALTA REFERENCIA A SD9.\\\hline
\texttt{1.\allowbreak{}11.\allowbreak{}5.\allowbreak{}1.\allowbreak{}2}\newline{}\textbf{Expediente validado y acta de cierre E1}\newline{}Davor Serey\newline{}1 instancia.\newline{}\textbf{Ventana:} 16, antes de CORTE-E1. & \textbf{Contenido y frontera.} Validación final del expediente E1, resolución de observaciones y presentación de acta a la contraparte. Reúne evidencia completa y certificación; la decisión corresponde a la contraparte autorizada. No vuelve a ejecutar ni estimar las pruebas o informes ya producidos.\newline{}\textbf{Aceptación.} Se cumplen conjuntamente BA17.3: últimas cuatro semanas del período con volumen real y disponibilidad/tiempos comprometidos, conciliación sin diferencias no explicadas, personal capacitado/certificado y ningún incidente crítico/alto abierto atribuible a la solución. Evidencias finales tienen versión y revisión. La contraparte firma después de completar la evidencia, con duración positiva para revisión y decisión antes del paso oficial; un borrador o aprobación interna no recibe la marcha blanca.\newline{}\textbf{Insumos.} 1.11.5.1.1 preparado; 1.11.3.1.2 recibido; registros finales de convivencia, resultados de pruebas y revisión de pendientes. La certificación y evidencia final se reciben antes de la decisión.\newline{}\textbf{Condiciones.} Predecesoras: 1.11.5.1.1,1.11.3.1.2 y evidencia completa de las cuatro últimas semanas de marcha blanca. Precede CORTE-E1; no exige CORTE-E1 como insumo ni autoriza paso anticipado. T-15/T-18 deben demostrar revisión/decisión con duración positiva en la ventana admisible. FALTA REFERENCIA A SD9; FALTA REFERENCIA A SD7.\\\hline
\texttt{1.\allowbreak{}11.\allowbreak{}5.\allowbreak{}2.\allowbreak{}1}\newline{}\textbf{Expediente de cierre E2 preparado}\newline{}Ignacio Silva\newline{}1 instancia.\newline{}\textbf{Ventana:} 19. & \textbf{Contenido y frontera.} Expediente de E2 preparado para consolidar la marcha blanca: estructura de evidencia, responsables, indicadores, conciliación, formación y observaciones. Integra referencias a pruebas y registros propietarios, sin repetir su ejecución ni anticipar el acta final.\newline{}\textbf{Aceptación.} Estructura, fuentes, versiones, responsables y pendientes completos; la evidencia todavía en producción tiene fecha requerida y dueño. La preparación es recibible como resultado documental y no requiere que ya concluyan las cuatro últimas semanas ni la certificación final. No declara cerrada la marcha blanca ni contiene una aceptación contractual anticipada.\newline{}\textbf{Insumos.} MB-INI-E2, protocolo de recepción, indicadores y registros disponibles de la etapa; plan de certificación y revisión de observaciones.\newline{}\textbf{Condiciones.} Predecesora de preparación: MB-INI-E2. Certificación 1.11.3.2.2 y evidencia final son condiciones del cierre en 1.11.5.2.2, no predecesoras que deban existir antes de preparar este expediente. BA17.3; RT-20.04; FALTA REFERENCIA A SD9.\\\hline
\texttt{1.\allowbreak{}11.\allowbreak{}5.\allowbreak{}2.\allowbreak{}2}\newline{}\textbf{Expediente validado y acta de cierre E2}\newline{}Davor Serey\newline{}1 instancia.\newline{}\textbf{Ventana:} 20, antes de CORTE-E2. & \textbf{Contenido y frontera.} Validación final del expediente E2, resolución de observaciones y presentación de acta a la contraparte. Reúne evidencia completa y certificación; la decisión corresponde a la contraparte autorizada. No vuelve a ejecutar ni estimar las pruebas o informes ya producidos.\newline{}\textbf{Aceptación.} Se cumplen conjuntamente BA17.3: últimas cuatro semanas del período con volumen real y disponibilidad/tiempos comprometidos, conciliación sin diferencias no explicadas, personal capacitado/certificado y ningún incidente crítico/alto abierto atribuible a la solución. Evidencias finales tienen versión y revisión. La contraparte firma después de completar la evidencia, con duración positiva para revisión y decisión antes del paso oficial; un borrador o aprobación interna no recibe la marcha blanca.\newline{}\textbf{Insumos.} 1.11.5.2.1 preparado; 1.11.3.2.2 recibido; registros finales de convivencia, resultados de pruebas y revisión de pendientes. La certificación y evidencia final se reciben antes de la decisión.\newline{}\textbf{Condiciones.} Predecesoras: 1.11.5.2.1,1.11.3.2.2 y evidencia completa de las cuatro últimas semanas de marcha blanca. Precede CORTE-E2; no exige CORTE-E2 como insumo ni autoriza paso anticipado. T-15/T-18 deben demostrar revisión/decisión con duración positiva en la ventana admisible. FALTA REFERENCIA A SD9; FALTA REFERENCIA A SD7.\\\hline
\texttt{1.\allowbreak{}11.\allowbreak{}6.\allowbreak{}1}\newline{}\textbf{Estabilización presencial E1 recibida}\newline{}José Lara Arce\newline{}1 paquete por etapa\newline{}\textbf{Ventana:} 16. & \textbf{Contenido y frontera.} Estabilización presencial posterior al paso oficial de E1. Resultado agregado de 28 días consecutivos con dos puestos diarios de ocho horas, uno en pantalanes y otro en varadero, incluidos fines de semana y relevos. Los 56 registros de día/zona son evidencia operativa; excluye acompañamiento de olas, estabilización técnica de datos y turnos de servicio continuo.\newline{}\textbf{Aceptación.} Los 28 días y ambos puestos tienen cobertura y relevo verificables; se registran atención, incidencias, resolución/escalamiento, pendientes, responsables y transferencia. El cierre integra todos los días y zonas, con revisión independiente; no basta recibir el primer turno. Dos puestos no implican dos personas sin descanso. Cobertura y duración conservan el diseño existente: agrupar no disminuye las horas ni la atención comprometida.\newline{}\textbf{Insumos.} CORTE-E1 recibido, runbooks, accesos, calendario natural, personas y suplentes disponibles; incidentes y pendientes de marcha blanca. T-15 determina recursos efectivos y relevo.\newline{}\textbf{Condiciones.} Comienza después de CORTE-E1; MARCHA-E1 no equivale a paso oficial. RT-20.06; E1 se imputa a implementación y E2 a Operación, una sola vez. FALTA REFERENCIA A SD7; FALTA REFERENCIA A SD10.\\\hline
\texttt{1.\allowbreak{}11.\allowbreak{}6.\allowbreak{}2}\newline{}\textbf{Estabilización presencial E2 recibida}\newline{}José Lara Arce\newline{}1 paquete por etapa\newline{}\textbf{Ventana:} 21. & \textbf{Contenido y frontera.} Estabilización presencial posterior al paso oficial de E2. Resultado agregado de 28 días consecutivos con dos puestos diarios de ocho horas, uno en pantalanes y otro en varadero, incluidos fines de semana y relevos. Los 56 registros de día/zona son evidencia operativa; excluye acompañamiento de olas, estabilización técnica de datos y turnos de servicio continuo.\newline{}\textbf{Aceptación.} Los 28 días y ambos puestos tienen cobertura y relevo verificables; se registran atención, incidencias, resolución/escalamiento, pendientes, responsables y transferencia. El cierre integra todos los días y zonas, con revisión independiente; no basta recibir el primer turno. Dos puestos no implican dos personas sin descanso. Cobertura y duración conservan el diseño existente: agrupar no disminuye las horas ni la atención comprometida.\newline{}\textbf{Insumos.} CORTE-E2 recibido, runbooks, accesos, calendario natural, personas y suplentes disponibles; incidentes y pendientes de marcha blanca. T-15 determina recursos efectivos y relevo.\newline{}\textbf{Condiciones.} Comienza después de CORTE-E2; MARCHA-E2 no equivale a paso oficial. RT-20.06; E1 se imputa a implementación y E2 a Operación, una sola vez. FALTA REFERENCIA A SD7; FALTA REFERENCIA A SD10.\\\hline
\texttt{1.\allowbreak{}11.\allowbreak{}7.\allowbreak{}1.\allowbreak{}1.\allowbreak{}1.\allowbreak{}\{1--100\}}\newline{}\textbf{F06 · Formación y certificación E1, perfil 1, cohortes codificadas}\newline{}José Lara Arce\newline{}100 identificadores reservados; cantidad activa según manifiesto.\newline{}\textbf{Ventana:} 12. & \textbf{Contenido y frontera.} Formación de una cohorte activada del perfil 1, etapa E1, conforme al manifiesto aprobado. Se registran asistentes, turno, sitio, modalidad, sesiones, ejercicio, evaluación y recuperación. Identificador reservado sin activación no constituye trabajo, costo, avance ni recepción.\newline{}\textbf{Aceptación.} Dos sesiones de dos horas, ejercicio, registro individual y recuperación; máximo doce asistentes y sin interrumpir turnos.\newline{}\textbf{Insumos.} Manifiesto de cantidad y población de 1.11.3.1.1, materiales 1.11.2.1.1, versión recibida y disponibilidad acordada de participantes e instructor. Los datos privados de asistentes permanecen en el registro autorizado, no en este formulario.\newline{}\textbf{Condiciones.} Predecesoras: 1.11.3.1.1,1.11.2.1.1,VERSION-E1. Códigos y cantidades siguen la tabla de control de cohortes; no se presume que todos los sufijos estén activos. Si la demanda supera la reserva, se amplía el registro sin limitar la formación comprometida. FALTA REFERENCIA A SD3; FALTA REFERENCIA A SD10.\\\hline
\texttt{1.\allowbreak{}11.\allowbreak{}7.\allowbreak{}1.\allowbreak{}2.\allowbreak{}1.\allowbreak{}\{1,\allowbreak{}2\}}\newline{}\textbf{F07 · Formación y certificación E1, perfil 2, cohortes codificadas}\newline{}José Lara Arce\newline{}2 identificadores reservados; cantidad activa según manifiesto.\newline{}\textbf{Ventana:} 12. & \textbf{Contenido y frontera.} Formación de una cohorte activada del perfil 2, etapa E1, conforme al manifiesto aprobado. Se registran asistentes, turno, sitio, modalidad, sesiones, ejercicio, evaluación y recuperación. Identificador reservado sin activación no constituye trabajo, costo, avance ni recepción.\newline{}\textbf{Aceptación.} Dos sesiones de dos horas, ejercicio, registro individual y recuperación; máximo doce asistentes y sin interrumpir turnos.\newline{}\textbf{Insumos.} Manifiesto de cantidad y población de 1.11.3.1.1, materiales 1.11.2.1.2, versión recibida y disponibilidad acordada de participantes e instructor. Los datos privados de asistentes permanecen en el registro autorizado, no en este formulario.\newline{}\textbf{Condiciones.} Predecesoras: 1.11.3.1.1,1.11.2.1.2,VERSION-E1. Códigos y cantidades siguen la tabla de control de cohortes; no se presume que todos los sufijos estén activos. Si la demanda supera la reserva, se amplía el registro sin limitar la formación comprometida. FALTA REFERENCIA A SD3; FALTA REFERENCIA A SD10.\\\hline
\texttt{1.\allowbreak{}11.\allowbreak{}7.\allowbreak{}1.\allowbreak{}3.\allowbreak{}1.\allowbreak{}1}\newline{}\textbf{F08 · Formación y certificación E1, perfil 3, cohortes codificadas}\newline{}José Lara Arce\newline{}1 identificadores reservados; cantidad activa según manifiesto.\newline{}\textbf{Ventana:} 12. & \textbf{Contenido y frontera.} Formación de una cohorte activada del perfil 3, etapa E1, conforme al manifiesto aprobado. Se registran asistentes, turno, sitio, modalidad, sesiones, ejercicio, evaluación y recuperación. Identificador reservado sin activación no constituye trabajo, costo, avance ni recepción.\newline{}\textbf{Aceptación.} Dos sesiones de dos horas, ejercicio, registro individual y recuperación; máximo doce asistentes y sin interrumpir turnos.\newline{}\textbf{Insumos.} Manifiesto de cantidad y población de 1.11.3.1.1, materiales 1.11.2.1.3, versión recibida y disponibilidad acordada de participantes e instructor. Los datos privados de asistentes permanecen en el registro autorizado, no en este formulario.\newline{}\textbf{Condiciones.} Predecesoras: 1.11.3.1.1,1.11.2.1.3,VERSION-E1. Códigos y cantidades siguen la tabla de control de cohortes; no se presume que todos los sufijos estén activos. Si la demanda supera la reserva, se amplía el registro sin limitar la formación comprometida. FALTA REFERENCIA A SD3; FALTA REFERENCIA A SD10.\\\hline
\texttt{1.\allowbreak{}11.\allowbreak{}7.\allowbreak{}1.\allowbreak{}4.\allowbreak{}1.\allowbreak{}1}\newline{}\textbf{F09 · Formación y certificación E1, perfil 4, cohortes codificadas}\newline{}José Lara Arce\newline{}1 identificadores reservados; cantidad activa según manifiesto.\newline{}\textbf{Ventana:} 12. & \textbf{Contenido y frontera.} Formación de una cohorte activada del perfil 4, etapa E1, conforme al manifiesto aprobado. Se registran asistentes, turno, sitio, modalidad, sesiones, ejercicio, evaluación y recuperación. Identificador reservado sin activación no constituye trabajo, costo, avance ni recepción.\newline{}\textbf{Aceptación.} Dos sesiones de dos horas, ejercicio, registro individual y recuperación; máximo doce asistentes y sin interrumpir turnos.\newline{}\textbf{Insumos.} Manifiesto de cantidad y población de 1.11.3.1.1, materiales 1.11.2.1.4, versión recibida y disponibilidad acordada de participantes e instructor. Los datos privados de asistentes permanecen en el registro autorizado, no en este formulario.\newline{}\textbf{Condiciones.} Predecesoras: 1.11.3.1.1,1.11.2.1.4,VERSION-E1. Códigos y cantidades siguen la tabla de control de cohortes; no se presume que todos los sufijos estén activos. Si la demanda supera la reserva, se amplía el registro sin limitar la formación comprometida. FALTA REFERENCIA A SD3; FALTA REFERENCIA A SD10.\\\hline
\texttt{1.\allowbreak{}11.\allowbreak{}7.\allowbreak{}1.\allowbreak{}5.\allowbreak{}1.\allowbreak{}\{1--3\}}\newline{}\textbf{F10 · Formación y certificación E1, perfil 5, cohortes codificadas}\newline{}José Lara Arce\newline{}3 identificadores reservados; cantidad activa según manifiesto.\newline{}\textbf{Ventana:} 12. & \textbf{Contenido y frontera.} Formación de una cohorte activada del perfil 5, etapa E1, conforme al manifiesto aprobado. Se registran asistentes, turno, sitio, modalidad, sesiones, ejercicio, evaluación y recuperación. Identificador reservado sin activación no constituye trabajo, costo, avance ni recepción.\newline{}\textbf{Aceptación.} Dos sesiones de dos horas, ejercicio, registro individual y recuperación; máximo doce asistentes y sin interrumpir turnos.\newline{}\textbf{Insumos.} Manifiesto de cantidad y población de 1.11.3.1.1, materiales 1.11.2.1.5, versión recibida y disponibilidad acordada de participantes e instructor. Los datos privados de asistentes permanecen en el registro autorizado, no en este formulario.\newline{}\textbf{Condiciones.} Predecesoras: 1.11.3.1.1,1.11.2.1.5,VERSION-E1. Códigos y cantidades siguen la tabla de control de cohortes; no se presume que todos los sufijos estén activos. Si la demanda supera la reserva, se amplía el registro sin limitar la formación comprometida. FALTA REFERENCIA A SD3; FALTA REFERENCIA A SD10.\\\hline
\texttt{1.\allowbreak{}11.\allowbreak{}7.\allowbreak{}2.\allowbreak{}1.\allowbreak{}1.\allowbreak{}\{1--100\}}\newline{}\textbf{F13 · Formación y certificación E2, perfil 1, cohortes codificadas}\newline{}José Lara Arce\newline{}100 identificadores reservados; cantidad activa según manifiesto.\newline{}\textbf{Ventana:} 18. & \textbf{Contenido y frontera.} Formación de una cohorte activada del perfil 1, etapa E2, conforme al manifiesto aprobado. Se registran asistentes, turno, sitio, modalidad, sesiones, ejercicio, evaluación y recuperación. Identificador reservado sin activación no constituye trabajo, costo, avance ni recepción.\newline{}\textbf{Aceptación.} Dos sesiones de dos horas, ejercicio, registro individual y recuperación; máximo doce asistentes y sin interrumpir turnos.\newline{}\textbf{Insumos.} Manifiesto de cantidad y población de 1.11.3.2.1, materiales 1.11.2.2.1, versión recibida y disponibilidad acordada de participantes e instructor. Los datos privados de asistentes permanecen en el registro autorizado, no en este formulario.\newline{}\textbf{Condiciones.} Predecesoras: 1.11.3.2.1,1.11.2.2.1,VERSION-E2. Códigos y cantidades siguen la tabla de control de cohortes; no se presume que todos los sufijos estén activos. Si la demanda supera la reserva, se amplía el registro sin limitar la formación comprometida. FALTA REFERENCIA A SD3; FALTA REFERENCIA A SD10.\\\hline
\texttt{1.\allowbreak{}11.\allowbreak{}7.\allowbreak{}2.\allowbreak{}2.\allowbreak{}1.\allowbreak{}\{1,\allowbreak{}2\}}\newline{}\textbf{F14 · Formación y certificación E2, perfil 2, cohortes codificadas}\newline{}José Lara Arce\newline{}2 identificadores reservados; cantidad activa según manifiesto.\newline{}\textbf{Ventana:} 18. & \textbf{Contenido y frontera.} Formación de una cohorte activada del perfil 2, etapa E2, conforme al manifiesto aprobado. Se registran asistentes, turno, sitio, modalidad, sesiones, ejercicio, evaluación y recuperación. Identificador reservado sin activación no constituye trabajo, costo, avance ni recepción.\newline{}\textbf{Aceptación.} Dos sesiones de dos horas, ejercicio, registro individual y recuperación; máximo doce asistentes y sin interrumpir turnos.\newline{}\textbf{Insumos.} Manifiesto de cantidad y población de 1.11.3.2.1, materiales 1.11.2.2.2, versión recibida y disponibilidad acordada de participantes e instructor. Los datos privados de asistentes permanecen en el registro autorizado, no en este formulario.\newline{}\textbf{Condiciones.} Predecesoras: 1.11.3.2.1,1.11.2.2.2,VERSION-E2. Códigos y cantidades siguen la tabla de control de cohortes; no se presume que todos los sufijos estén activos. Si la demanda supera la reserva, se amplía el registro sin limitar la formación comprometida. FALTA REFERENCIA A SD3; FALTA REFERENCIA A SD10.\\\hline
\texttt{1.\allowbreak{}11.\allowbreak{}7.\allowbreak{}2.\allowbreak{}3.\allowbreak{}1.\allowbreak{}1}\newline{}\textbf{F15 · Formación y certificación E2, perfil 3, cohortes codificadas}\newline{}José Lara Arce\newline{}1 identificadores reservados; cantidad activa según manifiesto.\newline{}\textbf{Ventana:} 18. & \textbf{Contenido y frontera.} Formación de una cohorte activada del perfil 3, etapa E2, conforme al manifiesto aprobado. Se registran asistentes, turno, sitio, modalidad, sesiones, ejercicio, evaluación y recuperación. Identificador reservado sin activación no constituye trabajo, costo, avance ni recepción.\newline{}\textbf{Aceptación.} Dos sesiones de dos horas, ejercicio, registro individual y recuperación; máximo doce asistentes y sin interrumpir turnos.\newline{}\textbf{Insumos.} Manifiesto de cantidad y población de 1.11.3.2.1, materiales 1.11.2.2.3, versión recibida y disponibilidad acordada de participantes e instructor. Los datos privados de asistentes permanecen en el registro autorizado, no en este formulario.\newline{}\textbf{Condiciones.} Predecesoras: 1.11.3.2.1,1.11.2.2.3,VERSION-E2. Códigos y cantidades siguen la tabla de control de cohortes; no se presume que todos los sufijos estén activos. Si la demanda supera la reserva, se amplía el registro sin limitar la formación comprometida. FALTA REFERENCIA A SD3; FALTA REFERENCIA A SD10.\\\hline
\texttt{1.\allowbreak{}11.\allowbreak{}7.\allowbreak{}2.\allowbreak{}4.\allowbreak{}1.\allowbreak{}1}\newline{}\textbf{F16 · Formación y certificación E2, perfil 4, cohortes codificadas}\newline{}José Lara Arce\newline{}1 identificadores reservados; cantidad activa según manifiesto.\newline{}\textbf{Ventana:} 18. & \textbf{Contenido y frontera.} Formación de una cohorte activada del perfil 4, etapa E2, conforme al manifiesto aprobado. Se registran asistentes, turno, sitio, modalidad, sesiones, ejercicio, evaluación y recuperación. Identificador reservado sin activación no constituye trabajo, costo, avance ni recepción.\newline{}\textbf{Aceptación.} Dos sesiones de dos horas, ejercicio, registro individual y recuperación; máximo doce asistentes y sin interrumpir turnos.\newline{}\textbf{Insumos.} Manifiesto de cantidad y población de 1.11.3.2.1, materiales 1.11.2.2.4, versión recibida y disponibilidad acordada de participantes e instructor. Los datos privados de asistentes permanecen en el registro autorizado, no en este formulario.\newline{}\textbf{Condiciones.} Predecesoras: 1.11.3.2.1,1.11.2.2.4,VERSION-E2. Códigos y cantidades siguen la tabla de control de cohortes; no se presume que todos los sufijos estén activos. Si la demanda supera la reserva, se amplía el registro sin limitar la formación comprometida. FALTA REFERENCIA A SD3; FALTA REFERENCIA A SD10.\\\hline
\texttt{1.\allowbreak{}11.\allowbreak{}7.\allowbreak{}2.\allowbreak{}5.\allowbreak{}1.\allowbreak{}\{1--3\}}\newline{}\textbf{F17 · Formación y certificación E2, perfil 5, cohortes codificadas}\newline{}José Lara Arce\newline{}3 identificadores reservados; cantidad activa según manifiesto.\newline{}\textbf{Ventana:} 18. & \textbf{Contenido y frontera.} Formación de una cohorte activada del perfil 5, etapa E2, conforme al manifiesto aprobado. Se registran asistentes, turno, sitio, modalidad, sesiones, ejercicio, evaluación y recuperación. Identificador reservado sin activación no constituye trabajo, costo, avance ni recepción.\newline{}\textbf{Aceptación.} Dos sesiones de dos horas, ejercicio, registro individual y recuperación; máximo doce asistentes y sin interrumpir turnos.\newline{}\textbf{Insumos.} Manifiesto de cantidad y población de 1.11.3.2.1, materiales 1.11.2.2.5, versión recibida y disponibilidad acordada de participantes e instructor. Los datos privados de asistentes permanecen en el registro autorizado, no en este formulario.\newline{}\textbf{Condiciones.} Predecesoras: 1.11.3.2.1,1.11.2.2.5,VERSION-E2. Códigos y cantidades siguen la tabla de control de cohortes; no se presume que todos los sufijos estén activos. Si la demanda supera la reserva, se amplía el registro sin limitar la formación comprometida. FALTA REFERENCIA A SD3; FALTA REFERENCIA A SD10.\\\hline
\texttt{1.\allowbreak{}11.\allowbreak{}8.\allowbreak{}1.\allowbreak{}\{13--20\}}\newline{}\textbf{Atención y soporte prestados: implementación}\newline{}José Lara Arce\newline{}8 instancias mensuales.\newline{}\textbf{Ventana:} 13--20. & \textbf{Contenido y frontera.} Prestación mensual efectiva de mesa bilingüe, orientación funcional/técnica, registro, clasificación, resolución o derivación y seguimiento hasta cierre. Incluye mantenimiento de la base de conocimiento, autoformación y evidencia de aprendizaje. Excluye respuesta especializada de guardia, ingeniería de cambios, acompañamiento presencial de olas y el informe agregado de SLA.\newline{}\textbf{Aceptación.} Cada mes conserva cuadrantes, personal/suplencia, canales disponibles, contactos y tickets trazables. Cobertura C07: 06:00 a 24:00 todos los días en temporada alta y 08:00 a 20:00 el resto del año, español/inglés. Se verifica 80\,\% de atención antes de 20 segundos, resolución al primer contacto de al menos 70\,\% y abandono no superior a 5\,\%, con denominadores y fallos registrados; cero contactos no prueba desempeño. Ticket tiene severidad, responsable, escalamiento y cierre confirmado o plazo automático declarado. Base de conocimiento actualizada/versionada, accesible, editable y de propiedad del cliente; cambios, uso y utilidad percibida registrados. Se acredita la prestación del período; informe favorable no sustituye cobertura ni corrige incumplimientos.\newline{}\textbf{Insumos.} Línea base/delta aplicable de 1.12.1.1 y 1.12.1.2 recibidos antes de activar cobertura, versión supervisada, herramientas/canales y dotación disponible; calendarios naturales, suplencias y permisos.\newline{}\textbf{Condiciones.} Cada mes es una instancia de prestación y se recibe con su universo completo. Un incidente conserva ticket único: mesa registra/deriva y guardia responde; cada intervención se imputa una vez. Registros diarios y turnos son actividades/evidencia de T-15, sin paquetes adicionales por persona/día. Soporte de implementación/marcha blanca, meses 13--20, según versión supervisada; no anticipa la fase contractual de Operación. Es transferido sin interrupción a 1.12.14 en mes 21. RT-21.01/05--10/13/15--18; C07 RT-21.06/16; FALTA REFERENCIA A SD10; FALTA REFERENCIA A SD12.\\\hline
\texttt{1.\allowbreak{}11.\allowbreak{}8.\allowbreak{}2.\allowbreak{}\{13--20\}}\newline{}\textbf{Guardia, monitoreo y respuesta prestados: implementación}\newline{}José Lara Arce\newline{}8 instancias mensuales.\newline{}\textbf{Ventana:} 13--20. & \textbf{Contenido y frontera.} Prestación mensual efectiva de NOC/SOC, monitoreo proactivo, guardia crítica y respuesta técnica, con escalamiento a especialistas. Atiende canales de emergencia y alertas de planes vencidos 24x7x365, además de los incidentes críticos de plataforma. Excluye otorgar zarpes o ejecutar salvamento, mesa ordinaria, acompañamiento de olas y producción de cambios ya imputados a ingeniería.\newline{}\textbf{Aceptación.} Guardias y relevos cubren todos los intervalos del mes, con especialistas/suplentes y registros de disponibilidad, alertas, contacto, diagnóstico, contención, escalamiento y recuperación. Se comprueba recepción/respuesta 24x7x365 del canal de emergencia y planes vencidos, criterios de severidad y tiempos aplicables de BA78/T-12, sin dar alertas técnicas por confirmación de hechos físicos. Traslado requerido a Bahía Ilque identifica embarcación, ventana marítima segura, especialista y suplencia, sin desatender otros sitios. La bitácora conserva intervalos sin cobertura y desvíos; la recepción exige evidencia del servicio real, no solo cuadrante diseñado.\newline{}\textbf{Insumos.} Línea base/delta aplicable de 1.12.1.1 y 1.12.1.2 recibidos antes de activar cobertura, versión supervisada, herramientas/canales y dotación disponible; calendarios naturales, suplencias y permisos.\newline{}\textbf{Condiciones.} Cada mes es una instancia de prestación y se recibe con su universo completo. Un incidente conserva ticket único: mesa registra/deriva y guardia responde; cada intervención se imputa una vez. Registros diarios y turnos son actividades/evidencia de T-15, sin paquetes adicionales por persona/día. Soporte de implementación/marcha blanca, meses 13--20, según versión supervisada; no anticipa la fase contractual de Operación. Es transferido sin interrupción a 1.12.14 en mes 21. RT-21.01/05--10/13/15--18; C07 RT-21.06/16; FALTA REFERENCIA A SD10; FALTA REFERENCIA A SD12.\\\hline
\end{longtable}\endgroup
\textbf{Frontera de servicio, mantenimiento e ingeniería.} \texttt{1.11.8}/\texttt{1.12.14} recibe la prestación efectiva; \texttt{1.12.3.1} produce su control mensual. \texttt{1.12.3.4} recibe mantenimiento preventivo y liberaciones de base, con referencia de 24 HH mensuales y 4,8 HH de deuda internas. \texttt{1.12.12} recibe servicio de ingeniería y cambios autorizados: correctiva 80, normativa 48, evolutiva 80 y deuda 32 HH mensuales de referencia, según SD6. Se conservan 240 HH de ingeniería y $32/240\simeq0{,}1333$ (13,33\%) de deuda; esas cantidades no son una nueva estimación ni sustituyen T-15. Las referencias correctiva/normativa no limitan obligaciones incluidas, garantías o plazos; la bolsa evolutiva ofrecida conserva 960 HH anuales y saldo acumulable conforme a BT21.4.

Cada orden clasifica causa/objetivo, intervención, resultado y cuenta de imputación. Un defecto pertenece a correctiva; un cambio legal aplicable a normativa; una mejora solicitada a evolutiva; mejora interna sin función nueva a deuda; revisión preventiva/actualización de base al mantenimiento. Trabajo mixto separa aportes disjuntos y las pruebas se imputan al resultado que las produce, una sola vez. Guardias, atención y diagnóstico no se vuelven a cargar como construcción de cambios. Cierre mensual concilia disponibilidad, ejecución, saldo y pendientes; ni un registro de reserva acredita un cambio ejecutado, ni recibir el período cierra automáticamente trabajos abiertos. La recepción técnica del cambio exige versión, prueba, revisión y decisión particulares. FALTA REFERENCIA A SD6; FALTA REFERENCIA A SD10.

\begingroup\setlength{\tabcolsep}{3pt}
\begin{longtable}{L{\dimexpr0.320000\textwidth-2\tabcolsep\relax}L{\dimexpr0.680000\textwidth-2\tabcolsep\relax}}
\caption{Diccionario: Servicio y salida}\label{tab:t14-dic-1-12}\\\hline
\EncabezadoTabla{Paquete y resultado} & \EncabezadoTabla{Frontera, aceptación e insumos}\\\hline\endfirsthead
\multicolumn{2}{l}{\textit{Continuación de la tabla \ref{tab:t14-dic-1-12}}}\\\hline
\EncabezadoTabla{Paquete y resultado} & \EncabezadoTabla{Frontera, aceptación e insumos}\\\hline\endhead
\hline\endfoot
\hline\endlastfoot
\texttt{1.\allowbreak{}12.\allowbreak{}1.\allowbreak{}1}\newline{}\textbf{Modelo y cobertura de soporte habilitados}\newline{}José Lara Arce\newline{}1 resultado, con línea base y actualización.\newline{}\textbf{Ventana:} 11--12,18. & \textbf{Contenido y frontera.} Modelo, proveedores, escalamiento y suplencias del soporte de implementación y del servicio de Operación. Línea base para meses 13--20 recibida antes del mes 13; actualización de Operación en mes 18 antes del mes 21. No ejecuta turnos ni sustituye la prestación mensual.\newline{}\textbf{Aceptación.} Mesa, NOC, SOC y especialistas tienen ubicación, responsabilidades, dotación/turnos, disponibilidad y reemplazos comprobables. El dimensionamiento conserva demanda, tiempo de atención, Erlang C/sensibilidad y restricciones; T-15 distribuye recursos. Uso de terceros requiere habilitación y acuerdo efectivo antes de cubrir un turno, no solo contrato proyectado. Antes de Operación constan designación del Gerente de Servicio, especialistas nominados y dedicaciones. Línea base y delta tienen revisión y recepción propias.\newline{}\textbf{Insumos.} H3, alcance por fase, requisitos de atención C07, demanda estimada, proveedores/recursos y arquitectura vigente. La actualización de mes 18 usa registros reales de soporte y criterios previstos de recepción E2; no exige un acta futura. Cambios derivados de evidencia posterior se incorporan al delta antes de activar Operación.\newline{}\textbf{Condiciones.} La línea base inicial precede 1.11.8; el delta precede 1.12.14. José integra la preparación técnica; Gerente de Servicio coordina Operación conforme a SD3/SD6. Preparar cobertura no acredita turnos prestados. RT-21.01--03/05/14; FALTA REFERENCIA A SD10; FALTA REFERENCIA A SD12.\\\hline
\texttt{1.\allowbreak{}12.\allowbreak{}1.\allowbreak{}2}\newline{}\textbf{Catálogo, runbooks y transferencia al soporte habilitados}\newline{}José Lara Arce\newline{}1 resultado, con línea base y actualización.\newline{}\textbf{Ventana:} 12,20. & \textbf{Contenido y frontera.} Catálogo y procedimientos de continuidad, atención, diagnóstico, mantenimiento y acceso. Línea base de soporte antes del mes 13 y delta/traspaso al servicio en mes 20. Reutiliza runbooks existentes y fuentes/pendientes de las etapas; excluye nuevas capacidades de software y ejecución mensual de atención.\newline{}\textbf{Aceptación.} Equipo de soporte reproduce los procedimientos, reconoce autoridad, registra tickets/escalamiento y verifica acceso y retorno. Contraparte recibe documentación de su ámbito. Cada versión tiene responsable, revisión, pendientes y decisión; los bloqueantes impiden activar la cobertura afectada. Se recibe la transferencia inicial y la actualización sin atribuirles turnos futuros.\newline{}\textbf{Insumos.} Línea base o delta aplicable de 1.12.1.1, runbooks técnicos, base de conocimiento de soporte, versiones y accesos autorizados.\newline{}\textbf{Condiciones.} Línea base recibida antes de 1.11.8; actualización recibida antes de 1.12.14. RT-21.04; BA77; FALTA REFERENCIA A SD10; FALTA REFERENCIA A SD6.\\\hline
\texttt{1.\allowbreak{}12.\allowbreak{}2.\allowbreak{}1.\allowbreak{}\{21--56\}.\allowbreak{}\{1,\allowbreak{}2\}}\newline{}\textbf{Seguimiento de innovación IN1.7}\newline{}José Lara Arce\newline{}72 instancias.\newline{}\textbf{Ventana:} 21--56. & \textbf{Contenido y frontera.} Seguimiento incremental de la innovación referida. Excluye desarrollo base, informe general de SLA y doble imputación del mismo seguimiento.\newline{}\textbf{Aceptación.} Informe quincenal IN1 con solicitudes, estados, primera respuesta, casos abiertos, indicadores y decisiones; no repite el informe general de servicio.\newline{}\textbf{Insumos.} Versión/registro recibido, acceso autorizado, herramientas, participantes y ventana del período.\newline{}\textbf{Condiciones.} Predecesoras de habilitación: herramientas, datos y decisiones de la etapa recibidos.\\\hline
\texttt{1.\allowbreak{}12.\allowbreak{}2.\allowbreak{}2.\allowbreak{}\{31--56\}}\newline{}\textbf{F25 · Seguimiento mensual incremental IN2}\newline{}José Lara Arce\newline{}26 instancias.\newline{}\textbf{Ventana:} 31--56. & \textbf{Contenido y frontera.} Registro mensual incremental IN2 con planes, solicitudes, trabajos aceptados/pendientes, ventanas y decisiones. Reutiliza información del servicio sin duplicar su informe general.\newline{}\textbf{Aceptación.} Estado de planes/trabajos y ventanas recibidos; seguimiento sin duplicar informes generales.\newline{}\textbf{Insumos.} Universo nominal/manifiesto por instancia, versión y política/protocolo aplicable recibidos, responsables, acceso autorizado y ventana admisible. Registro de instancia: código, población/activos, fecha, versión, evidencia y decisión; el conjunto no se acepta por cerrar solo su primera instancia.\newline{}\textbf{Condiciones.} Condiciones/predecesoras comunes del catálogo: 1.9.2.5.2. Referencia de detalle: FALTA REFERENCIA A SD13.\\\hline
\texttt{1.\allowbreak{}12.\allowbreak{}2.\allowbreak{}3.\allowbreak{}\{21--56\}.\allowbreak{}\{1,\allowbreak{}2\}}\newline{}\textbf{Seguimiento de innovación IN3.7}\newline{}José Lara Arce\newline{}72 instancias.\newline{}\textbf{Ventana:} 21--56. & \textbf{Contenido y frontera.} Seguimiento incremental de la innovación referida. Excluye desarrollo base, informe general de SLA y doble imputación del mismo seguimiento.\newline{}\textbf{Aceptación.} Seguimiento quincenal IN3 con muestra manual/asistida, validaciones humanas, errores y evolución; la sugerencia no confirma por sí sola un dato de negocio.\newline{}\textbf{Insumos.} Versión/registro recibido, acceso autorizado, herramientas, participantes y ventana del período.\newline{}\textbf{Condiciones.} Predecesoras de habilitación: herramientas, datos y decisiones de la etapa recibidos.\\\hline
\texttt{1.\allowbreak{}12.\allowbreak{}2.\allowbreak{}4.\allowbreak{}\{21--56\}}\newline{}\textbf{Seguimiento de innovación IN4.4}\newline{}José Lara Arce\newline{}36 instancias.\newline{}\textbf{Ventana:} 21--56. & \textbf{Contenido y frontera.} Seguimiento incremental de la innovación referida. Excluye desarrollo base, informe general de SLA y doble imputación del mismo seguimiento.\newline{}\textbf{Aceptación.} Control mensual del servicio gestionado por resultados, SLA, responsabilidad y acciones; no equivale a afirmar certificación del proveedor.\newline{}\textbf{Insumos.} Versión/registro recibido, acceso autorizado, herramientas, participantes y ventana del período.\newline{}\textbf{Condiciones.} Predecesoras de habilitación: herramientas, datos y decisiones de la etapa recibidos.\\\hline
\texttt{1.\allowbreak{}12.\allowbreak{}2.\allowbreak{}5.\allowbreak{}\{21--56\}}\newline{}\textbf{Seguimiento de innovación IN5.4}\newline{}José Lara Arce\newline{}36 instancias.\newline{}\textbf{Ventana:} 21--56. & \textbf{Contenido y frontera.} Seguimiento incremental de la innovación referida. Excluye desarrollo base, informe general de SLA y doble imputación del mismo seguimiento.\newline{}\textbf{Aceptación.} Actualización mensual del Pasaporte Ambiental Náutico con origen, integridad, consumos y residuos; excepciones y evidencia para auditoría identificadas.\newline{}\textbf{Insumos.} Versión/registro recibido, acceso autorizado, herramientas, participantes y ventana del período.\newline{}\textbf{Condiciones.} Predecesoras de habilitación: herramientas, datos y decisiones de la etapa recibidos.\\\hline
\texttt{1.\allowbreak{}12.\allowbreak{}3.\allowbreak{}1.\allowbreak{}21}\newline{}\textbf{Control mensual de servicio y SLA}\newline{}José Lara Arce\newline{}1 instancia.\newline{}\textbf{Ventana:} 21. & \textbf{Contenido y frontera.} SLA, incidentes, costos, compromisos y causas; excluye turnos de recepción continua.\newline{}\textbf{Aceptación.} Denominadores, tiempos, multas y acciones son trazables; una ausencia de registros no acredita cumplimiento. José Lara Arce responde por el paquete.\newline{}\textbf{Condiciones.} Predecesoras del catálogo: SERVICIO-ACTIVO. Referencia: BT21; BA78. Las habilitaciones técnicas se contrastan con arquitectura e inventario vigentes.\\\hline
\texttt{1.\allowbreak{}12.\allowbreak{}3.\allowbreak{}1.\allowbreak{}\{22--56\}}\newline{}\textbf{Control mensual de servicio y SLA}\newline{}José Lara Arce\newline{}35 instancias.\newline{}\textbf{Ventana:} 22--56. & \textbf{Contenido y frontera.} Resultado individual del período o conjunto indicado. Excluye resultados de otras cuentas y horas de contraparte, obras o adquisiciones del cliente.\newline{}\textbf{Aceptación.} Denominadores, tiempos, multas y acciones son trazables; una ausencia de registros no acredita cumplimiento. José Lara Arce responde por el paquete.\newline{}\textbf{Insumos.} Versión/registro recibido, acceso autorizado, herramientas, participantes y ventana del período.\newline{}\textbf{Condiciones.} Predecesoras de habilitación: SERVICIO-ACTIVO.\\\hline
\texttt{1.\allowbreak{}12.\allowbreak{}3.\allowbreak{}2.\allowbreak{}21}\newline{}\textbf{Comprobación de respaldos y restauración}\newline{}José Lara Arce\newline{}1 instancia.\newline{}\textbf{Ventana:} 21. & \textbf{Contenido y frontera.} Revisión de copias, prueba de muestra y evidencia; excluye conmutación semestral.\newline{}\textbf{Aceptación.} Restauración tiene integridad, accesos y tiempo registrados y brechas con responsable. José Lara Arce responde por el paquete.\newline{}\textbf{Condiciones.} Predecesoras del catálogo: SERVICIO-ACTIVO. Referencia: BT21; BA78. Las habilitaciones técnicas se contrastan con arquitectura e inventario vigentes.\\\hline
\texttt{1.\allowbreak{}12.\allowbreak{}3.\allowbreak{}2.\allowbreak{}\{22--56\}}\newline{}\textbf{Comprobación de respaldos y restauración}\newline{}José Lara Arce\newline{}35 instancias.\newline{}\textbf{Ventana:} 22--56. & \textbf{Contenido y frontera.} Resultado individual del período o conjunto indicado. Excluye resultados de otras cuentas y horas de contraparte, obras o adquisiciones del cliente.\newline{}\textbf{Aceptación.} Restauración tiene integridad, accesos y tiempo registrados y brechas con responsable. José Lara Arce responde por el paquete.\newline{}\textbf{Insumos.} Versión/registro recibido, acceso autorizado, herramientas, participantes y ventana del período.\newline{}\textbf{Condiciones.} Predecesoras de habilitación: SERVICIO-ACTIVO.\\\hline
\texttt{1.\allowbreak{}12.\allowbreak{}3.\allowbreak{}3.\allowbreak{}21}\newline{}\textbf{Revisión de seguridad y acceso}\newline{}Bruno Toro\newline{}1 instancia.\newline{}\textbf{Ventana:} 21. & \textbf{Contenido y frontera.} Accesos, alertas, agentes y cambios de permisos; excluye guardia SOC y ofensiva anual.\newline{}\textbf{Aceptación.} Cobertura por activo y cuentas efectivas quedan verificadas y desvíos tratados. Bruno Toro responde por el paquete.\newline{}\textbf{Condiciones.} Predecesoras del catálogo: SERVICIO-ACTIVO. Referencia: BT21; BA78. Las habilitaciones técnicas se contrastan con arquitectura e inventario vigentes.\\\hline
\texttt{1.\allowbreak{}12.\allowbreak{}3.\allowbreak{}3.\allowbreak{}\{22--56\}}\newline{}\textbf{Revisión de seguridad y acceso}\newline{}Bruno Toro\newline{}35 instancias.\newline{}\textbf{Ventana:} 22--56. & \textbf{Contenido y frontera.} Resultado individual del período o conjunto indicado. Excluye resultados de otras cuentas y horas de contraparte, obras o adquisiciones del cliente.\newline{}\textbf{Aceptación.} Cobertura por activo y cuentas efectivas quedan verificadas y desvíos tratados. Bruno Toro responde por el paquete.\newline{}\textbf{Insumos.} Versión/registro recibido, acceso autorizado, herramientas, participantes y ventana del período.\newline{}\textbf{Condiciones.} Predecesoras de habilitación: SERVICIO-ACTIVO.\\\hline
\texttt{1.\allowbreak{}12.\allowbreak{}3.\allowbreak{}4.\allowbreak{}21}\newline{}\textbf{Mantenimiento preventivo y liberación controlada recibidos}\newline{}Nikolai Navea\newline{}1 instancia.\newline{}\textbf{Ventana:} 21. & \textbf{Contenido y frontera.} Mantenimiento preventivo programado de componentes, revisiones/configuración y liberación acotada de versiones de base. Conserva referencia de 24 HH mensuales y 20\,\% interno (4,8 HH) para sus propias acciones de deuda; no representa correcciones/evoluciones ilimitadas ni cambios ya producidos en 1.12.12.\newline{}\textbf{Aceptación.} Calendario anual de mantenimiento y versiones acordado, revisiones ejecutadas, auditorías técnicas al menos trimestrales, informe y pendientes trazables. Cada liberación tiene aviso, ventana, prueba, versión y retorno; se verifica la acción preventiva y su efecto. Deuda interna se cuantifica dentro de las 24 HH y no se vuelve a reconocer en la bolsa de 32 HH. Un cambio de aplicación autorizado en ingeniería puede usar este procedimiento de liberación; solo la intervención distinta se imputa aquí, no su construcción/pruebas ya recibidas.\newline{}\textbf{Insumos.} SERVICIO-ACTIVO, inventario/versiones, plan anual recibido, herramientas, permisos y ventana del período; órdenes e imputaciones de ingeniería para conciliar fronteras.\newline{}\textbf{Condiciones.} Una orden identifica intervención, resultado y cuenta de imputación única. Correctiva corrige defecto; normativa responde al cambio aplicable; evolutiva incorpora mejora solicitada; deuda mejora estructura interna; mantenimiento previene/revisa componentes existentes según calendario. En trabajo mixto se separan aportes sin superposición. RT-21.21/22; BT21.4; FALTA REFERENCIA A SD6; FALTA REFERENCIA A SD11.\\\hline
\texttt{1.\allowbreak{}12.\allowbreak{}3.\allowbreak{}4.\allowbreak{}\{22--56\}}\newline{}\textbf{Mantenimiento preventivo y liberación controlada recibidos}\newline{}Nikolai Navea\newline{}35 instancias mensuales.\newline{}\textbf{Ventana:} 22--56. & \textbf{Contenido y frontera.} Mantenimiento preventivo programado de componentes, revisiones/configuración y liberación acotada de versiones de base. Conserva referencia de 24 HH mensuales y 20\,\% interno (4,8 HH) para sus propias acciones de deuda; no representa correcciones/evoluciones ilimitadas ni cambios ya producidos en 1.12.12.\newline{}\textbf{Aceptación.} Calendario anual de mantenimiento y versiones acordado, revisiones ejecutadas, auditorías técnicas al menos trimestrales, informe y pendientes trazables. Cada liberación tiene aviso, ventana, prueba, versión y retorno; se verifica la acción preventiva y su efecto. Deuda interna se cuantifica dentro de las 24 HH y no se vuelve a reconocer en la bolsa de 32 HH. Un cambio de aplicación autorizado en ingeniería puede usar este procedimiento de liberación; solo la intervención distinta se imputa aquí, no su construcción/pruebas ya recibidas.\newline{}\textbf{Insumos.} SERVICIO-ACTIVO, inventario/versiones, plan anual recibido, herramientas, permisos y ventana del período; órdenes e imputaciones de ingeniería para conciliar fronteras.\newline{}\textbf{Condiciones.} Una orden identifica intervención, resultado y cuenta de imputación única. Correctiva corrige defecto; normativa responde al cambio aplicable; evolutiva incorpora mejora solicitada; deuda mejora estructura interna; mantenimiento previene/revisa componentes existentes según calendario. En trabajo mixto se separan aportes sin superposición. RT-21.21/22; BT21.4; FALTA REFERENCIA A SD6; FALTA REFERENCIA A SD11.\\\hline
\texttt{1.\allowbreak{}12.\allowbreak{}3.\allowbreak{}5.\allowbreak{}21}\newline{}\textbf{Proyección de capacidad e infraestructura}\newline{}Fernando Guerrero\newline{}1 instancia.\newline{}\textbf{Ventana:} 21. & \textbf{Contenido y frontera.} Seguimiento mensual y consolidación trimestral sin segunda imputación. Incluye consumo/costo desglosado y utilización, FinOps y actualización de estimaciones de energía/carbono de la infraestructura según método aprobado; no confunde carbono de la plataforma con el Pasaporte Ambiental de embarcaciones. Seguimiento y consolidación del mismo período se reciben una sola vez.\newline{}\textbf{Aceptación.} Crecimiento, cuello de botella, alertas y propuesta de costo se vinculan a datos observados. Fernando Guerrero responde por el paquete. Datos medidos o estimados se distinguen con fuente, período y límites; se contrastan capacidad, etiquetado/presupuesto, ahorro de ambientes no productivos y desvíos de consumo. PUE e intensidad regional sostienen la estimación de huella de la plataforma; no se inventan mediciones ni certificaciones.\newline{}\textbf{Condiciones.} Predecesoras del catálogo: SERVICIO-ACTIVO. Referencia: BT21; BA78. Las habilitaciones técnicas se contrastan con arquitectura e inventario vigentes.\\\hline
\texttt{1.\allowbreak{}12.\allowbreak{}3.\allowbreak{}5.\allowbreak{}\{22--56\}}\newline{}\textbf{Proyección de capacidad e infraestructura}\newline{}Fernando Guerrero\newline{}35 instancias.\newline{}\textbf{Ventana:} 22--56. & \textbf{Contenido y frontera.} Resultado individual del período o conjunto indicado. Excluye resultados de otras cuentas y horas de contraparte, obras o adquisiciones del cliente. Incluye consumo/costo desglosado y utilización, FinOps y actualización de estimaciones de energía/carbono de la infraestructura según método aprobado; no confunde carbono de la plataforma con el Pasaporte Ambiental de embarcaciones. Seguimiento y consolidación del mismo período se reciben una sola vez.\newline{}\textbf{Aceptación.} Crecimiento, cuello de botella, alertas y propuesta de costo se vinculan a datos observados. Fernando Guerrero responde por el paquete. Datos medidos o estimados se distinguen con fuente, período y límites; se contrastan capacidad, etiquetado/presupuesto, ahorro de ambientes no productivos y desvíos de consumo. PUE e intensidad regional sostienen la estimación de huella de la plataforma; no se inventan mediciones ni certificaciones.\newline{}\textbf{Insumos.} Versión/registro recibido, acceso autorizado, herramientas, participantes y ventana del período.\newline{}\textbf{Condiciones.} Predecesoras de habilitación: SERVICIO-ACTIVO.\\\hline
\texttt{1.\allowbreak{}12.\allowbreak{}4.\allowbreak{}1.\allowbreak{}6}\newline{}\textbf{Resiliencia de componentes y recuperación: control semestral}\newline{}José Lara Arce\newline{}1 instancia.\newline{}\textbf{Ventana:} 24. & \textbf{Contenido y frontera.} Falla controlada de enlace, nodo y dependencia, con registro y recuperación.\newline{}\textbf{Aceptación.} Se conserva operación/degradación controlada y recuperación, con tiempos, transacciones y brechas documentados. José Lara Arce responde por el paquete.\newline{}\textbf{Condiciones.} Predecesoras del catálogo: SERVICIO-ACTIVO. Referencia: BT20.1; BA continuidad. Las habilitaciones técnicas se contrastan con arquitectura e inventario vigentes.\\\hline
\texttt{1.\allowbreak{}12.\allowbreak{}4.\allowbreak{}1.\allowbreak{}7}\newline{}\textbf{Conmutación real de desastre y retorno: control semestral}\newline{}José Lara Arce\newline{}1 instancia.\newline{}\textbf{Ventana:} 24. & \textbf{Contenido y frontera.} Conmutación regional real, verificación de datos/identidad y retorno; excluye usar existencia de réplica como evidencia.\newline{}\textbf{Aceptación.} RTO/RPO comprometidos se alcanzan y evidencia identifica último dato, tiempos, permisos y conciliación. José Lara Arce responde por el paquete.\newline{}\textbf{Condiciones.} Predecesoras del catálogo: SERVICIO-ACTIVO. Referencia: BT20.1; BA continuidad. Las habilitaciones técnicas se contrastan con arquitectura e inventario vigentes.\\\hline
\texttt{1.\allowbreak{}12.\allowbreak{}4.\allowbreak{}\{2--6\}.\allowbreak{}6}\newline{}\textbf{Resiliencia de componentes y recuperación: control semestral}\newline{}José Lara Arce\newline{}5 instancias.\newline{}\textbf{Ventana:} 28,33,36,40,45. & \textbf{Contenido y frontera.} Resultado individual del período o conjunto indicado. Excluye resultados de otras cuentas y horas de contraparte, obras o adquisiciones del cliente.\newline{}\textbf{Aceptación.} Se conserva operación/degradación controlada y recuperación, con tiempos, transacciones y brechas documentados. José Lara Arce responde por el paquete.\newline{}\textbf{Insumos.} Versión/registro recibido, acceso autorizado, herramientas, participantes y ventana del período.\newline{}\textbf{Condiciones.} Predecesoras de habilitación: SERVICIO-ACTIVO.\\\hline
\texttt{1.\allowbreak{}12.\allowbreak{}4.\allowbreak{}\{2--6\}.\allowbreak{}7}\newline{}\textbf{Conmutación real de desastre y retorno: control semestral}\newline{}José Lara Arce\newline{}5 instancias.\newline{}\textbf{Ventana:} 28,33,36,40,45. & \textbf{Contenido y frontera.} Resultado individual del período o conjunto indicado. Excluye resultados de otras cuentas y horas de contraparte, obras o adquisiciones del cliente.\newline{}\textbf{Aceptación.} RTO/RPO comprometidos se alcanzan y evidencia identifica último dato, tiempos, permisos y conciliación. José Lara Arce responde por el paquete.\newline{}\textbf{Insumos.} Versión/registro recibido, acceso autorizado, herramientas, participantes y ventana del período.\newline{}\textbf{Condiciones.} Predecesoras de habilitación: SERVICIO-ACTIVO.\\\hline
\texttt{1.\allowbreak{}12.\allowbreak{}4.\allowbreak{}\{7--9\}.\allowbreak{}6}\newline{}\textbf{F37 · Controles adicionales identificados de resiliencia}\newline{}Davor Serey\newline{}3 instancias.\newline{}\textbf{Ventana:} 48,52,56. & \textbf{Contenido y frontera.} Resultado de una prueba adicional de resiliencia identificado por control, versión, escenario, ventana, tiempos y recuperación; no repite la prueba básica ya recibida.\newline{}\textbf{Aceptación.} Prueba sobre versión identificada y ventana autorizada, sin degradar E1; recuperación y retorno medidos.\newline{}\textbf{Insumos.} Universo nominal/manifiesto por instancia, versión y política/protocolo aplicable recibidos, responsables, acceso autorizado y ventana admisible. Registro de instancia: código, población/activos, fecha, versión, evidencia y decisión; el conjunto no se acepta por cerrar solo su primera instancia.\newline{}\textbf{Condiciones.} Condiciones/predecesoras comunes del catálogo: SERVICIO-ACTIVO. Referencia de detalle: FALTA REFERENCIA A SD9.\\\hline
\texttt{1.\allowbreak{}12.\allowbreak{}4.\allowbreak{}\{7--9\}.\allowbreak{}7}\newline{}\textbf{F38 · Controles adicionales identificados de DR real y retorno}\newline{}José Lara Arce\newline{}3 instancias.\newline{}\textbf{Ventana:} 48,52,56. & \textbf{Contenido y frontera.} Resultado de una prueba adicional de DR real y retorno, con escenario, versión, autoridad de escritura, datos recuperados y tiempos medidos; depende de la pareja regional aprobada.\newline{}\textbf{Aceptación.} Prueba sobre versión identificada y ventana autorizada, sin degradar E1; recuperación y retorno medidos.\newline{}\textbf{Insumos.} Universo nominal/manifiesto por instancia, versión y política/protocolo aplicable recibidos, responsables, acceso autorizado y ventana admisible. Registro de instancia: código, población/activos, fecha, versión, evidencia y decisión; el conjunto no se acepta por cerrar solo su primera instancia.\newline{}\textbf{Condiciones.} Condiciones/predecesoras comunes del catálogo: SERVICIO-ACTIVO. Referencia de detalle: FALTA REFERENCIA A SD9.\\\hline
\texttt{1.\allowbreak{}12.\allowbreak{}5.\allowbreak{}1}\newline{}\textbf{Prueba ofensiva y cierre de hallazgos: control anual}\newline{}Bruno Toro\newline{}1 instancia.\newline{}\textbf{Ventana:} 28. & \textbf{Contenido y frontera.} Ejecución ofensiva de superficie completa y comprobación de resolución; correcciones de código se imputan a construcción.\newline{}\textbf{Aceptación.} No quedan hallazgos críticos ni altos abiertos y cada resultado corresponde a versión recibida. Bruno Toro responde por el paquete.\newline{}\textbf{Condiciones.} Predecesoras del catálogo: SERVICIO-ACTIVO. Referencia: BT20.1. Las habilitaciones técnicas se contrastan con arquitectura e inventario vigentes.\\\hline
\texttt{1.\allowbreak{}12.\allowbreak{}5.\allowbreak{}\{2,\allowbreak{}3\}}\newline{}\textbf{Prueba ofensiva y cierre de hallazgos: control anual}\newline{}Bruno Toro\newline{}2 instancias.\newline{}\textbf{Ventana:} 40,52. & \textbf{Contenido y frontera.} Resultado individual del período o conjunto indicado. Excluye resultados de otras cuentas y horas de contraparte, obras o adquisiciones del cliente.\newline{}\textbf{Aceptación.} No quedan hallazgos críticos ni altos abiertos y cada resultado corresponde a versión recibida. Bruno Toro responde por el paquete.\newline{}\textbf{Insumos.} Versión/registro recibido, acceso autorizado, herramientas, participantes y ventana del período.\newline{}\textbf{Condiciones.} Predecesoras de habilitación: SERVICIO-ACTIVO.\\\hline
\texttt{1.\allowbreak{}12.\allowbreak{}6.\allowbreak{}1}\newline{}\textbf{Jornada de actualización y transferencia}\newline{}José Lara Arce\newline{}1 instancia.\newline{}\textbf{Ventana:} 24. & \textbf{Contenido y frontera.} Actualización, ejercicios y evaluación en ventana acordada; excluye incorporación adicional y capacitación en peak protegido.\newline{}\textbf{Aceptación.} Se conserva asistencia, evaluación, materiales actualizados y brechas; se cumple mínimo de dos jornadas por año de operación. José Lara Arce responde por el paquete.\newline{}\textbf{Condiciones.} Predecesoras del catálogo: SERVICIO-ACTIVO. Referencia: RT-22.06/07. Las habilitaciones técnicas se contrastan con arquitectura e inventario vigentes.\\\hline
\texttt{1.\allowbreak{}12.\allowbreak{}6.\allowbreak{}\{2--6\}}\newline{}\textbf{Jornada de actualización y transferencia}\newline{}José Lara Arce\newline{}5 instancias.\newline{}\textbf{Ventana:} 30,36,42,48,54. & \textbf{Contenido y frontera.} Resultado individual del período o conjunto indicado. Excluye resultados de otras cuentas y horas de contraparte, obras o adquisiciones del cliente.\newline{}\textbf{Aceptación.} Se conserva asistencia, evaluación, materiales actualizados y brechas; se cumple mínimo de dos jornadas por año de operación. José Lara Arce responde por el paquete.\newline{}\textbf{Insumos.} Versión/registro recibido, acceso autorizado, herramientas, participantes y ventana del período.\newline{}\textbf{Condiciones.} Predecesoras de habilitación: SERVICIO-ACTIVO.\\\hline
\texttt{1.\allowbreak{}12.\allowbreak{}7.\allowbreak{}1}\newline{}\textbf{Plan de reversibilidad actualizado}\newline{}Ignacio Silva\newline{}1 instancia.\newline{}\textbf{Ventana:} 24. & \textbf{Contenido y frontera.} Actualiza activos, dependencias, formatos, esfuerzo y custodia; excluye ejecutar sustitución total de proveedor.\newline{}\textbf{Aceptación.} Inventario/credenciales/exportaciones y estrategia de traslado reflejan versión efectiva sin copiar claves no exportables. Ignacio Silva responde por el paquete.\newline{}\textbf{Condiciones.} Predecesoras del catálogo: 1.1.3.1. Referencia: BA77; SD4 RT-03.07. Las habilitaciones técnicas se contrastan con arquitectura e inventario vigentes.\\\hline
\texttt{1.\allowbreak{}12.\allowbreak{}7.\allowbreak{}\{2,\allowbreak{}3\}}\newline{}\textbf{Plan de reversibilidad actualizado}\newline{}Ignacio Silva\newline{}2 instancias.\newline{}\textbf{Ventana:} 36,48. & \textbf{Contenido y frontera.} Resultado individual del período o conjunto indicado. Excluye resultados de otras cuentas y horas de contraparte, obras o adquisiciones del cliente.\newline{}\textbf{Aceptación.} Inventario/credenciales/exportaciones y estrategia de traslado reflejan versión efectiva sin copiar claves no exportables. Ignacio Silva responde por el paquete.\newline{}\textbf{Insumos.} Versión/registro recibido, acceso autorizado, herramientas, participantes y ventana del período.\newline{}\textbf{Condiciones.} Predecesoras de habilitación: 1.1.3.1.\\\hline
\texttt{1.\allowbreak{}12.\allowbreak{}8.\allowbreak{}1}\newline{}\textbf{Inventario de salida y responsables}\newline{}Ignacio Silva\newline{}1 instancia.\newline{}\textbf{Ventana:} 53. & \textbf{Contenido y frontera.} Transferencia de salida incluida; excluye confundir entrega contractual con reconstrucción completa de otra plataforma.\newline{}\textbf{Aceptación.} Cliente valida activos, interfaces, custodios y destino de transferencia. Ignacio Silva responde por el paquete.\newline{}\textbf{Condiciones.} Predecesoras del catálogo: PLAN-SALIDA. Referencia: BA77; SD4 reversibilidad. Las habilitaciones técnicas se contrastan con arquitectura e inventario vigentes.\\\hline
\texttt{1.\allowbreak{}12.\allowbreak{}8.\allowbreak{}2}\newline{}\textbf{Exportación abierta y manifiestos de datos}\newline{}Nicolás Torfan\newline{}1 instancia.\newline{}\textbf{Ventana:} 53--56. & \textbf{Contenido y frontera.} Exportación abierta y manifiestos de transferencia: entrega inicial en mes 53 y actualización/conciliación de cambios hasta el traspaso autorizado de mes 56. Se conservan identificadores y evidencia de cada versión; el delta no es una segunda exportación imputada del mismo estado. Excluye reconstruir la plataforma del receptor.\newline{}\textbf{Aceptación.} Datos, objetos, relaciones, versiones y huellas se exportan y concilian sin depender del software legado. La entrega inicial queda identificada como inicial; la recepción final exige manifiesto vigente, fecha/hora y autoridad del corte acordado, cambios posteriores conciliados y ausencia de diferencias no explicadas. La reproducción por el receptor usa ese estado final; una copia del mes 53 no acredita por sí sola integridad al término.\newline{}\textbf{Condiciones.} PLAN-SALIDA recibido, inventario/destino autorizados y acceso controlado. Servicio e ingeniería continúan hasta el traspaso; cada cambio posterior a la entrega inicial actualiza el manifiesto y la conciliación de este mismo paquete. 1.12.8.5 requiere su versión final conforme. BA77.2/77.3; procedimiento de reversibilidad/transferencia en T-18 y programación en T-15. FALTA REFERENCIA A SD4.\\\hline
\texttt{1.\allowbreak{}12.\allowbreak{}8.\allowbreak{}3}\newline{}\textbf{Fuentes, IaC, artefactos y documentación transferidos}\newline{}Fernando Guerrero\newline{}1 instancia.\newline{}\textbf{Ventana:} 54--56. & \textbf{Contenido y frontera.} Transferencia inicial de fuentes, IaC, artefactos y documentación en mes 54; actualización de los mismos resultados hasta la última versión desplegada al traspaso autorizado de mes 56. Excluye reconstruir otra plataforma o reconocer de nuevo cambios ya imputados a ingeniería.\newline{}\textbf{Aceptación.} Repositorio del cliente contiene fuentes, artefactos/builds, configuración y procedimientos reproducibles, con versión y huellas correspondientes a la última versión desplegada al traspaso. Documentación técnica, funcional, arquitectura, operación, seguridad y base de conocimiento quedan actualizadas a esa versión; escrow se verifica. Diferencias entre entrega inicial y final están conciliadas y la reproducción del receptor identifica la versión efectivamente recibida.\newline{}\textbf{Condiciones.} PLAN-SALIDA recibido y repositorios/accesos del receptor autorizados. Cambios/despliegues posteriores a mes 54 actualizan el manifiesto, los artefactos y su documentación sin otra imputación de construcción. 1.12.8.5 exige la entrega final conforme antes de aceptación. BA77.1/77.3; FALTA REFERENCIA A SD4; FALTA REFERENCIA A SD11.\\\hline
\texttt{1.\allowbreak{}12.\allowbreak{}8.\allowbreak{}4}\newline{}\textbf{Custodia, credenciales y recifrado autorizado}\newline{}Bruno Toro\newline{}1 instancia.\newline{}\textbf{Ventana:} 54--56. & \textbf{Contenido y frontera.} Preparación y transferencia autorizada de custodia, credenciales y recifrado desde mes 54, con comprobación final al traspaso de mes 56. Excluye divulgar secretos o retirar anticipadamente el acceso necesario para prestar el servicio vigente.\newline{}\textbf{Aceptación.} Claves no exportables se sustituyen por procedimiento autorizado; custodia, permisos y capacidad de acceso del receptor se comprueban sin divulgar secretos. Se identifica autoridad y vigencia de los accesos de cada parte durante la transición. Revocación/rotación de accesos salientes ocurre conforme al traspaso autorizado, con evidencia, sin interrumpir obligaciones previas.\newline{}\textbf{Condiciones.} PLAN-SALIDA y decisiones de autoridad/acceso recibidos; preparación en mes 54 y actualización de cambios de claves/permisos hasta el traspaso. 1.12.8.5 requiere configuración y acceso final conformes. BA77.2/77.3; FALTA REFERENCIA A SD4; FALTA REFERENCIA A SD11.\\\hline
\texttt{1.\allowbreak{}12.\allowbreak{}8.\allowbreak{}5}\newline{}\textbf{Aceptación de transferencia y eliminación autorizada}\newline{}Ignacio Silva\newline{}1 instancia.\newline{}\textbf{Ventana:} 56. & \textbf{Contenido y frontera.} Recepción final conjunta de datos, fuentes, artefactos, documentación, custodia y accesos actualizados al traspaso autorizado. Excluye confundir entregas iniciales de meses 53/54 con el estado final recibido y reconstruir otra plataforma.\newline{}\textbf{Aceptación.} El acta identifica fecha/hora y autoridad del corte final, manifiesto de datos y delta conciliado, última versión desplegada con fuentes/artefactos/configuración/documentación correspondientes y decisión de autoridad/accesos. Receptor reproduce operación sobre ese estado final y se demuestra continuidad de la transferencia; diferencias bloqueantes impiden recepción. Pendientes no bloqueantes quedan expresamente aceptados con dueño/plazo. Acompañamiento cumple al menos 90 días reales. Solo se elimina origen tras aceptación, autorización y evidencia, respetando conservación/custodia obligatorias; no se retira información necesaria para el servicio o la auditoría.\newline{}\textbf{Condiciones.} Predecesoras: versiones finales conformes de 1.12.8.2, 1.12.8.3 y 1.12.8.4; acompañamiento agregado 1.12.9 con al menos 90 días reales comprobados. T-18/plan de reversibilidad define la secuencia y T-15 los recursos y revisión positiva antes del cierre contractual. Toda modificación entre entrega inicial y traspaso debe estar incluida o resuelta por decisión autorizada que preserve integridad/continuidad. BA77.1--77.3; FALTA REFERENCIA A SD4; FALTA REFERENCIA A SD11.\\\hline
\texttt{1.\allowbreak{}12.\allowbreak{}9.\allowbreak{}53.\allowbreak{}1}\newline{}\textbf{Acompañamiento quincenal de transición}\newline{}José Lara Arce\newline{}1 instancia.\newline{}\textbf{Ventana:} 53. & \textbf{Contenido y frontera.} Soporte de transición, reproducción y atención al receptor; excluye guardias continuas y doble imputación de exportación.\newline{}\textbf{Aceptación.} La instancia identifica su quincena, cobertura prestada, consultas, ejecución observada, brechas y responsables; la evidencia del período se recibe individualmente. El conjunto de 1.12.9 debe acreditar al menos 90 días reales de acompañamiento sin cobro adicional antes de 1.12.8.5; no se exigen 90 días a cada registro quincenal. José Lara Arce integra la evidencia del período, sin sustituir la recepción contractual.\newline{}\textbf{Condiciones.} Predecesoras del catálogo: PLAN-SALIDA. Referencia: BA77. Las habilitaciones técnicas se contrastan con arquitectura e inventario vigentes.\\\hline
\texttt{1.\allowbreak{}12.\allowbreak{}9.\allowbreak{}\{53.\allowbreak{}2;\allowbreak{}\{54--56\}.\allowbreak{}\{1,\allowbreak{}2\}\}}\newline{}\textbf{Acompañamiento quincenal de transición}\newline{}José Lara Arce\newline{}7 instancias.\newline{}\textbf{Ventana:} 53--56. & \textbf{Contenido y frontera.} Resultado individual del período o conjunto indicado. Excluye resultados de otras cuentas y horas de contraparte, obras o adquisiciones del cliente.\newline{}\textbf{Aceptación.} La instancia identifica su quincena, cobertura prestada, consultas, ejecución observada, brechas y responsables; la evidencia del período se recibe individualmente. El conjunto de 1.12.9 debe acreditar al menos 90 días reales de acompañamiento sin cobro adicional antes de 1.12.8.5; no se exigen 90 días a cada registro quincenal. José Lara Arce integra la evidencia del período, sin sustituir la recepción contractual.\newline{}\textbf{Insumos.} Versión/registro recibido, acceso autorizado, herramientas, participantes y ventana del período.\newline{}\textbf{Condiciones.} Predecesoras de habilitación: PLAN-SALIDA.\\\hline
\texttt{1.\allowbreak{}12.\allowbreak{}10.\allowbreak{}\{24,\allowbreak{}36,\allowbreak{}48\}.\allowbreak{}\{1.\allowbreak{}1.\allowbreak{}\{1--18\};\allowbreak{}2.\allowbreak{}1.\allowbreak{}\{1,\allowbreak{}2\};\allowbreak{}\{3,\allowbreak{}4\}.\allowbreak{}1.\allowbreak{}1;\allowbreak{}5.\allowbreak{}1.\allowbreak{}\{1--3\}\}}\newline{}\textbf{F20 · Formación estacional por campaña y cohorte activada}\newline{}José Lara Arce\newline{}75 identificadores reservados; cantidad activa según manifiesto.\newline{}\textbf{Ventana:} 24,36,48. & \textbf{Contenido y frontera.} Formación, evaluación, recuperación y certificación de cada cohorte estacional activada. Manifiesto anual/campaña reúne población nueva y monitores, perfiles, turnos, sitios y modalidades antes del 15 de diciembre. Las 25 posiciones por campaña (18/2/1/1/3) son sufijos históricos disponibles, no una población ni un reparto de ejecución aprobados.\newline{}\textbf{Aceptación.} Nómina autorizada y población objetivo tienen fuente/versión; cohortes se calculan por grupos compatibles y capacidad efectiva, con límite de doce asistentes, sesiones y objetivos según plan. Cada persona objetivo tiene formación/evaluación y recuperación identificadas, sin duplicación injustificada. Se conserva la capacidad ofrecida de 300 plazas anuales y ampliación sin costo si se requiere; plazas disponibles no acreditan cohortes impartidas. Reserva sin activación no constituye costo de ejecución, avance ni recepción. Formación finaliza antes del congelamiento.\newline{}\textbf{Insumos.} Manifiesto estacional validado por responsables del cliente, nuevas incorporaciones/monitores y turnos, material/versión vigentes, instructores y ventanas admisibles. José solicita/actualiza población; contraparte valida calendario. Un dato material pendiente bloquea cerrar cantidad/HH, no elimina el derecho a formación.\newline{}\textbf{Condiciones.} SERVICIO-ACTIVO y manifiesto de la campaña recibidos antes de ejecutar. El control anual aplica la fórmula y reglas de la tabla de cohortes; campañas 24/36/48 no limitan formación de otras altas del contrato. Si surge demanda fuera de esos sufijos, se amplía el registro y calendario versionados sin cargo adicional. Dos jornadas anuales de actualización no sustituyen incorporación de personal nuevo. RT-22.04/06; FALTA REFERENCIA A SD10; FALTA REFERENCIA A SD7.\\\hline
\texttt{1.\allowbreak{}12.\allowbreak{}11.\allowbreak{}1}\newline{}\textbf{Actualización semestral de custodia de fuentes}\newline{}Ignacio Silva\newline{}1 instancia.\newline{}\textbf{Ventana:} 24. & \textbf{Contenido y frontera.} Depósito de fuentes y manifiesto ante custodio independiente; excluye reconstruir fuentes no versionadas.\newline{}\textbf{Aceptación.} Custodio confirma versión, integridad y condiciones de acceso/liberación y se conserva comprobante. Ignacio Silva responde por el paquete.\newline{}\textbf{Condiciones.} Predecesoras del catálogo: 1.1.3.2. Referencia: BA custodia independiente; actualización semestral. Las habilitaciones técnicas se contrastan con arquitectura e inventario vigentes.\\\hline
\texttt{1.\allowbreak{}12.\allowbreak{}11.\allowbreak{}\{2--6\}}\newline{}\textbf{Actualización semestral de custodia de fuentes}\newline{}Ignacio Silva\newline{}5 instancias.\newline{}\textbf{Ventana:} 30,36,42,48,54. & \textbf{Contenido y frontera.} Resultado individual del período o conjunto indicado. Excluye resultados de otras cuentas y horas de contraparte, obras o adquisiciones del cliente.\newline{}\textbf{Aceptación.} Custodio confirma versión, integridad y condiciones de acceso/liberación y se conserva comprobante. Ignacio Silva responde por el paquete.\newline{}\textbf{Insumos.} Versión/registro recibido, acceso autorizado, herramientas, participantes y ventana del período.\newline{}\textbf{Condiciones.} Predecesoras de habilitación: 1.1.3.2.\\\hline
\texttt{1.\allowbreak{}12.\allowbreak{}12.\allowbreak{}correctiva.\allowbreak{}\{21--56\}}\newline{}\textbf{F21 · Servicio mensual de ingeniería correctiva recibido}\newline{}Nikolai Navea\newline{}36 instancias mensuales.\newline{}\textbf{Ventana:} 21--56. & \textbf{Contenido y frontera.} Servicio de ingeniería correctiva del período: disponibilidad de capacidad de referencia (80 HH mensuales según SD6), atención del backlog autorizado, producción de cambios y conciliación de consumo/saldo. El registro mensual es evidencia; el resultado comprende servicio disponible y trabajos efectivamente ejecutados. Excluye turnos de mesa/NOC/SOC, mantenimiento preventivo ya incluido y duplicar otras bolsas.\newline{}\textbf{Aceptación.} Corrección de defectos de la solución: reproducción, causa, cambio y regresión; se verifica resolución según severidad y obligaciones de BA78, sin costo adicional. La capacidad de referencia no limita garantía/corrección ni habilita cobrar por un defecto incluido. Cada trabajo conserva autorización, código de imputación único, versión, pruebas ejecutadas, revisión, retorno y decisión. Cierre mensual informa capacidad ofrecida/disponible, uso, saldo y pendientes; no recibe cambios sin terminar ni declara ejecutadas horas no usadas. Trabajo mixto se separa en componentes disjuntos aprobados, sin cargar una misma intervención a dos cuentas.\newline{}\textbf{Insumos.} SERVICIO-ACTIVO, backlog y órdenes aprobadas, criterio de clasificación, recursos/herramientas, versiones y ventanas admisibles; saldo anterior conciliado y pruebas/revisores disponibles.\newline{}\textbf{Condiciones.} Aplicar frontera mantenimiento/ingeniería definida en este formulario y trazada en SD6. T-15 distribuye actividades/esfuerzo y controla imputación por orden; cambios que alteren alcance se aprueban antes de iniciar. Recepción del período y recepción de cada cambio se distinguen. Mantener servicio y obligaciones aun sin cambios autorizados, con evidencia de capacidad disponible; no se inventan tareas para consumir reserva. FALTA REFERENCIA A SD6; FALTA REFERENCIA A SD11.\\\hline
\texttt{1.\allowbreak{}12.\allowbreak{}12.\allowbreak{}deuda.\allowbreak{}\{21--56\}}\newline{}\textbf{F24 · Servicio mensual de ingeniería deuda recibido}\newline{}Nikolai Navea\newline{}36 instancias mensuales.\newline{}\textbf{Ventana:} 21--56. & \textbf{Contenido y frontera.} Servicio de ingeniería deuda del período: disponibilidad de capacidad de referencia (32 HH mensuales según SD6), atención del backlog autorizado, producción de cambios y conciliación de consumo/saldo. El registro mensual es evidencia; el resultado comprende servicio disponible y trabajos efectivamente ejecutados. Excluye turnos de mesa/NOC/SOC, mantenimiento preventivo ya incluido y duplicar otras bolsas.\newline{}\textbf{Aceptación.} Reducción de deuda interna sin agregar función: causa/riesgo, horas, componente, objetivo medible y verificación de reducción. Solo trabajo de esta bolsa integra su denominador de deuda; las acciones internas del mantenimiento de 24 HH se registran separadas. Cada trabajo conserva autorización, código de imputación único, versión, pruebas ejecutadas, revisión, retorno y decisión. Cierre mensual informa capacidad ofrecida/disponible, uso, saldo y pendientes; no recibe cambios sin terminar ni declara ejecutadas horas no usadas. Trabajo mixto se separa en componentes disjuntos aprobados, sin cargar una misma intervención a dos cuentas.\newline{}\textbf{Insumos.} SERVICIO-ACTIVO, backlog y órdenes aprobadas, criterio de clasificación, recursos/herramientas, versiones y ventanas admisibles; saldo anterior conciliado y pruebas/revisores disponibles.\newline{}\textbf{Condiciones.} Aplicar frontera mantenimiento/ingeniería definida en este formulario y trazada en SD6. T-15 distribuye actividades/esfuerzo y controla imputación por orden; cambios que alteren alcance se aprueban antes de iniciar. Recepción del período y recepción de cada cambio se distinguen. Mantener servicio y obligaciones aun sin cambios autorizados, con evidencia de capacidad disponible; no se inventan tareas para consumir reserva. FALTA REFERENCIA A SD6; FALTA REFERENCIA A SD11.\\\hline
\texttt{1.\allowbreak{}12.\allowbreak{}12.\allowbreak{}evolutiva.\allowbreak{}\{21--56\}}\newline{}\textbf{F23 · Servicio mensual de ingeniería evolutiva recibido}\newline{}Nikolai Navea\newline{}36 instancias mensuales.\newline{}\textbf{Ventana:} 21--56. & \textbf{Contenido y frontera.} Servicio de ingeniería evolutiva del período: disponibilidad de capacidad de referencia (80 HH mensuales según SD6), atención del backlog autorizado, producción de cambios y conciliación de consumo/saldo. El registro mensual es evidencia; el resultado comprende servicio disponible y trabajos efectivamente ejecutados. Excluye turnos de mesa/NOC/SOC, mantenimiento preventivo ya incluido y duplicar otras bolsas.\newline{}\textbf{Aceptación.} Mejora funcional solicitada y aprobada: requisito/delta, diseño, pruebas, seguridad, documentación y aceptación con estándar del desarrollo original. Se concilia la bolsa anual ofrecida de 960 HH, consumo y saldo acumulable entre años; el cierre mensual no extingue saldo no utilizado. Excedentes requieren autorización conforme al expediente económico, sin tarifa inventada en T-14. Cada trabajo conserva autorización, código de imputación único, versión, pruebas ejecutadas, revisión, retorno y decisión. Cierre mensual informa capacidad ofrecida/disponible, uso, saldo y pendientes; no recibe cambios sin terminar ni declara ejecutadas horas no usadas. Trabajo mixto se separa en componentes disjuntos aprobados, sin cargar una misma intervención a dos cuentas.\newline{}\textbf{Insumos.} SERVICIO-ACTIVO, backlog y órdenes aprobadas, criterio de clasificación, recursos/herramientas, versiones y ventanas admisibles; saldo anterior conciliado y pruebas/revisores disponibles.\newline{}\textbf{Condiciones.} Aplicar frontera mantenimiento/ingeniería definida en este formulario y trazada en SD6. T-15 distribuye actividades/esfuerzo y controla imputación por orden; cambios que alteren alcance se aprueban antes de iniciar. Recepción del período y recepción de cada cambio se distinguen. Mantener servicio y obligaciones aun sin cambios autorizados, con evidencia de capacidad disponible; no se inventan tareas para consumir reserva. FALTA REFERENCIA A SD6; FALTA REFERENCIA A SD11.\\\hline
\texttt{1.\allowbreak{}12.\allowbreak{}12.\allowbreak{}normativa.\allowbreak{}\{21--56\}}\newline{}\textbf{F22 · Servicio mensual de ingeniería normativa recibido}\newline{}Nikolai Navea\newline{}36 instancias mensuales.\newline{}\textbf{Ventana:} 21--56. & \textbf{Contenido y frontera.} Servicio de ingeniería normativa del período: disponibilidad de capacidad de referencia (48 HH mensuales según SD6), atención del backlog autorizado, producción de cambios y conciliación de consumo/saldo. El registro mensual es evidencia; el resultado comprende servicio disponible y trabajos efectivamente ejecutados. Excluye turnos de mesa/NOC/SOC, mantenimiento preventivo ya incluido y duplicar otras bolsas.\newline{}\textbf{Aceptación.} Adecuación a cambio legal/regulatorio aplicable: fuente, entrada en vigencia, impacto, prueba y aceptación antes del plazo exigible. Se incluye en Operación sin costo adicional; las 48 HH de referencia no excusan una obligación ni amplían automáticamente el alcance funcional. Cada trabajo conserva autorización, código de imputación único, versión, pruebas ejecutadas, revisión, retorno y decisión. Cierre mensual informa capacidad ofrecida/disponible, uso, saldo y pendientes; no recibe cambios sin terminar ni declara ejecutadas horas no usadas. Trabajo mixto se separa en componentes disjuntos aprobados, sin cargar una misma intervención a dos cuentas.\newline{}\textbf{Insumos.} SERVICIO-ACTIVO, backlog y órdenes aprobadas, criterio de clasificación, recursos/herramientas, versiones y ventanas admisibles; saldo anterior conciliado y pruebas/revisores disponibles.\newline{}\textbf{Condiciones.} Aplicar frontera mantenimiento/ingeniería definida en este formulario y trazada en SD6. T-15 distribuye actividades/esfuerzo y controla imputación por orden; cambios que alteren alcance se aprueban antes de iniciar. Recepción del período y recepción de cada cambio se distinguen. Mantener servicio y obligaciones aun sin cambios autorizados, con evidencia de capacidad disponible; no se inventan tareas para consumir reserva. FALTA REFERENCIA A SD6; FALTA REFERENCIA A SD11.\\\hline
\texttt{1.\allowbreak{}12.\allowbreak{}13.\allowbreak{}21}\newline{}\textbf{Referentes y protocolo técnico transferidos}\newline{}José Lara Arce\newline{}1 instancia.\newline{}\textbf{Ventana:} 21. & \textbf{Contenido y frontera.} Mentoría mensual a referentes técnicos y suplentes; casos, transferencia y continuidad del conocimiento.\newline{}\textbf{Aceptación.} Referente y suplente; agenda/casos; acta, documentación y verificación independiente; pendientes con dueño/plazo. No sustituye mesa, estabilización ni salida.\newline{}\textbf{Insumos.} Versión/registro recibido, acceso autorizado, herramientas, participantes y ventana del período.\newline{}\textbf{Condiciones.} Predecesoras de habilitación: CORTE-E2.\\\hline
\texttt{1.\allowbreak{}12.\allowbreak{}13.\allowbreak{}22}\newline{}\textbf{Casos de recuperación y conciliación acompañados}\newline{}José Lara Arce\newline{}1 instancia.\newline{}\textbf{Ventana:} 22. & \textbf{Contenido y frontera.} Mentoría mensual a referentes técnicos y suplentes; casos, transferencia y continuidad del conocimiento.\newline{}\textbf{Aceptación.} Referente y suplente; agenda/casos; acta, documentación y verificación independiente; pendientes con dueño/plazo. No sustituye mesa, estabilización ni salida.\newline{}\textbf{Insumos.} Versión/registro recibido, acceso autorizado, herramientas, participantes y ventana del período.\newline{}\textbf{Condiciones.} Predecesoras de habilitación: 1.12.13.21.\\\hline
\texttt{1.\allowbreak{}12.\allowbreak{}13.\allowbreak{}23}\newline{}\textbf{Cambios e indicadores acompañados}\newline{}José Lara Arce\newline{}1 instancia.\newline{}\textbf{Ventana:} 23. & \textbf{Contenido y frontera.} Mentoría mensual a referentes técnicos y suplentes; casos, transferencia y continuidad del conocimiento.\newline{}\textbf{Aceptación.} Referente y suplente; agenda/casos; acta, documentación y verificación independiente; pendientes con dueño/plazo. No sustituye mesa, estabilización ni salida.\newline{}\textbf{Insumos.} Versión/registro recibido, acceso autorizado, herramientas, participantes y ventana del período.\newline{}\textbf{Condiciones.} Predecesoras de habilitación: 1.12.13.22.\\\hline
\texttt{1.\allowbreak{}12.\allowbreak{}13.\allowbreak{}24}\newline{}\textbf{Procedimientos estacionales transferidos}\newline{}José Lara Arce\newline{}1 instancia.\newline{}\textbf{Ventana:} 24. & \textbf{Contenido y frontera.} Mentoría mensual a referentes técnicos y suplentes; casos, transferencia y continuidad del conocimiento.\newline{}\textbf{Aceptación.} Referente y suplente; agenda/casos; acta, documentación y verificación independiente; pendientes con dueño/plazo. No sustituye mesa, estabilización ni salida.\newline{}\textbf{Insumos.} Versión/registro recibido, acceso autorizado, herramientas, participantes y ventana del período.\newline{}\textbf{Condiciones.} Predecesoras de habilitación: 1.12.13.23.\\\hline
\texttt{1.\allowbreak{}12.\allowbreak{}13.\allowbreak{}25}\newline{}\textbf{Casos y suplencia acompañados en congelamiento}\newline{}José Lara Arce\newline{}1 instancia.\newline{}\textbf{Ventana:} 25. & \textbf{Contenido y frontera.} Mentoría mensual a referentes técnicos y suplentes; casos, transferencia y continuidad del conocimiento.\newline{}\textbf{Aceptación.} Referente y suplente; agenda/casos; acta, documentación y verificación independiente; pendientes con dueño/plazo. No sustituye mesa, estabilización ni salida.\newline{}\textbf{Insumos.} Versión/registro recibido, acceso autorizado, herramientas, participantes y ventana del período.\newline{}\textbf{Condiciones.} Predecesoras de habilitación: 1.12.13.24.\\\hline
\texttt{1.\allowbreak{}12.\allowbreak{}13.\allowbreak{}26}\newline{}\textbf{Transferencia técnica y continuidad del conocimiento recibidas}\newline{}José Lara Arce\newline{}1 instancia.\newline{}\textbf{Ventana:} 26. & \textbf{Contenido y frontera.} Mentoría mensual a referentes técnicos y suplentes; casos, transferencia y continuidad del conocimiento.\newline{}\textbf{Aceptación.} Referente y suplente; agenda/casos; acta, documentación y verificación independiente; pendientes con dueño/plazo. No sustituye mesa, estabilización ni salida.\newline{}\textbf{Insumos.} Versión/registro recibido, acceso autorizado, herramientas, participantes y ventana del período.\newline{}\textbf{Condiciones.} Predecesoras de habilitación: 1.12.13.25.\\\hline
\texttt{1.\allowbreak{}12.\allowbreak{}14.\allowbreak{}1.\allowbreak{}\{21--56\}}\newline{}\textbf{Atención y soporte prestados: Operación}\newline{}José Lara Arce\newline{}36 instancias mensuales.\newline{}\textbf{Ventana:} 21--56. & \textbf{Contenido y frontera.} Prestación mensual efectiva de mesa bilingüe, orientación funcional/técnica, registro, clasificación, resolución o derivación y seguimiento hasta cierre. Incluye mantenimiento de la base de conocimiento, autoformación y evidencia de aprendizaje. Excluye respuesta especializada de guardia, ingeniería de cambios, acompañamiento presencial de olas y el informe agregado de SLA.\newline{}\textbf{Aceptación.} Cada mes conserva cuadrantes, personal/suplencia, canales disponibles, contactos y tickets trazables. Cobertura C07: 06:00 a 24:00 todos los días en temporada alta y 08:00 a 20:00 el resto del año, español/inglés. Se verifica 80\,\% de atención antes de 20 segundos, resolución al primer contacto de al menos 70\,\% y abandono no superior a 5\,\%, con denominadores y fallos registrados; cero contactos no prueba desempeño. Ticket tiene severidad, responsable, escalamiento y cierre confirmado o plazo automático declarado. Base de conocimiento actualizada/versionada, accesible, editable y de propiedad del cliente; cambios, uso y utilidad percibida registrados. Se acredita la prestación del período; informe favorable no sustituye cobertura ni corrige incumplimientos.\newline{}\textbf{Insumos.} Línea base/delta aplicable de 1.12.1.1 y 1.12.1.2 recibidos antes de activar cobertura, versión supervisada, herramientas/canales y dotación disponible; calendarios naturales, suplencias y permisos.\newline{}\textbf{Condiciones.} Cada mes es una instancia de prestación y se recibe con su universo completo. Un incidente conserva ticket único: mesa registra/deriva y guardia responde; cada intervención se imputa una vez. Registros diarios y turnos son actividades/evidencia de T-15, sin paquetes adicionales por persona/día. SERVICIO-ACTIVO habilita Operación de ambos alcances, meses 21--56. Gerente de Servicio coordina y José Lara Arce integra el resultado técnico; especialistas ejecutan sus ámbitos. Control mensual 1.12.3.1 evalúa, pero no reemplaza, esta prestación. RT-21.01/05--10/13/15--18; C07 RT-21.06/16; FALTA REFERENCIA A SD10; FALTA REFERENCIA A SD12.\\\hline
\texttt{1.\allowbreak{}12.\allowbreak{}14.\allowbreak{}2.\allowbreak{}\{21--56\}}\newline{}\textbf{Guardia, monitoreo y respuesta prestados: Operación}\newline{}José Lara Arce\newline{}36 instancias mensuales.\newline{}\textbf{Ventana:} 21--56. & \textbf{Contenido y frontera.} Prestación mensual efectiva de NOC/SOC, monitoreo proactivo, guardia crítica y respuesta técnica, con escalamiento a especialistas. Atiende canales de emergencia y alertas de planes vencidos 24x7x365, además de los incidentes críticos de plataforma. Excluye otorgar zarpes o ejecutar salvamento, mesa ordinaria, acompañamiento de olas y producción de cambios ya imputados a ingeniería.\newline{}\textbf{Aceptación.} Guardias y relevos cubren todos los intervalos del mes, con especialistas/suplentes y registros de disponibilidad, alertas, contacto, diagnóstico, contención, escalamiento y recuperación. Se comprueba recepción/respuesta 24x7x365 del canal de emergencia y planes vencidos, criterios de severidad y tiempos aplicables de BA78/T-12, sin dar alertas técnicas por confirmación de hechos físicos. Traslado requerido a Bahía Ilque identifica embarcación, ventana marítima segura, especialista y suplencia, sin desatender otros sitios. La bitácora conserva intervalos sin cobertura y desvíos; la recepción exige evidencia del servicio real, no solo cuadrante diseñado.\newline{}\textbf{Insumos.} Línea base/delta aplicable de 1.12.1.1 y 1.12.1.2 recibidos antes de activar cobertura, versión supervisada, herramientas/canales y dotación disponible; calendarios naturales, suplencias y permisos.\newline{}\textbf{Condiciones.} Cada mes es una instancia de prestación y se recibe con su universo completo. Un incidente conserva ticket único: mesa registra/deriva y guardia responde; cada intervención se imputa una vez. Registros diarios y turnos son actividades/evidencia de T-15, sin paquetes adicionales por persona/día. SERVICIO-ACTIVO habilita Operación de ambos alcances, meses 21--56. Gerente de Servicio coordina y José Lara Arce integra el resultado técnico; especialistas ejecutan sus ámbitos. Control mensual 1.12.3.1 evalúa, pero no reemplaza, esta prestación. RT-21.01/05--10/13/15--18; C07 RT-21.06/16; FALTA REFERENCIA A SD10; FALTA REFERENCIA A SD12.\\\hline
\end{longtable}\endgroup
\section{Trazabilidad y cobertura de requisitos}\label{sec:t14-traza}
\subsection{Universo y asignación}
El universo de T-12 comprende 576 identificadores: 146 RF, 56 RNF y 374 RT. Los 27 complementos particulares C07 se asocian a sus RT y no se suman como identificadores independientes. Enunciados, obligatoriedad, etapa, umbrales y declaración de oferta permanecen en T-12. La asignación identifica trabajo planificado; no significa cumplimiento demostrado mediante pruebas ejecutadas.

El cotejo de las celdas del T-12 v0.9 arroja 546 identificadores ofrecidos y 30 no ofrecidos: 146 RF y 56 RNF ofrecidos, y 344 RT ofrecidos de 374. RT-26.08 permanece deseable no ofrecido; no se compromete medir una innovación antes del mes 16. El registro antecedente 547/29 se sustituye por este conteo, sin modificar las cinco innovaciones ni sus ventanas. Se mantienen complementos C07 obligatorios de RT-03.24, RT-05.10 y RT-16.14 aunque la parte deseable BT tenga otro tratamiento. Una exclusión deseable no retira la obligación del caso.

\subsection{Bloques y relación con verificación}
Los bloques identifican productores o verificaciones concretas; no crean entregables ni HH adicionales. La relación distingue diseño, construcción, prueba, control periódico y evidencia académica externa. Los códigos reservados de formación solo pertenecen al universo activo cuando constan en el manifiesto recibido; los conteos de bloques son nominales. Para cada requisito se comprueba la cláusula específica y el alcance BT/C07 vigente en T-12, el resultado, su universo/versiones y su aceptación. Una prueba complementaria no sustituye al productor. FALTA REFERENCIA A SD9 para los localizadores de prueba particulares. Las variantes de un bloque conservan su ID en esta versión, pero sustituyen expresamente su composición anterior. B90 identifica prestación de atención, B91 guardia/monitoreo/respuesta y B92 creación/mantenimiento de conocimiento; sus registros verifican trabajo real. Los 75 códigos estacionales son reservas y los cuatro paquetes agregados sustituyen únicamente el nivel terminal diario; ninguna relación de bloque revive los registros históricos como entregables.
\begingroup\setlength{\tabcolsep}{3pt}
\begin{longtable}{L{\dimexpr0.100000\textwidth-2\tabcolsep\relax}L{\dimexpr0.800000\textwidth-2\tabcolsep\relax}L{\dimexpr0.100000\textwidth-2\tabcolsep\relax}}
\caption{Bloques de paquetes para trazabilidad}\label{tab:t14-bloques}\\\hline
\EncabezadoTabla{Bloque} & \EncabezadoTabla{Códigos nominales} & \EncabezadoTabla{N.º}\\\hline\endfirsthead
\multicolumn{3}{l}{\textit{Continuación de la tabla \ref{tab:t14-bloques}}}\\\hline
\EncabezadoTabla{Bloque} & \EncabezadoTabla{Códigos nominales} & \EncabezadoTabla{N.º}\\\hline\endhead
\hline\endfoot
\hline\endlastfoot
\texttt{B01} & \texttt{1.\allowbreak{}4.\allowbreak{}1.\allowbreak{}\{1--3\}.\allowbreak{}1} & 3\\\hline
\texttt{B02} & \texttt{1.\allowbreak{}7.\allowbreak{}4.\allowbreak{}1} & 1\\\hline
\texttt{B03} & \texttt{1.\allowbreak{}4.\allowbreak{}2.\allowbreak{}\{1--3\}.\allowbreak{}1} & 3\\\hline
\texttt{B04} & \texttt{1.\allowbreak{}6.\allowbreak{}1.\allowbreak{}\{1--4\}.\allowbreak{}1} & 4\\\hline
\texttt{B05} & \texttt{1.\allowbreak{}5.\allowbreak{}1.\allowbreak{}\{1--5\}.\allowbreak{}1} & 5\\\hline
\texttt{B06} & \texttt{1.\allowbreak{}5.\allowbreak{}1.\allowbreak{}5.\allowbreak{}1} & 1\\\hline
\texttt{B07} & \texttt{1.\allowbreak{}5.\allowbreak{}3.\allowbreak{}\{1--5\}.\allowbreak{}1} & 5\\\hline
\texttt{B08} & \texttt{1.\allowbreak{}5.\allowbreak{}3.\allowbreak{}\{1--5\}.\allowbreak{}1} & 5\\\hline
\texttt{B09} & \texttt{1.\allowbreak{}4.\allowbreak{}3.\allowbreak{}\{1--3\}.\allowbreak{}1} & 3\\\hline
\texttt{B10} & \texttt{1.\allowbreak{}6.\allowbreak{}2.\allowbreak{}\{1--4\}.\allowbreak{}1} & 4\\\hline
\texttt{B11} & \texttt{1.\allowbreak{}6.\allowbreak{}\{1,\allowbreak{}2\}.\allowbreak{}\{1--4\}.\allowbreak{}1} & 8\\\hline
\texttt{B12} & \texttt{1.\allowbreak{}7.\allowbreak{}3.\allowbreak{}\{1--3\}} & 3\\\hline
\texttt{B13} & \texttt{1.\allowbreak{}10.\allowbreak{}2.\allowbreak{}\{1,\allowbreak{}2\}.\allowbreak{}\{1--10\}} & 20\\\hline
\texttt{B14} & \texttt{1.\allowbreak{}12.\allowbreak{}4.\allowbreak{}\{1--9\}.\allowbreak{}\{6,\allowbreak{}7\}} & 18\\\hline
\texttt{B15} & \texttt{1.\allowbreak{}10.\allowbreak{}2.\allowbreak{}\{1,\allowbreak{}2\}.\allowbreak{}\{1--10\}} & 20\\\hline
\texttt{B16} & \texttt{1.\allowbreak{}10.\allowbreak{}2.\allowbreak{}\{1,\allowbreak{}2\}.\allowbreak{}\{1--10\}} & 20\\\hline
\texttt{B17} & \texttt{1.\allowbreak{}3.\allowbreak{}1.\allowbreak{}\{1--3\}} & 3\\\hline
\texttt{B18} & \texttt{1.\allowbreak{}10.\allowbreak{}2.\allowbreak{}\{1,\allowbreak{}2\}.\allowbreak{}\{1--10\}} & 20\\\hline
\texttt{B19} & \texttt{1.\allowbreak{}7.\allowbreak{}\{1.\allowbreak{}\{1--3\};\allowbreak{}4.\allowbreak{}\{1--4\}\}} & 7\\\hline
\texttt{B20} & \texttt{1.\allowbreak{}12.\allowbreak{}\{1.\allowbreak{}\{1,\allowbreak{}2\};\allowbreak{}2.\allowbreak{}\{\{1,\allowbreak{}3\}.\allowbreak{}\{21--56\}.\allowbreak{}\{1,\allowbreak{}2\};\allowbreak{}2.\allowbreak{}\{31--56\};\allowbreak{}\{4,\allowbreak{}5\}.\allowbreak{}\{21--56\}\};\allowbreak{}3.\allowbreak{}\{1--5\}.\allowbreak{}\{21--56\};\allowbreak{}4.\allowbreak{}\{1--9\}.\allowbreak{}\{6,\allowbreak{}7\};\allowbreak{}\{5,\allowbreak{}7\}.\allowbreak{}\{1--3\};\allowbreak{}\{6,\allowbreak{}11\}.\allowbreak{}\{1--6\};\allowbreak{}8.\allowbreak{}\{1--5\};\allowbreak{}9.\allowbreak{}\{53--56\}.\allowbreak{}\{1,\allowbreak{}2\};\allowbreak{}10.\allowbreak{}\{24,\allowbreak{}36,\allowbreak{}48\}.\allowbreak{}\{1.\allowbreak{}1.\allowbreak{}\{1--18\};\allowbreak{}2.\allowbreak{}1.\allowbreak{}\{1,\allowbreak{}2\};\allowbreak{}\{3,\allowbreak{}4\}.\allowbreak{}1.\allowbreak{}1;\allowbreak{}5.\allowbreak{}1.\allowbreak{}\{1--3\}\};\allowbreak{}12.\allowbreak{}\{correctiva,\allowbreak{}deuda,\allowbreak{}evolutiva,\allowbreak{}normativa\}.\allowbreak{}\{21--56\}\}} & 692\\\hline
\texttt{B21} & \texttt{1.\allowbreak{}8.\allowbreak{}5.\allowbreak{}\{1--3\}} & 3\\\hline
\texttt{B22} & \texttt{1.\allowbreak{}12.\allowbreak{}4.\allowbreak{}\{1--9\}.\allowbreak{}\{6,\allowbreak{}7\}} & 18\\\hline
\texttt{B23} & \texttt{1.\allowbreak{}10.\allowbreak{}1.\allowbreak{}\{1,\allowbreak{}2\}.\allowbreak{}\{1,\allowbreak{}2\}.\allowbreak{}\{1--5\}} & 20\\\hline
\texttt{B24} & \texttt{1.\allowbreak{}2.\allowbreak{}3.\allowbreak{}\{\{1--5\};\allowbreak{}7.\allowbreak{}\{1--10\}\}} & 15\\\hline
\texttt{B25} & \texttt{1.\allowbreak{}12.\allowbreak{}9.\allowbreak{}\{53--56\}.\allowbreak{}\{1,\allowbreak{}2\}} & 8\\\hline
\texttt{B26} & \texttt{1.\allowbreak{}11.\allowbreak{}\{1.\allowbreak{}\{\{1--3\}.\allowbreak{}\{1--3\};\allowbreak{}EST.\allowbreak{}\{1,\allowbreak{}2\}.\allowbreak{}\{1,\allowbreak{}2\};\allowbreak{}MB.\allowbreak{}\{1.\allowbreak{}\{13--15\}.\allowbreak{}\{1,\allowbreak{}2\};\allowbreak{}2.\allowbreak{}\{19,\allowbreak{}20\}.\allowbreak{}\{1,\allowbreak{}2\}\}\};\allowbreak{}2.\allowbreak{}\{1,\allowbreak{}2\}.\allowbreak{}\{1--5\};\allowbreak{}\{3,\allowbreak{}5\}.\allowbreak{}\{1,\allowbreak{}2\}.\allowbreak{}\{1,\allowbreak{}2\};\allowbreak{}\{4,\allowbreak{}6\}.\allowbreak{}\{1,\allowbreak{}2\};\allowbreak{}7.\allowbreak{}\{1,\allowbreak{}2\}.\allowbreak{}\{1.\allowbreak{}1.\allowbreak{}\{1--100\};\allowbreak{}2.\allowbreak{}1.\allowbreak{}\{1,\allowbreak{}2\};\allowbreak{}\{3,\allowbreak{}4\}.\allowbreak{}1.\allowbreak{}1;\allowbreak{}5.\allowbreak{}1.\allowbreak{}\{1--3\}\}\}} & 259\\\hline
\texttt{B27} & \texttt{1.\allowbreak{}12.\allowbreak{}10.\allowbreak{}\{24,\allowbreak{}36,\allowbreak{}48\}.\allowbreak{}\{1.\allowbreak{}1.\allowbreak{}\{1--18\};\allowbreak{}2.\allowbreak{}1.\allowbreak{}\{1,\allowbreak{}2\};\allowbreak{}\{3,\allowbreak{}4\}.\allowbreak{}1.\allowbreak{}1;\allowbreak{}5.\allowbreak{}1.\allowbreak{}\{1--3\}\}} & 75\\\hline
\texttt{B28} & \texttt{1.\allowbreak{}3.\allowbreak{}\{\{1,\allowbreak{}4\}.\allowbreak{}\{1--3\};\allowbreak{}\{2,\allowbreak{}3\}.\allowbreak{}\{1,\allowbreak{}2\}\}} & 10\\\hline
\texttt{B29} & \texttt{1.\allowbreak{}2.\allowbreak{}\{\{1,\allowbreak{}2\}.\allowbreak{}\{1--5\};\allowbreak{}3.\allowbreak{}\{\{1--5\};\allowbreak{}7.\allowbreak{}\{1--10\}\};\allowbreak{}4.\allowbreak{}\{1--3\}\}} & 28\\\hline
\texttt{B30} & \texttt{1.\allowbreak{}2.\allowbreak{}\{2.\allowbreak{}\{1--5\};\allowbreak{}4.\allowbreak{}\{1--3\}\}} & 8\\\hline
\texttt{B31} & \texttt{1.\allowbreak{}8.\allowbreak{}\{\{1,\allowbreak{}8\}.\allowbreak{}\{1--4\};\allowbreak{}\{2--4\}.\allowbreak{}\{1--6\};\allowbreak{}\{5,\allowbreak{}6\}.\allowbreak{}\{1--3\};\allowbreak{}7.\allowbreak{}\{1.\allowbreak{}\{1--14\};\allowbreak{}2.\allowbreak{}\{1--56\}\}\}} & 102\\\hline
\texttt{B32} & \texttt{1.\allowbreak{}\{2.\allowbreak{}3.\allowbreak{}\{\{1--5\};\allowbreak{}7.\allowbreak{}\{1--10\}\};\allowbreak{}12.\allowbreak{}3.\allowbreak{}5.\allowbreak{}\{21--56\}\}} & 51\\\hline
\texttt{B33} & \texttt{1.\allowbreak{}\{2.\allowbreak{}\{2.\allowbreak{}5;\allowbreak{}4.\allowbreak{}\{2,\allowbreak{}3\}\};\allowbreak{}3.\allowbreak{}4.\allowbreak{}1;\allowbreak{}10.\allowbreak{}2.\allowbreak{}\{1,\allowbreak{}2\}.\allowbreak{}7;\allowbreak{}12.\allowbreak{}\{3.\allowbreak{}2.\allowbreak{}\{21--56\};\allowbreak{}8.\allowbreak{}\{1--5\}\}\}} & 47\\\hline
\texttt{B34} & \texttt{1.\allowbreak{}\{1.\allowbreak{}5.\allowbreak{}4;\allowbreak{}2.\allowbreak{}3.\allowbreak{}\{\{1--5\};\allowbreak{}7.\allowbreak{}\{1--10\}\};\allowbreak{}12.\allowbreak{}8.\allowbreak{}\{1--5\}\}} & 21\\\hline
\texttt{B35} & \texttt{1.\allowbreak{}12.\allowbreak{}\{3.\allowbreak{}\{1--5\}.\allowbreak{}\{21--56\};\allowbreak{}4.\allowbreak{}\{1--9\}.\allowbreak{}\{6,\allowbreak{}7\};\allowbreak{}8.\allowbreak{}\{1--5\}\}} & 203\\\hline
\texttt{B36} & \texttt{1.\allowbreak{}\{2.\allowbreak{}1.\allowbreak{}\{2,\allowbreak{}4,\allowbreak{}5\};\allowbreak{}3.\allowbreak{}\{1.\allowbreak{}\{1--3\};\allowbreak{}\{2,\allowbreak{}3\}.\allowbreak{}\{1,\allowbreak{}2\};\allowbreak{}4.\allowbreak{}\{1--4\}\};\allowbreak{}10.\allowbreak{}2.\allowbreak{}\{1,\allowbreak{}2\}.\allowbreak{}8;\allowbreak{}12.\allowbreak{}\{1.\allowbreak{}\{1,\allowbreak{}2\};\allowbreak{}3.\allowbreak{}3.\allowbreak{}\{21--56\};\allowbreak{}5.\allowbreak{}\{1--3\}\}\}} & 57\\\hline
\texttt{B37} & \texttt{1.\allowbreak{}\{3.\allowbreak{}\{1.\allowbreak{}\{2,\allowbreak{}3\};\allowbreak{}2.\allowbreak{}\{1,\allowbreak{}2\};\allowbreak{}3.\allowbreak{}1\};\allowbreak{}7.\allowbreak{}\{3.\allowbreak{}\{1,\allowbreak{}3\};\allowbreak{}4.\allowbreak{}2\}\}} & 8\\\hline
\texttt{B38} & \texttt{1.\allowbreak{}\{2.\allowbreak{}4.\allowbreak{}1;\allowbreak{}3.\allowbreak{}3.\allowbreak{}1;\allowbreak{}12.\allowbreak{}\{1.\allowbreak{}2;\allowbreak{}3.\allowbreak{}1.\allowbreak{}\{21--56\}\}\}} & 39\\\hline
\texttt{B39} & \texttt{1.\allowbreak{}\{2.\allowbreak{}\{2.\allowbreak{}\{1--3\};\allowbreak{}3.\allowbreak{}1\};\allowbreak{}12.\allowbreak{}3.\allowbreak{}5.\allowbreak{}\{21--56\}\}} & 40\\\hline
\texttt{B40} & \texttt{1.\allowbreak{}7.\allowbreak{}3.\allowbreak{}\{1--3\}} & 3\\\hline
\texttt{B41} & \texttt{1.\allowbreak{}7.\allowbreak{}\{2.\allowbreak{}\{\{1--3\};\allowbreak{}5.\allowbreak{}\{1--16\}\};\allowbreak{}3.\allowbreak{}\{1--3\}\}} & 22\\\hline
\texttt{B42} & \texttt{1.\allowbreak{}9.\allowbreak{}3.\allowbreak{}\{\{1,\allowbreak{}2,\allowbreak{}4\}.\allowbreak{}\{1,\allowbreak{}2\};\allowbreak{}3.\allowbreak{}\{1--3\};\allowbreak{}5.\allowbreak{}\{\{1,\allowbreak{}3\};\allowbreak{}2.\allowbreak{}\{19,\allowbreak{}20\}.\allowbreak{}\{1,\allowbreak{}2\}\};\allowbreak{}6.\allowbreak{}1\}} & 16\\\hline
\texttt{B43} & \texttt{1.\allowbreak{}1.\allowbreak{}\{1.\allowbreak{}\{1--3\};\allowbreak{}2.\allowbreak{}\{1,\allowbreak{}2\};\allowbreak{}3.\allowbreak{}\{\{1,\allowbreak{}2,\allowbreak{}4\};\allowbreak{}3.\allowbreak{}\{12,\allowbreak{}18\}\};\allowbreak{}4.\allowbreak{}\{1--20\}.\allowbreak{}\{1,\allowbreak{}2\};\allowbreak{}5.\allowbreak{}\{1--6\}\}} & 56\\\hline
\texttt{B44} & \texttt{1.\allowbreak{}11.\allowbreak{}\{1.\allowbreak{}\{\{1--3\}.\allowbreak{}\{1--3\};\allowbreak{}EST.\allowbreak{}\{1,\allowbreak{}2\}.\allowbreak{}\{1,\allowbreak{}2\};\allowbreak{}MB.\allowbreak{}\{1.\allowbreak{}\{13--15\}.\allowbreak{}\{1,\allowbreak{}2\};\allowbreak{}2.\allowbreak{}\{19,\allowbreak{}20\}.\allowbreak{}\{1,\allowbreak{}2\}\}\};\allowbreak{}2.\allowbreak{}\{1,\allowbreak{}2\}.\allowbreak{}\{1--5\};\allowbreak{}\{3,\allowbreak{}5\}.\allowbreak{}\{1,\allowbreak{}2\}.\allowbreak{}\{1,\allowbreak{}2\};\allowbreak{}\{4,\allowbreak{}6\}.\allowbreak{}\{1,\allowbreak{}2\};\allowbreak{}7.\allowbreak{}\{1,\allowbreak{}2\}.\allowbreak{}\{1.\allowbreak{}1.\allowbreak{}\{1--100\};\allowbreak{}2.\allowbreak{}1.\allowbreak{}\{1,\allowbreak{}2\};\allowbreak{}\{3,\allowbreak{}4\}.\allowbreak{}1.\allowbreak{}1;\allowbreak{}5.\allowbreak{}1.\allowbreak{}\{1--3\}\}\}} & 259\\\hline
\texttt{B45} & \texttt{1.\allowbreak{}12.\allowbreak{}\{1.\allowbreak{}\{1,\allowbreak{}2\};\allowbreak{}3.\allowbreak{}\{1--5\}.\allowbreak{}\{21--56\}\}} & 182\\\hline
\texttt{B46} & \texttt{1.\allowbreak{}12.\allowbreak{}1.\allowbreak{}\{1,\allowbreak{}2\}} & 2\\\hline
\texttt{B47} & \texttt{1.\allowbreak{}12.\allowbreak{}\{1.\allowbreak{}\{1,\allowbreak{}2\};\allowbreak{}12.\allowbreak{}\{correctiva,\allowbreak{}deuda,\allowbreak{}evolutiva,\allowbreak{}normativa\}.\allowbreak{}\{21--56\}\}} & 146\\\hline
\texttt{B48} & \texttt{1.\allowbreak{}12.\allowbreak{}1.\allowbreak{}\{1,\allowbreak{}2\}} & 2\\\hline
\texttt{B49} & \texttt{1.\allowbreak{}12.\allowbreak{}12.\allowbreak{}evolutiva.\allowbreak{}\{21--56\}} & 36\\\hline
\texttt{B50} & \texttt{1.\allowbreak{}12.\allowbreak{}12.\allowbreak{}evolutiva.\allowbreak{}\{21--56\}} & 36\\\hline
\texttt{B51} & \texttt{1.\allowbreak{}12.\allowbreak{}12.\allowbreak{}deuda.\allowbreak{}\{21--56\}} & 36\\\hline
\texttt{B52} & \texttt{1.\allowbreak{}12.\allowbreak{}\{1.\allowbreak{}2;\allowbreak{}3.\allowbreak{}4.\allowbreak{}\{21--56\}\}} & 37\\\hline
\texttt{B53} & \texttt{1.\allowbreak{}11.\allowbreak{}\{2.\allowbreak{}\{1,\allowbreak{}2\}.\allowbreak{}\{1--5\};\allowbreak{}7.\allowbreak{}\{1,\allowbreak{}2\}.\allowbreak{}\{1.\allowbreak{}1.\allowbreak{}\{1--100\};\allowbreak{}2.\allowbreak{}1.\allowbreak{}\{1,\allowbreak{}2\};\allowbreak{}\{3,\allowbreak{}4\}.\allowbreak{}1.\allowbreak{}1;\allowbreak{}5.\allowbreak{}1.\allowbreak{}\{1--3\}\}\}} & 224\\\hline
\texttt{B54} & \texttt{1.\allowbreak{}11.\allowbreak{}\{2.\allowbreak{}\{1,\allowbreak{}2\}.\allowbreak{}\{1--5\};\allowbreak{}4.\allowbreak{}\{1,\allowbreak{}2\};\allowbreak{}7.\allowbreak{}\{1,\allowbreak{}2\}.\allowbreak{}\{1.\allowbreak{}1.\allowbreak{}\{1--100\};\allowbreak{}2.\allowbreak{}1.\allowbreak{}\{1,\allowbreak{}2\};\allowbreak{}\{3,\allowbreak{}4\}.\allowbreak{}1.\allowbreak{}1;\allowbreak{}5.\allowbreak{}1.\allowbreak{}\{1--3\}\}\}} & 226\\\hline
\texttt{B55} & \texttt{1.\allowbreak{}11.\allowbreak{}2.\allowbreak{}\{1,\allowbreak{}2\}.\allowbreak{}\{1--5\}} & 10\\\hline
\texttt{B56} & \texttt{1.\allowbreak{}11.\allowbreak{}\{3.\allowbreak{}\{1,\allowbreak{}2\}.\allowbreak{}\{1,\allowbreak{}2\};\allowbreak{}7.\allowbreak{}\{1,\allowbreak{}2\}.\allowbreak{}\{1.\allowbreak{}1.\allowbreak{}\{1--100\};\allowbreak{}2.\allowbreak{}1.\allowbreak{}\{1,\allowbreak{}2\};\allowbreak{}\{3,\allowbreak{}4\}.\allowbreak{}1.\allowbreak{}1;\allowbreak{}5.\allowbreak{}1.\allowbreak{}\{1--3\}\}\}} & 218\\\hline
\texttt{B57} & \texttt{1.\allowbreak{}12.\allowbreak{}\{6.\allowbreak{}\{1--6\};\allowbreak{}10.\allowbreak{}\{24,\allowbreak{}36,\allowbreak{}48\}.\allowbreak{}\{1.\allowbreak{}1.\allowbreak{}\{1--18\};\allowbreak{}2.\allowbreak{}1.\allowbreak{}\{1,\allowbreak{}2\};\allowbreak{}\{3,\allowbreak{}4\}.\allowbreak{}1.\allowbreak{}1;\allowbreak{}5.\allowbreak{}1.\allowbreak{}\{1--3\}\}\}} & 81\\\hline
\texttt{B58} & \texttt{1.\allowbreak{}12.\allowbreak{}13.\allowbreak{}\{21--26\}} & 6\\\hline
\texttt{B59} & \texttt{1.\allowbreak{}12.\allowbreak{}\{3.\allowbreak{}\{1--5\}.\allowbreak{}\{21--56\};\allowbreak{}6.\allowbreak{}\{1--6\}\}} & 186\\\hline
\texttt{B60} & \texttt{1.\allowbreak{}1.\allowbreak{}5.\allowbreak{}\{1--6\}} & 6\\\hline
\texttt{B61} & \texttt{1.\allowbreak{}7.\allowbreak{}7.\allowbreak{}\{1,\allowbreak{}2\}} & 2\\\hline
\texttt{B62} & \texttt{1.\allowbreak{}9.\allowbreak{}\{1.\allowbreak{}\{\{1--3\}.\allowbreak{}\{1--3\};\allowbreak{}4.\allowbreak{}\{1,\allowbreak{}2\};\allowbreak{}5.\allowbreak{}\{1--6\};\allowbreak{}6.\allowbreak{}1\};\allowbreak{}2.\allowbreak{}\{\{1,\allowbreak{}2,\allowbreak{}4\}.\allowbreak{}\{1,\allowbreak{}2\};\allowbreak{}3.\allowbreak{}\{1--3\};\allowbreak{}5.\allowbreak{}\{\{1,\allowbreak{}2,\allowbreak{}6\};\allowbreak{}3.\allowbreak{}\{0--2\}.\allowbreak{}\{1,\allowbreak{}2\};\allowbreak{}4.\allowbreak{}\{1--5\};\allowbreak{}5.\allowbreak{}\{3--5\}.\allowbreak{}\{1,\allowbreak{}2\}\};\allowbreak{}6.\allowbreak{}1\};\allowbreak{}3.\allowbreak{}\{\{1,\allowbreak{}2,\allowbreak{}4\}.\allowbreak{}\{1,\allowbreak{}2\};\allowbreak{}3.\allowbreak{}\{1--3\};\allowbreak{}5.\allowbreak{}\{\{1,\allowbreak{}3\};\allowbreak{}2.\allowbreak{}\{19,\allowbreak{}20\}.\allowbreak{}\{1,\allowbreak{}2\}\};\allowbreak{}6.\allowbreak{}1\};\allowbreak{}4.\allowbreak{}\{1.\allowbreak{}\{1,\allowbreak{}2\};\allowbreak{}2.\allowbreak{}\{1--4\};\allowbreak{}3.\allowbreak{}\{\{1,\allowbreak{}2\}.\allowbreak{}\{1,\allowbreak{}2\};\allowbreak{}3\};\allowbreak{}4.\allowbreak{}1\};\allowbreak{}5.\allowbreak{}\{1.\allowbreak{}\{1,\allowbreak{}2\};\allowbreak{}2.\allowbreak{}\{1--4\};\allowbreak{}3.\allowbreak{}\{\{1,\allowbreak{}4\};\allowbreak{}\{2,\allowbreak{}3\}.\allowbreak{}\{1,\allowbreak{}2\}\};\allowbreak{}4.\allowbreak{}1\}\}} & 89\\\hline
\texttt{B63} & \texttt{1.\allowbreak{}7.\allowbreak{}6.\allowbreak{}1} & 1\\\hline
\texttt{B64} & \texttt{1.\allowbreak{}7.\allowbreak{}6.\allowbreak{}2} & 1\\\hline
\texttt{B65} & \texttt{1.\allowbreak{}7.\allowbreak{}6.\allowbreak{}3} & 1\\\hline
\texttt{B66} & \texttt{1.\allowbreak{}7.\allowbreak{}6.\allowbreak{}4} & 1\\\hline
\texttt{B67} & \texttt{1.\allowbreak{}8.\allowbreak{}1.\allowbreak{}5} & 1\\\hline
\texttt{B68} & \texttt{1.\allowbreak{}\{4.\allowbreak{}\{1.\allowbreak{}\{1--4\}.\allowbreak{}1;\allowbreak{}\{2,\allowbreak{}3\}.\allowbreak{}\{1--3\}.\allowbreak{}1\};\allowbreak{}5.\allowbreak{}\{\{1,\allowbreak{}3\}.\allowbreak{}\{1--5\}.\allowbreak{}1;\allowbreak{}2.\allowbreak{}\{1,\allowbreak{}2\}.\allowbreak{}1\};\allowbreak{}6.\allowbreak{}\{1,\allowbreak{}2\}.\allowbreak{}\{1--4\}.\allowbreak{}1\}} & 30\\\hline
\texttt{B69} & \texttt{1.\allowbreak{}\{2.\allowbreak{}1.\allowbreak{}3;\allowbreak{}7.\allowbreak{}6.\allowbreak{}7\}} & 2\\\hline
\texttt{B70} & \texttt{1.\allowbreak{}3.\allowbreak{}3.\allowbreak{}1} & 1\\\hline
\texttt{B71} & \texttt{1.\allowbreak{}3.\allowbreak{}\{1.\allowbreak{}1;\allowbreak{}3.\allowbreak{}2;\allowbreak{}4.\allowbreak{}3\}} & 3\\\hline
\texttt{B72} & \texttt{1.\allowbreak{}7.\allowbreak{}1.\allowbreak{}\{1--3\}} & 3\\\hline
\texttt{B73} & \texttt{1.\allowbreak{}7.\allowbreak{}4.\allowbreak{}\{1--4\}} & 4\\\hline
\texttt{B74} & \texttt{1.\allowbreak{}10.\allowbreak{}1.\allowbreak{}\{1,\allowbreak{}2\}.\allowbreak{}\{1,\allowbreak{}2\}.\allowbreak{}\{1--5\}} & 20\\\hline
\texttt{B75} & \texttt{1.\allowbreak{}\{8.\allowbreak{}4.\allowbreak{}\{1--6\};\allowbreak{}11.\allowbreak{}1.\allowbreak{}\{1--3\}.\allowbreak{}3\}} & 9\\\hline
\texttt{B76} & \texttt{1.\allowbreak{}7.\allowbreak{}6.\allowbreak{}5} & 1\\\hline
\texttt{B77} & \texttt{1.\allowbreak{}7.\allowbreak{}6.\allowbreak{}6} & 1\\\hline
\texttt{B78} & \texttt{1.\allowbreak{}3.\allowbreak{}4.\allowbreak{}4} & 1\\\hline
\texttt{B79} & \texttt{1.\allowbreak{}2.\allowbreak{}3.\allowbreak{}3} & 1\\\hline
\texttt{B80} & \texttt{1.\allowbreak{}7.\allowbreak{}\{1,\allowbreak{}4\}.\allowbreak{}1} & 2\\\hline
\texttt{B81} & \texttt{1.\allowbreak{}7.\allowbreak{}8.\allowbreak{}1} & 1\\\hline
\texttt{B82} & \texttt{1.\allowbreak{}\{7.\allowbreak{}\{3.\allowbreak{}\{1,\allowbreak{}3\};\allowbreak{}7.\allowbreak{}\{1,\allowbreak{}2\}\};\allowbreak{}10.\allowbreak{}2.\allowbreak{}\{1,\allowbreak{}2\}.\allowbreak{}9\}} & 6\\\hline
\texttt{B83} & \texttt{1.\allowbreak{}10.\allowbreak{}2.\allowbreak{}\{1,\allowbreak{}2\}.\allowbreak{}10} & 2\\\hline
\texttt{B84} & \texttt{1.\allowbreak{}7.\allowbreak{}4.\allowbreak{}1} & 1\\\hline
\texttt{B85} & \texttt{1.\allowbreak{}\{3.\allowbreak{}3.\allowbreak{}1;\allowbreak{}4.\allowbreak{}\{1.\allowbreak{}\{1--4\}.\allowbreak{}1;\allowbreak{}\{2,\allowbreak{}3\}.\allowbreak{}\{1--3\}.\allowbreak{}1\};\allowbreak{}5.\allowbreak{}\{\{1,\allowbreak{}3\}.\allowbreak{}\{1--5\}.\allowbreak{}1;\allowbreak{}2.\allowbreak{}\{1,\allowbreak{}2\}.\allowbreak{}1\};\allowbreak{}6.\allowbreak{}\{1,\allowbreak{}2\}.\allowbreak{}\{1--4\}.\allowbreak{}1\}} & 31\\\hline
\texttt{B86} & \texttt{1.\allowbreak{}\{2.\allowbreak{}4.\allowbreak{}2;\allowbreak{}12.\allowbreak{}3.\allowbreak{}2.\allowbreak{}\{21--56\}\}} & 37\\\hline
\texttt{B87} & \texttt{1.\allowbreak{}11.\allowbreak{}3.\allowbreak{}\{1,\allowbreak{}2\}.\allowbreak{}1} & 2\\\hline
\texttt{B88} & \texttt{1.\allowbreak{}1.\allowbreak{}5.\allowbreak{}4} & 1\\\hline
\texttt{B89} & \texttt{1.\allowbreak{}\{2.\allowbreak{}3.\allowbreak{}\{\{1--5\};\allowbreak{}7.\allowbreak{}\{1--10\}\};\allowbreak{}12.\allowbreak{}8.\allowbreak{}\{1--5\}\}} & 20\\\hline
\texttt{B90} & \texttt{1.\allowbreak{}\{11.\allowbreak{}8.\allowbreak{}1.\allowbreak{}\{13--20\};\allowbreak{}12.\allowbreak{}14.\allowbreak{}1.\allowbreak{}\{21--56\}\}} & 44\\\hline
\texttt{B91} & \texttt{1.\allowbreak{}\{11.\allowbreak{}8.\allowbreak{}2.\allowbreak{}\{13--20\};\allowbreak{}12.\allowbreak{}14.\allowbreak{}2.\allowbreak{}\{21--56\}\}} & 44\\\hline
\texttt{B92} & \texttt{1.\allowbreak{}\{11.\allowbreak{}\{2.\allowbreak{}\{1,\allowbreak{}2\}.\allowbreak{}5;\allowbreak{}8.\allowbreak{}1.\allowbreak{}\{13--20\}\};\allowbreak{}12.\allowbreak{}14.\allowbreak{}1.\allowbreak{}\{21--56\}\}} & 46\\\hline
\end{longtable}\endgroup
\begingroup\setlength{\tabcolsep}{3pt}
\begin{longtable}{L{\dimexpr0.340000\textwidth-2\tabcolsep\relax}L{\dimexpr0.240000\textwidth-2\tabcolsep\relax}L{\dimexpr0.420000\textwidth-2\tabcolsep\relax}}
\caption{Asignación de requisitos a resultados y recepción}\label{tab:t14-asignacion}\\\hline
\EncabezadoTabla{Identificadores T-12} & \EncabezadoTabla{Bloques productores} & \EncabezadoTabla{Verificación y condición}\\\hline\endfirsthead
\multicolumn{3}{l}{\textit{Continuación de la tabla \ref{tab:t14-asignacion}}}\\\hline
\EncabezadoTabla{Identificadores T-12} & \EncabezadoTabla{Bloques productores} & \EncabezadoTabla{Verificación y condición}\\\hline\endhead
\hline\endfoot
\hline\endlastfoot
\texttt{RF-NAV-\{01--13\}} & \texttt{B01}\newline{}Alcance ofrecido: construcción/verificación/servicio & 1.10.2.1.10 / 1.10.2.2.10\newline{}Obligatoriedad y etapa por cláusula en T-12. Trabajo de implementación/servicio según paquete.\\\hline
\texttt{RF-ALE-\{01--12\}} & \texttt{B02}\newline{}Alcance ofrecido: construcción/verificación/servicio & 1.10.2.1.10 / 1.10.2.2.10\newline{}Obligatoriedad y etapa por cláusula en T-12. Trabajo de implementación/servicio según paquete.\\\hline
\texttt{RF-ESC-\{01--14\}} & \texttt{B03}\newline{}Alcance ofrecido: construcción/verificación/servicio & 1.10.2.1.10 / 1.10.2.2.10\newline{}Obligatoriedad y etapa por cláusula en T-12. Trabajo de implementación/servicio según paquete.\\\hline
\texttt{RF-AMA-\{01--12\}} & \texttt{B04}\newline{}Alcance ofrecido: construcción/verificación/servicio & 1.10.2.1.10 / 1.10.2.2.10\newline{}Obligatoriedad y etapa por cláusula en T-12. Trabajo de implementación/servicio según paquete.\\\hline
\texttt{RF-VAR-\{01--08\}} & \texttt{B05}\newline{}Alcance ofrecido: construcción/verificación/servicio & 1.10.2.1.10 / 1.10.2.2.10\newline{}Obligatoriedad y etapa por cláusula en T-12. Trabajo de implementación/servicio según paquete.\\\hline
\texttt{RF-CTR-\{01--07\}} & \texttt{B06}\newline{}Alcance ofrecido: construcción/verificación/servicio & 1.10.2.1.10 / 1.10.2.2.10\newline{}Obligatoriedad y etapa por cláusula en T-12. Trabajo de implementación/servicio según paquete.\\\hline
\texttt{RF-AMB-\{01--07\}} & \texttt{B07}\newline{}Alcance ofrecido: construcción/verificación/servicio & 1.10.2.1.10 / 1.10.2.2.10\newline{}Obligatoriedad y etapa por cláusula en T-12. Trabajo de implementación/servicio según paquete.\\\hline
\texttt{RF-CON-\{01--09\}} & \texttt{B08}\newline{}Alcance ofrecido: construcción/verificación/servicio & 1.10.2.1.10 / 1.10.2.2.10\newline{}Obligatoriedad y etapa por cláusula en T-12. Trabajo de implementación/servicio según paquete.\\\hline
\texttt{RF-BOY-\{01--07\}} & \texttt{B09}\newline{}Alcance ofrecido: construcción/verificación/servicio & 1.10.2.1.10 / 1.10.2.2.10\newline{}Obligatoriedad y etapa por cláusula en T-12. Trabajo de implementación/servicio según paquete.\\\hline
\texttt{RF-FAC-\{01--09\}} & \texttt{B10}\newline{}Alcance ofrecido: construcción/verificación/servicio & 1.10.2.1.10 / 1.10.2.2.10\newline{}Obligatoriedad y etapa por cláusula en T-12. Trabajo de implementación/servicio según paquete.\\\hline
\texttt{RF-ADM-\{01--37\}} & \texttt{B11}\newline{}Alcance ofrecido: construcción/verificación/servicio & 1.10.2.1.10 / 1.10.2.2.10\newline{}Obligatoriedad y etapa por cláusula en T-12. Trabajo de implementación/servicio según paquete.\\\hline
\texttt{RF-AUT-\{01--11\}} & \texttt{B12}\newline{}Alcance ofrecido: construcción/verificación/servicio & 1.10.2.1.10 / 1.10.2.2.10\newline{}Obligatoriedad y etapa por cláusula en T-12. Trabajo de implementación/servicio según paquete.\\\hline
\texttt{RNF-DES-\{01--12\}} & \texttt{B13}\newline{}Alcance ofrecido: construcción/verificación/servicio & 1.10.2.1.10 / 1.10.2.2.10\newline{}Obligatoriedad y etapa por cláusula en T-12. Trabajo de implementación/servicio según paquete.\\\hline
\texttt{RNF-CONT-\{01--06\}} & \texttt{B14}\newline{}Alcance ofrecido: construcción/verificación/servicio & 1.10.2.1.10 / 1.10.2.2.10\newline{}Obligatoriedad y etapa por cláusula en T-12. Trabajo de implementación/servicio según paquete.\\\hline
\texttt{RNF-SEG-\{01--07\}} & \texttt{B15}\newline{}Alcance ofrecido: construcción/verificación/servicio & 1.10.2.1.10 / 1.10.2.2.10\newline{}Obligatoriedad y etapa por cláusula en T-12. Trabajo de implementación/servicio según paquete.\\\hline
\texttt{RNF-USA-\{01--04\}} & \texttt{B16}\newline{}Alcance ofrecido: construcción/verificación/servicio & 1.10.2.1.10 / 1.10.2.2.10\newline{}Obligatoriedad y etapa por cláusula en T-12. Trabajo de implementación/servicio según paquete.\\\hline
\texttt{RNF-IDM-\{01,\allowbreak{}02\}} & \texttt{B17}\newline{}Alcance ofrecido: construcción/verificación/servicio & 1.10.2.1.10 / 1.10.2.2.10\newline{}Obligatoriedad y etapa por cláusula en T-12. Trabajo de implementación/servicio según paquete.\\\hline
\texttt{RNF-ESCAL-\{01,\allowbreak{}02\}} & \texttt{B18}\newline{}Alcance ofrecido: construcción/verificación/servicio & 1.10.2.1.10 / 1.10.2.2.10\newline{}Obligatoriedad y etapa por cláusula en T-12. Trabajo de implementación/servicio según paquete.\\\hline
\texttt{RNF-INT-\{01,\allowbreak{}02\}} & \texttt{B19}\newline{}Alcance ofrecido: construcción/verificación/servicio & 1.10.2.1.10 / 1.10.2.2.10\newline{}Obligatoriedad y etapa por cláusula en T-12. Trabajo de implementación/servicio según paquete.\\\hline
\texttt{RNF-OPE-\{01--03\};\allowbreak{} RNF-SOP-\{01--03\}} & \texttt{B20}\newline{}Alcance ofrecido: construcción/verificación/servicio & 1.10.2.1.10 / 1.10.2.2.10\newline{}Obligatoriedad y etapa por cláusula en T-12. Trabajo de implementación/servicio según paquete.\\\hline
\texttt{RNF-RET-\{01--08\}} & \texttt{B21}\newline{}Alcance ofrecido: construcción/verificación/servicio & 1.10.2.1.10 / 1.10.2.2.10\newline{}Obligatoriedad y etapa por cláusula en T-12. Trabajo de implementación/servicio según paquete.\\\hline
\texttt{RNF-DISP-01} & \texttt{B22}\newline{}Alcance ofrecido: construcción/verificación/servicio & 1.10.2.1.10 / 1.10.2.2.10\newline{}Obligatoriedad y etapa por cláusula en T-12. Trabajo de implementación/servicio según paquete.\\\hline
\texttt{RNF-MIG-01} & \texttt{B23}\newline{}Alcance ofrecido: construcción/verificación/servicio & 1.10.2.1.10 / 1.10.2.2.10\newline{}Obligatoriedad y etapa por cláusula en T-12. Trabajo de implementación/servicio según paquete.\\\hline
\texttt{RNF-RED-01} & \texttt{B24}\newline{}Alcance ofrecido: construcción/verificación/servicio & 1.10.2.1.10 / 1.10.2.2.10\newline{}Obligatoriedad y etapa por cláusula en T-12. Trabajo de implementación/servicio según paquete.\\\hline
\texttt{RNF-TER-01} & \texttt{B25}\newline{}Alcance ofrecido: construcción/verificación/servicio & 1.10.2.1.10 / 1.10.2.2.10\newline{}Obligatoriedad y etapa por cláusula en T-12. Trabajo de implementación/servicio según paquete.\\\hline
\texttt{RNF-VEN-01} & \texttt{B26}\newline{}Alcance ofrecido: construcción/verificación/servicio & 1.10.2.1.10 / 1.10.2.2.10\newline{}Obligatoriedad y etapa por cláusula en T-12. Trabajo de implementación/servicio según paquete.\\\hline
\texttt{RNF-CAP-01} & \texttt{B27}\newline{}Alcance ofrecido: construcción/verificación/servicio & 1.10.2.1.10 / 1.10.2.2.10\newline{}Obligatoriedad y etapa por cláusula en T-12. Trabajo de implementación/servicio según paquete.\\\hline
\texttt{RNF-CUMP-01} & \texttt{B28}\newline{}Alcance ofrecido: construcción/verificación/servicio & 1.10.2.1.10 / 1.10.2.2.10\newline{}Obligatoriedad y etapa por cláusula en T-12. Trabajo de implementación/servicio según paquete.\\\hline
\texttt{RT-02.\allowbreak{}\{01--04,\allowbreak{}11,\allowbreak{}13\}} & \texttt{B81}\newline{}Productores contractuales ofrecidos & Vistas y ADR de la arquitectura integrada 1.7.8.1; no se acreditan mediante controles de seguridad aislados.\newline{}Línea base H2 y delta H8; decisiones de oferta localizadas en SD4.\\\hline
\texttt{RT-02.\allowbreak{}\{05--10,\allowbreak{}12\}} & \texttt{B81;\allowbreak{} B72;\allowbreak{} B68;\allowbreak{} B29}\newline{}Productores contractuales ofrecidos & Diseño recibido, infraestructura e integración configuradas y adaptaciones funcionales; idempotencia, resiliencia, estado externo y escalamiento se verifican en los productores pertinentes.\newline{}Diseño/construcción antes de la recepción de su etapa; verificación en T-13/T-17.\\\hline
\texttt{RT-13.\allowbreak{}\{01,\allowbreak{}02,\allowbreak{}05--08,\allowbreak{}10--12\}} & \texttt{B82;\allowbreak{} B68}\newline{}Productores contractuales ofrecidos & Diseño recibido, canales/vistas propietarias y verificación accesible de la versión; requisitos por perfil, tarea y dispositivo.\newline{}Diseño antes de construir; comprobaciones manuales/automáticas y UAT antes de producción.\\\hline
\texttt{RT-13.\allowbreak{}\{03,\allowbreak{}04,\allowbreak{}09\}} & \texttt{B61}\newline{}Productores contractuales ofrecidos & Investigación, prototipado de trabajo, validación con usuarios, indicadores y sistema de diseño en 1.7.7.1/2.\newline{}Diseño E1 antes de construcción definitiva; delta E2 recibido en H8 antes de construirlo. T-13 desarrolla las pruebas, no reemplaza este resultado previo.\\\hline
\texttt{RT-02.\allowbreak{}14;\allowbreak{} RT-03.\allowbreak{}\{08,\allowbreak{}09,\allowbreak{}19\};\allowbreak{} RT-04.\allowbreak{}14;\allowbreak{} RT-05.\allowbreak{}\{24,\allowbreak{}30\};\allowbreak{} RT-06.\allowbreak{}\{12,\allowbreak{}15,\allowbreak{}34\};\allowbreak{} RT-08.\allowbreak{}\{15,\allowbreak{}19\};\allowbreak{} RT-09.\allowbreak{}10;\allowbreak{} RT-11.\allowbreak{}21;\allowbreak{} RT-14.\allowbreak{}09;\allowbreak{} RT-15.\allowbreak{}\{05,\allowbreak{}06\};\allowbreak{} RT-16.\allowbreak{}\{05,\allowbreak{}26\};\allowbreak{} RT-17.\allowbreak{}08;\allowbreak{} RT-18.\allowbreak{}10;\allowbreak{} RT-21.\allowbreak{}12;\allowbreak{} RT-22.\allowbreak{}09;\allowbreak{} RT-23.\allowbreak{}\{07--10\};\allowbreak{} RT-24.\allowbreak{}03;\allowbreak{} RT-25.\allowbreak{}10} & \texttt{---}\newline{}Deseable BT no ofrecido; exclusión conservada & No procede; conservar exclusión\newline{}Alcance Deseable BT no ofrecido; no genera trabajo contractual por sí solo.\\\hline
\texttt{RT-03.\allowbreak{}\{01--05,\allowbreak{}14,\allowbreak{}15,\allowbreak{}17,\allowbreak{}18,\allowbreak{}20,\allowbreak{}21,\allowbreak{}23\}} & \texttt{B29;\allowbreak{} B81}\newline{}Productores contractuales ofrecidos & Plataforma/recinto, red y terminales: 1.2; especificaciones, titular de provisión y cálculos en SD4/T-11. Diseño recibido precede habilitación H3.\newline{}Diseño H2, habilitación H3 y control posterior del activo; no sustituye obligaciones externas por pruebas de software.\\\hline
\texttt{RT-03.\allowbreak{}06} & \texttt{B29;\allowbreak{} B39}\newline{}Productores contractuales ofrecidos & Etiquetado y presupuestos de recursos; seguimiento de consumo/costo en 1.12.3.5.\newline{}Habilitación y reporte de cada período; montos en Oferta Económica.\\\hline
\texttt{RT-03.\allowbreak{}07} & \texttt{B81;\allowbreak{} B43;\allowbreak{} B25}\newline{}Productores contractuales ofrecidos & Decisiones de portabilidad y plan de reversibilidad; transferencia final de fuentes/datos.\newline{}Diseño H2; actualizaciones y salida en sus paquetes respectivos.\\\hline
\texttt{RT-03.\allowbreak{}\{10--13\}} & \texttt{B29;\allowbreak{} B41;\allowbreak{} B68;\allowbreak{} B37}\newline{}Productores contractuales ofrecidos & Autoridad local, registro crítico, sincronización/conflictos y catálogo de funciones degradadas; prueba integral de desconexión antes de producción.\newline{}Diseño/construcción antes de la recepción de su etapa; verificación en T-13/T-17.\\\hline
\texttt{RT-03.\allowbreak{}16} & \texttt{B29;\allowbreak{} B38}\newline{}Productores contractuales ofrecidos & Monitoreo integrado nube/recinto, alertas y evidencia por activo.\newline{}Diseño/construcción antes de la recepción de su etapa; verificación en T-13/T-17.\\\hline
\texttt{RT-03.\allowbreak{}22} & \texttt{B29;\allowbreak{} B37;\allowbreak{} B36}\newline{}Productores contractuales ofrecidos & Acceso remoto con identidad, postura y controles de exposición; comprobación de denegaciones.\newline{}Diseño/construcción antes de la recepción de su etapa; verificación en T-13/T-17.\\\hline
\texttt{RT-04.\allowbreak{}\{01,\allowbreak{}02\}} & \texttt{B30}\newline{}Productores contractuales ofrecidos & Cinco ambientes y diferencias PREPROD/PROD documentadas en su expediente de recepción.\newline{}Antes de H3, según RT-04.01; equivalencia comprobada antes de ensayos.\\\hline
\texttt{RT-04.\allowbreak{}\{03--08,\allowbreak{}10,\allowbreak{}11\}} & \texttt{B29;\allowbreak{} B68;\allowbreak{} B13}\newline{}Productores contractuales ofrecidos & Repositorio y cadena 1.2.1.1/5; promoción, pruebas bloqueantes y migraciones de esquema; consumidores y ensayo en PREPROD.\newline{}Antes del paso de producción correspondiente; desarrollo y evidencia no se reemplazan por disponer de ambientes.\\\hline
\texttt{RT-04.\allowbreak{}09} & \texttt{B17}\newline{}Productores contractuales ofrecidos & Gestión de secretos/claves 1.3.1.1 y uso de los consumidores; rotación y ausencia de secretos embebidos.\newline{}Diseño/construcción antes de la recepción de su etapa; verificación en T-13/T-17.\\\hline
\texttt{RT-04.\allowbreak{}12} & \texttt{B29;\allowbreak{} B35}\newline{}Productores contractuales ofrecidos & Objetivos de liberación y medición durante servicio; versión y período identificados.\newline{}Objetivos de diseño; medición en Operación.\\\hline
\texttt{RT-04.\allowbreak{}13} & \texttt{B39}\newline{}Productores contractuales ofrecidos & Reducción/apagado de ambientes no productivos y ahorro sustentado; tratamiento económico fuera de T-14.\newline{}Diseño/construcción antes de la recepción de su etapa; verificación en T-13/T-17.\\\hline
\texttt{RT-06.\allowbreak{}\{01--09,\allowbreak{}11,\allowbreak{}13,\allowbreak{}16--18,\allowbreak{}20--25,\allowbreak{}29--33\}} & \texttt{B32}\newline{}Productores contractuales ofrecidos & Especificación/coordinación y expediente de recepción del recinto en 1.2.3.1; configuración propietaria de nodos/red y seguimiento. Cada habilitación indica responsable externo o propio y evidencia, según SD4/T-11.\newline{}Diseño/habilitación antes de H3; controles periódicos según cláusula. La obra a cargo del cliente conserva ese titular.\\\hline
\texttt{RT-06.\allowbreak{}10} & \texttt{B32;\allowbreak{} B20}\newline{}Productores contractuales ofrecidos & Revisión semestral de instalación eléctrica y constancia de medición; actividad/ventana/responsable del mantenimiento físico declarados en el plan de operación.\newline{}Habilitación y controles semestrales durante servicio; FALTA REFERENCIA A SD11.\\\hline
\texttt{RT-06.\allowbreak{}\{14,\allowbreak{}19\}} & \texttt{B32;\allowbreak{} B38}\newline{}Productores contractuales ofrecidos & Sensores y alarmas físicas integrados al monitoreo; aviso a NOC y contraparte.\newline{}Diseño/construcción antes de la recepción de su etapa; verificación en T-13/T-17.\\\hline
\texttt{RT-06.\allowbreak{}\{26--28\}} & \texttt{B32;\allowbreak{} B86}\newline{}Productores contractuales ofrecidos & Custodia e inventario de medios, condiciones del recinto, rotación y legibilidad; la política de respaldo no acredita sola la custodia física.\newline{}Diseño/construcción antes de la recepción de su etapa; verificación en T-13/T-17.\\\hline
\texttt{RT-07.\allowbreak{}\{01--08\}} & \texttt{B33;\allowbreak{} B81}\newline{}Productores contractuales ofrecidos & Diseño DR, ambiente secundario, replicación y procedimientos; conmutación/retorno medidos y controles periódicos.\newline{}Diseño H2; habilitación H3; prueba antes de producción y periodicidad RT-07, sin presentar ensayos futuros como realizados.\\\hline
\texttt{RT-07.\allowbreak{}\{09--14\}} & \texttt{B86;\allowbreak{} B76}\newline{}Productores contractuales ofrecidos & Política por dominio, copias protegidas y restauración verificable; habilitación 1.2.4.2 y controles 1.12.3.2.\newline{}Habilitación previa al servicio y prueba mensual; la retención conserva el alcance específico C07.\\\hline
\texttt{RT-08.\allowbreak{}\{01--14\}} & \texttt{B34;\allowbreak{} B81}\newline{}Productores contractuales ofrecidos & Especificaciones, recepción y configuración del hardware/terminales 1.2.3; inventario, garantías y ciclo de vida según SD4/T-11. Se conserva la responsabilidad de compra/provisión de cada activo.\newline{}Diseño y habilitación; operación/retiro en la ventana del activo. No se acredita equipamiento mediante una matriz de permisos.\\\hline
\texttt{RT-08.\allowbreak{}\{16--18\}} & \texttt{B89}\newline{}Productores contractuales ofrecidos & Ciclo de vida, retiro, sanitización y disposición certificada de activos; evidencia por medio y titular del equipamiento.\newline{}Recepción, mantenimiento y retiro; procedimiento de disposición del propietario enlazado, no compra nueva.\\\hline
\texttt{RT-09.\allowbreak{}\{01--05,\allowbreak{}08\}} & \texttt{B81;\allowbreak{} B29;\allowbreak{} B68;\allowbreak{} B13}\newline{}Productores contractuales ofrecidos & Dimensionamiento y adaptación de los componentes, con ensayo de rendimiento/degradación; el informe de prueba no sustituye al diseño.\newline{}Diseño/construcción antes de la recepción de su etapa; verificación en T-13/T-17.\\\hline
\texttt{RT-09.\allowbreak{}\{06,\allowbreak{}07\}} & \texttt{B13}\newline{}Productores contractuales ofrecidos & Ensayo de carga/estrés y curvas de resultados en 1.10.2.1.4/1.10.2.2.4.\newline{}PREPROD antes de cada paso de producción; criterios específicos en T-13.\\\hline
\texttt{RT-09.\allowbreak{}09} & \texttt{B39}\newline{}Productores contractuales ofrecidos & Proyección de capacidad e infraestructura 1.12.3.5; datos observados, cuellos de botella y consolidación trimestral.\newline{}Operación, sin duplicar el control mensual con su consolidación.\\\hline
\texttt{RT-10.\allowbreak{}\{01--06,\allowbreak{}08,\allowbreak{}09\}} & \texttt{B81;\allowbreak{} B29;\allowbreak{} B35;\allowbreak{} B68}\newline{}Productores contractuales ofrecidos & Criticidad y continuidad del diseño, comportamiento de productores e indicadores del servicio; dependencias y restricciones operacionales identificadas.\newline{}Diseño/construcción antes de la recepción de su etapa; verificación en T-13/T-17.\\\hline
\texttt{RT-10.\allowbreak{}07} & \texttt{B13;\allowbreak{} B35}\newline{}Productores contractuales ofrecidos & Ensayos de resiliencia 1.10.2.1.6/1.10.2.2.6 y controles semestrales del servicio.\newline{}Antes de producción y una vez por semestre en Operación, con evidencias diferentes.\\\hline
\texttt{RT-11.\allowbreak{}\{01--09,\allowbreak{}11--19,\allowbreak{}22--28\}} & \texttt{B36}\newline{}Productores contractuales ofrecidos & Productores de seguridad de 1.3, controles de borde/cadena en 1.2 y expedientes de servicio; cada control se contrasta con su componente y evidencia específica.\newline{}Diseño H2, habilitación/recepción técnica de su etapa y revisión periódica del control.\\\hline
\texttt{RT-11.\allowbreak{}10} & \texttt{B78;\allowbreak{} B68}\newline{}Productores contractuales ofrecidos & Protección adicional de campos sensibles 1.3.4.4 y adaptaciones propietarias; cifrado de disco o KMS solos no acreditan este control.\newline{}Diseño/construcción antes de la recepción de su etapa; verificación en T-13/T-17.\\\hline
\texttt{RT-11.\allowbreak{}20} & \texttt{B13;\allowbreak{} B36}\newline{}Productores contractuales ofrecidos & Prueba ofensiva independiente antes de cada producción 1.10.2.1.8/1.10.2.2.8 y controles anuales 1.12.5.1--3.\newline{}Antes de producción y anualmente; cierre de hallazgos según protocolo, sin reutilizar una prueba periódica para recibir construcción inicial.\\\hline
\texttt{RT-12.\allowbreak{}\{01--04,\allowbreak{}06--08,\allowbreak{}12,\allowbreak{}13\}} & \texttt{B37}\newline{}Productores contractuales ofrecidos & Identidad 1.3.1.2: directorio/federación, SSO, MFA/FIDO2, privilegios, sesiones, autoservicio y emergencia; prueba específica del ciclo de acceso.\newline{}Diseño/construcción antes de la recepción de su etapa; verificación en T-13/T-17.\\\hline
\texttt{RT-12.\allowbreak{}05} & \texttt{B37;\allowbreak{} B68}\newline{}Productores contractuales ofrecidos & Matriz 1.3.1.3 y autorización efectiva de cada dominio; atributos y segregación cuando corresponda.\newline{}Diseño/construcción antes de la recepción de su etapa; verificación en T-13/T-17.\\\hline
\texttt{RT-12.\allowbreak{}09} & \texttt{B37;\allowbreak{} B70}\newline{}Productores contractuales ofrecidos & Emisión de eventos de identidad y conservación protegida en 1.3.3.1; creación, modificación, elevación, bloqueo y baja registrados.\newline{}Diseño/construcción antes de la recepción de su etapa; verificación en T-13/T-17.\\\hline
\texttt{RT-12.\allowbreak{}10} & \texttt{B37;\allowbreak{} B73}\newline{}Productores contractuales ofrecidos & Fuente laboral, adaptador 1.7.4.2, revocación central/local y baja dentro de 24 horas; se verifican todos los consumidores aplicables.\newline{}Diseño/construcción antes de la recepción de su etapa; verificación en T-13/T-17.\\\hline
\texttt{RT-12.\allowbreak{}11} & \texttt{B37;\allowbreak{} B40;\allowbreak{} B61}\newline{}Productores contractuales ofrecidos & Diseño y autenticación adecuados a perfiles de terreno, guantes, turnos y dispositivos compartidos; identidad nominativa y permisos conservados.\newline{}Diseño/construcción antes de la recepción de su etapa; verificación en T-13/T-17.\\\hline
\texttt{RT-14.\allowbreak{}\{01--08\}} & \texttt{B38}\newline{}Productores contractuales ofrecidos & Observabilidad 1.2.4.1, auditoría protegida y tableros/runbooks del servicio; retención y control de incidentes específicos.\newline{}Diseño/construcción antes de la recepción de su etapa; verificación en T-13/T-17.\\\hline
\texttt{RT-15.\allowbreak{}\{01--04\}} & \texttt{B39;\allowbreak{} B81}\newline{}Productores contractuales ofrecidos & Dimensionamiento/utilización, reducción de ambientes y estimación de energía/carbono de la plataforma; método, límites y fuente de PUE/intensidad declarados.\newline{}Diseño/habilitación y actualización con datos de Operación; costos fuera de T-14.\\\hline
\texttt{RT-15.\allowbreak{}\{07--09\}} & \texttt{B88}\newline{}Productores contractuales ofrecidos & Expediente de control enlazado a T-5/SD12: certificados verificables, asignación efectiva, dedicaciones y reposición.\newline{}Evidencia de oferta y vigencia contractual diferenciadas; FALTA REFERENCIA A SD12.\\\hline
\texttt{RT-16.\allowbreak{}\{01--04\}} & \texttt{B40;\allowbreak{} B37;\allowbreak{} B68;\allowbreak{} B70}\newline{}Productores contractuales ofrecidos & Consola común, roles y adaptaciones de parámetros por dominio; catálogo sin desarrollo, aprobación dual y auditoría de cambios.\newline{}Diseño/construcción antes de la recepción de su etapa; verificación en T-13/T-17.\\\hline
\texttt{RT-16.\allowbreak{}\{06--09\}} & \texttt{B85}\newline{}Productores contractuales ofrecidos & Generación de cambios/consultas por dominios y conservación protegida en 1.3.3.1; consulta/exportación autorizadas. Para RT-16.09 se comprueba el universo sensible C07, además de las modificaciones.\newline{}Diseño/construcción antes de la recepción de su etapa; verificación en T-13/T-17.\\\hline
\texttt{RT-16.\allowbreak{}\{11--13\}} & \texttt{B40;\allowbreak{} B68;\allowbreak{} B84;\allowbreak{} B70}\newline{}Productores contractuales ofrecidos & Flujos de los dominios, configuración de responsables/plazos/aprobación, bandeja común, alertas y auditoría; la base de canal sola no produce el flujo.\newline{}Diseño/construcción antes de la recepción de su etapa; verificación en T-13/T-17.\\\hline
\texttt{RT-16.\allowbreak{}\{15,\allowbreak{}16,\allowbreak{}19\}} & \texttt{B68;\allowbreak{} B40;\allowbreak{} B69;\allowbreak{} B76}\newline{}Productores contractuales ofrecidos & Documentos de los procesos propietarios: versiones, metadatos, búsqueda, plantillas y acceso; persistencia/cifrado y conservación por política vigente.\newline{}Diseño/construcción antes de la recepción de su etapa; verificación en T-13/T-17.\\\hline
\texttt{RT-16.\allowbreak{}\{17,\allowbreak{}18\}} & \texttt{B77;\allowbreak{} B68}\newline{}Productores contractuales ofrecidos & Capacidad de firma, verificación y evidencia 1.7.6.6 y cinco procesos C07 afectados; modalidad admisible y conservación por acto/versión.\newline{}Diseño/construcción antes de la recepción de su etapa; verificación en T-13/T-17.\\\hline
\texttt{RT-16.\allowbreak{}\{20--25\}} & \texttt{B84;\allowbreak{} B68}\newline{}Productores contractuales ofrecidos & Notificaciones 1.7.4.1 y contenido/reglas del dominio; plantillas, preferencias, proveedor, envío/duplicados/estados y baja comprobados.\newline{}Construcción antes de recepción; costos/proyección en Oferta Económica, no una tarifa supuesta de T-14.\\\hline
\texttt{RT-16.\allowbreak{}\{27--29\}} & \texttt{B40;\allowbreak{} B68;\allowbreak{} B63;\allowbreak{} B64}\newline{}Productores contractuales ofrecidos & Búsqueda/listados de dominios, autoservicio y exportación masiva/analítica; permisos, filtros, formato y procesamiento asíncrono cuando procede.\newline{}Diseño/construcción antes de la recepción de su etapa; verificación en T-13/T-17.\\\hline
\texttt{RT-16.\allowbreak{}30} & \texttt{B85;\allowbreak{} B64}\newline{}Productores contractuales ofrecidos & Exportación sensible: productor registra persona, universo, instante y versión; auditoría recibe y protege el evento.\newline{}Diseño/construcción antes de la recepción de su etapa; verificación en T-13/T-17.\\\hline
\texttt{RT-16.\allowbreak{}\{31--34\}} & \texttt{B40;\allowbreak{} B68;\allowbreak{} B29;\allowbreak{} B61}\newline{}Productores contractuales ofrecidos & Portal y autoatención de funciones ofertadas; diseño por perfiles, indicador de atención asistida y aislamiento de recursos del servicio interno.\newline{}Diseño/construcción antes de la recepción de su etapa; verificación en T-13/T-17.\\\hline
\texttt{RT-17.\allowbreak{}\{01--07\}} & \texttt{B41}\newline{}Alcance ofrecido: construcción/verificación/servicio & Protocolo integral 1.10.2.1.1 / 1.10.2.2.1 y evidencia específica de la cláusula\newline{}Obligatoriedad y etapa por cláusula en T-12. Trabajo de implementación/servicio según paquete.\\\hline
\texttt{RT-18.\allowbreak{}\{01--09\}} & \texttt{B42;\allowbreak{} B62;\allowbreak{} B60;\allowbreak{} B70}\newline{}Productores contractuales ofrecidos & Componentes IA efectivamente declarados, control de uso/datos/proveedor, intervención humana y auditoría; modelos predictivos solo si pertenecen al alcance vigente, sin incorporar uno nuevo.\newline{}Declaración de oferta, recepción de innovación y control posterior; expediente IA de 1.1.5.3 enlaza su evidencia.\\\hline
\texttt{RT-19.\allowbreak{}\{01--10\}} & \texttt{B43}\newline{}Alcance ofrecido: construcción/verificación/servicio & Protocolo integral 1.10.2.1.1 / 1.10.2.2.1 y evidencia específica de la cláusula\newline{}Obligatoriedad y etapa por cláusula en T-12. Trabajo de implementación/servicio según paquete.\\\hline
\texttt{RT-20.\allowbreak{}\{01--08\}} & \texttt{B44;\allowbreak{} B13;\allowbreak{} B74}\newline{}Productores contractuales ofrecidos & Preparación de implantación/convivencia, ensayos y expedientes de aceptación; definición de terminado 1.1.1.3 y protocolo por hito.\newline{}Antes de recibir y activar su etapa; T-18/T-17 detallan el procedimiento, sin declarar recepción por horas consumidas.\\\hline
\texttt{RT-21.\allowbreak{}01} & \texttt{B46;\allowbreak{} B91}\newline{}Productores y evidencia ofrecidos & Modelo habilitado y prestación real de NOC/SOC/guardia en 1.11.8.2 y 1.12.14.2; preparación y cuadrante no equivalen a turnos prestados.\newline{}Soporte de implementación 13--20; Operación 21--56, sin adelantar su inicio contractual.\\\hline
\texttt{RT-21.\allowbreak{}02} & \texttt{B46}\newline{}Alcance ofrecido: construcción/verificación/servicio & Designación gerente, dedicación y registros de comité; SD1\newline{}Obligatoriedad y etapa por cláusula en T-12. Trabajo de implementación/servicio según paquete.\\\hline
\texttt{RT-21.\allowbreak{}03} & \texttt{B47;\allowbreak{} B90;\allowbreak{} B91}\newline{}Productores y evidencia ofrecidos & Especialistas nominados y dedicaciones en el modelo, con disponibilidad y suplencia efectiva en prestación de soporte/guardia.\newline{}Antes de activar la cobertura y durante cada período; T-15/SD12 detallan dotación.\\\hline
\texttt{RT-21.\allowbreak{}04} & \texttt{B46}\newline{}Productores y evidencia ofrecidos & Catálogo/runbooks y transferencia de 1.12.1.2, versión inicial y delta de Operación; procedimientos reproducidos antes de uso.\newline{}Línea base antes de mes13 y delta antes de mes21; no se exige recibir futuro servicio para prepararlo.\\\hline
\texttt{RT-21.\allowbreak{}07} & \texttt{B90}\newline{}Productores y evidencia ofrecidos & Prestación de atención/soporte bilingüe por período, ticket único, seguimiento/escalamiento y cierre conforme.\newline{}Implementación 13--20 y Operación 21--56; no sustituida por informe de SLA.\\\hline
\texttt{RT-21.\allowbreak{}08} & \texttt{B90}\newline{}Productores y evidencia ofrecidos & Prestación de atención/soporte bilingüe por período, ticket único, seguimiento/escalamiento y cierre conforme.\newline{}Implementación 13--20 y Operación 21--56; no sustituida por informe de SLA.\\\hline
\texttt{RT-21.\allowbreak{}09} & \texttt{B90;\allowbreak{} B45}\newline{}Productores y evidencia ofrecidos & Histórico de atención, análisis mensual, indicadores por autorización e informe de niveles de servicio; el control usa registros de prestación, sin reemplazarla.\newline{}Evidencia de cada período; no se acredita cumplimiento por ausencia de registros.\\\hline
\texttt{RT-21.\allowbreak{}10} & \texttt{B92}\newline{}Productores y evidencia ofrecidos & Base de conocimiento/autoformación creada en materiales, mantenida por soporte con versiones y evidencia de aprendizaje, avance y utilidad.\newline{}Desde habilitación inicial y durante prestación; no sustituida por jornadas semestrales.\\\hline
\texttt{RT-21.\allowbreak{}11} & \texttt{B92}\newline{}Productores y evidencia ofrecidos & Base de conocimiento/autoformación creada en materiales, mantenida por soporte con versiones y evidencia de aprendizaje, avance y utilidad.\newline{}Desde habilitación inicial y durante prestación; no sustituida por jornadas semestrales.\\\hline
\texttt{RT-21.\allowbreak{}13} & \texttt{B90;\allowbreak{} B45}\newline{}Productores y evidencia ofrecidos & Histórico de atención, análisis mensual, indicadores por autorización e informe de niveles de servicio; el control usa registros de prestación, sin reemplazarla.\newline{}Evidencia de cada período; no se acredita cumplimiento por ausencia de registros.\\\hline
\texttt{RT-21.\allowbreak{}15} & \texttt{B90}\newline{}Productores y evidencia ofrecidos & Prestación de atención/soporte bilingüe por período, ticket único, seguimiento/escalamiento y cierre conforme.\newline{}Implementación 13--20 y Operación 21--56; no sustituida por informe de SLA.\\\hline
\texttt{RT-21.\allowbreak{}17} & \texttt{B90}\newline{}Productores y evidencia ofrecidos & Prestación de atención/soporte bilingüe por período, ticket único, seguimiento/escalamiento y cierre conforme.\newline{}Implementación 13--20 y Operación 21--56; no sustituida por informe de SLA.\\\hline
\texttt{RT-21.\allowbreak{}18} & \texttt{B90;\allowbreak{} B45}\newline{}Productores y evidencia ofrecidos & Histórico de atención, análisis mensual, indicadores por autorización e informe de niveles de servicio; el control usa registros de prestación, sin reemplazarla.\newline{}Evidencia de cada período; no se acredita cumplimiento por ausencia de registros.\\\hline
\texttt{RT-21.\allowbreak{}05} & \texttt{B90}\newline{}Productores y evidencia ofrecidos & Prestación de atención/soporte bilingüe por período, ticket único, seguimiento/escalamiento y cierre conforme.\newline{}Implementación 13--20 y Operación 21--56; no sustituida por informe de SLA.\\\hline
\texttt{RT-21.\allowbreak{}06} & \texttt{B90;\allowbreak{} B91;\allowbreak{} B45}\newline{}Productores y evidencia ofrecidos & Atención con umbrales/denominadores y horarios C07; canales de emergencia y planes vencidos 24x7x365 en guardia. Prestación y control mensual tienen fronteras distintas.\newline{}Recepción efectiva por período; T-13/T-17 verifican métricas y evidencia de cobertura.\\\hline
\texttt{RT-21.\allowbreak{}14} & \texttt{B48;\allowbreak{} B90;\allowbreak{} B91}\newline{}Productores y evidencia ofrecidos & Dimensionamiento de cobertura con demanda/TMA/Erlang C/sensibilidad; revisión con registros y disponibilidad real.\newline{}Modelo antes de activar turnos; actualización durante el contrato; recursos en T-15.\\\hline
\texttt{RT-21.\allowbreak{}16} & \texttt{B91}\newline{}Productores y evidencia ofrecidos & Especialista/suplente, embarcación y ventana segura para intervención en Bahía Ilque, con continuidad de otros sitios y evidencia de atención.\newline{}Cuando la resolución lo requiere; programación/costos en sus documentos, no tarifa inventada en T-14.\\\hline
\texttt{RT-21.\allowbreak{}19} & \texttt{B49}\newline{}Alcance ofrecido: construcción/verificación/servicio & Bolsa anual 960 HH, composición por perfil, aprobación/liquidación y acumulación de saldo\newline{}Obligatoriedad y etapa por cláusula en T-12. Trabajo de implementación/servicio según paquete.\\\hline
\texttt{RT-21.\allowbreak{}20} & \texttt{B50}\newline{}Alcance ofrecido: construcción/verificación/servicio & Cambio evolutivo recibido con pruebas, regresión y seguridad equivalentes al desarrollo original\newline{}Obligatoriedad y etapa por cláusula en T-12. Trabajo de implementación/servicio según paquete.\\\hline
\texttt{RT-21.\allowbreak{}21} & \texttt{B51}\newline{}Alcance ofrecido: construcción/verificación/servicio & Registro cuantificado y 32/240 HH mensuales de capacidad de ingeniería a deuda técnica\newline{}Obligatoriedad y etapa por cláusula en T-12. Trabajo de implementación/servicio según paquete.\\\hline
\texttt{RT-21.\allowbreak{}22} & \texttt{B52}\newline{}Productores y evidencia ofrecidos & Plan anual de versiones, autorización de ventana, pruebas, liberación y retorno en mantenimiento 1.12.3.4; no se acredita mediante pentest o bolsa normativa.\newline{}Operación, con calendario recibido y evidencia de cada intervención.\\\hline
\texttt{RT-22.\allowbreak{}01} & \texttt{B53;\allowbreak{} B87}\newline{}Productores contractuales ofrecidos & Materiales, manifiesto activado por perfil/etapa y formación individual; reserva de códigos no acredita población formada.\newline{}Manifiesto antes de fijar cantidad/HH; materiales y versión antes de impartir formación; cierre individual comprobado.\\\hline
\texttt{RT-22.\allowbreak{}02} & \texttt{B53;\allowbreak{} B87}\newline{}Productores contractuales ofrecidos & Materiales, manifiesto activado por perfil/etapa y formación individual; reserva de códigos no acredita población formada.\newline{}Manifiesto antes de fijar cantidad/HH; materiales y versión antes de impartir formación; cierre individual comprobado.\\\hline
\texttt{RT-22.\allowbreak{}03} & \texttt{B55;\allowbreak{} B92}\newline{}Productores y evidencia ofrecidos & Materiales en español editables y de propiedad del cliente: manuales, guías, FAQ, videos y base de conocimiento; mantenimiento enlazado al soporte.\newline{}Materiales antes de formación; base de conocimiento actualizada desde su primera versión hasta cierre.\\\hline
\texttt{RT-22.\allowbreak{}04} & \texttt{B27}\newline{}Alcance ofrecido: construcción/verificación/servicio & Turnos y formación previa al congelamiento; nuevos estacionales y monitores, sin formación interna delegada al CLIENTE\newline{}Obligatoriedad y etapa por cláusula en T-12. Trabajo de implementación/servicio según paquete.\\\hline
\texttt{RT-22.\allowbreak{}05} & \texttt{B53;\allowbreak{} B87}\newline{}Productores contractuales ofrecidos & Materiales, manifiesto activado por perfil/etapa y formación individual; reserva de códigos no acredita población formada.\newline{}Manifiesto antes de fijar cantidad/HH; materiales y versión antes de impartir formación; cierre individual comprobado.\\\hline
\texttt{RT-22.\allowbreak{}06} & \texttt{B57}\newline{}Alcance ofrecido: construcción/verificación/servicio & Dos jornadas anuales y alta de personal nuevo, evaluación y recuperación, sin cargo adicional\newline{}Obligatoriedad y etapa por cláusula en T-12. Trabajo de implementación/servicio según paquete.\\\hline
\texttt{RT-22.\allowbreak{}07} & \texttt{B58}\newline{}Alcance ofrecido: construcción/verificación/servicio & Seis paquetes mensuales 21--26, referentes/suplentes, actas de casos, documentación y revisión independiente; 144 HH\newline{}Obligatoriedad y etapa por cláusula en T-12. Trabajo de implementación/servicio según paquete.\\\hline
\texttt{RT-22.\allowbreak{}08} & \texttt{B92}\newline{}Productores y evidencia ofrecidos & Base de conocimiento creada/versionada y mantenida en prestación de soporte; cambios, uso y utilidad percibida verificables.\newline{}Desde primera versión recibida durante el contrato, con traspaso sin discontinuidad entre implementación y Operación.\\\hline
\texttt{RT-23.\allowbreak{}\{01--06\};\allowbreak{} RT-24.\allowbreak{}\{01,\allowbreak{}02,\allowbreak{}04--06\}} & \texttt{B60}\newline{}Obligación de oferta previa al contrato; vigencia gestionada en 1.1.5 & Control previo de expediente de oferta; contrato: gestión de vigencia 1.1.5\newline{}Evidencia de oferta previa a contratación; vigencia posterior, cuando corresponda.\\\hline
\texttt{RT-25.\allowbreak{}\{01--09\}} & \texttt{---}\newline{}Evidencia académica ofrecida; externa al contrato & Prototipo académico y arquitectura de información: enlace navegable, componentes, acceso de seis meses y evidencia UX según cap. 25. Evidencia externa del Informe 3, pendiente de localización en T-12; no se recibe por entregar un portal productivo.\newline{}Obligación académica ofrecida con Informe 3. No crea un paquete contractual ni exige el prototipo terminado en Informe 2. Diseño contractual se traza separadamente por RT-13. FALTA REFERENCIA A SD4 para el diseño asociado.\\\hline
\texttt{RT-26.\allowbreak{}\{01--07\}} & \texttt{B62}\newline{}Alcance ofrecido: construcción/verificación/servicio & Recepción periódica del paquete y control contractual SLA\newline{}Obligatoriedad y etapa por cláusula en T-12. Trabajo de implementación/servicio según paquete.\\\hline
\texttt{RT-05.\allowbreak{}01} & \texttt{B68;\allowbreak{} B63}\newline{} Alcance ofrecido & Criterio de la ficha productora y prueba específica de la cláusula; no se acredita por ejecutar únicamente migración.\newline{} FALTA REFERENCIA A SD9.\\\hline
\texttt{RT-05.\allowbreak{}03} & \texttt{B70;\allowbreak{} B68}\newline{} Alcance ofrecido & Criterio de la ficha productora y prueba específica de la cláusula; no se acredita por ejecutar únicamente migración.\newline{} FALTA REFERENCIA A SD9.\\\hline
\texttt{RT-05.\allowbreak{}04} & \texttt{B68;\allowbreak{} B64;\allowbreak{} B31}\newline{} Alcance ofrecido & Criterio de la ficha productora y prueba específica de la cláusula; no se acredita por ejecutar únicamente migración.\newline{} FALTA REFERENCIA A SD9.\\\hline
\texttt{RT-05.\allowbreak{}05} & \texttt{B69;\allowbreak{} B63}\newline{} Alcance ofrecido & Criterio de la ficha productora y prueba específica de la cláusula; no se acredita por ejecutar únicamente migración.\newline{} FALTA REFERENCIA A SD9.\\\hline
\texttt{RT-05.\allowbreak{}06} & \texttt{B64;\allowbreak{} B31}\newline{} Alcance ofrecido & Criterio de la ficha productora y prueba específica de la cláusula; no se acredita por ejecutar únicamente migración.\newline{} FALTA REFERENCIA A SD9.\\\hline
\texttt{RT-05.\allowbreak{}09} & \texttt{B68}\newline{} Alcance ofrecido & Criterio de la ficha productora y prueba específica de la cláusula; no se acredita por ejecutar únicamente migración.\newline{} FALTA REFERENCIA A SD9.\\\hline
\texttt{RT-05.\allowbreak{}11} & \texttt{B31}\newline{} Alcance ofrecido & Criterio de la ficha productora y prueba específica de la cláusula; no se acredita por ejecutar únicamente migración.\newline{} FALTA REFERENCIA A SD9.\\\hline
\texttt{RT-05.\allowbreak{}12} & \texttt{B31}\newline{} Alcance ofrecido & Criterio de la ficha productora y prueba específica de la cláusula; no se acredita por ejecutar únicamente migración.\newline{} FALTA REFERENCIA A SD9.\\\hline
\texttt{RT-05.\allowbreak{}13} & \texttt{B74}\newline{} Alcance ofrecido & Ensayos/conciliación de 1.10 y 1.11 por corte y versión; SD5 tabla 5.13.\newline{} FALTA REFERENCIA A SD9.\\\hline
\texttt{RT-05.\allowbreak{}14} & \texttt{B31;\allowbreak{} B75}\newline{} Alcance ofrecido & Ensayos/conciliación de 1.10 y 1.11 por corte y versión; SD5 tabla 5.13.\newline{} FALTA REFERENCIA A SD9.\\\hline
\texttt{RT-05.\allowbreak{}15} & \texttt{B31;\allowbreak{} B67}\newline{} Alcance ofrecido & Historia consultable y medida preventiva mes 6; no son el mismo resultado.\newline{} FALTA REFERENCIA A SD9.\\\hline
\texttt{RT-05.\allowbreak{}16} & \texttt{B72}\newline{} Alcance ofrecido & Criterio de la ficha productora y prueba específica de la cláusula; no se acredita por ejecutar únicamente migración.\newline{} FALTA REFERENCIA A SD9.\\\hline
\texttt{RT-05.\allowbreak{}17} & \texttt{B72}\newline{} Alcance ofrecido & Criterio de la ficha productora y prueba específica de la cláusula; no se acredita por ejecutar únicamente migración.\newline{} FALTA REFERENCIA A SD9.\\\hline
\texttt{RT-05.\allowbreak{}18} & \texttt{B72;\allowbreak{} B71}\newline{} Alcance ofrecido & Criterio de la ficha productora y prueba específica de la cláusula; no se acredita por ejecutar únicamente migración.\newline{} FALTA REFERENCIA A SD9.\\\hline
\texttt{RT-05.\allowbreak{}19} & \texttt{B72;\allowbreak{} B70}\newline{} Alcance ofrecido & Criterio de la ficha productora y prueba específica de la cláusula; no se acredita por ejecutar únicamente migración.\newline{} FALTA REFERENCIA A SD9.\\\hline
\texttt{RT-05.\allowbreak{}20} & \texttt{B73}\newline{} Alcance ofrecido & Criterio de la ficha productora y prueba específica de la cláusula; no se acredita por ejecutar únicamente migración.\newline{} FALTA REFERENCIA A SD9.\\\hline
\texttt{RT-05.\allowbreak{}21} & \texttt{B73}\newline{} Alcance ofrecido & Criterio de la ficha productora y prueba específica de la cláusula; no se acredita por ejecutar únicamente migración.\newline{} FALTA REFERENCIA A SD9.\\\hline
\texttt{RT-05.\allowbreak{}22} & \texttt{B64}\newline{} Alcance ofrecido & Criterio de la ficha productora y prueba específica de la cláusula; no se acredita por ejecutar únicamente migración.\newline{} FALTA REFERENCIA A SD9.\\\hline
\texttt{RT-05.\allowbreak{}25} & \texttt{B63}\newline{} Alcance ofrecido & Criterio de la ficha productora y prueba específica de la cláusula; no se acredita por ejecutar únicamente migración.\newline{} FALTA REFERENCIA A SD9.\\\hline
\texttt{RT-05.\allowbreak{}26} & \texttt{B63}\newline{} Alcance ofrecido & Criterio de la ficha productora y prueba específica de la cláusula; no se acredita por ejecutar únicamente migración.\newline{} FALTA REFERENCIA A SD9.\\\hline
\texttt{RT-05.\allowbreak{}27} & \texttt{B63}\newline{} Alcance ofrecido & Criterio de la ficha productora y prueba específica de la cláusula; no se acredita por ejecutar únicamente migración.\newline{} FALTA REFERENCIA A SD9.\\\hline
\texttt{RT-05.\allowbreak{}28} & \texttt{B65}\newline{} Alcance ofrecido & Criterio de la ficha productora y prueba específica de la cláusula; no se acredita por ejecutar únicamente migración.\newline{} FALTA REFERENCIA A SD9.\\\hline
\texttt{RT-05.\allowbreak{}29} & \texttt{B66}\newline{} Alcance ofrecido & Criterio de la ficha productora y prueba específica de la cláusula; no se acredita por ejecutar únicamente migración.\newline{} FALTA REFERENCIA A SD9.\\\hline
\texttt{RT-03.\allowbreak{}24} & \texttt{B79}\newline{}Alcance obligatorio ofrecido & Criterio específico de la ficha productora, universo y versión E1/E2; cláusula BT/C07 diferenciada. FALTA REFERENCIA A SD9.\\\hline
\texttt{RT-05.\allowbreak{}02} & \texttt{B69}\newline{}Alcance obligatorio ofrecido & Criterio específico de la ficha productora, universo y versión E1/E2; cláusula BT/C07 diferenciada. FALTA REFERENCIA A SD9.\\\hline
\texttt{RT-05.\allowbreak{}07} & \texttt{B76}\newline{}Alcance obligatorio ofrecido & Criterio específico de la ficha productora, universo y versión E1/E2; cláusula BT/C07 diferenciada. FALTA REFERENCIA A SD9.\\\hline
\texttt{RT-05.\allowbreak{}08} & \texttt{B78;\allowbreak{} B68}\newline{}Alcance obligatorio ofrecido & Criterio específico de la ficha productora, universo y versión E1/E2; cláusula BT/C07 diferenciada. FALTA REFERENCIA A SD9.\\\hline
\texttt{RT-05.\allowbreak{}10} & \texttt{B76}\newline{}Alcance obligatorio ofrecido & Criterio específico de la ficha productora, universo y versión E1/E2; cláusula BT/C07 diferenciada. FALTA REFERENCIA A SD9.\\\hline
\texttt{RT-05.\allowbreak{}23} & \texttt{B80;\allowbreak{} B68}\newline{}Alcance obligatorio ofrecido & Criterio específico de la ficha productora, universo y versión E1/E2; cláusula BT/C07 diferenciada. FALTA REFERENCIA A SD9.\\\hline
\texttt{RT-16.\allowbreak{}14} & \texttt{B77;\allowbreak{} B68}\newline{}Alcance obligatorio ofrecido & Criterio específico de la ficha productora, universo y versión E1/E2; cláusula BT/C07 diferenciada. FALTA REFERENCIA A SD9.\\\hline
\texttt{RT-16.\allowbreak{}10} & \texttt{B76;\allowbreak{} B70}\newline{}Alcance obligatorio ofrecido & Criterio específico de la ficha productora, universo y versión E1/E2; cláusula BT/C07 diferenciada. FALTA REFERENCIA A SD9.\\\hline
\texttt{RT-26.\allowbreak{}08} & \texttt{---}\newline{}Deseable BT no ofrecido & No procede: se conserva exclusión T-12; IN1 línea base mes 19, piloto 20 y despliegue 21. No genera un compromiso de beneficio antes del mes 16.\\\hline
\end{longtable}\endgroup
\section{Programa de trabajo y Gantt de 56 meses}\label{sec:t14-programa}
\subsection{Secuencia, dependencias y calendario}
El programa parte de la línea base de alcance y criterios (H1), seguida de arquitectura, seguridad y datos (H2), habilitación técnica (H3), versiones de software, certificación, marcha blanca y producción por etapa. H3 exige además DR conforme a RT-04.01, aunque E-25 enumere cuatro ambientes. Las decisiones, permisos y recepciones externas son puertas de dependencia y no horas productivas del equipo.

Durante meses 13--15 se separan el frente de marcha blanca E1 y el frente de desarrollo E2. Durante 19--20 se protege E1 productiva mientras E2 se observa con datos y usuarios reales. Las capacidades y responsables concurrentes se demuestran en T-15. Hasta el acta y cambio oficial, cada dato compartido conserva una sola autoridad de escritura y los efectos externos se ejecutan una sola vez; T-18 detalla cómo se aplica la regla durante la convivencia.

\begingroup\setlength{\tabcolsep}{3pt}
\begin{longtable}{L{\dimexpr0.219355\textwidth-2\tabcolsep\relax}L{\dimexpr0.225806\textwidth-2\tabcolsep\relax}L{\dimexpr0.477419\textwidth-2\tabcolsep\relax}L{\dimexpr0.077420\textwidth-2\tabcolsep\relax}}
\caption{Dependencias de los principales resultados del programa}\label{tab:t14-secuencia}\\\hline
\EncabezadoTabla{Precedencia} & \EncabezadoTabla{Cuentas de trabajo} & \EncabezadoTabla{Puerta de habilitación} & \EncabezadoTabla{Meses}\\\hline\endfirsthead
\multicolumn{4}{l}{\textit{Continuación de la tabla \ref{tab:t14-secuencia}}}\\\hline
\EncabezadoTabla{Precedencia} & \EncabezadoTabla{Cuentas de trabajo} & \EncabezadoTabla{Puerta de habilitación} & \EncabezadoTabla{Meses}\\\hline\endhead
\hline\endfoot
\hline\endlastfoot
Inicio $\rightarrow$ H1 & 1.1; delimitación del alcance funcional; diagnóstico 1.8.1 & Accesos y contraparte designados; alcance, requisitos y criterios individualizados. & 1--2\\\hline
H1 $\rightarrow$ H2 & 1.2; 1.3; diseño 1.4--1.7; reglas/modelos 1.8.2 & Diseño y seguridad recibidos; dependencias externas e interfaces delimitadas. & 2--4\\\hline
Preparación desde H1; H2 $\rightarrow$ H3 & 1.2; 1.3; 1.7 & Cinco ambientes, borde, red, terminales y observabilidad comprobados; compra/obra externa recibida. & 3--6\\\hline
H2/H3 $\rightarrow$ versiones E1 $\rightarrow$ H4/H5 & 1.4--1.8; 1.10; 1.11 & Software, herramientas de migración, pruebas y usuarios preparados; revisión/subsanación preceden recepción. & 4--12\\\hline
H5/H6 $\rightarrow$ marcha E1 $\rightarrow$ H7 & 1.10; 1.11, incluido soporte 1.11.8; preparación 1.12.1 & Cuatro semanas finales con operación real y condiciones copulativas; acta antes del cambio oficial. & 13--16\\\hline
Base E1/H8 $\rightarrow$ versiones E2 $\rightarrow$ H9/H10 & 1.4--1.9; 1.10; 1.11 & Reutilización de capacidades recibidas; E1 protegida; pruebas y ensayos sobre versión/delta E2. & 13--18\\\hline
H10/H11 $\rightarrow$ marcha E2 $\rightarrow$ H12/Operación & 1.10; 1.11; 1.12 & Evidencia completa, decisión expresa y activación ordenadas; no se presume firma instantánea. & 19--21\\\hline
Servicio activo $\rightarrow$ controles recurrentes $\rightarrow$ salida & 1.12; campaña IN2 y seguimiento de innovaciones recibidas & SLA, respaldo, DR, seguridad, mantenimiento, formación, custodia y transferencia por período. & 21--56\\\hline
\end{longtable}\endgroup

\begingroup\setlength{\tabcolsep}{3pt}
\begin{longtable}{L{\dimexpr0.264151\textwidth-2\tabcolsep\relax}L{\dimexpr0.320755\textwidth-2\tabcolsep\relax}L{\dimexpr0.415094\textwidth-2\tabcolsep\relax}}
\caption{Habilitaciones externas del escenario}\label{tab:t14-externas}\\\hline
\EncabezadoTabla{ID} & \EncabezadoTabla{Insumo externo} & \EncabezadoTabla{Resultado bloqueado si falta}\\\hline\endfirsthead
\multicolumn{3}{l}{\textit{Continuación de la tabla \ref{tab:t14-externas}}}\\\hline
\EncabezadoTabla{ID} & \EncabezadoTabla{Insumo externo} & \EncabezadoTabla{Resultado bloqueado si falta}\\\hline\endhead
\hline\endfoot
\hline\endlastfoot
\texttt{EXT-RECINTO} & Obras y condiciones del recinto & Infraestructura del sitio y pruebas técnicas previas a H3.\\\hline
\texttt{EXT-EQUIPOS /\allowbreak{} EXT-LICENCIAS} & Activos conformes al T-11 y licencias efectivas & Instalación, enrolamiento, ambientes y soporte antes de su recepción.\\\hline
\texttt{EXT-RED /\allowbreak{} EXT-TELEMETRIA} & Enlaces, tendidos, buses, concentradores y acceso & Pruebas de segregación, cobertura, respaldo e integración de medición antes de H3.\\\hline
\texttt{EXT-IDENTIDAD} & Fuente laboral, responsables y permisos & Integración de identidad y autoridad local antes de su comprobación.\\\hline
\texttt{EXT-CONTABLE /\allowbreak{} EXT-PAGOS} & Contratos de interfaz y acceso a terceros & Adaptadores y validación E2 antes de H9/H10.\\\hline
\end{longtable}\endgroup


Fernando Guerrero coordina las interfaces técnicas; Ignacio Silva controla compromisos y decisiones con la contraparte designada. Cada habilitación debe registrar responsable externo, fecha requerida/ofrecida, condición de recepción, evidencia y efecto en el programa. La disponibilidad de una cuenta, licencia, interfaz u obra no se presume por incluirla en esta tabla.

Los meses contractuales gobiernan las fases. La ingeniería se programa en jornadas hábiles y las guardias/observación se programan en tiempo natural, sin transformar noches o fines de semana en horas de ingeniería diurna. El escenario conserva un inicio de referencia en diciembre de 2026. El calendario de feriados futuro es una hipótesis a confirmar antes de fijar fechas diarias; T-15 calcula duraciones, esperas y holguras con recursos efectivos.

Se conserva el congelamiento estival del 15 de diciembre al 15 de marzo recogido en SD3. Las habilitaciones y sesiones nuevas E1 se anticipan al 15 de diciembre; durante el congelamiento continúa el acompañamiento de procedimientos ya activos y la observación permitida. PREPROD aislado no autoriza intervenir producción, usuarios de sitio o maniobras. Ante aviso de marejadas se suspende inmediatamente toda actividad programada del proyecto, incluida construcción aislada o remota, conforme a SD6. El Jefe de Proyecto o Gerente de Servicio registra trabajo, capacidad y reserva afectada; reprograma solo al cesar la condición y verificar una ventana admisible. Atención de una emergencia real sigue el procedimiento de incidente y no habilita trabajo programado. Eventos y cargas suspendidas mantienen sus restricciones particulares. Las reservas se cuantifican por exposición y se contabilizan una sola vez en T-15. La primera campaña IN2 se conserva en junio de 2029, mes 31 del escenario, y la evaluación ambiental debe proteger la temporada 2029.

\subsection{Fases e hitos de recepción}
Las ventanas de la tabla~\ref{tab:t14-fases} son exigencias de las Bases Administrativas, art. 17, y no se sustituyen por las fechas de un escenario de atraso. Las fechas absolutas de la columna final derivan exclusivamente del supuesto de inicio ya señalado.

\begingroup\setlength{\tabcolsep}{3pt}
\begin{longtable}{L{\dimexpr0.227273\textwidth-2\tabcolsep\relax}L{\dimexpr0.103896\textwidth-2\tabcolsep\relax}L{\dimexpr0.461039\textwidth-2\tabcolsep\relax}L{\dimexpr0.207792\textwidth-2\tabcolsep\relax}}
\caption{Ventanas contractuales del programa}\label{tab:t14-fases}\\\hline
\EncabezadoTabla{Fase} & \EncabezadoTabla{Meses} & \EncabezadoTabla{Contenido} & \EncabezadoTabla{Referencia absoluta}\\\hline\endfirsthead
\multicolumn{4}{l}{\textit{Continuación de la tabla \ref{tab:t14-fases}}}\\\hline
\EncabezadoTabla{Fase} & \EncabezadoTabla{Meses} & \EncabezadoTabla{Contenido} & \EncabezadoTabla{Referencia absoluta}\\\hline\endhead
\hline\endfoot
\hline\endlastfoot
Desarrollo E1 & 1--12 & Ingeniería, infraestructura, integración, migración, pruebas y certificación. & Dic. 2026--nov. 2027\\\hline
Marcha blanca E1 & 13--15 & Tres meses con datos/usuarios reales, convivencia, retorno activo e indicadores diarios. & Dic. 2027--feb. 2028\\\hline
Producción E1 & 16 & Cambio a registro oficial del alcance E1, previa recepción. & Mar. 2028\\\hline
Desarrollo E2 & 13--18 & Segundo alcance, sin degradar E1; cierre inclusive en mes 18. & Dic. 2027--may. 2028\\\hline
Marcha blanca E2 & 19--20 & Dos meses supervisados, con E1 productiva y datos íntegros. & Jun.--jul. 2028\\\hline
Producción E2 e inicio de Operación & 21 & Ambos alcances bajo servicio desde el primer día; acta previa al cambio oficial. & Ago. 2028\\\hline
Operación & 21--56 & 36 meses de servicio continuo para ambos alcances. & Ago. 2028--jul. 2031\\\hline
\end{longtable}\endgroup

\begingroup\setlength{\tabcolsep}{3pt}
\begin{longtable}{L{\dimexpr0.076923\textwidth-2\tabcolsep\relax}L{\dimexpr0.160256\textwidth-2\tabcolsep\relax}L{\dimexpr0.461538\textwidth-2\tabcolsep\relax}L{\dimexpr0.301283\textwidth-2\tabcolsep\relax}}
\caption{Hitos de implementación vinculados a E-25}\label{tab:t14-hitos}\\\hline
\EncabezadoTabla{Hito} & \EncabezadoTabla{Mes / fecha} & \EncabezadoTabla{Entregable recibido} & \EncabezadoTabla{Dependencia de recepción}\\\hline\endfirsthead
\multicolumn{4}{l}{\textit{Continuación de la tabla \ref{tab:t14-hitos}}}\\\hline
\EncabezadoTabla{Hito} & \EncabezadoTabla{Mes / fecha} & \EncabezadoTabla{Entregable recibido} & \EncabezadoTabla{Dependencia de recepción}\\\hline\endhead
\hline\endfoot
\hline\endlastfoot
\texttt{H1} & 2\newline{}ene. 2027 & Alcance y matriz de trazabilidad aprobados & 1.1; fichas de alcance; T-12\\\hline
\texttt{H2} & 4\newline{}mar. 2027 & Arquitectura, seguridad y modelo de datos aprobados & 1.2; 1.3; 1.7; 1.8.2; SD4/SD5\\\hline
\texttt{H3} & 6\newline{}may. 2027 & Infraestructura híbrida y cuatro ambientes con observabilidad; DR adicional por RT-04.01 & 1.2; 1.3; 1.7; habilitaciones externas\\\hline
\texttt{H4} & 10\newline{}sep. 2027 & Software E1 con QA superado y evidencia de cobertura & Versión E1; 1.4--1.8; 1.10.2.1\\\hline
\texttt{H5} & 12\newline{}nov. 2027 & E1 certificada: UAT, carga, resiliencia y seguridad ofensiva & H4; expediente 1.10; observaciones cerradas\\\hline
\texttt{H6} & 13\newline{}dic. 2027 & Inicio de marcha blanca E1, retorno activo y usuarios capacitados & H5; 1.11; datos y cobertura habilitados\\\hline
\texttt{H7} & 16\newline{}mar. 2028 & Producción E1 y cierre conforme de marcha blanca & H6; cuatro semanas finales; acta art. 17.3\\\hline
\texttt{H8} & 14\newline{}ene. 2028 & Alcance E2 y diseño detallado aprobados & Base E1; alcance E2; H2 actualizado\\\hline
\texttt{H9} & 17\newline{}abr. 2028 & Software E2 con QA superado & H8; versión E2; 1.10.2.2\\\hline
\texttt{H10} & 18\newline{}may. 2028 & Certificación E2 y cierre de desarrollo & H9; ensayos y pruebas E2; observaciones cerradas\\\hline
\texttt{H11} & 19\newline{}jun. 2028 & Inicio de marcha blanca E2 con E1 en producción & H10; 1.11; cobertura y reversión habilitadas\\\hline
\texttt{H12} & 21\newline{}ago. 2028 & Producción E2 y aceptación final de implementación & H11; condiciones art. 17.3; acta y activación\\\hline
\end{longtable}\endgroup


Los porcentajes, montos y condiciones económicas se mantienen en E-25 y la Oferta Económica. La tabla conserva sus doce hitos y condiciones de entregable para programar las recepciones. Cada hito requiere acta suscrita por la Contraparte Técnica, conforme al art. 18.1. Las ventanas de revisión y subsanación del art. 18.3 se insertan en la red; una fecha de presentación no equivale a la de aprobación.

El cierre de marcha blanca exige conjuntamente ausencia de incidentes críticos/altos atribuibles, volumen real comprometido durante al menos las cuatro últimas semanas, disponibilidad y respuesta sostenidas durante ese período, conciliación sin diferencias no explicadas, capacitación/certificación y acta. T-18 desarrolla indicadores, denominadores y el procedimiento; T-17 conserva el protocolo de recepción. Si faltan condiciones, se bloquea el cambio y se registra la desviación: extender la marcha blanca a costo del adjudicatario no desplaza las fases siguientes.

Para E2, la evidencia completa debe preceder la decisión y esta debe preceder la activación. No se ofrece como demostrada la secuencia que cierra medición al final del mes 20 y supone firma/corte instantáneos al inicio del 21. El objetivo contractual del primer día se conserva; su solución temporal requiere duraciones positivas, coordinación expresa y evidencia de cabida en T-15/T-18, sin reducir las cuatro semanas ni los derechos de revisión. Mientras eso no esté conciliado, permanece una dependencia abierta del programa.

\subsection{Innovaciones y ensayos vinculados al corte}
IN1 conserva diseño en 14--15, construcción/integración en 16--17, pruebas en 18 y línea base en 19 con flujo adicional desactivado. La capacitación comienza después de cerrar la línea base y termina antes del piloto del mes 20; T-15/T-18 precisan su intervalo con disponibilidad de usuarios. Piloto y evaluación se realizan en 20; despliegue operativo en 21, previa aceptación. No se mezcla capacitación con la medición sin flujo para aparentar una línea base comparable.

IN2 conserva diseño 14--15, implementación 16--17, integración 17--18 y pruebas 19--20. Piloto desde el primer junio admisible con IN1 e interfaces recibidas, no anterior al mes 21. En el escenario actual se programa de mes 31 a 36; evaluación y transferencia posteriores deben conciliarse con SD13/T-19. FALTA REFERENCIA A SD13. Los controles de botadura preservan anticipación y restricciones de faena; no se considera recibido el piloto dentro de marcha blanca E2 por haber concluido su diseño.

\textbf{Dependencia temporal de IN4.} Las interfaces de \texttt{1.9.4.2.1} se especifican y reciben en 14--15 después de recibir el flujo, estados y roles de IN1 \texttt{1.9.1.1.1}. Captura \texttt{1.9.4.2.2} en 16; cálculo \texttt{1.9.4.2.3} y pruebas \texttt{1.9.4.2.4} en 17, en ese orden y con correcciones antes de iniciar la línea base de 18. Compartir mes no elimina la precedencia. Se conservan evaluación en 19 sin efectos de pago y habilitación contractual en 21. T-15 debe demostrar duraciones y disponibilidad para esa secuencia.

\begingroup\setlength{\tabcolsep}{3pt}
\begin{longtable}{L{\dimexpr.16\textwidth-2\tabcolsep\relax}L{\dimexpr.40\textwidth-2\tabcolsep\relax}L{\dimexpr.44\textwidth-2\tabcolsep\relax}}
\caption{Secuencia de innovaciones tomada del diccionario vigente}\label{tab:t14-innovaciones-programa}\\\hline
\EncabezadoTabla{Innovación} & \EncabezadoTabla{Preparación y verificación} & \EncabezadoTabla{Medición, evaluación y transferencia}\\\hline\endfirsthead
\multicolumn{3}{l}{\textit{Continuación de la tabla \ref{tab:t14-innovaciones-programa}}}\\\hline
\EncabezadoTabla{Innovación} & \EncabezadoTabla{Preparación y verificación} & \EncabezadoTabla{Medición, evaluación y transferencia}\\\hline\endhead
\hline\endfoot\hline\endlastfoot
IN1 & Diseño 14--15; construcción/integración 16--17; pruebas 18. & Base 19; capacitación posterior a base y previa a piloto 20; evaluación 20; despliegue 21.\\\hline
IN2 & Diseño 14--15; construcción 16--17; integración 17--18; pruebas 19--20; inscripción/planes/formación 30. & Piloto y resultados 31--36; revisiones de botadura 33--35; evaluación/transferencia 37.\\\hline
IN3 & Inventario 13; interfaces y construcción 14; integración 15; pruebas y preparación 16. & Registros 18--19; evaluación 19; transferencia 21. La ventana de la ficha gobierna; un sufijo histórico no sustituye su mes.\\\hline
IN4 & Ficha 2 y regla 4; especificación 14--15 tras diseño IN1; captura 16; cálculo seguido de pruebas 17. & Base 18; evaluación 19 sin efectos de pago; habilitación contractual 21.\\\hline
IN5 & Plan 13; diseño/cálculo 14; portal/integración 15; pruebas/preparación 16. & Base 18; evaluación 19; transferencia 21; seguimiento protege la temporada ambiental 2029 conforme al alcance vigente.\\\hline
\end{longtable}\endgroup

En migración, C identifica el cambio de autoridad del grupo. T-18 individualiza grupo, versión, autorizador y fecha admisible; T-15 convierte reservas de SD5 en actividades, recursos y dependencias. Los antiguos meses fijos 10/11 y 16/17 no demuestran que una evidencia sea vigente al corte. Si se usan como comprobaciones tempranas, se identifica trabajo incremental de actualización sin volver a imputar la misma ejecución. Dos ensayos completos deben estar aprobados antes del primer cambio de autoridad; cambios posteriores actualizan evidencia y pueden exigir nueva simulación definitiva. F27/F28 se activan desde el hallazgo, corrigen y reensayan antes de autorizar C. La reserva de dos horas para la revisión operativa de habilitación no sustituye los derechos de revisión contractual del art. 18.3; T-15 debe ubicar presentación, revisión y subsanación con tiempos positivos antes del corte.
\begingroup\setlength{\tabcolsep}{3pt}
\begin{longtable}{L{\dimexpr0.230000\textwidth-2\tabcolsep\relax}L{\dimexpr0.400000\textwidth-2\tabcolsep\relax}L{\dimexpr0.370000\textwidth-2\tabcolsep\relax}}
\caption{Preparación relativa al corte C, según SD5}\label{tab:t14-cortes}\\\hline
\EncabezadoTabla{Ventana} & \EncabezadoTabla{Resultado y reserva} & \EncabezadoTabla{Condición}\\\hline\endfirsthead
\multicolumn{3}{l}{\textit{Continuación de la tabla \ref{tab:t14-cortes}}}\\\hline
\EncabezadoTabla{Ventana} & \EncabezadoTabla{Resultado y reserva} & \EncabezadoTabla{Condición}\\\hline\endhead
\hline\endfoot
\hline\endlastfoot
C menos 6 a 4 semanas & Inventario, extracción y perfilado; tres semanas & Manifiestos y defectos clasificados.\\\hline
C menos 3 semanas & Ensayo completo 1; dos jornadas de 8 h & Carga conciliada y retorno probado.\\\hline
C menos 2 semanas & Corrección y precarga; una semana & Reglas y versiones estabilizadas.\\\hline
C menos 1 semana & Ensayo completo 2; 8 h de precarga y 4 h de simulación & Delta y ventana comprobados; duraciones observadas actualizan la reserva.\\\hline
C menos 2 días hábiles & Revisión de habilitación; dos horas & Autorización o reprogramación, sin presumir permiso concedido.\\\hline
Desde C & Conciliación diaria y de cierre mensual & Diferencias y decisiones registradas.\\\hline
\end{longtable}\endgroup
\subsection{Carta Gantt integrada}
La carta cubre los 56 meses y conserva las fases del art. 17 y los doce hitos E-25. Es el programa de referencia de T-14 para construir la red y nivelar recursos en T-15. Las barras por cuenta son la unión de las ventanas numéricas de sus fichas: no representan una actividad continua, ni duración calculada, ni HH asignadas a los padres. Los espacios entre ventanas no eliminan las actividades de construcción que T-15 debe secuenciar entre diseño y entrega. Los detalles de innovación, mesa/guardia y salida ya están incluidos en las cuentas; no se vuelven a sumar como trabajo.

\textbf{Lectura de las figuras.} Oscuro identifica fases contractuales; azul, ventanas del diccionario; claro, detalle de resultados ya incluidos; rombos, objetivos de hitos cuya recepción requiere acta. La ubicación mensual del rombo no fija día ni acredita aprobación. H12 y Operación conservan el objetivo del primer día de 21; la cabida de revisión y decisión previa permanece por demostrar. Ninguna barra se marca como crítica antes de calcular CPM y nivelar recursos.

\textbf{Resultados relativos a C.} Las filas 1.8 y 1.10 muestran sus ventanas numéricas y se complementan con la tabla~\ref{tab:t14-cortes}. Los ensayos y F27/F28 se anclan al cambio de autoridad de cada grupo identificado en T-18, no automáticamente a los meses 16/21 ni a los antiguos 10/11 y 16/17. Los grupos activados en 13/19 requieren preparación y ensayos aprobados antes de su propio C. T-15 debe incorporar esas instancias y actualizar la carta al fijar los cortes; su ausencia como barra mensual fija no implica excluirlos del alcance.

\textbf{Cobertura y salida.} Preparación del soporte: \texttt{1.12.1.1} en 11--12/18 y \texttt{1.12.1.2} en 12/20; prestación de mesa y guardia \texttt{1.11.8} en 13--20; Operación \texttt{1.12.14} en 21--56. El detalle de servicio distingue esas dos prestaciones sin adelantar Operación. La salida conserva inventario y exportación inicial en 53; fuentes/configuración/custodia desde 54 y actualización hasta el traspaso de 56. \texttt{1.12.9} acredita al menos 90 días reales de acompañamiento; \texttt{1.12.8.5} recibe el estado final conciliado y autoriza la eliminación solo cuando corresponde. El servicio continúa durante la transferencia.

\textbf{Cantidades y calendarios pendientes.} Las barras de formación identifican ventanas admisibles, no cohortes confirmadas. Los manifiestos determinan $N_F$ y $N_E$; las olas no se fijan por el número de códigos históricos. Congelamiento, marejadas, faenas, feriados y suplencias se aplican al convertir las ventanas en actividades diarias. Holguras, reservas y cabida de ejecución/revisión se calculan en T-15; T-18 define su aplicación operacional. La tabla~\ref{tab:t14-innovaciones-programa} conserva las precedencias que una escala mensual no puede mostrar dentro de cada mes.

\begin{paginaHorizontalSynaptix}
\begin{figure}[p]\centering
\begin{tikzpicture}[x=.52cm,y=1cm,font=\fontsize{9}{10}\selectfont]
\node at (0.50,.48) {1};\draw[SynLine!55] (0,.20)--(0,-12.75);
\node at (1.50,.48) {2};\draw[SynLine!55] (1,.20)--(1,-12.75);
\node at (2.50,.48) {3};\draw[SynLine!55] (2,.20)--(2,-12.75);
\node at (3.50,.48) {4};\draw[SynLine!55] (3,.20)--(3,-12.75);
\node at (4.50,.48) {5};\draw[SynLine!55] (4,.20)--(4,-12.75);
\node at (5.50,.48) {6};\draw[SynLine!55] (5,.20)--(5,-12.75);
\node at (6.50,.48) {7};\draw[SynLine!55] (6,.20)--(6,-12.75);
\node at (7.50,.48) {8};\draw[SynLine!55] (7,.20)--(7,-12.75);
\node at (8.50,.48) {9};\draw[SynLine!55] (8,.20)--(8,-12.75);
\node at (9.50,.48) {10};\draw[SynLine!55] (9,.20)--(9,-12.75);
\node at (10.50,.48) {11};\draw[SynLine!55] (10,.20)--(10,-12.75);
\node at (11.50,.48) {12};\draw[SynLine!55] (11,.20)--(11,-12.75);
\node at (12.50,.48) {13};\draw[SynLine!55] (12,.20)--(12,-12.75);
\node at (13.50,.48) {14};\draw[SynLine!55] (13,.20)--(13,-12.75);
\node at (14.50,.48) {15};\draw[SynLine!55] (14,.20)--(14,-12.75);
\node at (15.50,.48) {16};\draw[SynLine!55] (15,.20)--(15,-12.75);
\node at (16.50,.48) {17};\draw[SynLine!55] (16,.20)--(16,-12.75);
\node at (17.50,.48) {18};\draw[SynLine!55] (17,.20)--(17,-12.75);
\node at (18.50,.48) {19};\draw[SynLine!55] (18,.20)--(18,-12.75);
\node at (19.50,.48) {20};\draw[SynLine!55] (19,.20)--(19,-12.75);
\node at (20.50,.48) {21};\draw[SynLine!55] (20,.20)--(20,-12.75);
\node at (21.50,.48) {22};\draw[SynLine!55] (21,.20)--(21,-12.75);
\node at (22.50,.48) {23};\draw[SynLine!55] (22,.20)--(22,-12.75);
\node at (23.50,.48) {24};\draw[SynLine!55] (23,.20)--(23,-12.75);
\node at (24.50,.48) {25};\draw[SynLine!55] (24,.20)--(24,-12.75);
\node at (25.50,.48) {26};\draw[SynLine!55] (25,.20)--(25,-12.75);
\node at (26.50,.48) {27};\draw[SynLine!55] (26,.20)--(26,-12.75);
\node at (27.50,.48) {28};\draw[SynLine!55] (27,.20)--(27,-12.75);
\draw[SynLine!55] (28,.20)--(28,-12.75);
\node[anchor=east,align=right,text width=5.5cm,inner sep=0pt] at (-.25,-0.00) {Desarrollo E1};
\fill[SynNavy!75] (0.10,-0.13) rectangle (11.90,0.13);
\node[anchor=east,align=right,text width=5.5cm,inner sep=0pt] at (-.25,-0.48) {Marcha blanca E1};
\fill[SynNavy!75] (12.10,-0.61) rectangle (14.90,-0.35);
\node[anchor=east,align=right,text width=5.5cm,inner sep=0pt] at (-.25,-0.96) {Paso a producción E1};
\node[fill=SynNavy,draw=SynNavy,diamond,inner sep=2pt] at (15.50,-0.96) {};
\node[anchor=east,align=right,text width=5.5cm,inner sep=0pt] at (-.25,-1.44) {Desarrollo E2};
\fill[SynNavy!75] (12.10,-1.57) rectangle (17.90,-1.31);
\node[anchor=east,align=right,text width=5.5cm,inner sep=0pt] at (-.25,-1.92) {Marcha blanca E2};
\fill[SynNavy!75] (18.10,-2.05) rectangle (19.90,-1.79);
\node[anchor=east,align=right,text width=5.5cm,inner sep=0pt] at (-.25,-2.40) {Producción E2 / Operación};
\node[fill=SynNavy,draw=SynNavy,diamond,inner sep=2pt] at (20.50,-2.40) {};
\node[anchor=east,align=right,text width=5.5cm,inner sep=0pt] at (-.25,-2.88) {Operación de ambos alcances};
\fill[SynNavy!75] (20.10,-3.01) rectangle (27.90,-2.75);
\node[anchor=east,align=right,text width=5.5cm,inner sep=0pt] at (-.25,-3.36) {Hitos E-25};
\node[fill=SynNavy,draw=SynNavy,diamond,inner sep=1.5pt] at (1.50,-3.36) {};
\node at (1.50,-3.20) {H1};
\node[fill=SynNavy,draw=SynNavy,diamond,inner sep=1.5pt] at (3.50,-3.36) {};
\node at (3.50,-3.52) {H2};
\node[fill=SynNavy,draw=SynNavy,diamond,inner sep=1.5pt] at (5.50,-3.36) {};
\node at (5.50,-3.20) {H3};
\node[fill=SynNavy,draw=SynNavy,diamond,inner sep=1.5pt] at (9.50,-3.36) {};
\node at (9.50,-3.52) {H4};
\node[fill=SynNavy,draw=SynNavy,diamond,inner sep=1.5pt] at (11.50,-3.36) {};
\node at (11.50,-3.20) {H5};
\node[fill=SynNavy,draw=SynNavy,diamond,inner sep=1.5pt] at (12.50,-3.36) {};
\node at (12.50,-3.52) {H6};
\node[fill=SynNavy,draw=SynNavy,diamond,inner sep=1.5pt] at (13.50,-3.36) {};
\node at (13.50,-3.20) {H8};
\node[fill=SynNavy,draw=SynNavy,diamond,inner sep=1.5pt] at (15.50,-3.36) {};
\node at (15.50,-3.52) {H7};
\node[fill=SynNavy,draw=SynNavy,diamond,inner sep=1.5pt] at (16.50,-3.36) {};
\node at (16.50,-3.20) {H9};
\node[fill=SynNavy,draw=SynNavy,diamond,inner sep=1.5pt] at (17.50,-3.36) {};
\node at (17.50,-3.52) {H10};
\node[fill=SynNavy,draw=SynNavy,diamond,inner sep=1.5pt] at (18.50,-3.36) {};
\node at (18.50,-3.20) {H11};
\node[fill=SynNavy,draw=SynNavy,diamond,inner sep=1.5pt] at (20.50,-3.36) {};
\node at (20.50,-3.52) {H12};
\node[anchor=east,align=right,text width=5.5cm,inner sep=0pt] at (-.25,-3.84) {1.1 Dirección y control};
\fill[SynBlue!65] (0.10,-3.97) rectangle (19.90,-3.71);
\node[anchor=east,align=right,text width=5.5cm,inner sep=0pt] at (-.25,-4.32) {1.2 Plataforma híbrida};
\fill[SynBlue!65] (2.10,-4.45) rectangle (5.90,-4.19);
\node[anchor=east,align=right,text width=5.5cm,inner sep=0pt] at (-.25,-4.80) {1.3 Auditoría y seguridad};
\fill[SynBlue!65] (3.10,-4.93) rectangle (8.90,-4.67);
\fill[SynBlue!65] (12.10,-4.93) rectangle (14.90,-4.67);
\node[anchor=east,align=right,text width=5.5cm,inner sep=0pt] at (-.25,-5.28) {1.4 Núcleo Operacional};
\fill[SynBlue!65] (2.10,-5.41) rectangle (2.90,-5.15);
\fill[SynBlue!65] (7.10,-5.41) rectangle (7.90,-5.15);
\node[anchor=east,align=right,text width=5.5cm,inner sep=0pt] at (-.25,-5.76) {1.5 Núcleo de Servicios};
\fill[SynBlue!65] (2.10,-5.89) rectangle (2.90,-5.63);
\fill[SynBlue!65] (7.10,-5.89) rectangle (7.90,-5.63);
\fill[SynBlue!65] (12.10,-5.89) rectangle (14.90,-5.63);
\node[anchor=east,align=right,text width=5.5cm,inner sep=0pt] at (-.25,-6.24) {1.6 Núcleo Administrativo};
\fill[SynBlue!65] (2.10,-6.37) rectangle (2.90,-6.11);
\fill[SynBlue!65] (7.10,-6.37) rectangle (7.90,-6.11);
\fill[SynBlue!65] (12.10,-6.37) rectangle (14.90,-6.11);
\node[anchor=east,align=right,text width=5.5cm,inner sep=0pt] at (-.25,-6.72) {1.7 Canales e integración};
\fill[SynBlue!65] (1.10,-6.85) rectangle (17.90,-6.59);
\node[anchor=east,align=right,text width=5.5cm,inner sep=0pt] at (-.25,-7.20) {1.8 Migración e historia (+ C)};
\fill[SynBlue!65] (0.10,-7.33) rectangle (2.90,-7.07);
\fill[SynBlue!65] (5.10,-7.33) rectangle (7.90,-7.07);
\fill[SynBlue!65] (13.10,-7.33) rectangle (13.90,-7.07);
\node[anchor=east,align=right,text width=5.5cm,inner sep=0pt] at (-.25,-7.68) {1.9 Innovaciones};
\fill[SynBlue!65] (1.10,-7.81) rectangle (1.90,-7.55);
\fill[SynBlue!65] (3.10,-7.81) rectangle (3.90,-7.55);
\fill[SynBlue!65] (12.10,-7.81) rectangle (20.90,-7.55);
\node[anchor=east,align=right,text width=5.5cm,inner sep=0pt] at (-.25,-8.16) {1.10 Verificación (+ C)};
\fill[SynBlue!65] (8.10,-8.29) rectangle (11.90,-8.03);
\fill[SynBlue!65] (14.10,-8.29) rectangle (14.90,-8.03);
\fill[SynBlue!65] (16.10,-8.29) rectangle (17.90,-8.03);
\node[anchor=east,align=right,text width=5.5cm,inner sep=0pt] at (-.25,-8.64) {1.11 Formación e implantación};
\fill[SynBlue!65] (10.10,-8.77) rectangle (20.90,-8.51);
\node[anchor=east,align=right,text width=5.5cm,inner sep=0pt] at (-.25,-9.12) {1.12 Servicio y salida};
\fill[SynBlue!65] (10.10,-9.25) rectangle (11.90,-8.99);
\fill[SynBlue!65] (17.10,-9.25) rectangle (17.90,-8.99);
\fill[SynBlue!65] (19.10,-9.25) rectangle (27.90,-8.99);
\node[anchor=east,align=right,text width=5.5cm,inner sep=0pt] at (-.25,-9.60) {IN1: solicitud / piloto / despliegue};
\fill[SynBlue!30] (13.10,-9.73) rectangle (20.90,-9.47);
\node[anchor=east,align=right,text width=5.5cm,inner sep=0pt] at (-.25,-10.08) {IN2: preparación / campaña};
\fill[SynBlue!30] (13.10,-10.21) rectangle (19.90,-9.95);
\node[anchor=east,align=right,text width=5.5cm,inner sep=0pt] at (-.25,-10.56) {IN3: extracción asistida};
\fill[SynBlue!30] (12.10,-10.69) rectangle (15.90,-10.43);
\fill[SynBlue!30] (17.10,-10.69) rectangle (18.90,-10.43);
\fill[SynBlue!30] (20.10,-10.69) rectangle (20.90,-10.43);
\node[anchor=east,align=right,text width=5.5cm,inner sep=0pt] at (-.25,-11.04) {IN4: medición e incentivo};
\fill[SynBlue!30] (1.10,-11.17) rectangle (1.90,-10.91);
\fill[SynBlue!30] (3.10,-11.17) rectangle (3.90,-10.91);
\fill[SynBlue!30] (13.10,-11.17) rectangle (18.90,-10.91);
\fill[SynBlue!30] (20.10,-11.17) rectangle (20.90,-10.91);
\node[anchor=east,align=right,text width=5.5cm,inner sep=0pt] at (-.25,-11.52) {IN5: pasaporte ambiental};
\fill[SynBlue!30] (12.10,-11.65) rectangle (15.90,-11.39);
\fill[SynBlue!30] (17.10,-11.65) rectangle (18.90,-11.39);
\fill[SynBlue!30] (20.10,-11.65) rectangle (20.90,-11.39);
\node[anchor=east,align=right,text width=5.5cm,inner sep=0pt] at (-.25,-12.00) {Mesa y guardia prestadas};
\fill[SynBlue!30] (12.10,-12.13) rectangle (27.90,-11.87);
\draw[SynNavy,line width=.6pt] (20.00,-12.16)--(20.00,-11.84);
\node[anchor=east,align=right,text width=5.5cm,inner sep=0pt] at (-.25,-12.48) {Salida y acompañamiento};
\end{tikzpicture}
\par\smallskip{\fontsize{9}{11}\selectfont Oscuro: fase contractual; azul: ventanas numéricas del diccionario; claro: detalle de filas ya incluidas; rombo: objetivo de hito.}
\caption{Gantt de referencia: meses 1--28; entregables, fases y cobertura}\label{fig:t14-gantt-1}
\end{figure}
\end{paginaHorizontalSynaptix}
\begin{paginaHorizontalSynaptix}
\begin{figure}[p]\centering
\begin{tikzpicture}[x=.52cm,y=1cm,font=\fontsize{9}{10}\selectfont]
\node at (0.50,.48) {29};\draw[SynLine!55] (0,.20)--(0,-12.75);
\node at (1.50,.48) {30};\draw[SynLine!55] (1,.20)--(1,-12.75);
\node at (2.50,.48) {31};\draw[SynLine!55] (2,.20)--(2,-12.75);
\node at (3.50,.48) {32};\draw[SynLine!55] (3,.20)--(3,-12.75);
\node at (4.50,.48) {33};\draw[SynLine!55] (4,.20)--(4,-12.75);
\node at (5.50,.48) {34};\draw[SynLine!55] (5,.20)--(5,-12.75);
\node at (6.50,.48) {35};\draw[SynLine!55] (6,.20)--(6,-12.75);
\node at (7.50,.48) {36};\draw[SynLine!55] (7,.20)--(7,-12.75);
\node at (8.50,.48) {37};\draw[SynLine!55] (8,.20)--(8,-12.75);
\node at (9.50,.48) {38};\draw[SynLine!55] (9,.20)--(9,-12.75);
\node at (10.50,.48) {39};\draw[SynLine!55] (10,.20)--(10,-12.75);
\node at (11.50,.48) {40};\draw[SynLine!55] (11,.20)--(11,-12.75);
\node at (12.50,.48) {41};\draw[SynLine!55] (12,.20)--(12,-12.75);
\node at (13.50,.48) {42};\draw[SynLine!55] (13,.20)--(13,-12.75);
\node at (14.50,.48) {43};\draw[SynLine!55] (14,.20)--(14,-12.75);
\node at (15.50,.48) {44};\draw[SynLine!55] (15,.20)--(15,-12.75);
\node at (16.50,.48) {45};\draw[SynLine!55] (16,.20)--(16,-12.75);
\node at (17.50,.48) {46};\draw[SynLine!55] (17,.20)--(17,-12.75);
\node at (18.50,.48) {47};\draw[SynLine!55] (18,.20)--(18,-12.75);
\node at (19.50,.48) {48};\draw[SynLine!55] (19,.20)--(19,-12.75);
\node at (20.50,.48) {49};\draw[SynLine!55] (20,.20)--(20,-12.75);
\node at (21.50,.48) {50};\draw[SynLine!55] (21,.20)--(21,-12.75);
\node at (22.50,.48) {51};\draw[SynLine!55] (22,.20)--(22,-12.75);
\node at (23.50,.48) {52};\draw[SynLine!55] (23,.20)--(23,-12.75);
\node at (24.50,.48) {53};\draw[SynLine!55] (24,.20)--(24,-12.75);
\node at (25.50,.48) {54};\draw[SynLine!55] (25,.20)--(25,-12.75);
\node at (26.50,.48) {55};\draw[SynLine!55] (26,.20)--(26,-12.75);
\node at (27.50,.48) {56};\draw[SynLine!55] (27,.20)--(27,-12.75);
\draw[SynLine!55] (28,.20)--(28,-12.75);
\node[anchor=east,align=right,text width=5.5cm,inner sep=0pt] at (-.25,-0.00) {Desarrollo E1};
\node[anchor=east,align=right,text width=5.5cm,inner sep=0pt] at (-.25,-0.48) {Marcha blanca E1};
\node[anchor=east,align=right,text width=5.5cm,inner sep=0pt] at (-.25,-0.96) {Paso a producción E1};
\node[anchor=east,align=right,text width=5.5cm,inner sep=0pt] at (-.25,-1.44) {Desarrollo E2};
\node[anchor=east,align=right,text width=5.5cm,inner sep=0pt] at (-.25,-1.92) {Marcha blanca E2};
\node[anchor=east,align=right,text width=5.5cm,inner sep=0pt] at (-.25,-2.40) {Producción E2 / Operación};
\node[anchor=east,align=right,text width=5.5cm,inner sep=0pt] at (-.25,-2.88) {Operación de ambos alcances};
\fill[SynNavy!75] (0.10,-3.01) rectangle (27.90,-2.75);
\node[anchor=east,align=right,text width=5.5cm,inner sep=0pt] at (-.25,-3.36) {Hitos E-25};
\node[anchor=west] at (.2,-3.36) {H1--H12: meses 2--21 (figura anterior)};
\node[anchor=east,align=right,text width=5.5cm,inner sep=0pt] at (-.25,-3.84) {1.1 Dirección y control};
\node[anchor=east,align=right,text width=5.5cm,inner sep=0pt] at (-.25,-4.32) {1.2 Plataforma híbrida};
\node[anchor=east,align=right,text width=5.5cm,inner sep=0pt] at (-.25,-4.80) {1.3 Auditoría y seguridad};
\node[anchor=east,align=right,text width=5.5cm,inner sep=0pt] at (-.25,-5.28) {1.4 Núcleo Operacional};
\node[anchor=east,align=right,text width=5.5cm,inner sep=0pt] at (-.25,-5.76) {1.5 Núcleo de Servicios};
\node[anchor=east,align=right,text width=5.5cm,inner sep=0pt] at (-.25,-6.24) {1.6 Núcleo Administrativo};
\node[anchor=east,align=right,text width=5.5cm,inner sep=0pt] at (-.25,-6.72) {1.7 Canales e integración};
\node[anchor=east,align=right,text width=5.5cm,inner sep=0pt] at (-.25,-7.20) {1.8 Migración e historia (+ C)};
\node[anchor=east,align=right,text width=5.5cm,inner sep=0pt] at (-.25,-7.68) {1.9 Innovaciones};
\fill[SynBlue!65] (1.10,-7.81) rectangle (8.90,-7.55);
\node[anchor=east,align=right,text width=5.5cm,inner sep=0pt] at (-.25,-8.16) {1.10 Verificación (+ C)};
\node[anchor=east,align=right,text width=5.5cm,inner sep=0pt] at (-.25,-8.64) {1.11 Formación e implantación};
\node[anchor=east,align=right,text width=5.5cm,inner sep=0pt] at (-.25,-9.12) {1.12 Servicio y salida};
\fill[SynBlue!65] (0.10,-9.25) rectangle (27.90,-8.99);
\node[anchor=east,align=right,text width=5.5cm,inner sep=0pt] at (-.25,-9.60) {IN1: solicitud / piloto / despliegue};
\node[anchor=east,align=right,text width=5.5cm,inner sep=0pt] at (-.25,-10.08) {IN2: preparación / campaña};
\fill[SynBlue!30] (1.10,-10.21) rectangle (8.90,-9.95);
\node[anchor=east,align=right,text width=5.5cm,inner sep=0pt] at (-.25,-10.56) {IN3: extracción asistida};
\node[anchor=east,align=right,text width=5.5cm,inner sep=0pt] at (-.25,-11.04) {IN4: medición e incentivo};
\node[anchor=east,align=right,text width=5.5cm,inner sep=0pt] at (-.25,-11.52) {IN5: pasaporte ambiental};
\node[anchor=east,align=right,text width=5.5cm,inner sep=0pt] at (-.25,-12.00) {Mesa y guardia prestadas};
\fill[SynBlue!30] (0.10,-12.13) rectangle (27.90,-11.87);
\node[anchor=east,align=right,text width=5.5cm,inner sep=0pt] at (-.25,-12.48) {Salida y acompañamiento};
\fill[SynBlue!30] (24.10,-12.61) rectangle (27.90,-12.35);
\end{tikzpicture}
\par\smallskip{\fontsize{9}{11}\selectfont Oscuro: fase contractual; azul: ventanas numéricas del diccionario; claro: detalle de filas ya incluidas; rombo: objetivo de hito.}
\caption{Gantt de referencia: meses 29--56; entregables, fases y cobertura}\label{fig:t14-gantt-2}
\end{figure}
\end{paginaHorizontalSynaptix}
\subsection{Seguimiento y control de integración}
El control quincenal registra resultados recibidos, pendientes, cambios y desviaciones. Avance físico corresponde a evidencia y aceptación por paquete, no a HH consumidas o porcentajes subjetivos. Modificar, subdividir o retirar resultados exige actualizar diccionario, trazabilidad, recursos, red, Gantt y recepción conjuntamente. Los resultados compartidos y los trabajos activados contra reserva se contabilizan una sola vez.

Antes de emitir, se concilia T-12 con los códigos agrupados; se completa la selección regional y su impacto en SD4/T-11; se cierra método de esfuerzo, capacidad, CPM/PERT y recepción E2 en T-15/T-18; y se sustituyen las marcas de referencia por localizadores efectivos. Los compromisos de SD3/SD5/SD6 se conservan y se reflejan en los paquetes que los realizan. Ninguna de estas comprobaciones atribuye ejecución, recepción contractual o revisión humana aún no realizada.